\documentclass[11pt,a4paper,oneside,openany]{book}
\usepackage[T1]{fontenc}
\usepackage[utf8]{inputenc}
\usepackage{lmodern,amsmath,amssymb,amsthm,mathtools,microtype}
\usepackage[margin=28mm]{geometry}
\usepackage{booktabs,array,longtable,enumitem,xcolor}
\usepackage[colorlinks=true,linkcolor=blue!45!black,citecolor=blue!45!black,urlcolor=blue!45!black]{hyperref}
\hypersetup{pdftitle={Understanding the linear degree gap},pdfauthor={Soren Riis (mathematical direction, supervision and editing)},pdfsubject={AI-generated tutorial textbook on uniform polynomial consequences and the lifting principle},pdfkeywords={proof complexity, uniform consequences, degree profiles, linear lower bounds, AI-generated textbook}}
\usepackage{fancyhdr}
\pagestyle{fancy}\fancyhf{}\fancyhead[L]{\small Uniform consequences and the lifting principle}\fancyhead[R]{\small\thepage}
\renewcommand{\headrulewidth}{0.3pt}\setlength{\headheight}{14pt}
\setlength{\emergencystretch}{2em}
\setlist{nosep,leftmargin=*,topsep=5pt}
\newtheorem{theorem}{Theorem}[chapter]
\newtheorem{lemma}[theorem]{Lemma}
\newtheorem{proposition}[theorem]{Proposition}
\newtheorem{corollary}[theorem]{Corollary}
\theoremstyle{definition}\newtheorem{definition}[theorem]{Definition}
\newtheorem{example}[theorem]{Example}
\theoremstyle{remark}\newtheorem{remark}[theorem]{Remark}
\newcommand{\NS}{\mathrm{NS}}\newcommand{\PC}{\mathrm{PC}}
\newcommand{\NN}{\mathbb N}\newcommand{\N}{\mathbb N}\newcommand{\Z}{\mathbb Z}\newcommand{\Q}{\mathbb Q}\newcommand{\Fp}{\mathbb F_p}
\DeclareMathOperator{\supp}{supp}\DeclareMathOperator{\im}{im}\DeclareMathOperator{\Span}{span}
\DeclareMathOperator{\rec}{Rec}\DeclareMathOperator{\Soc}{Soc}\DeclareMathOperator{\Top}{Top}\DeclareMathOperator{\LIS}{LIS}
\newcommand{\join}{\mathbin{\vee_*}}
\newcommand{\Ftwo}{\mathbb F_2}
\newcommand{\D}{\mathcal D}\newcommand{\Del}{\mathcal D}\newcommand{\Lift}{\mathsf L}
\DeclareMathOperator{\Inj}{Inj}
\widowpenalty=10000\clubpenalty=10000
\newtheorem{exercise}[theorem]{Exercise}
\newenvironment{solution}{\begin{proof}[Solution]}{\end{proof}}
\newenvironment{checkpoint}{\begin{quote}\small\textbf{Checkpoint.}\ }{\end{quote}}
\numberwithin{equation}{chapter}
\usepackage{amscd}
\title{Understanding the linear degree gap\\[8pt]
\Large Uniform polynomial consequences and the lifting principle\\[10pt]
\large An AI-generated tutorial textbook\\[4pt]
\normalsize Making machine-generated proofs readable}
\author{\normalsize Mathematical direction, supervision and editing by\\[4pt]
\large S\o ren Riis}
\date{Review edition aligned with the paper of 4 November 2026}
\begin{document}
\frontmatter
\hypersetup{pageanchor=false}
\maketitle
\hypersetup{pageanchor=true}

\chapter*{Preface: mathematics, machines and exposition}
\addcontentsline{toc}{chapter}{Preface: mathematics, machines and exposition}
\label{tx:preface:provenance}
\begingroup\fontsize{10.5}{12.6}\selectfont
AI is changing how mathematics can be explored, checked and explained.
This project suggests a corresponding way to present mathematical research:
an interactive paper, with further explanations attached to its arguments,
and a tutorial textbook which develops the ideas at a human pace.
The two formats serve different purposes. The paper states the results
and their proofs concisely; the textbook supplies the definitions,
examples and intermediate reasoning needed to understand the whole route.

Machine-generated proofs can be hard to follow. They may combine ideas
from several research communities without explaining why a construction
was chosen or how its pieces fit together. The initial proofs generated
by Astra and Claude also omitted several key references to prior human
work. Their mathematical knowledge reflects a body of human research,
although the outputs do not establish which particular sources were
used in training or how they influenced the constructions. A formally
checked argument therefore also needs a readable account of the
mathematics and its antecedents. This book provides that account for
the linear degree gap.

This is an \emph{AI-generated tutorial textbook}, developed under my
mathematical direction, supervision and editing. Astra supplied much of
the drafting and restructuring. Astra and Claude also contributed
substantially to the proof search and formalisation. I selected the
programme, supplied my earlier work and mathematical input, challenged
arguments, read successive chapters and requested clearer explanations
and worked examples. This continuing exchange shaped both the proof
development and its presentation.

The published mathematical starting point is my LICS~2008
paper~\cite{Riis}; the 53-page 2009 report supplies a longer account
\cite{RiisLong}. The conference paper established a four-way
classification, or \emph{quadrichotomy}, of uniform algebraic
refutation complexity: $\mathrm T_1,\ldots,\mathrm T_4$, with
periodicity, NS--PC type agreement and undecidability for every
nontrivial union of types, including the no-finite-model promise.
The logarithmic and double-logarithmic scales used
Kraj\'i\v cek's Theorem~3.5~\cite{Krajicek}, based on Ajtai's
symmetric-system work~\cite{Ajtai}.
My joint journal papers with Meera Sitharam develop uniform
translations~\cite{RiisSitharamTautologies} and uniformly generated
permutation submodules over fields of characteristic
zero~\cite{RiisSitharamModules}. A separate 1997 announcement proposed
a nonconstant-to-linear degree gap. Its inducing-down argument has
an unresolved gap, as I recorded in the 2009 report
\cite[p.~50, footnote~3; p.~53, reference~33]{RiisLong}; this concerns
the proof and does not refute the asserted theorem. My later
unpublished notes pursued stronger bounds and suggested lifting
separating vectors between ranks, with important steps incomplete.
The aim here was to obtain complete proofs of linear degree bounds in
the universe size, in every fixed positive characteristic, and to treat
uniform polynomial consequences as well as refutations.

The central proof search included, in my recollection, roughly two to
three days of concentrated work with Astra and Claude, alongside my
interactions. Checking, repairs, extensions and exposition continued
afterwards. When the construction reached an $\Omega(\sqrt n)$ bound in
characteristic two, I pressed for an $\Omega(n)$ bound and an argument
covering all positive characteristics. The difficult work was to build
one lifting mechanism compatible with all the required operations.

\newpage
The development and subsequent comparison with the literature brought
together proof complexity, model theory, modular representation theory
and algebraic combinatorics. Some of these connections were unfamiliar
to me; the assistants helped bring them into the discussion.
Kraj\'i\v cek's earlier work already treated uniform polynomial targets,
Euler structures and bounded-degree NS--PC comparison, including the
NS-span plateau argument~\cite{KrajEuler}\cite[Theorem~5.5]{Krajicek}. The wider context
includes Church--Ellenberg--Farb's FI framework for representations as
the number of labels grows, Harman's modular periodicity in a linear
range, and Hone and coauthors' explicit periodic functors through
semisimplification~\cite{CEF,Harman,HKPSYZ}. Harman's equivalences
already respect induction, restriction, tensor products and
duality~\cite[Propositions~2.7 and~2.12]{Harman}. The construction
here proves the required closure preservation and target detection
on the actual polynomial maps with its own quantitative bounds.
Integral centraliser and
rook combinatorics provide another connection~\cite{BDM}: realising
compatible linear maps as combinations of label permutations. These
communities study stability as the number of labels grows and the
construction of maps between representations.
Chapter~\ref{tx:chap:contract} explains these antecedents; the finite
constructions required for our polynomial proof rules are proved here.

The principal construction developed in this project is in
Chapters~\ref{tx:ch:linear-layers}--\ref{tx:ch:prime}: compatible
corrections reuse one retained pool of labels, and the odd-characteristic
argument completes the same linear budget.
Chapter~\ref{tx:chap:assemble-transfer} transports the actual axiom,
multiplication and indexed-target maps.
Chapter~\ref{tx:sec:comparison} independently supplies the comparison's
uniform chain bound. The four qualitative types are inherited from my
earlier work; the present quantitative bounds are linear in
$\mathrm T_2$, linear on the hard residues in $\mathrm T_3$, and give
finite linear peaks in $\mathrm T_4$.
Chapter~\ref{tx:sec:undecidability} gives explicit reductions separating
any two types under the no-finite-model promise already stated in
LICS~2008, refining its every-union formulation. This treatment was
developed interactively by me and Astra from my earlier undecidability
arguments. The NS proof system is due to Beame, Impagliazzo,
Kraj\'i\v cek, Pitassi and Pudl\'ak~\cite{NSFoundations}; PC is due
to Clegg, Edmonds and Impagliazzo~\cite{PCFoundations}.
The induction-chain separation used in
Chapter~\ref{tx:ch:running-order} is due to Buss and
Pitassi~\cite{BussPitassiInduction}; the design method is developed
in Buss~\cite{BussDesign}, and the modular bijection antecedent is
Beame--Riis~\cite{BeameRiis}. The compression discussion belongs to
separate companion work.

I developed this tutorial interactively by reading drafts, identifying
steps I could not follow, proposing abstractions and requesting examples
that would carry the reader through the general argument. Alongside
supplying my earlier manuscripts and guiding the proof search, an
important part of my contribution was to identify relevant prior human
work, check the connections against the literature, and ensure that
these contributions received proper credit. My intention is to make
the textbook available on GitHub alongside the interactive
version of the paper, and potentially use it for lectures and
presentations. The contribution to exposition is this coordinated
combination: a scientific paper with expandable explanations and a
longer, progressively developed tutorial.

I stated and proved the following undecidability theorem: for every
nonempty proper subset $A\subsetneq\{1,2,3,4\}$, it is undecidable
whether a uniform family arising from a first-order sentence belongs
to one of the types $\mathrm T_i$ with $i\in A$.
This theorem appears in my LICS~2008 paper, with further details in
the 2009 report~\cite{Riis,RiisLong}.
The no-finite-model promise is explicitly part of its published
statement~\cite[Theorem~1]{Riis}. Its Lemma~8 gives a
$\mathrm T_1$-versus-prescribed-type switch; the controlled-join
construction in Chapter~\ref{tx:sec:undecidability} supplies the
remaining pairwise reductions. This comparison concerns that
published formulation and inspected construction.
The core degree, comparison and undecidability results now have Lean
formalisations and independent source, statement and axiom audits.
Newly added teaching examples and the topological reformulation are
written deductions; they are not additional Lean theorem endpoints. The book explains
the mathematics and its provenance, while making explicit how its text
was generated and how I directed and revised it.

\par\endgroup

\chapter*{How to read this book}
\addcontentsline{toc}{chapter}{How to read this book}
This book develops the mathematics behind \emph{Degree profiles of
uniform polynomial consequences}. It begins with actual algebraic proof
lines and ends with a self-contained account of the linear degree gap.
The intended reader is comfortable with polynomials, coefficient
vectors, linear combinations and spans. We recall kernels and quotient
spaces and introduce tensor products through explicit calculations.
Nullstellensatz proofs, Polynomial Calculus and the required label
actions are introduced here. No prior knowledge of modular
representation theory, homological algebra or of the earlier
complexity-gap paper is assumed. Understanding the conclusion is easier
than checking the complete proof: the common-witness construction in
Chapters~\ref{tx:ch:linear-layers}--\ref{tx:ch:prime} requires sustained
work with the displayed maps.

The central problem is
\[
                    \Gamma_n\vdash G_n,
\]
where both the axioms and the polynomial expression to be derived come
from fixed rules as the universe size $n$ changes. A target with a
bounded-support index can be included as well. Deriving $1$ is the
important refutation special case, but it is not the only case.
Here $\Gamma_n$ is the set of input polynomials, each interpreted as
an equation equal to zero; $G_n$ is the target polynomial, whose equation
we want to derive. The symbol $\vdash$ means that an algebraic proof
exists. The two proof systems and their degree conventions are defined
in Chapter~\ref{tx:ch:consequences}.

\section*{The question which organises the whole proof}
Suppose a uniform system at size $n$ has a proof with degree at most
$d$. Can that degree budget also work at $n+q$, for a suitable
prime-power block $q$? For an invariant uniform target $G_n$, the
picture is
\[
 \Gamma_n\vdash_{\le d}G_n
 \quad\Longrightarrow\quad
 \Gamma_{n+q}\vdash_{\le d}G_{n+q}
 \quad\Longrightarrow\quad
 \Gamma_{n+2q}\vdash_{\le d}G_{n+2q}\quad\Longrightarrow\cdots.
\]
Here the proof system is fixed as either NS or ordinary PC. The
axioms and target follow their prescribed uniform rules at each size;
the larger proof may have a different list of lines. The resource
preserved by lifting is the degree of the whole proof.
In this picture, \emph{invariant} means that every permutation of the
universe labels fixes the target polynomial. A target with distinguished
indices has a different invariance condition, defined in
Chapter~\ref{tx:ch:uniformity}.

This is the main question of the book. First we describe exactly
what has to be transported. Then we construct progressively stronger
lifts, measuring how large $n$ must be in relation to $d$. Finally
we reverse that numerical condition to obtain degree lower bounds.
General indexed targets use the same membership principle, with an
additional argument to detect their images; their individual vectors
need not be literal images under the invariant-target lift.

\section*{The linear-algebra picture behind lifting}
At a fixed degree budget $d$, ordinary PC is a calculation in a finite
vector space. Start from usable axioms, close under linear combinations,
and multiply a vector by a variable whenever that vector has degree at
most $d-1$. This gives a subspace $C_d(n)$. NS uses a different subspace
$I_d(n)$, spanned by the permitted monomial--axiom products. The target
question is membership in one of these spaces.

The difficult theorem relates the membership calculations at $n$ and
$n+q$, where $q=p^a$ is a suitably large power of the characteristic.
One exact operation deletes the extra block of labels. A related linear
operator lifts smaller consequences into the larger universe. The
construction has to preserve actual maps, intersections and images;
a count of dimensions would not suffice.

For a first reading of this transition, begin with the worked
multiplication calculation in Section~\ref{tx:sec:multiplication-registers}.
Then follow the four operations making one closure stage in
Section~\ref{tx:sec:grammar}. Section~\ref{tx:sec:target-reflection}
explains why preserving this calculation still leaves a separate
target question. Chapter~\ref{tx:chap:deletion} introduces the objects
in Theorem~\ref{tx:thm:deletion} one by one and gives a guide to each
clause. At that point the theorem is a precise statement of what
the later construction will prove, not a fact to accept without proof.

There are two different operations which should never be confused.
Multiplication by a variable takes degree-$d-1$ vectors to degree-$d$
vectors \emph{within one universe}. The lifting operator takes
polynomials from the smaller universe to the larger one \emph{without
increasing degree}. It is linear, but is not generally a homomorphism
of polynomial rings.

\section*{Why one proof gives a constant bound on a progression}
Fix one degree $d$ and one eligible starting size $n$. Lifting gives
proofs with that same budget at $n+q,n+2q,\ldots$. The absolute budget
stays $d$. Once that starting proof has been selected, $d$ is a
\emph{constant} for the whole progression, even if it was originally
described as, say, a small multiple of $\log n$. At the later ranks
we keep its numerical value; we do not recompute the logarithm.
Its ratio to universe size tends to zero:
\[
                         \frac{d}{n+\ell q}\longrightarrow0.
\]
The constant bound initially belongs to one arithmetic progression.
It is not yet a constant bound over all universe sizes. Two-way transfer
at every lower threshold shows exactly which finite minimum degree
recurs on that progression. The four types describe how those
progressions fit together.

The lower-bound argument reverses this observation. If a finite degree
cannot recur, it must have occurred before transfer became available.
The theorem makes that availability threshold linear in $d$. Solving
this inequality for $d$ gives the linear lower bound. In the fluctuating
case, a separate bound on the length of the period locates finite linear
peaks. These are different uses of the same transfer theorem.

\section*{The stages are improvements in how far we can lift}
The earlier refutation classification gives logarithmic NS and
double-logarithmic PC lower bounds~\cite{Riis,RiisLong}. The underlying
transfer estimates summarised in the LICS paper have exponential and
double-exponential starting ranks as functions of the fixed degree.
These are historical refutation estimates, obtained there from
substantial earlier work; they are not elementary lemmas reconstructed
in this book.

The constructions developed here make the starting rank much smaller:
\begin{center}
\begin{tabular}{@{}>{\raggedright\arraybackslash}p{.43\linewidth}>{\raggedright\arraybackslash}p{.25\linewidth}>{\raggedright\arraybackslash}p{.22\linewidth}@{}}
\toprule
Lifting stage & Starting rank for degree $d$ & Lower-bound scale\\
\midrule
Earlier NS refutation result & at most $2^{cd}$ & $\log n$\\
Earlier PC refutation result & at most $2^{2^{cd}}$ & $\log\log n$\\
Characteristic-two construction with fresh labels & $O(d^2)$ & $\sqrt n$\\
One retained pool, every positive characteristic & $O(d)$ & $n$\\
\bottomrule
\end{tabular}
\end{center}
The constants depend on the fixed family and field. A small degree
means small enough to meet the relevant starting-rank inequality.
Chapter~\ref{tx:chap:lift} gives an elementary counting explanation
for the prime-power steps. Chapter~\ref{tx:ch:quadratic} supplies the
first fully worked lifting construction. The central
advance in Chapters~\ref{tx:ch:linear-layers}--\ref{tx:ch:prime} is to
replace its quadratic label cost by a linear one, and to make that
construction work also in odd characteristic. The T4 conclusion uses
an additional linear bound on the period and gives finite linear peaks.

This development connects proof complexity with modular representation
theory and algebraic combinatorics. Those communities already study
representations built from bounded numbers of labels and periodic
behaviour as the universe grows~\cite{CEF,Harman,HKPSYZ,BDM}.
Section~\ref{tx:sec:stability-context} identifies the close antecedents.
The task here is to prove transport for the actual polynomial proof
rules and targets, with all the required finite constructions supplied
in the book.

\section*{A layered route through the construction}
The first part replaces proof syntax by consequence spaces, then works
through a Boolean first-order example over a fixed finite field: a
strict order in which every element chooses a smaller predecessor.
Chapter~\ref{tx:ch:running-order} proves the elementary growing-degree
obstruction and shows exactly how comparison and lifting strengthen it
to a linear bound. The predecessor may be many-to-one; the example
uses no matching or injection constraint. The second
part states the precise transfer contract and derives the four profiles
from it; it also identifies the desired lifting map. The third part
constructs the common algebraic projector which fulfils that contract,
first in characteristic two and then in every positive characteristic.
The fourth part puts the target and consequence maps into the same
argument and proves the independent NS--PC comparison. The final part
works through examples, dual witnesses, matching obstructions and the
logical undecidability consequence.

Within a technical chapter the order is always the same: say which
property is needed, explain why it is enough, prove the structural
lemma, and then inspect the calculation that realises it. The full
normalisation, zero-marginal, rook and odd-characteristic arguments
remain in the main chapters. Exercises have worked solutions at the
end, but the proofs of the main results do not depend on doing them.

A first reading can follow the contracts, examples and checkpoints.
A second reading can fill in each finite algebraic construction in turn.
This is a way of learning the complete proof, not permission to regard
its hard inputs as unproved assumptions.

When returning to a technical section, keep track of the space on
which an operator acts and which parameters are fixed. Powers of a
linear operator mean repeated composition. A polynomial degree, a
tuple length and a support budget measure different things, even
when the argument relates them by an inequality. Local reminders at
the main changes of construction, together with the notation guide
in the final chapter, identify these roles.

The most direct first route through the opening is to read the basic
template definition in Chapter~\ref{tx:ch:uniformity}, then the target
request and support budget in
Sections~\ref{tx:sec:target-request}--\ref{tx:sec:target-budgets}.
The orbit-sum and compression discussion can be read later. In
Chapter~\ref{tx:ch:closure}, follow the closure recurrence and its maps;
the coefficient-list size interpretation is a separate application.
Then follow the running order example in
Chapter~\ref{tx:ch:running-order}. Its final steps specify the two
general inputs still to be constructed: uniform comparison and
degree-preserving lifting.
Chapter~\ref{tx:chap:contract} then names the three obligations which
organise the rest of the proof: identify the data, preserve the
calculation, and reflect target vanishing. The chapters below refer
back to these obligations as each construction is supplied.

\section*{Specialist tools and a route through the proof}
The proof uses several specialist languages. Its representation-theoretic
part concerns permutation modules, group algebras, annihilator ideals
and nilpotent chains. Its transport language consists of tensor products,
quotients, naturality and exactness. Its finite constructions use common
partial inverses, ordered projectors, marginals, sections and rook
coordinates. Dual functionals and matching designs supply obstructions;
definability and effective reductions explain the logical consequences.
The background tools are developed at their first use and recalled
when their role changes. The following map identifies the main
transitions between mathematical subjects. It also helps a reader
returning to the proof locate the relevant definitions.

\begin{longtable}{@{}p{.23\linewidth}p{.47\linewidth}p{.22\linewidth}@{}}
\toprule
Tool & Concrete use in this book & First development\\
\midrule
Ideals and degree cutoffs & Polynomial combinations, with the
allowed products constrained by degree & Chapters~\ref{tx:ch:consequences}--\ref{tx:ch:closure}\\
Orbits and stabilisers & Grouping coefficient patterns under relabelling
& Chapter~\ref{tx:ch:uniformity}\\
Scalar extension & Solving the same coordinate problem with auxiliary
field coefficients & Chapter~\ref{tx:ch:closure}\\
Tensor products & Coordinates for paired inputs, with a diagonal
label action & Chapter~\ref{tx:chap:contract}\\
Naturality and exactness & Commuting maps and preserving the actual
kernel, image and quotient calculations & Chapter~\ref{tx:chap:contract}\\
Nilpotent chains & Reading which directions a block operator kills
or leaves in its quotient & Chapter~\ref{tx:chap:deletion}\\
Frobenius and sections & Prime-power counting and a chosen lift of
each retained vector & Chapter~\ref{tx:chap:lift}\\
Common partial inverses & One group-algebra expression certifying
the required projector on several spaces & Chapters~\ref{tx:chap:lift}--\ref{tx:ch:square-zero}\\
Marginals and layer sections & Compatible coordinate sums and
corrections using one retained label pool & Chapter~\ref{tx:ch:linear-layers}\\
Rook coordinates & Encoding overlap patterns and solving a
triangular coefficient system & Chapter~\ref{tx:ch:rooks}\\
Dual functionals and pairings & Coefficient tests separating targets
from consequence spaces & Chapters~\ref{tx:ch:running-order}--\ref{tx:chap:contract}\\
Matching designs & Bounded-degree coefficient tests and their
extension rules & Chapter~\ref{tx:app:matching}\\
Definability and reductions & The model-theoretic interpretation and
the transfer of an undecidable computation problem & Chapters~\ref{tx:sec:model-theory}--\ref{tx:sec:undecidability}\\
\bottomrule
\end{longtable}

\paragraph{Background reading.}
For ordinary algebra, Lang~\cite{LangAlgebra} has chapters on
modules, algebraic extensions, a single endomorphism, bilinear forms
and tensor products. Alperin~\cite{AlperinLocal} introduces the modular
representation theory of finite groups. The exactness and reflection
vocabulary is described in the Stacks Project~\cite{StacksExact}; here
we develop it through kernels and images of the displayed linear maps.
Cutland~\cite{Cutland}, especially the chapters on computable functions,
universal programs, undecidability and reducibility, gives background
for the counter-program reductions. The references in
Sections~\ref{tx:sec:stability-context} and~\ref{tx:sec:model-theory}
locate the representation-stability and Euler-structure connections.
These are reading routes for the terminology. The common witness,
its natural map compatibility, the rook correction and the ordinary-PC
extension invariant are established explicitly in this book.

The opening route above is for learning the problem: it reaches the
running example and the transport contract quickly. The book's chapter
order is for learning the construction: it explains the destination
before constructing the engine, including the simpler quadratic stage.
The following dependency order is for checking that every input of the
linear theorem has been discharged; it returns to classification only
after transfer is proved. These purposes account for their different
orders. Each item
states a concrete conclusion to check, rather than asking the reader to
accept a general stability slogan.
\begin{enumerate}
\item \textbf{Replace proofs by linear algebra.}
In Chapters~\ref{tx:ch:consequences}--\ref{tx:ch:closure}, check that
$I_d$ and $C_d$ describe \emph{exactly} the legal degree-$d$ proofs.
In particular, PC intersects with $P_{\le d-1}$ after cancellation.
Check stable coefficients and the separate degree-zero convention.
\item \textbf{Identify the certificate for exactness.}
In Chapter~\ref{tx:ch:square-zero}, verify that one common identity
$zwz=z$ gives the projector $(1-zw)(1-wz)$ and exact deletion on
the entire common annihilator class. The two-point failure of exactness
explains why the moved-basis calculation alone is insufficient.
\item \textbf{Construct that common certificate in characteristic two.}
In Chapter~\ref{tx:ch:linear-layers}, check the partition of unity,
the signed layer sections and the correction interface
(Lemma~\ref{tx:lem:layer-interface}). In Chapter~\ref{tx:ch:rooks},
check strengthened forcing and triangular realisation. These prove that
one correction repairs length $b$ while annihilating every shorter
length, on the same pool of retained labels.
\item \textbf{Check the changes for an arbitrary prime.}
Chapter~\ref{tx:ch:prime} replaces $z$ by its top power
$u=z^{p-1}$ in the witness identity. Check purity on tuple generators,
the ordered projector factors, the $p$-level sections and the actual
transpose relabelling. Purity is not assumed for arbitrary nilpotent
representations.
\item \textbf{Return to the prescribed polynomials.}
In Chapter~\ref{tx:chap:assemble-transfer}, check the axiom and
multiplication maps entry by entry, then the full target map $J$.
For indexed targets, check the nonzero-subquotient pairing
(Lemma~\ref{tx:lem:detect}); exactness alone does not reflect vanishing.
The pairing doubles the original monomial support $kd$.
\item \textbf{Read off the lower bounds.}
Now check Chapter~\ref{tx:ch:profiles}. The linear onset controls
nonrecurring values in T2 and on T3 hard residues. The period bound
has the additional job of locating \emph{finite} T4 peaks. This step
is an argument about integer sequences after transfer has been proved.
\item \textbf{Check the separate comparison theorem.}
Chapter~\ref{tx:sec:comparison} uses a strict-chain bound and an NS
plateau to contain the whole PC closure. It proves same-type agreement;
it is not the source of the linear PC lower bound.
\end{enumerate}

Chapter~\ref{tx:ch:quadratic} is a complete alternative construction in
characteristic two, useful for seeing why spending fresh labels gives
only a square-root bound. The linear route does not depend on that
construction or its normalisation lemma. Chapters~\ref{tx:chap:deletion}
and~\ref{tx:chap:lift} state and explain the operation before its
coefficients are constructed; their promised input is discharged by
Chapters~\ref{tx:ch:square-zero}--\ref{tx:ch:prime}.
The later example, matching and undecidability chapters have their own
additional arguments.

The core formal results have independently checked Lean source builds,
theorem statements, axioms and dependencies, as recorded with the paper.
That evidence does not certify every new pedagogical formulation in this
book. A human check can follow the finite calculations above without
treating either the formalisation or an external representation-stability
theorem as a substitute for a missing step.

\section*{Notation which keeps the roles separate}
\begin{center}
\begin{tabular}{@{}ll@{}}
\toprule
Symbol & Meaning\\
\midrule
$p$ & Fixed prime characteristic of the coefficient field\\
$n$ & Retained, smaller universe size\\
$q=p^a$ & Size of the added/deleted block; $a$ is an exponent\\
$N=n+q$ & Larger universe size\\
$d$ & Degree budget for the whole proof\\
$D_G$ & Fixed upper bound on the target polynomial degree\\
$k,s,t$ & Supports of a variable, axiom index, target index\\
$r=R(d)$ & One support budget once $d$ has been fixed\\
$N(d)$ & Onset beyond which a degree-$d$ threshold transfers\\
$Q(d)$ & A suitable prime-power period for that threshold\\
$I_d(n),C_d(n)$ & NS and ordinary-PC consequence subspaces\\
$\mathsf L$ & Degree-preserving lift from $n$ to $N$\\
$\Del$ & Exact deletion operation from larger modules to smaller ones\\
\bottomrule
\end{tabular}
\end{center}
In particular, the support budget $r$ is \emph{not} the exponent $a$
in $p^a$, and the variable support $k$ is neither of them. The number
of polynomial variables can be of order $n^k$ and is not itself $n$.

\section*{Reading alongside the current paper}
This edition follows the paper revised on 4 November 2026. The paper
is the concise statement of the research; this book supplies the slower
route through its arguments. In particular $N(d)$ now denotes the
same periodicity onset in both texts. Ordinary $N=n+q$ still denotes
a larger universe size. The older refutation-only tutorials are
historical guides, rather than additional prerequisites.

\begin{center}\small
\begin{tabular}{@{}>{\raggedright\arraybackslash}p{.35\textwidth}>{\raggedright\arraybackslash}p{.59\textwidth}@{}}
\toprule
Question while reading the paper & Where to turn here\\
\midrule
What is uniform, and what does degree count? & Chapters 1--3; distinguish a full polynomial from its bounded support of distinguished labels.\\[3pt]
How does transfer imply the four cases? & Chapter~\ref{tx:ch:profiles}; compare \emph{thresholds}, then recurring exact values.\\[3pt]
Why does deletion preserve the proof? & Chapter~\ref{tx:chap:contract}; transport maps and intersections, not just dimensions.\\[3pt]
Where is the linear improvement? & Chapters~\ref{tx:ch:quadratic}--\ref{tx:ch:prime}; compare fresh labels with reuse of one common pool.\\[3pt]
Why is the prescribed target preserved? & Chapter~\ref{tx:chap:assemble-transfer}; the quotient obstruction must survive deletion.\\[3pt]
How does topology separate T1--T4? & Section~\ref{tx:sec:topological-spectrum}; the residue tree, compactness, and the binary worked example.\\[3pt]
What is formalised? & The accompanying Lean packet's scope and endpoint tables; the new topological reformulation has a written proof.\\
\bottomrule
\end{tabular}\end{center}

A short first pass is Chapters 1 and 3, the contract in
Chapter~\ref{tx:chap:contract}, Chapter~\ref{tx:ch:profiles}, and
Section~\ref{tx:sec:topological-spectrum}. Then return to the common
witness and rook constructions. Their purpose is to justify the
contract with a \emph{linear} label budget; the classification proof
itself is considerably simpler.

\tableofcontents
\mainmatter

\part{From proof lines to a finite linear closure}

\chapter{Polynomial consequences and proof degree}\label{tx:ch:consequences}
A polynomial equation can be used in two different ways. One can evaluate
it at a proposed solution, or one can combine it algebraically with other
equations. This book studies the second activity. Given equations
$f=0$ for $f$ in a set $\Gamma$, how complicated is a derivation of a
specified equation $G=0$? Our measure of complexity is the largest total
degree needed in the derivation.

The distinction between an equation and its left-hand polynomial is
worth keeping in view. We usually write $f$ instead of $f=0$, and write
$\Gamma\vdash G$ for deriving $G=0$. Thus a derivation of the polynomial
$1$ is a contradiction, since it asserts $1=0$. A derivation of $x+y$
need not be a contradiction: it may simply express a consequence shared
by all solutions of the original equations.

\section{The ambient polynomial ring}
Let $K$ be a field and let $V$ be a finite set of variable names. We work
in the ordinary polynomial ring $K[x_v:v\in V]$. A monomial is a product
$\prod_vx_v^{a_v}$ with nonnegative integer exponents. Its total degree
is $\sum_va_v$. The total degree of a nonzero polynomial is the largest
degree of a monomial with nonzero coefficient, after equal monomials
have been collected. We use the convention $\deg0=0$.

For example, over $\mathbb F_2$ the polynomial
\[
                   (xy+x)+(xy+1)=x+1
\]
has degree one. The two degree-two terms have cancelled. In a polynomial
ring over a field, nonzero products satisfy
$\deg(fg)=\deg f+\deg g$. We will use this fact when we define the
degree of NS certificates and PC derivations in the next two sections.

A system $\Gamma$ consists of polynomials regarded as equations equal
to zero. It may include Boolean equations $x_v^2-x_v$, but their presence
does not change the ambient ring. In particular, $x_v^2$ is still a
polynomial of degree two. Replacing it by $x_v$ requires a justified
algebraic step using $x_v^2-x_v$.

\section{Nullstellensatz certificates}
The algebraic proof system was introduced by Beame, Impagliazzo,
Kraj\'i\v cek, Pitassi and Pudl\'ak~\cite{NSFoundations}. We use its
polynomial-identity form, with degree measured on the complete products.
\begin{definition}[NS consequence]\label{tx:def:NS}
An NS derivation of $G$ from $\Gamma$ is a finite polynomial identity
\begin{equation}\label{tx:eq:NS}
                         G=\sum_{f\in F}a_ff,
                   \qquad F\subseteq\Gamma\text{ finite}.
\end{equation}
Its degree is the largest degree of a nonzero product $a_ff$.
Thus it has degree at most $d$ if
$\deg a_f+\deg f\le d$ for every nonzero summand. Zero summands can be
removed. The empty identity is a degree-zero derivation of $G=0$.
\end{definition}

The letters NS refer to Nullstellensatz, but the definition is simply
an identity which can be checked by multiplication and addition. There
is no appeal to an existence theorem hidden in a certificate. The
multiplier $a_f$ records how the input equation $f=0$ is used.

\begin{example}[A first certificate]\label{tx:eg:first-NS}
From $f_1=x+y$ and $f_2=y-z$, over any field, we derive $G=x+z$ by
\[
                       x+z=f_1-f_2.
\]
Both products have degree one. From $f_1=x$ and $f_2=1-x$ we instead
obtain $1=f_1+f_2$. The first identity expresses a consequence; the
second is a refutation.
\end{example}

The relevant degree is the degree of each product before the products
are added. A large-degree contribution cannot be ignored merely because
it cancels against another summand in the final identity. Otherwise
every refutation would appear to have degree zero, since its final
polynomial is the constant $1$.

\section{Ordinary Polynomial Calculus}
Polynomial Calculus was introduced by Clegg, Edmonds and
Impagliazzo~\cite{PCFoundations}. The rules below fix the ordinary
polynomial convention used throughout this book.
\begin{definition}[Ordinary-PC consequence]\label{tx:def:PC}
A PC derivation of $G$ is a finite sequence $P_1,\ldots,P_M$ ending in
$G$. A line may be an input polynomial from $\Gamma$, a linear
combination $\alpha P_i+\beta P_j$ of earlier lines, or a variable
multiple $x_vP_i$ of an earlier line. Coefficients $\alpha,\beta$ belong
to $K$. The zero polynomial is allowed as the empty linear combination.
The degree of the derivation is $\max_i\deg P_i$, with each line
collected as an actual polynomial.
\end{definition}

The word \emph{ordinary} stresses that there is no automatic Boolean
reduction and no rule which multiplies by an arbitrary polynomial in a
single step. A monomial multiple can be obtained by repeated variable
multiplication, but every intermediate line counts towards the degree.
The number of lines is unrestricted.

\begin{example}[A cancellation which enables the next step]\label{tx:eg:cancellation}
Over $\mathbb F_2$, consider the inputs
\[
                   f_1=xy+x,\qquad f_2=xy+1,\qquad f_3=xz.
\]
We want a PC derivation of $z$ in which every line has degree at most
two. This is a restriction on the whole proof, including its inputs and
intermediate lines. After listing these inputs, PC may continue with
\begin{align*}
 f_1+f_2&=x+1,\\
 z(x+1)&=xz+z,\\
 (xz+z)+f_3&=z.
\end{align*}
Every line stays within the chosen degree-two budget. Multiplication
by $z$ is always a PC rule, but here we must also check that its result
fits that budget. The preceding line $x+1$ has degree one, so
$z(x+1)=xz+z$ has degree two and is an admissible line in this
degree-two derivation. Multiplying one of the original quadratic inputs
by $z$ would instead produce a degree-three line. Expanding the same calculation into
one NS identity gives
\[
                          z=zf_1+zf_2+f_3,
\]
whose first two products have degree three. This is a concrete reason to
keep track of cancellations inside a PC proof.

The axioms are consistent: $x=y=1$ and $z=0$ is a solution. We have
proved the nonconstant consequence $z=0$, not a contradiction.
Exercise~\ref{tx:ex:small-gap} checks that the two minimum degrees in
this example really are different.
\end{example}

More generally, suppose we require every line to have degree at most
$d$, and the proof has derived a nonzero polynomial $P$. A variable
multiple stays within this degree budget only if $\deg P\le d-1$. This follows from
$\deg(x_vP)=\deg P+1$. Even if $P$ is represented as a linear combination
of large-degree earlier lines, it is its actual degree \emph{after
cancellation} that controls this next step. This observation is the
source of the intersection operation in our later linear-algebra model.

\section{Least degrees, zero, and infinity}
\begin{definition}[Least target degree]\label{tx:def:least-degree}
For $S\in\{\NS,\PC\}$, let $h_S(\Gamma\vdash G)$ be the least degree
of an $S$ derivation of $G$, or $\infty$ if there is no derivation.
We order $\NN\cup\{\infty\}$ by placing $\infty$ above every integer.
For refutations write
\[
                      d_S(\Gamma)=h_S(\Gamma\vdash1).
\]
\end{definition}

Thus $h_{\PC}(\Gamma\vdash G)=2$ says both that a PC derivation with
every line of degree at most two exists and that no derivation confined
to degree at most one exists. The symbol $S$ selects the proof system;
$h_S$ then records the smallest degree achievable in that system.

These conventions separate three situations. The zero target has degree
zero, even from empty input. A nonzero constant target is equivalent to
the target $1$, by rescaling with a field element. A nonconstant target
has a positive degree floor, but its minimum proof degree can be much
larger or infinite.

\begin{proposition}[Degree floor and simulation]\label{tx:prop:basic-comparison}
For every target $G$,
\begin{equation}\label{tx:eq:basic-comparison}
 \deg G\le h_{\PC}(\Gamma\vdash G)
          \le h_{\NS}(\Gamma\vdash G).
\end{equation}
If a finite proof exists, the first inequality is an ordinary integer
inequality. With an infinite value it is read in the extended order.
\end{proposition}
\begin{proof}
The last PC line is $G$, so its degree is at least $\deg G$. For NS,
a sum of products of degree at most $d$ is also of degree at most $d$.
To simulate an NS certificate, expand each multiplier into monomials.
Derive a permitted product $\mu f$ by listing $f$ and multiplying by
one variable factor of $\mu$ at a time. None of these intermediate
products exceeds $\deg\mu+\deg f$. Rescale and add all products to
obtain $G$. Every step remains within the original NS degree bound.
\end{proof}

The algebraic object which organises unrestricted derivability is an
\emph{ideal}. An ideal in a polynomial ring is a subspace closed under
multiplication by every polynomial. The ideal generated by $\Gamma$,
written $\langle\Gamma\rangle$, consists of all finite sums
$\sum_f a_ff$ with $f\in\Gamma$ and polynomial coefficients $a_f$.
It is the smallest ideal containing the inputs. Thus unrestricted NS
derivability is precisely membership in $\langle\Gamma\rangle$.
At a fixed degree, the restriction on each product is an additional
condition; an arbitrary ideal representation need not respect it.

There is a converse without a degree bound. Annotate each PC input $f$
with its expression as a single input multiple. Linear combination of
two lines combines their annotations, and multiplication by $x_v$
multiplies every coefficient in the annotation by $x_v$. At the last
line this produces an NS identity. It may have a higher product degree
than the largest actual PC line, because cancellations have been
expanded rather than reused. Consequently NS and PC derive precisely
the same polynomials when degree is unrestricted: the ideal generated
by $\Gamma$. The difficult comparison later in the book concerns their
\emph{bounded-degree} behaviour uniformly in the universe size.

\section{Algebraic consequence and solutions}
Every derivation is sound under evaluation. If a tuple $a$ satisfies
$f(a)=0$ for every input, an NS identity gives $G(a)=0$. The PC rules
preserve this assertion line by line. The same is true when $a$ lies
in an extension field of $K$.

The converse need not hold for the proof systems defined here. From
$x^2=0$, every field-valued solution has $x=0$, but $x$ is not a
polynomial multiple of $x^2$. It therefore has no NS or PC derivation.
So there are two different questions: does $G$ vanish at every solution,
and does an NS or PC derivation of $G$ exist? We call the first
\emph{semantic implication} and the second \emph{algebraic derivability}.
For the proof systems defined here, algebraic derivability is ideal
membership. The example shows that semantic implication alone does not
guarantee it.

For Boolean systems we can prove a converse for refutation. Here a
\emph{Boolean assignment} assigns each variable either zero or one, and
an \emph{inconsistent Boolean system} means that no such assignment
satisfies all its equations. We now prove completeness: an inconsistent
Boolean system with the Boolean axioms explicitly present has an
algebraic refutation. The argument constructs the certificate directly.

\begin{proposition}[Completeness for Boolean refutation]\label{tx:prop:boolean-completeness}
Let $\Gamma$ be a finite set of polynomials in finitely many variables,
over any field $K$, containing $x_v^2-x_v$ for every variable $x_v$.
If no Boolean assignment satisfies $\Gamma$, there is an NS identity
deriving $1$ from $\Gamma$, and hence a PC refutation.
\end{proposition}
\begin{proof}
The Boolean axiom gives $x_v(x_v-1)=0$ at every solution. Since a field
has no zero divisors, each coordinate is zero or one; this remains true
over an extension field. For each Boolean assignment $a$, put
\[
           \delta_a=\prod_{a_v=1}x_v\prod_{a_v=0}(1-x_v).
\]
This is its indicator polynomial: $\delta_a(a)=1$, while
$\delta_a(b)=0$ for every other Boolean assignment $b$. Expanding
$\prod_v(x_v+(1-x_v))$ gives $\sum_a\delta_a=1$.

There are only $2^{|V|}$ Boolean assignments. Because none satisfies the
system, for each $a$ choose one input $f_a$ with $f_a(a)\ne0$. Its
nonzero value has an inverse in $K$. Form the explicit polynomial
\[
                  H=1-\sum_a\frac{\delta_a}{f_a(a)}f_a.
\]
At any Boolean assignment $b$, all summands except the one indexed by
$b$ vanish, and that summand has value one. Thus $H(b)=0$.

We claim that a polynomial vanishing on every Boolean assignment is
an explicit sum of multiples of the Boolean axioms. To see this,
repeatedly reduce powers using the identity
\[
 x_v^e-x_v=(x_v^2-x_v)(1+x_v+\cdots+x_v^{e-2})\qquad(e\ge2).
\]
Each replacement records a polynomial multiple of $x_v^2-x_v$.
After finitely many replacements every exponent is at most one, so
\[
                  H=R+\sum_v b_v(x)(x_v^2-x_v),
\]
where $R$ is multilinear. The replacements do not change values on
Boolean assignments, so $R$ also vanishes at all of them. A multilinear
polynomial is determined by these values. For a direct induction, write
$R=A+x_v B$: evaluating at $x_v=0$ gives that $A$ vanishes on the
smaller Boolean cube, and evaluating at $x_v=1$ then gives the same for
$B$. Starting with constant polynomials, induction proves $R=0$.
Therefore rearranging the definition of $H$ gives the actual identity
\[
       1=\sum_a\frac{\delta_a}{f_a(a)}f_a
                            +\sum_v b_v(x)(x_v^2-x_v).
\]
Every polynomial on the right which is used as an input belongs to
$\Gamma$. This is an NS refutation. The simulation proved earlier
converts it into a PC refutation.
\end{proof}

Completeness here means that absence of a satisfying assignment
guarantees existence of a refutation. The construction uses an
exhaustive list of assignments and power reductions, so it gives an
ordinary proof with finitely many terms or lines. It does not promise
that the proof is short or has low degree. The degree lower bounds
later in the book concern how much degree may be unavoidable despite
this existence result.

\section{The next question: can a proof survive a change of universe?}
The definitions so far concern one fixed polynomial system. In the
next chapter the axioms and targets will come from fixed rules indexed
by the universe size $n$. We shall then ask whether a degree-$d$ proof
at one size can be lifted to a larger size with the same degree budget.
The earlier proof may use many lines; our aim is to preserve the
degree bound without imposing a bound on their number.

If this can be repeated at $n,n+q,n+2q,\ldots$, the one chosen proof
degree becomes a constant bound on that progression. If the family
is known to have no such bounded progression, an eligible starting
proof cannot exist. Measuring which degrees are eligible is how the
lifting construction will become a lower-bound argument.

\section{Exercises}
\begin{exercise}\label{tx:ex:zero-constant}
Find the least NS and PC degrees of the targets $0$, $1$, and $x$ from
$\Gamma=\varnothing$ and from $\Gamma=\{1\}$. Explain the role of the
last-line degree floor.
\end{exercise}
\noindent{\small\hyperref[tx:ans:zero-constant]{Model answer (p.~\pageref*{tx:ans:zero-constant}).}}\par\medskip

\begin{exercise}\label{tx:ex:small-gap}
For Example~\ref{tx:eg:cancellation}, prove
$h_{\PC}(\Gamma\vdash z)=2$ and
$h_{\NS}(\Gamma\vdash z)=3$. For the NS lower bound, list the monomials
which can occur in a degree-two certificate.
\end{exercise}
\noindent{\small\hyperref[tx:ans:small-gap]{Model answer (p.~\pageref*{tx:ans:small-gap}).}}\par\medskip

\begin{exercise}\label{tx:ex:Boolean-rule}
A proposed degree-two PC calculation multiplies the derived line $xy$
by $x$ and replaces $x^2y$ by $xy$, using a Boolean equation. Explain
why this is not an ordinary degree-two step, even if $x^2-x$ is an input.
\end{exercise}
\noindent{\small\hyperref[tx:ans:Boolean-rule]{Model answer (p.~\pageref*{tx:ans:Boolean-rule}).}}\par\medskip


\chapter{What it means for a family to be uniform}\label{tx:ch:uniformity}\label{tx:sec:systems}
The number of variables in our examples grows, but the rule describing
them does not. At universe size $n$ there may be a variable $x_{ij}$ for
every pair of labels $i,j<n$, an equation for every row, and a target
for every distinguished label. The classification theorem requires a
precise version of this fixed-rule condition. Symmetry at each size,
by itself, does not imply the four complexity profiles: the equations
could be chosen differently at each size, according to an arbitrary
rule. We need symmetry together with one coefficient rule which remains
the same as the universe grows.

For the main proof, this chapter has three jobs: specify the labelled
variables, fix the coefficient rule for the axioms, and do the same for
the requested targets. The intervening orbit-sum and compression results
explain the geometry of supported expressions, but the linear lower
bound does not use compression. On a first reading, one may pass from
Definition~\ref{tx:def:template} and its examples directly to
Sections~\ref{tx:sec:target-request}--\ref{tx:sec:target-budgets}, then
return to the structural results later.

\section{Labels, support, and relabelling}
We also call the universe size $n$ the \emph{rank} of the labelled
family. It counts available labels. The dimension of a polynomial
space at that rank is a separate number, usually much larger.

Write $[n]=\{0,\ldots,n-1\}$. Think of its members as names rather
than numerical values. A variable $x_{ij}$ is indexed by the ordered
pair $(i,j)$, and relabelling by a permutation $\pi$ sends it to
$x_{\pi(i),\pi(j)}$. The labels needed to specify that index are
$\{i,j\}$; this is the support of the pair, and therefore of the
variable indexed by it. When $i=j$, its support has size one, not two.
Permutations fixing those labels leave that variable unchanged. We
start with individual indices because monomials and whole polynomial
expressions can have different kinds of support, as the row-sum example
below will show.

\begin{definition}[Bounded basis system]\label{tx:def:basis}
A bounded basis system is a collection of objects labelled from one
infinite pool $\NN$, with an action of label permutations. There is
one fixed integer $b$, independent of $n$ and of the object, and each
object $v$ is assigned a set $\supp(v)\subseteq\NN$ with
$|\supp(v)|\le b$. Supports are equivariant:
$\supp(\pi v)=\pi(\supp(v))$. Every permutation fixing
$\supp(v)$ pointwise fixes $v$. At rank $n$, we keep precisely the
objects whose supports lie in $[n]$; we require that this set be finite.
We call $b$ a support bound for the system.
\end{definition}

Examples include ordered pairs, tuples with specified equality patterns,
and finitely many tagged sorts of such tuples. There may also be objects
with empty support: a distinguished variable $y$ fixed by every label
permutation is an example. A support bound is a bound on the labels
needed by one object, not on the number of objects present at rank $n$.
The important requirement is the \emph{same bound at every rank}.
Finiteness at one rank would not provide that: a bound of $n$ would
grow with the universe and would not be a bounded-support system in
this sense.

A variable basis $\mathcal V$ gives the ring
$K[x_v:v\in\mathcal V_n]$. We write $P_{n,\le D}$ for its vector
space of polynomials of total degree at most $D$. If each variable has support at most $k$,
a monomial of degree at most $d$ involves at most $kd$ labels: it has
at most $d$ variable factors, each involving at most $k$ labels. Powers
and overlapping supports may make the true number smaller.

\begin{example}[The row sum mentions every label]\label{tx:eg:row-support}
With variables $x_{ij}$, consider
\[
                         f_{i,n}=\sum_{j<n}x_{ij}-1.
\]
Its index is $i$, of support one. Every permutation of $[n]$ fixing $i$ fixes the
whole polynomial, because it merely permutes the summands. Nevertheless
the union of labels mentioned in the polynomial is generally all of
$[n]$. Each monomial uses at most two labels, and the index uses one;
neither assertion says that the whole polynomial uses only one or two
labels. This distinction permits totality sums in the theory.
\end{example}

\section{Stable coefficient rules}
\begin{definition}[Uniform axiom template]\label{tx:def:template}
A uniform template consists of a variable basis $\mathcal V$ with
support bound $k\ge1$, an axiom-index basis $\mathcal A$ with support
bound $s$, and coefficients $c(a,\mu)\in K$. Here $a$ is an axiom
index and $\mu$ is a monomial in the labelled variables. The rule is
independent of rank and invariant under simultaneous relabelling:
\[
                 c(\pi a,\pi\mu)=c(a,\mu).
\]
For a fixed $D$, it vanishes whenever $\deg\mu>D$. At rank $n$ the
inputs are
\begin{equation}\label{tx:eq:template-instance}
 \Gamma_n=\{f_{a,n}:a\in\mathcal A_n\},\qquad
 f_{a,n}=\sum_{\mu\text{ supported in }[n]}c(a,\mu)\mu.
\end{equation}
Fix $n_0\ge\max\{1,kD+s\}$ as part of the template data.
\end{definition}

Because there are finitely many variables at each rank and the degree
is bounded, each displayed sum is finite. The same coefficients are
used when the universe grows; there are simply more monomials available.
For the row sum, the coefficient of $x_{ij}$ is one when its row index
is $i$, the constant coefficient is minus one, and other coefficients
are zero. None of these rules asks how large $n$ is.

For a group acting on indices, the \emph{orbit} of $a$ is the set
of all its relabellings $\pi a$. Its \emph{stabiliser} consists of the
permutations which fix $a$. These familiar group-action notions have
a specific use here: indices in the same orbit describe the same
coefficient pattern with different label names. Fixing a support
\emph{pointwise} means fixing each of its labels, not merely preserving
the set of labels as a whole. Later, $H_S$ denotes the subgroup fixing
the distinguished set $S$ pointwise.

The stable range allows us to recognise the degree of an axiom from
its index orbit. A monomial witnessing the largest nonzero coefficient
uses at most $kD$ labels. Together with the axiom index it uses at most
$kD+s$. A permutation fixing the index can move all the additional
labels into $[n]$ when $n\ge kD+s$. Hence, beyond $n_0$, each index
orbit has one fixed polynomial degree. This matters when deciding which
multiplier--axiom products are legal at a proposed proof degree.

\begin{example}[Symmetric but not uniform]\label{tx:eg:not-uniform}
Choose an arbitrary set $A$ of positive integers. At rank $n\in A$
let the inputs be $\{1\}$, and at other ranks let there be no inputs.
Each instance is invariant under every label permutation. Yet the
refutation degree is zero on $A$ and infinite elsewhere. Arbitrary
choices of $A$ cannot satisfy the periodic conclusions of this book.
They are excluded because the axiom rule changes with rank.
\end{example}

\section{Supported expressions are orbit-sum expressions}\label{tx:sec:supported-expressions}
The row sum suggests a converse question. Suppose we are given a
polynomial and know that every permutation outside a short list of
labels fixes it. Can we describe the polynomial by a uniform rule?
At one rank the answer is yes. To represent a whole family by one
rule, we also need compatibility as rank changes.

Here is the support notion for a whole expression. At rank $n$, a set
$S\subseteq[n]$ \emph{supports a polynomial $F$} if every permutation
of $[n]$ fixing $S$ pointwise fixes $F$. The permutation may move every
other label. This is different from collecting every label mentioned
in its monomials. For example, $\sum_{j<n}x_j$ has empty support as an
expression, although its monomials collectively mention all labels.
The row sum $f_{i,n}$ is supported by $\{i\}$ in this sense.

\begin{proposition}[Orbit-sum description]\label{tx:prop:orbit-sum}
Let $F$ have degree at most $D$, and let $H_S$ be the permutations of
$[n]$ fixing $S$ pointwise. The polynomial $F$ is supported by $S$
exactly when it has a unique expression
\[
             F=\sum_O a_O\,\Omega_{O,n},\qquad
             \Omega_{O,n}=\sum_{\mu\in O}\mu,
\]
where $O$ runs over the $H_S$-orbits of degree-at-most-$D$ monomials.
\end{proposition}
\begin{proof}
Relabelling permutes the monomial basis. Thus $hF=F$ for all $h\in H_S$
exactly when monomials in the same orbit have the same coefficient.
Collect these equal coefficients to obtain the displayed sum.
Different orbits have disjoint monomial sets, so the orbit sums are
linearly independent and their coefficients are unique.
\end{proof}

This is a grouping of distinct monomials, not an average over a group.
In particular, it is valid in characteristic $p$ even when the order
of the stabiliser is divisible by $p$.

\begin{example}[All supported linear expressions in one-variable indices]\label{tx:eg:supported-linear}
For variables $x_j$ and a distinguished label $i$, every linear
polynomial fixed by the permutations fixing $i$ has the form
\[
                    a x_i+b\sum_{j\ne i}x_j+c.
\]
For $n\ge2$ there are three monomial orbits: the constant monomial, $x_i$, and all
other variables. At $n=1$ the last orbit is absent, so its coefficient
cannot yet be recovered from the polynomial. At $n\ge2$ all three
orbits are present. Keeping $a,b,c$ fixed gives a uniform expression
at every larger rank.
\end{example}

\subsection{From one supported expression to one uniform template}
Return to a variable basis of support at most $k$. Each monomial of
degree at most $D$ uses at most $r=kD$ labels. Let $s=|S|$ and choose
an ordering $a_0$ of $S$. Relative to those $s$ distinguished labels,
every global monomial orbit can be represented using at most $s+r$
labels in total. Hence all its orbit types are visible at any rank
$n\ge s+r$ containing $S$.

\begin{proposition}[Uniform extension of one supported polynomial]\label{tx:prop:supported-extension}
Suppose $n\ge s+kD$ and $F\in P_{n,\le D}$ is supported by an
$s$-element set $S$. Using ordered distinct $s$-tuples as indices, there
is a unique rank-independent, simultaneously relabelling-invariant
coefficient rule of degree at most $D$ whose polynomial at index
$a_0$ and rank $n$ is $F$. It gives a uniform axiom or target orbit
with index-support bound $s$ and no increase in polynomial degree.
\end{proposition}
\begin{proof}
Consider a pair $(a,\mu)$ consisting of an ordered distinct $s$-tuple
and a monomial of degree at most $D$. First relabel $a$ to $a_0$.
The remaining monomial labels can then be moved into $[n]$ by a
permutation fixing $S$, because there are at most $kD$ of them and
$n-s\ge kD$. Thus every global orbit of pairs has a representative
$(a_0,\mu_0)$ at the chosen rank.

Assign to that orbit the coefficient of $\mu_0$ in $F$. This value
does not depend on the representative. Indeed, a global relabelling
between two such representatives fixes $a_0$. Its restriction to the
labels used by the tuple and monomial is a finite injection within
$[n]$, which extends to a permutation of $[n]$ fixing $S$. It therefore
relates the two monomials by $H_S$, and their coefficients agree by
Proposition~\ref{tx:prop:orbit-sum}.

The assigned value defines $c(a,\mu)$ on all pairs, with zero
coefficients above degree $D$. Being constant on simultaneous pair
orbits makes it relabelling invariant. Summing its coefficients over
rank-$m$ monomials gives the required polynomial for every index
available at rank $m$. At the original rank and index, every coefficient
agrees with $F$. All pair orbits already occurred there, so the
coefficient rule is unique.
\end{proof}

Thus one supported polynomial at a sufficiently large rank determines
its uniform continuation. This does not say that a separately chosen
polynomial at every rank necessarily follows that continuation.

\subsection{The additional condition for a whole family}
For $m\le n$, let $\rho_{n,m}$ retain the coefficients of monomials
supported in $[m]$ and discard the other monomials. This is the
coefficient restriction map, rather than a transport rule for proof
steps. A stable uniform family satisfies
\begin{equation}\label{tx:eq:coefficient-coherence}
        \rho_{n,m}(F_{a,n})=F_{a,m}
        \quad\text{whenever the index }a\text{ is supported in }[m].
\end{equation}

\begin{proposition}[Characterisation of a stable coefficient rule]\label{tx:prop:coherent-family}
Fix bounded variable and index bases and a degree bound $D$. A family
$F_{a,n}\in P_{n,\le D}$, defined for every $n\ge1$ and every index
$a$ available at that rank, comes from one uniform coefficient rule if
and only if it is equivariant under permutations at each rank and
satisfies \eqref{tx:eq:coefficient-coherence}.
\end{proposition}
\begin{proof}
A uniform rule gives equivariance by simultaneous relabelling of
indices and monomials. Restriction preserves exactly its coefficients
on the retained monomials, proving coherence.

Conversely, for a global pair $(a,\mu)$ choose a rank containing all
its labels and define its coefficient to be the coefficient of $\mu$
in $F_{a,n}$. Coherence makes this independent of the chosen rank.
To compare it with the coefficient of $(\pi a,\pi\mu)$, choose a
larger rank containing the labels of both pairs. The action of $\pi$
on the finite support of the pair extends to a permutation of that
rank. Equivariance there gives equality of the two coefficients.
The resulting rule has all the required properties, and reproduces
the family coefficient by coefficient.
\end{proof}

For example, $F_n=a_n\sum_{i<n}x_i$ is invariant at every rank for
\emph{any} sequence of scalars $a_n$. If the scalars change, the
coefficient of a retained variable changes, and coherence fails.
Small support and bounded degree control the form of an expression;
coherence controls its behaviour across different universe sizes.

\section{Smaller-support generators: what can be compressed?}\label{tx:sec:support-compression}
The bound $kD$ applies to one monomial. It need not bound the number
of distinguished labels needed to specify a polynomial. A further,
different possibility is to replace a polynomial's entire permutation
orbit by generators with smaller support. First we state a sufficient
condition for replacing axioms without changing proof degrees.

\begin{proposition}[Degree-preserving axiom replacement]\label{tx:prop:filtered-replacement}
Suppose two axiom sets $\Gamma$ and $\Gamma'$ in the same polynomial
ring satisfy, for every integer $e\ge0$,
\[
 \Span_K\{f\in\Gamma:\deg f\le e\}
       =\Span_K\{g\in\Gamma':\deg g\le e\}.
\]
Then every target has the same least NS degree and the same least PC
degree from the two sets, including degree zero and infinity.
\end{proposition}
\begin{proof}
An input $f\in\Gamma$ of degree $e$ is a constant-coefficient linear
combination of inputs from $\Gamma'$ of degree at most $e$. In an NS
certificate replace $f$ by that combination. Every new multiplier--input
product has degree no larger than the product it replaces, since its
new input degree is at most $e$. This transforms any certificate
without increasing its degree. In a PC derivation, derive each old
input by these linear combinations of new inputs, within its old
degree, and then reuse the original inference steps. This also
preserves the degree budget. Interchanging the two sets gives the
converse transformations and proves equality of their least degrees.
\end{proof}

The increasing spaces
$A_e=\Span_K\{f\in\Gamma:\deg f\le e\}$ form a \emph{filtration}:
$A_0\subseteq A_1\subseteq\cdots$, organised by the usable input degree.
The replacement criterion compares this entire filtration.
The unrestricted span forgets which inputs were usable at each cutoff.

An arbitrary change of basis for the full axiom span need not meet this
condition. In Example~\ref{tx:eg:cancellation}, replacing $f_2$ by
$f_1+f_2=x+1$ leaves the axiom span unchanged, but now
$z=z(x+1)+f_3$ is an NS certificate of degree two. From the original
axioms the least NS degree is three, as Exercise~\ref{tx:ex:small-gap}
shows. Equality of the unrestricted span alone therefore does not
justify a claim about exact proof degrees. The proposition supplies
a sufficient condition by comparing the inputs usable at \emph{each}
degree.

The next example satisfies the criterion and supplies a successful
support compression.

\begin{proposition}[Compression of a linear subset sum]\label{tx:prop:linear-compression}
Let $\varnothing\ne S\subsetneq[n]$, put $s=|S|$, and consider
$F_S=\sum_{i\in S}x_i$. Its smallest support as an expression has size
$\min\{s,n-s\}$. Nevertheless its $S_n$-orbit span has generators
supported on at most two labels. More precisely, with $s$ interpreted
as a scalar in $K$,
\[
 \Span_K\{\pi F_S:\pi\in S_n\}=
 \begin{cases}
     \Span_K\{x_i:i<n\},&s\ne0\text{ in }K,\\
     \Span_K\{x_i-x_j:i,j<n\},&s=0\text{ in }K.
 \end{cases}
\]
\end{proposition}
\begin{proof}
The coefficient of $x_i$ is one on $S$ and zero off $S$. If labels
outside a proposed support include one label of each kind, swapping
those two labels changes the polynomial. Thus a support must contain
all of $S$ or all of its complement, and either choice is sufficient.
This proves the minimum-support formula.

Take $i\in S$ and $j\notin S$. Then
\[
                       F_S-(ij)F_S=x_i-x_j.
\]
Relabelling shows that the orbit span contains every difference.
Those differences span the subspace
$A=\{\sum_i a_ix_i:\sum_i a_i=0\}$: fix a label $j$ and write a
vector of $A$ as $\sum_{i\ne j}a_i(x_i-x_j)$.
Each orbit vector has coefficient sum $s$. If $s=0$ in $K$, the
orbit span is contained in $A$ and is therefore exactly $A$. If
$s\ne0$, it also contains $F_S\notin A$; since $A$ has codimension
one, it is the whole linear polynomial space.
\end{proof}

For example, $\sum_{i\in S}x_i$ can need many labels to specify even
though each monomial uses only one. The proposition replaces its orbit
span, not the polynomial itself, by one- or two-label generators.
All replacements are linear combinations of the original orbit
polynomials and conversely. Thus making this replacement in an axiom
set changes neither the NS nor the PC consequence spaces at any
positive degree budget: the starting linear spans agree, and all
replacement equations have degree one.

The two-label bound can really be necessary. Suppose the characteristic
$p$ divides both $n$ and $s$, with $n\ge3$. A vector of $A$ supported
by at most one label $i$ has the form
$a x_i+b\sum_{j\ne i}x_j$. Its coefficient sum is
$a+(n-1)b=a-b$, so $a=b$. Such vectors are scalar multiples of the
fully invariant sum $\sum_jx_j$. They span only one dimension, whereas
$A$ has dimension $n-1$. Consequently one-label generators cannot
replace the whole orbit span in this case.

\begin{example}[Four-label axioms replaced by two-label axioms]\label{tx:eg:four-to-two}
In characteristic two, let $\Gamma_n$ consist of all sums
$\sum_{i\in S}x_i$ over four-element subsets $S\subseteq[n]$.
For every $n\ge5$, their span is exactly the span of the differences
$x_i-x_j$. Thus we can replace all the four-label equations by the
two-label equations $x_i-x_j=0$, preserving every NS and PC consequence
degree. Both descriptions have fixed coefficient rules across ranks.
The starting-rank condition matters: at $n=4$ there is only one
four-element subset, so its single sum does not span all differences.

At $n=12$, an individual four-label sum has smallest support four,
whereas its orbit span has two-label generators. Since the characteristic
divides both 12 and 4, the preceding argument shows that one-label
generators would not suffice. This gives a concrete reason to distinguish
the polynomial's own support from the support of generators for its
orbit span.
\end{example}

The general compression question has been resolved in separate
companion work~\cite{CompressionCompanion}; its proof is not developed
in this book. In the stable-template setting over a field, put $r=kD$.
There is one coherent replacement template with the same variables
and degree bound, and with axiom-index support at most $2r$, such that
for every
\[
                  n\ge\max\{n_0,\max(s,r)+2r\}
\]
the old and new axiom spans agree at \emph{every input-degree cutoff}.
The replacement criterion above then preserves every target's least
NS degree and its least ordinary-PC degree separately, including zero
and infinity. This is an optional normal form, not an assertion about
arbitrary submodules without a bound on the starting supports, and not
an executable compression algorithm. No equivalence at the exceptional
small ranks is claimed.

The monomial bound $kD$ alone does not establish that compression;
the companion proof supplies an additional structural argument.
The lifting and classification proofs below retain the given index
support $s$ in their budget $2kd+s+t$; they do not assume that it has
already been compressed to $kD$ or $2kD$.

\section{Targets also need a fixed rule}\label{tx:sec:target-request}
\begin{definition}[Uniform target request]\label{tx:def:target}
A uniform target family for a fixed axiom template has a target-index
basis $\mathcal B$ with support bound $t$, a fixed degree bound $D_G$,
and coefficients $\gamma(b,\mu)\in K$. They are independent of rank,
invariant under simultaneous relabelling of $b$ and $\mu$, and zero
above degree $D_G$. The rank-$n$ target indexed by $b\in\mathcal B_n$ is
\[
                  G_{b,n}=\sum_{\mu\text{ supported in }[n]}
                                     \gamma(b,\mu)\mu.
\]
For $S=\NS$ or $\PC$, define the all-target profile
\begin{equation}\label{tx:eq:target-profile}
 h_S(n)=\max_{b\in\mathcal B_n}h_S(\Gamma_n\vdash G_{b,n}).
\end{equation}
The empty maximum is zero. If any requested target is underivable,
the maximum is infinity.
\end{definition}

The statement $h_S(n)\le d$ means that \emph{every} requested target
has some derivation within degree $d$. Those derivations need not share
lines, and their total length is not bounded. This is a maximum of
individual minimum degrees, not the degree of the sum of the targets.
Nor are the targets inserted into the axiom system.

There are three scopes which we will keep distinct.
\begin{enumerate}
\item A fully invariant target has a single index of empty support.
Its polynomial $G_n$ is fixed by every permutation of $[n]$. For
example, $G_n=\sum_{i<n}x_i$ is invariant and has degree one, although
it mentions all labels.
\item A fixed target orbit can contain many indices at a rank. For
$G_{i,n}=x_i$, the indices $i$ form one orbit. Any one representative
has the same least degree as every other representative, since
relabelling transports proofs.
\item An all-target request may contain several orbits. Its profile
is the maximum of their individual profiles.
\end{enumerate}

Every global index orbit has a representative at rank $n\ge t$.
Indeed, move its support into $[t]$. There are only finitely many global
orbits because $\mathcal B_t$ is finite. Within one orbit, two rank-$n$
indices can be related by a permutation of $[n]$: first use the given
relabelling on their supports and then extend that finite bijection.
Thus choosing a representative which varies with $n$ is harmless
\emph{provided it stays in one fixed orbit}. Switching between different
orbits according to an arbitrary rule is not a uniform request.

\section{Stable target degrees and the label budgets}\label{tx:sec:target-budgets}
The same support argument used for axioms applies to targets, now with
at most $kD_G+t$ labels. We therefore use
\begin{equation}\label{tx:eq:stable-rank}
                          n_*=\max\{n_0,1,kD_G+t\}.
\end{equation}
Beyond this rank, every target orbit has a fixed degree. Their largest
degree is a floor for both all-target profiles. A proof budget below
that floor fails; we do not replace the target by its lower-degree
part in order to make it fit.

For an axiom multiplier, the combined support of a monomial of degree
at most $d$ and an axiom index is at most $kd+s$. For multiplication
inside a PC proof, a variable and a monomial of degree at most $d-1$
use at most $k+k(d-1)=kd$ labels. These elementary counts explain why
proof degree enters the transfer theorem linearly.

General indexed targets require a further detection argument. The
resulting safe budget is
\begin{equation}\label{tx:eq:support-budget}
                R(d)=2kd+s+t=\kappa d+\sigma,
                    \qquad\kappa=2k,\quad\sigma=s+t.
\end{equation}
For one fully invariant target the sharper budget is $R(d)=kd+s$.
The factor of two in the general case is not an extra degree cost in a
proof. It is a support cost in the finite-dimensional argument which
shows that target membership is reflected, as well as preserved, when
labels are deleted.

\section{Where first-order systems enter}
A universal clause is a disjunction of literals required to hold for
every assignment of its logical variables to universe labels. A
\emph{literal} is an atom or its negation. A \emph{grounding} replaces
all those logical variables by particular labels; one clause pattern
therefore produces many polynomial equations at rank $n$.
For example, grounding $R(u)\vee\neg R(v)$ at $u=i,v=j$ gives the
clause whose polynomial translation appears below.

The normalised syntax is fixed before proof degrees are measured.
\emph{Flattened} function-graph atoms name an output explicitly,
as in $f(i)=j$, instead of putting nested function terms inside an
atom. They can therefore be indexed by bounded tuples of labels.
Equality between two grounded labels is evaluated as equality of
those labels; it is not an extra unknown relation.

A fixed first-order sentence can supply a uniform template. Relation
atoms receive Boolean variables. A function $f$ is represented by its
graph variables $x^f_{\bar i,j}$, with equations
\[
 \sum_{j<n}x^f_{\bar i,j}-1=0,\qquad
 x^f_{\bar i,j}x^f_{\bar i,j'}=0\quad(j\ne j').
\]
Boolean equations make each value zero or one; collision equations give
at most one output; the sum gives at least one. The collisions cannot
be dropped in characteristic $p$, where $p+1$ ones also sum to one.
A clause is translated as a product of its literal-falsity indicators.
For example $R(i)\vee\neg R(j)$ gives $(1-r_i)r_j=0$.

Fixed arities and a fixed finite list of clause patterns give bounded
variable and axiom-index supports, and expansion gives a stable
coefficient rule. Functions of arity zero represent constants. There
is no built-in numerical order on the labels: an order, if used in a
sentence, is a relation symbol with its own axioms. Later examples keep
this graph-and-sum encoding visible, since an arbitrary change of
encoding need not preserve minimum degree.

\begin{proposition}[Normalised first-order compiler]\label{tx:prop:compiler}
A fixed finite list of universal clauses in flattened relation, function
graph, and equality atoms gives a uniform template by the displayed
translation. At every $n\ge1$, its roots over $K$ correspond exactly
to models on $[n]$ of that normalised sentence.
\end{proposition}
\begin{proof}
Use atom tuples as variable indices, and grounding tuples as clause
indices. Totality and collision indices are fixed tuples as well.
Expansion of a clause product depends only on the fixed syntax and on
equalities among the grounding labels; it therefore gives stable,
simultaneously invariant coefficients, including repetitions and
cancellations. If $a_R,a_f$ are the maximum relation and function
arities, safe parameters are
\[
 k=a_R+a_f+1,\qquad s=k+1+\max_C\operatorname{nvars}(C),\qquad
 D=2+\max_C|C|,\qquad n_0=kD+s+1.
\]
Here $\operatorname{nvars}(C)$ is one more than the largest variable
index occurring in $C$, and is one if none occurs; an unused variable
slot is harmless. Maxima over empty symbol or clause lists may be taken
as zero where appropriate.

At a root, every relation and graph variable is zero or one. Totality
and collisions select exactly one output for each function input, and
a clause product vanishes precisely when one of its literals is true.
Equality atoms use their actual label-equality values. This constructs
a model. Conversely a model supplies the corresponding Boolean values
and satisfies all translated equations.
\end{proof}

Normalisation is part of the input convention. This proposition does
not say that different normalisations, or arbitrary introductions of
witness functions, preserve exact proof degree. The later examples use
specific normalised sentences so that their equations are unambiguous.

\section{Exercises}
\begin{exercise}\label{tx:ex:support-count}
For variables $x_{ij}$, find the label supports and total degrees of
$x_{ij}^3$ and $x_{ij}x_{jk}$ when $i,j,k$ are distinct. Compare the
actual supports with the general bound $2d$. Explain why the row sum
in Example~\ref{tx:eg:row-support} has an index of support one.
\end{exercise}
\noindent{\small\hyperref[tx:ans:support-count]{Model answer (p.~\pageref*{tx:ans:support-count}).}}\par\medskip

\begin{exercise}\label{tx:ex:orbit-switch}
Take no axioms, a single variable $x$ of empty support, and two target
indices of empty support requesting $0$ and $x$. Determine the two
fixed-orbit profiles and the all-target profile. What happens if one
chooses the first target at sizes in an arbitrary set and the second
elsewhere?
\end{exercise}
\noindent{\small\hyperref[tx:ans:orbit-switch]{Model answer (p.~\pageref*{tx:ans:orbit-switch}).}}\par\medskip

\begin{exercise}\label{tx:ex:stable-floor}
Take variables $x_i$, no axioms, and targets $x_i-x_j$ indexed by
ordered distinct pairs. Use $k=1,s=0,t=2,D_G=1,n_0=1$ to compute $n_*$.
At which positive ranks does this request actually contain a target?
Why is the stated stable rank allowed to be larger than necessary?
\end{exercise}
\noindent{\small\hyperref[tx:ans:stable-floor]{Model answer (p.~\pageref*{tx:ans:stable-floor}).}}\par\medskip


\chapter{All proofs inside one finite-dimensional space}\label{tx:ch:closure}\label{tx:sec:pc-space}
A proof can be arbitrarily long even when its degree is fixed. To study
all such proofs, we need not enumerate their possible lengths. There
are only finitely many monomials of degree at most $d$ in a fixed finite
set of variables. Every line of a degree-$d$ proof is a vector in the
space on those monomials. The proof rules describe which vectors can
be reached.

This chapter identifies the reachable space exactly. The distinction
between NS and PC becomes the distinction between one linear span and
an iterated closure. In particular, the closure must recognise a
low-degree vector created by cancellation among higher-degree vectors.

\section{Coordinates for bounded-degree polynomials}
Let $\mathcal X$ be the finite set of variable names and put
$v=|\mathcal X|$. Work over the fixed field $K$. The space $P_{\le d}$
consists of polynomials in $x_u$, $u\in\mathcal X$, of total degree at
most $d$. Its monomial basis has size
\[
                     \dim P_{\le d}=\binom{v+d}{d}.
\]
Indeed, a monomial has exponents $a_1,\ldots,a_v$ with
$\sum a_i\le d$. Add the slack coordinate
$a_{v+1}=d-\sum a_i$. This gives $v+1$ nonnegative integers summing
to $d$. Placing $v$ separators among $d+v$ positions counts those
choices by $\binom{v+d}{v}=\binom{v+d}{d}$.

For two variables and $d=2$, one possible ordering is
$1,x,y,x^2,xy,y^2$. A polynomial is its vector of coefficients in this
basis. Addition and scalar multiplication are the usual vector-space
operations; polynomial multiplication is a different operation which
may leave this space.

Fix a finite system $\Gamma$. We omit its zero inputs when forming
spanning sets, since zero contributes no new vector. For NS set
\begin{equation}\label{tx:eq:NSspan}
 I_d=\Span_K\{\mu f:f\in\Gamma,\ \mu\text{ a monomial},\quad
                              \deg\mu+\deg f\le d\}.
\end{equation}
There are only finitely many permitted pairs $(\mu,f)$. If desired,
put one column $\mu f$ in a matrix for each such pair. Deriving a target
by NS in degree $d$ will mean that its coefficient vector belongs to
the column span of this matrix.

\section{The closure for ordinary PC}
For $d\ge1$ define
\begin{equation}\label{tx:eq:closure}
 \begin{split}
 V_0&=\Span_K\{f\in\Gamma:\deg f\le d\},\\
 V_{j+1}&=V_j+\sum_{u\in\mathcal X}x_u(V_j\cap P_{\le d-1}),\\
 C_d&=V_\infty:=\bigcup_{j\ge0}V_j.
 \end{split}
\end{equation}
For $d=0$, define $C_0=V_0$ without a multiplication term.
The notation $x_uU$ means the image of a subspace $U$ under the linear
map $f\mapsto x_uf$. The sum of subspaces consists of finite sums of
vectors from them.

The intersection in this formula is crucial. It asks which vectors
already derivable in degree $d$ happen to have actual degree at most
$d-1$. It does not merely select input axioms of degree at most $d-1$,
nor basis vectors that were short when first listed. A linear combination
may cancel its highest terms and enter the intersection.

\begin{example}[The first closure step sees cancellation]
Return to $f_1=xy+x$, $f_2=xy+1$, $f_3=xz$ over $\mathbb F_2$, with
$d=2$. A combination
\[
                   af_1+bf_2+cf_3
                  =(a+b)xy+cxz+ax+b
\]
is linear precisely when $a=b$ and $c=0$. Thus
\[
                     V_0\cap P_{\le1}=\Span\{x+1\}.
\]
The vector $z(x+1)$ lies in the new multiplication space, and adding
$f_3$ produces $z\in V_1$. The degree-two NS space, in contrast, is
just the span of the three quadratic inputs and contains no $z$ term.
\end{example}

\begin{proposition}[Exact description of consequences]\label{tx:prop:closure}
For every polynomial $G$ and every $d\ge0$,
\[
 \begin{aligned}
 h_{\NS}(\Gamma\vdash G)\le d&\iff G\in I_d,\\
 h_{\PC}(\Gamma\vdash G)\le d&\iff G\in C_d.
 \end{aligned}
\]
The chain $V_j$ stabilises after finitely many strict increases, and
$I_d\subseteq C_d\subseteq P_{\le d}$.
\end{proposition}
\begin{proof}
Expand every coefficient $a_f$ in an NS identity into monomials.
Every resulting product belongs to the spanning set in
\eqref{tx:eq:NSspan}. Conversely, group a linear combination of such
products by input $f$; the resulting multiplier has degree at most
$d-\deg f$. This proves the NS equivalence, including the empty
identity for $G=0$.

For PC, first prove by induction on a derivation that every line is in
$C_d$. Usable input axioms belong to $V_0$. Two vectors in stages $V_i$
and $V_j$ both belong to $V_{\max(i,j)}$, which is closed under linear
combinations. If a nonzero line $f$ is multiplied by a variable while
staying within degree $d$, then $\deg f\le d-1$. If $f\in V_j$, its
multiple belongs to $V_{j+1}$ by definition. Multiples of zero are zero.

For the converse, induct on the closure stage. Every vector of $V_0$
is a finite linear combination of usable axioms, hence has a legal
proof. A vector of $V_{j+1}$ is a finite sum of an old vector and
terms $x_vf_v$ with $f_v\in V_j\cap P_{\le d-1}$. Concatenate proofs
of the finitely many $f_v$, append the permitted multiplications, and
add to obtain the desired vector as the last line. Every partial sum
still has degree at most $d$.

The chain is increasing inside $P_{\le d}$. A strict increase raises
its dimension, so there are at most $\dim P_{\le d}$ such increases.
If $V_{j+1}=V_j$, the recurrence gives equality at every later stage.
Thus one stage equals the whole union. Finally each NS generator
$\mu f$ has a PC derivation within the same bound by successive
variable multiplication, giving $I_d\subseteq C_d$.
\end{proof}

The proposition includes two different finiteness statements. Every
particular vector in the closure has a finite proof, by the constructive
induction. The space containing \emph{all} degree-$d$ proofs is also
reached after finitely many stages. Neither statement assumes in advance
an upper bound on the length of proofs presented to us.

\section{The linear maps behind the recurrence}
The closure formula can be expressed entirely through fixed ambient
spaces and linear maps. This is the form needed for transferring it
between universe sizes.

The immediate idea is elementary. To describe all permitted variable
multiplications, keep a separate low-degree input polynomial $f_v$
for each variable $x_v$, and send this list to $\sum_v x_vf_v$.
This is a linear map on lists. Later, in
Section~\ref{tx:sec:multiplication-registers}, we call the space of
these lists $X(n)\otimes L_d(n)$ and the map $m_n$.
The tensor notation packages the separate registers; it does not
introduce a new PC inference rule. All these operations occur at
one fixed universe size. Lifting to another size is a separate task.

The first map sends an axiom index to its coefficient vector. Restrict
to axioms of degree at most $d$; its image is $V_0$. The NS map instead
sends a permitted pair $(\mu,a)$ to $\mu f_a$; its image is $I_d$.
Both use the actual input coefficients.

Let $X$ be the vector space with one basis element for each polynomial
variable, and put $L=P_{\le d-1}$. Multiplication is represented by a
single linear map
\[
              m:X\otimes L\longrightarrow P_{\le d},
                         \qquad x_v\otimes f\longmapsto x_vf.
\]
Here the tensor product can be understood concretely: its basis is
made of pairs $(x_v,\mu)$ with $\deg\mu\le d-1$, and the map sends
such a pair to the monomial $x_v\mu$. Bilinearity is built into the
vector space on these pairs.

For any current stage $V_j$, its low-degree part is the intersection
$V_j\cap L$. It can be described as the kernel of
\[
                  V_j\oplus L\longrightarrow P_{\le d},
                              \qquad(f,g)\longmapsto f-g,
\]
identified with either equal coordinate. Multiplication then gives an
image, and adjoining old vectors gives a sum:
\[
 V_{j+1}=V_j+\im\bigl(m|_{X\otimes(V_j\cap L)}\bigr).
\]
Thus the construction uses a small vocabulary: coefficient maps,
kernels, intersections, images, sums, and a finite iteration. A
transport operation preserving this vocabulary, and agreeing with the
actual maps, transports the closure. Merely matching dimensions would
not identify which target vectors lie in it.

\section{Changing the coefficient field}
Extending scalars enlarges the allowed coefficients while keeping the
same coordinate equations. If $U\subseteq K^M$ has basis
$u_1,\ldots,u_a$ and $K\subseteq E$, its scalar extension is
$U_E=\Span_E\{u_1,\ldots,u_a\}\subseteq E^M$.
It is also written $E\otimes_K U$. The subscript $K$ records which
scalars may be moved between the two tensor factors.
For example, the kernel of $(a,b)\mapsto a+b$ is $K(1,-1)$ over
$K$ and $E(1,-1)$ over $E$. The same basis and relation remain,
although there are more choices of scalar multiples.
The next lemma applies this coordinate principle to every stage
of the proof calculation.

\begin{lemma}[Field extension preserves minimum degree]\label{tx:lem:field}
If $K\subseteq E$ are fields and $\Gamma,G$ have coefficients in $K$,
then the least NS and ordinary-PC degrees of $\Gamma\vdash G$ are the
same over $K$ and over $E$.
\end{lemma}
\begin{proof}
At a fixed degree, NS membership is solvability of a finite linear
system with coefficients in $K$. Row reduction over $K$ already
determines whether it is consistent; extending scalars does not change
that answer. Equivalently, a matrix and its augmented matrix have the
same ranks over the two fields.

For PC, extending scalars commutes with sums and images of finite
matrices, and also with their kernels: row reduction describes these
spaces over $K$ and gives bases after extension. Intersections are
kernels of the difference map just displayed. Induction on
\eqref{tx:eq:closure} therefore identifies each extended stage with the
stage computed over $E$. Once both chains have stabilised, membership
of the original vector $G$ is unchanged. This proves equivalence at
every integer degree, including degree zero, and hence equality of
the minimum degrees, including infinity.
\end{proof}

This lemma will let us use convenient auxiliary coefficients in the
finite-dimensional construction. It does not change the target to an
arbitrary polynomial over the larger field. The target whose membership
is tested remains the original $K$-coefficient vector.

\section{When degree gives a proof-size bound}
This section is a separate application of the consequence-space view.
It explains the size interpretation under full coefficient lists; the
construction of the lifting operator and the degree lower bounds do
not depend on this encoding. A reader following the degree proof can
continue to Chapter~\ref{tx:chap:contract} and return here for the size
consequences.

Degree alone is not a count of the nonzero monomials in a proof. We use
one particular, deliberately explicit representation in which a degree
bound has a direct size interpretation.

\begin{definition}[Full coefficient-list size]\label{tx:def:dense-size}
In an ambient ring with $v$ variables, store a nonzero polynomial of
degree $e$ by the full list of coefficients of all monomials of degree
at most $e$, including zeros. The number of entries is
$B_v(e)=\binom{v+e}{e}$. Store zero by one entry. PC size is the sum of
these lengths over all lines, including the final target. NS size is
the sum over its nonzero multipliers; the inputs and target are supplied
separately. The empty NS certificate for zero has size zero.
\end{definition}

Write $S_{\NS}$ and $S_{\PC}$ for the least coefficient-list sizes of
proofs of the fixed target $G$ from $\Gamma$. These are minimum sizes
over all proofs in the specified system; they are infinity when no
proof exists.

\begin{lemma}[Degree versus coefficient lists]\label{tx:lem:dense-size}
With $v$ variables and input degree at most $D$, a degree-$d$ NS proof
has a representative certificate of size at most $B_v(d)^2$. A
 degree-$d$ PC proof has a representative proof of size $O(B_v(d)^3)$.
Conversely, if the least degrees are finite, minimum sizes satisfy
\[
 S_{\PC}\ge B_v(h_{\PC}),\qquad
 S_{\NS}\ge B_v(\max\{0,h_{\NS}-D\})\quad(G\ne0).
\]
Furthermore $B_v(d)\ge2^{\min(v,d)}$.
\end{lemma}
\begin{proof}
Put $M=B_v(d)$. A vector in $I_d$ is a linear combination of at most
$M$ generator columns. Grouping by input gives at most $M$ nonzero
multipliers, each representable by at most $M$ entries.

For PC retain a basis of usable axioms. If its span $W$ is not closed
under the permitted multiplications, choose
$f\in W\cap P_{\le d-1}$ with $x_vf\notin W$. Derive $f$ as a
combination of the at most $M$ retained basis vectors, append $x_vf$,
and enlarge the basis. Each enlargement raises dimension, so there are
at most $M$ enlargements. At the end, the closed space contains every
degree-$d$ consequence. Append a combination giving the desired target.
This uses $O(M^2)$ lines, each with at most $M$ entries.

Any PC proof has a line of degree at least the least possible proof
degree. An NS certificate for a nonzero target has a nonzero product
of degree at least its least degree; its multiplier has degree at least
$h_{\NS}-D$. These give the lower bounds. Finally the table contains
all squarefree monomials in any $\min(v,d)$ variables, of which there
are $2^{\min(v,d)}$.
\end{proof}

All variables of the ambient ring count, including those absent from
nonzero terms. Over a fixed finite field, coefficient entries have
bounded bit length. The results do not automatically become lower
bounds for sparse representations, where zero coefficients are omitted.
This distinction will remain explicit in the size classification.

\section{Exercises}
\begin{exercise}\label{tx:ex:closure-zero}
Describe $I_0$ and $C_0$ for an arbitrary system. Prove directly that a
nonzero constant target has a degree-zero proof exactly when an input
is a nonzero constant.
\end{exercise}
\noindent{\small\hyperref[tx:ans:closure-zero]{Model answer (p.~\pageref*{tx:ans:closure-zero}).}}\par\medskip

\begin{exercise}\label{tx:ex:closure-witness}
With input $\Gamma=\{x^2\}$, show that every vector in every closure
$C_d$ has zero coefficient of $x$. Deduce that $x$ is underivable,
although it vanishes at every field-valued solution of $\Gamma$.
\end{exercise}
\noindent{\small\hyperref[tx:ans:closure-witness]{Model answer (p.~\pageref*{tx:ans:closure-witness}).}}\par\medskip

\begin{exercise}\label{tx:ex:dense-count}
List the monomials counted by $B_2(2)$, and compare the full coefficient
list of $x^2$ with its sparse representation. For $v\ge d$, identify
explicitly the $2^d$ entries used in the final lower bound of
Lemma~\ref{tx:lem:dense-size}.
\end{exercise}
\noindent{\small\hyperref[tx:ans:dense-count]{Model answer (p.~\pageref*{tx:ans:dense-count}).}}\par\medskip


\chapter{A running example: an order with no first element}
\label{tx:ch:running-order}
We now apply the consequence-space viewpoint to one concrete family.
The equations describe a finite total order in which every element has
a selected predecessor. A finite order has a first element, so each
instance is inconsistent. The question is how much polynomial degree
is needed to expose that contradiction.

For the running example fix a finite field $K$ of characteristic $p$.
The chain and restriction arguments below work over every field; the
final linear bound also holds over any fixed field of positive characteristic.

The proof has three stages. First, a restriction to an elementary
induction chain shows that NS degree tends to infinity. Second, the
uniform NS--PC comparison proved in Chapter~\ref{tx:sec:comparison}
shows that PC degree also tends to infinity. Third, the degree-preserving
lifting theorem, constructed in the intervening chapters and assembled
in Chapter~\ref{tx:chap:assemble-transfer}, upgrades this qualitative
information to a linear lower bound. The chain obstruction and the
restriction are proved here in full. The two general results are stated
at their point of use, so the reader can see exactly what their later
proofs must deliver.

\section{The equations and their finite proofs}
For this chapter write $[n]=\{1,\ldots,n\}$; this is just a
renaming of the book's universe. There are two sorts of variables:
\[
 L_{ij}\quad\text{means that }i\text{ precedes }j,\qquad
 F_{ij}\quad\text{means that the selected predecessor of }i\text{ is }j.
\]
Let $\mathcal O_n$ consist of the following polynomials, all interpreted
as equations equal to zero:
\begin{subequations}\label{tx:eq:running-order-axioms}
\begin{align}
 L_{ij}^2-L_{ij},\quad F_{ij}^2-F_{ij}
       &\qquad(i,j\in[n]),\label{tx:eq:order-Boolean}\\
 L_{ii}&\qquad(i\in[n]),\label{tx:eq:order-irreflexive}\\
 L_{ij}L_{ji},\quad(1-L_{ij})(1-L_{ji})
       &\qquad(i\ne j),\label{tx:eq:order-total}\\
 L_{ij}L_{jk}(1-L_{ik})
       &\qquad(i,j,k\in[n]),\label{tx:eq:order-transitive}\\
 \sum_{j=1}^nF_{ij}-1
       &\qquad(i\in[n]),\label{tx:eq:order-function-total}\\
 F_{ij}F_{ik}&\qquad(j\ne k),\label{tx:eq:order-function-collision}\\
 F_{ij}(1-L_{ji})&\qquad(i,j\in[n]).\label{tx:eq:order-predecessor}
\end{align}
\end{subequations}
Here a selected predecessor means any strictly earlier element; it
need not be the immediately preceding element. Different elements may
choose the same predecessor.

The first four lines describe a strict total order. The next two select
exactly one predecessor in each row. The last requires that this
predecessor really come earlier. This is the graph-and-sum encoding of
the first-order assertion that $L$ is a strict total order and
$L(F(i),i)$ holds for every $i$.

The functionality collisions matter in positive characteristic.
Boolean values and a sum equal to one alone could permit $p+1$ selected
entries. The collisions allow at most one, and the sum then forces one.

\begin{proposition}[Finite inconsistency and finite proof degree]
\label{tx:prop:order-finite}
For every $n\ge1$ and every field $K$, the system $\mathcal O_n$ has
an NS refutation and an ordinary-PC refutation. Their minimum degrees
are finite and positive.
\end{proposition}
\begin{proof}
At a root, all variables are zero or one. The order equations give a
strict total order on the labels. Starting from any label, move to an
earlier one while possible. Transitivity and irreflexivity forbid a
repeated label, so finiteness forces this process to stop at a first
element $i$.
The function equations select a label $j$ with $F_{ij}=1$.
Equation~\eqref{tx:eq:order-predecessor} then gives $L_{ji}=1$,
contradicting the choice of $i$. The same argument applies over every
extension field.

Boolean refutation completeness
(Proposition~\ref{tx:prop:boolean-completeness}) constructs a finite
NS certificate, and its simulation gives a finite PC proof. No input
is a nonzero constant, so the degree-zero consequence space is zero.
Neither refutation can have degree zero.
\end{proof}

At $n=1$ the contradiction is already visible in six ordinary-PC lines:
\begin{center}
\begin{tabular}{@{}cllc@{}}
\toprule
Line & Polynomial & Reason & Degree\\
\midrule
1 & $L_{11}$ & input & 1\\
2 & $F_{11}-1$ & input & 1\\
3 & $F_{11}(1-L_{11})$ & input & 2\\
4 & $F_{11}L_{11}$ & multiply line 1 by $F_{11}$ & 2\\
5 & $F_{11}$ & add lines 3 and 4 & 1\\
6 & $1$ & subtract line 2 from line 5 & 0\\
\bottomrule
\end{tabular}
\end{center}
For larger $n$ there is no designated first label. The order variables
must determine it, while every row insists on a predecessor.

\subsection{Why this is one uniform family}
Relabelling sends $L_{ij}$ and $F_{ij}$ to the variables with the
relabelled ordered pair. Every input has an index using at most three
labels, and all coefficients follow one fixed rule. For example, the
row with index $i$ has coefficient one on every $F_{ij}$ and constant
coefficient minus one, regardless of $n$.

Safe template parameters are
\begin{equation}\label{tx:eq:order-parameters}
                  k=2,\qquad s=3,\qquad D=3,\qquad n_0=9.
\end{equation}
There are $2n^2$ variables, but one variable uses at most two labels.
The target $1$ is fully invariant and has degree zero. At proof degree
$d$ the invariant-target support budget is therefore
\begin{equation}\label{tx:eq:order-support-budget}
                              r=kd+s=2d+3.
\end{equation}
This budget bounds the bases and maps of the consequence calculation;
it does not require each proof line to mention only $2d+3$ labels.

\section{A smaller chain problem with two different proof degrees}
Here the input comes with a \emph{fixed successor chain}: the numerical
chain $1,2,\ldots,n$, its successor relation and its two endpoints
are fixed in advance. They are part of the input rule, rather than
relations to be determined by a satisfying assignment. For variables
$x_1,\ldots,x_n$, define the induction system
\begin{equation}\label{tx:eq:order-induction}
 \mathrm{IND}_n=\{1-x_1,\ x_i(1-x_{i+1})\ (1\le i<n),\ x_n,
                                      \ x_i^2-x_i\ (1\le i\le n)\}.
\end{equation}
Its meaning is sequential: the first variable is one, a one forces the
next variable to be one, and the last variable is zero.

There is an ordinary-PC refutation of degree at most two. Starting
from a derived line $P=1-x_i$, first multiply it by the single variable
$x_{i+1}$. Then form
\[
       P-x_{i+1}P+x_i(1-x_{i+1})=1-x_{i+1}
\]
by two linear combinations. Every intermediate line has degree at
most two. Iterate and add the final input $x_n$ to obtain $1$.
For $n=1$, the two endpoint inputs already sum to one.

NS behaves differently: its minimum degree is $\Theta(\log n)$.
We prove both a logarithmic lower bound and a logarithmic upper bound
in the next section. Thus the fixed chain gives a constant versus
logarithmic degree separation. The linear NS and PC bounds belong
to the original uniform order family $\mathcal O_n$, whose order is
an unknown relation. The distinction is possible because this chain family is
\emph{not} a uniform template of the kind defined in Chapter~\ref{tx:ch:uniformity}.
It singles out numerical adjacency and the last label. An arbitrary
label permutation does not preserve its input rule. Thus the uniform
NS--PC comparison cannot be applied to $\mathrm{IND}_n$ itself.

The induction system is a useful restriction of the uniform order
family. We next prove its obstruction using a finite graph of
monomials. The classical NS degree behaviour of induction is discussed
by Buss and Pitassi~\cite{BussPitassiInduction}; the needed lower bound
is proved below rather than used as an external theorem.

\section{A finite graph that witnesses the NS obstruction}
Fix an integer degree budget $b\ge0$. A multilinear monomial
$x^S=\prod_{i\in S}x_i$ is indexed by a subset $S\subseteq[n]$.
Make a graph whose vertices are the sets with $|S|\le b$. Its edges
allow the following changes:
\begin{enumerate}
\item Insert or delete label $1$.
\item Insert or delete label $i>1$ when label $i-1$ is present.
\end{enumerate}
The number of occupied labels must remain at most $b$. Call an occupied
label a pebble. The moves are reversible: the predecessor is still
present after the move, so the reverse move is legal.

\subsection{The clearance invariant}
We need a bound on how far a sequence starting with no pebbles can go.
The following auxiliary test supplies it. Write a configuration as
$a_1<\cdots<a_k$, put $a_0=0$, and call pebble $a_j$ \emph{eligible}
for removal when
\begin{equation}\label{tx:eq:order-clearance-test}
                        a_j-a_{j-1}\le2^{b-k}.
\end{equation}
A configuration is \emph{clearable} if successive eligible removals
can empty it, recalculating the threshold after each removal.
These auxiliary removals need not be legal graph moves. They give an
invariant for the actual graph moves, whose predecessor gaps are one.
In this clearance argument, $k$ counts the currently occupied pebbles;
it is independent of the variable-support parameter $k=2$ for
$\mathcal O_n$.

\begin{lemma}[Any eligible removal can come first]
\label{tx:lem:order-clearance-exchange}
Removing any eligible pebble from a clearable configuration leaves a
clearable configuration.
\end{lemma}
\begin{proof}
Induct on the number $k$ of pebbles. The assertion for $k\le1$ is
immediate. Put $h=2^{b-k}$, choose an eligible pebble $a$, and let $c$
be the first removal in one clearing order. If $a=c$, use the remainder
of that order.

Otherwise both are initially eligible. Removing either doubles the
threshold to $2h$. The other's gap stays the same unless the two
pebbles are adjacent in the configuration; in that case its gap can
become the sum of two gaps at most $h$, hence is at most $2h$. Thus either order
can remove both. The configuration with $c$ removed is clearable;
by induction it can remove $a$ next and then clear the remaining
configuration. Starting instead by removing $a$, remove $c$ next and
use that same remainder. This gives a clearing order after removing $a$.
\end{proof}

\begin{lemma}[The reachable frontier]
\label{tx:lem:order-pebble-frontier}
Every configuration reachable from the empty set using at most $b$
pebbles is clearable. In every nonempty clearable configuration of
$k$ pebbles,
\begin{equation}\label{tx:eq:order-pebble-frontier}
                     a_k\le2^b-2^{b-k}\le2^b-1.
\end{equation}
\end{lemma}
\begin{proof}
The empty configuration is clearable. For a legal deletion, the
deleted pebble has gap one, so is eligible; the exchange lemma preserves
clearability. For a legal insertion, the new pebble has gap one and
can be removed immediately by the auxiliary test. This returns to the
previous clearable configuration, proving that the new one is also
clearable. This proves the invariant along every graph path.

For the bound, induct on $k$, assigning maximum zero to the empty
configuration. A clearing order first removes an eligible pebble with
gap at most $h=2^{b-k}$. The remaining configuration is clearable and
has $k-1$ pebbles, so its maximum is at most $2^b-2h$. If the deleted
pebble was not the largest, this bounds the original maximum as well.
If it was the largest, restoring it adds at most its gap $h$ to that
maximum. Either way, $a_k\le2^b-h$. Since $k\le b$, one has $h\ge1$,
giving the final inequality.
\end{proof}

For example, two pebbles can reach label three:
\[
 \varnothing,\ \{1\},\ \{1,2\},\ \{2\},\ \{2,3\}.
\]
They cannot reach label four. The clearance test explains the limitation
without bounding the number of moves: with two occupied labels the
threshold is one, and after an eligible removal it becomes two.

\subsection{From graph components to a linear functional}
A \emph{linear functional} assigns a scalar to each coefficient vector
and respects addition and scalar multiplication. Specifying its value
on every monomial specifies it on every polynomial by linearity.
To exclude an NS identity for $1$, it suffices to find such a functional
which is zero on every permitted multiplier--input product but nonzero
on $1$. Applying it to the identity would give a contradiction.
This is a dual certificate for nonmembership in a span.

The graph supplies its weights. A difference of two adjacent monomial
coordinates receives value zero if the weights are constant along
edges, hence on connected components. We give the empty component
weight one and the other components weight zero. The proof below
checks that this assignment also handles powers and endpoint inputs.

\begin{proposition}[An elementary logarithmic NS obstruction]
\label{tx:prop:order-induction-lower}
If $n\ge2^b$, the system $\mathrm{IND}_n$ has no total-product
degree-$b$ NS refutation, over any field.
\end{proposition}
\begin{proof}
Let $\mathcal E$ be the component of the empty vertex in the graph.
By Lemma~\ref{tx:lem:order-pebble-frontier}, no vertex of $\mathcal E$
contains $n$. For a monomial $\mu$ of degree at most $b$, let $S(\mu)$
be the labels with positive exponent and define
\[
 \Lambda(\mu)=
 \begin{cases}1,&S(\mu)\in\mathcal E,\\0,&S(\mu)\notin\mathcal E.
 \end{cases}
\]
Extend $\Lambda$ linearly to $P_{\le b}$. In particular $\Lambda(1)=1$.
We check that it vanishes on every permitted monomial--input product.

The two monomials of $\mu(1-x_1)$ reduce to the same subset or to
vertices joined by a label-one edge. Their $\Lambda$ values are equal,
so their difference has value zero. For
$\mu x_i(1-x_{i+1})$, the subsets are
$S(\mu)\cup\{i\}$ and $S(\mu)\cup\{i,i+1\}$. They are identical
or adjacent, because $i$ is present. They both have size at most $b$
by the product-degree restriction, so this difference also has value
zero. The product $\mu x_n$ corresponds to a vertex containing $n$,
hence has value zero. Finally $\mu(x_i^2-x_i)$ has the same set of
positive exponents in both monomials, so again has value zero.

Every allowed NS identity is a linear combination of these products.
Applying $\Lambda$ to a supposed identity for $1$ would give $0=1$.
This contradiction proves the proposition.
\end{proof}

The functional used for the lower bound is a test on coefficient vectors, not a truth
assignment. Replacing a positive exponent by one in its definition
does not grant a free Boolean-reduction rule to a PC proof.
Indeed, the functional need not respect multiplication of derived
low-degree lines; the degree-two PC refutation above demonstrates
precisely why an NS witness need not be an ordinary-PC witness.

For a small calculation, take $b=2$ and $n\ge4$. The empty component
contains $\{2\}$ and $\{2,3\}$ but not $\{3\}$: the singleton's gap
three fails its clearance threshold two. Consequently
\[
 \Lambda(1-x_2)=0,\qquad
 \Lambda\bigl(x_3(1-x_2)\bigr)=0-1=-1\ne0.
\]
The line $1-x_2$ has a degree-two PC derivation. Its permitted
multiplication by $x_3$ is exactly an operation this NS functional does
not preserve.

\subsection{Why the chain also has logarithmic NS certificates}
The chain obstruction cannot be strengthened to a linear bound.
There is a short divide-and-conquer construction that makes this
clear without appealing to an external upper-bound theorem.

\begin{proposition}[A logarithmic NS upper bound for the fixed chain]
\label{tx:prop:order-induction-upper}
For $n\ge2$, the system $\mathrm{IND}_n$ has a total-product NS
refutation of degree at most $2+\lceil\log_2(n-1)\rceil$, over every field.
\end{proposition}
\begin{proof}
For $a<b$ write $T_{ab}=x_a(1-x_b)$. When $b=a+1$, this is
an input. For any $a<c<b$, the polynomial identity
\[
 T_{ab}=(1-x_b)T_{ac}+x_aT_{cb}
\]
combines two shorter intervals. Choose $c$ so that their lengths
differ by at most one, and recursively expand both terms until only
the original adjacent-transition inputs remain. Each recursive level
multiplies their certificates by a polynomial of degree one. An
interval of length $b-a$ needs at most
$\lceil\log_2(b-a)\rceil$ levels, so every resulting multiplier--input
product has degree at most $2+\lceil\log_2(b-a)\rceil$.
This is an NS identity in the original inputs, obtained by expanding
all intermediate transitions.

Finally use the endpoint identity
\[
                1=(1-x_1)+T_{1n}+x_1x_n.
\]
The first term is an input, the last is a degree-two multiple of the
input $x_n$, and the middle has the certificate just constructed.
The displayed degree bound follows. For $n=1$, the endpoint inputs
already sum to one with total-product degree one.
\end{proof}

Together with Proposition~\ref{tx:prop:order-induction-lower}, this
proves the claimed $\Theta(\log n)$ NS degree. The fixed successor chain
allows PC to carry a single degree-one conclusion down the chain;
an NS certificate records the original inputs, and balanced
composition keeps its growing degree logarithmic.

\section{Restricting the uniform order problem to the chain}
In this argument, a \emph{substitution} assigns a polynomial to each
variable and extends by addition and multiplication to the whole
polynomial ring. It is a ring homomorphism:
$\sigma(af+bg)=a\sigma(f)+b\sigma(g)$ and
$\sigma(fg)=\sigma(f)\sigma(g)$, where $a,b\in K$.
Thus it sends an identity $1=\sum a_ff$ to
$1=\sum\sigma(a_f)\sigma(f)$.
To obtain a certificate from the chain inputs, we must additionally
express each substituted input $\sigma(f)$ using those inputs.
The degree accounting below controls both steps. Substitution here
is a fixed-size reduction between two systems, whereas the later
lifting operator is a linear map between universe sizes.

We now relate $\mathcal O_n$ to $\mathrm{IND}_n$, for $n\ge2$.
The substitution need not respect relabelling: it will be used only
to transform an NS identity at one fixed size. Set
\begin{equation}\label{tx:eq:order-cycle-restriction}
 \sigma(F_{ij})=
 \begin{cases}
 1,&i>1\text{ and }j=i-1,\text{ or }i=1\text{ and }j=n,\\
 0,&\text{otherwise},
 \end{cases}
\end{equation}
and put
\begin{equation}\label{tx:eq:order-relation-restriction}
 \sigma(L_{ij})=
 \begin{cases}
 0,&i=j,\\
 1-x_i(1-x_j),&i<j,\\
 (1-x_i)x_j,&i>j.
 \end{cases}
\end{equation}
Every substituted variable has degree at most two. On a Boolean
assignment, the relation in \eqref{tx:eq:order-relation-restriction}
orders labels first by their Boolean value, with zero before one,
and then by their numerical label. It is therefore always a strict
total order.

The selected predecessors run around one fixed cycle. For $i>1$,
the nonzero guard becomes exactly
\[
            1-\sigma(L_{i-1,i})=x_{i-1}(1-x_i).
\]
The wraparound guard is
\begin{equation}\label{tx:eq:order-wrap-guard}
       1-\sigma(L_{n,1})=(1-x_1)+x_1x_n.
\end{equation}
This has a degree-two NS identity from the two endpoint inputs of
$\mathrm{IND}_n$. All unselected guards become zero, and the function
Boolean, totality and collision equations become zero identically.

\begin{lemma}[Boolean reduction retains a product-degree bound]
\label{tx:lem:order-Boolean-reduction}
If a polynomial $P$ of degree at most $m$ vanishes on every Boolean
assignment, it is a sum of multiples of $x_i^2-x_i$ with each nonzero
product of degree at most $m$.
\end{lemma}
\begin{proof}
Reduce powers one monomial at a time using
\[
 x_i^a-x_i=(x_i^2-x_i)(1+x_i+\cdots+x_i^{a-2})\qquad(a\ge2).
\]
Multiplying this identity by the other factors of that monomial
does not exceed its original degree. Continue until every monomial
is multilinear, recording the Boolean-axiom products. The final
multilinear remainder still vanishes on the Boolean cube, so is zero
by the induction proved in Chapter~\ref{tx:ch:consequences}.
\end{proof}

Apply this lemma to the substituted order equations, which vanish on
every Boolean assignment. A Boolean or two-label order axiom has
degree two before substitution and degree at most four afterwards.
A transitivity axiom has degree three and becomes a polynomial of
degree at most six. Irreflexivity becomes zero. Thus each source
axiom $f$ of actual degree $e$ has a substituted NS identity from
$\mathrm{IND}_n$ with product degree at most $2e$; the guards have
the stronger bound two, as just computed.

\begin{lemma}[Degree accounting for the restriction]
\label{tx:lem:order-NS-restriction}
A degree-$d$ NS refutation of $\mathcal O_n$ gives a degree-at-most-$2d$
NS refutation of $\mathrm{IND}_n$.
\end{lemma}
\begin{proof}
Write the original identity as $1=\sum_f a_ff$, omitting zero
products. If $e=\deg f$, then $\deg a_f\le d-e$. Hence
$\deg\sigma(a_f)\le2(d-e)$. Replace $\sigma(f)$ by the identity
just established, whose individual products have degree at most $2e$.
After multiplication by $\sigma(a_f)$, every new product has degree
at most $2(d-e)+2e=2d$. Summing gives the required NS refutation.
\end{proof}

\begin{corollary}[The qualitative NS seed]
\label{tx:cor:order-NS-seed}
For every $n\ge2$ and every field,
\begin{equation}\label{tx:eq:order-log-seed}
                        d_{\NS}(\mathcal O_n)>
                                      \tfrac12\log_2 n.
\end{equation}
In particular, the NS minimum degrees tend to infinity.
\end{corollary}
\begin{proof}
If a degree-$d$ refutation existed with $n\ge2^{2d}$,
Lemma~\ref{tx:lem:order-NS-restriction} would produce a chain refutation
forbidden by Proposition~\ref{tx:prop:order-induction-lower} at budget
$b=2d$. Therefore $n<2^{2d}$ for the finite minimum degree $d$,
which is the displayed inequality.
\end{proof}

\section{Why PC degree also tends to infinity}
We now use the fixed positive characteristic. The uniform comparison
theorem, Theorem~\ref{tx:thm:comparison}, supplies the
following input for the original order family: for each fixed PC
budget $d$ there is a finite NS budget $E(d)$, independent of $n$,
such that
\begin{equation}\label{tx:eq:order-comparison-input}
 d_{\PC}(\mathcal O_n)\le d
               \quad\Longrightarrow\quad
 d_{\NS}(\mathcal O_n)\le E(d).
\end{equation}
The function may grow quickly; only independence from $n$ is needed.

For a fixed $d$, inequality~\eqref{tx:eq:order-log-seed} eventually
exceeds $E(d)$. Consequently a degree-$d$ PC refutation is impossible
at every sufficiently large size. Since this holds for each fixed
$d$, the PC minimum degree tends to infinity too. No finite degree
can recur infinitely often in either proof system.

This uses comparison on $\mathcal O_n$, where the uniform hypothesis
holds. Applying it to the restricted chain would be invalid: that
nonuniform family has the degree-two PC proofs already displayed.
The comparison transfers \emph{qualitative unboundedness}; the linear
PC bound will come from lifting with its original PC degree unchanged.

\section{The same lifting mechanism gives the linear bound}
For this family, the invariant-target case of
Theorem~\ref{tx:thm:transfer} gives a particularly simple statement.
For either $S=\NS$ or $S=\PC$, if
\begin{equation}\label{tx:eq:order-lift-range}
 q=p^a>2d+3,\qquad n\ge\max\{9,4d+5\},
\end{equation}
then a degree-$d$ refutation at $n$ gives one at $n+q$. These numbers
come from $r=2d+3$ and the retained-label requirement $n\ge2r-1$.
We use the actual range of the transfer theorem; the profile chapter
also defines a slightly more conservative onset for general bookkeeping.

\begin{theorem}[A finite linear lower bound for the order family]
\label{tx:thm:order-linear}
Over every fixed field of prime characteristic, the NS and ordinary-PC
refutation degrees of $\mathcal O_n$ are finite at every positive
size, tend to infinity, and satisfy
\begin{equation}\label{tx:eq:order-linear-bound}
 d_{\PC}(\mathcal O_n)>\frac{n-5}{4},\qquad
 d_{\NS}(\mathcal O_n)>\frac{n-5}{4}\qquad(n\ge9).
\end{equation}
Both profiles are of type T2 in Chapter~\ref{tx:ch:profiles}.
\end{theorem}
\begin{proof}
Finiteness and positivity were proved in
Proposition~\ref{tx:prop:order-finite}. The preceding two sections show
that both minimum degrees tend to infinity.

Suppose a degree-$d$ refutation existed at a size satisfying
\eqref{tx:eq:order-lift-range}. Keep this $d$ fixed and choose one
admissible prime power $q$. Iterated forward lifting gives refutations
of degree at most $d$ at
\[
                            n,n+q,n+2q,\ldots.
\]
Every step remains eligible: the retained rank grows, while $d,r,q$
stay fixed. This contradicts the fact that the minimum degree tends
to infinity. The contradiction works separately for NS and PC.

Now take any $n\ge9$ and let $d$ be its finite minimum degree in the
chosen system. If $n\ge4d+5$, the preceding contradiction applies.
Therefore $n<4d+5$, which rearranges to
$d>(n-5)/4$. Absence of any recurring finite degree is exactly the
T2 case of the later classification.
\end{proof}

The conclusion concerns unavoidable proof degree at every sufficiently
large size. It uses an infinite progression, not a drop in the degree:
the lifted proofs would all have the same budget $d$. A proof of
arbitrary length cannot evade the bound, because lifting transports
the complete degree-$d$ consequence space.

\section{What the later common projector must do for this example}
The application has now isolated the jobs of the algebraic construction.
The common projector will be a linear projection obtained from one
shared combination of label permutations. Write $\Lift$ for the
resulting linear map from smaller-rank polynomials to larger-rank ones.
The variable basis has two-label support; monomials of degree at most
$d$ use at most $2d$ labels. The usable-axiom source and the NS
multiplier--axiom source are covered by $r=2d+3$. Every closure stage
stays inside these fixed ambient spaces.

Consider the row input $\sum_{j\le n}F_{ij}-1$. Plain inclusion into
a larger universe leaves out its added summands. Setting all variables
involving deleted labels to zero instead turns a deleted row into
$-1$. The common projector must repair the actual axiom coefficient
map. Its axiom-lifting formula gives a linear combination of genuine
larger-instance inputs, rather than assuming that the same retained
polynomial remains an axiom.

For PC it must also preserve $V_j\cap P_{\le d-1}$ after cancellation,
and the joint tensor multiplication map. This ensures that every
lifted vector in the larger closure has a legal ordinary-PC derivation.
The target here is $1$, so the identity $\Lift(1)=1$ turns consequence
lifting into refutation lifting. The retained-label cost $2r-1=4d+5$
is what can be inverted to obtain the linear bound.

The separate comparison theorem uses a uniformly bounded chain of
invariant subspaces and an NS plateau to contain the whole PC closure.
In this example it rules out bounded PC degree before lifting supplies
the numerical lower bound. The restriction to the induction chain,
the clearance invariant, comparison, and the common projector each
have one task. Together they turn a sequential finite-order
contradiction into the same linear-degree mechanism used throughout
the book.

\part{The lifting principle and the four profiles}

\chapter{The lifting problem: transport a closure and its target}
\label{tx:chap:contract}
The syntax of an NS or PC proof has now been replaced by a finite
membership calculation. We can therefore ask a more general linear
question before constructing any operator: which parts of that calculation
must survive a change of universe?

We separate three tasks: identify the actual polynomial maps, preserve
the subspace calculation, and detect whether a target remains outside
the consequence space. The finite constructions later in the book
will fulfil this single transport contract.

\section{From one eligible proof to an infinite progression}
It helps to see the whole argument before its linear-algebra interfaces.
Recall that $k$ bounds the labels needed by one variable, $s$ bounds
the labels in one axiom index, and $t$ bounds those in one target index.
We work over a fixed field of positive characteristic $p$.
Fix one proof degree $d\ge1$, and choose a support budget $r$ covering
the degree-$d$ data. A fully invariant target uses $r=kd+s$; the
general indexed-target calculation uses $r=2kd+s+t$. The completed
lifting theorem will say that, for
\[
 q=p^a>r,\qquad n\ge\max\{n_*,2r-1\},
\]
a prescribed target request has proofs of degree at most $d$ at size
$n$ exactly when it does at size $n+q$.
Here $n_*$ is a fixed starting rank after which the template and target
data are stable; its precise value is given in
Chapter~\ref{tx:ch:profiles}. The number $p$ is the characteristic,
and $q$ is the size of the added block, not necessarily the size of
the coefficient field.

The first implication can be repeated with the same $d,r,q$, because
every larger starting size still satisfies the displayed inequality.
Thus one proof gives a constant degree bound on an infinite
progression. The reverse implication, applied also at lower degree
thresholds, identifies exact recurring minimum degrees. This is the
bridge to the four profiles.

Now suppose a finite value $d$ cannot recur. It cannot occur at an
eligible starting size, so its occurrence forces
\[
                    n<\max\{n_*,2r-1\}.
\]
Since $r$ is affine in $d$, this is a linear lower bound on $d$ once
$n$ is beyond the fixed starting rank. In T4 a value does recur; the
two-way comparison and a bound on $q$ instead move an occurrence to
a representative of size $O(d)$, creating a finite linear peak.
The remaining mathematical task is to prove the displayed lifting
statement with its linear retained-label budget. The definitions below
split that task into explicit obligations.

\section{Label actions in linear-algebra language}\label{tx:sec:action-primer}
We need some language from linear algebra and representation theory.
It is useful language, but none of the words in this section should
stand for an unexplained theorem. We introduce them through the
polynomial spaces already familiar from Chapter~\ref{tx:ch:closure}.

\subsection{A vector space together with relabelling}
Consider the degree-one space
\[
                  M=Kx_1\oplus Kx_2\oplus Kx_3.
\]
The permutation $(1\ 2)$ exchanges $x_1$ and $x_2$ and fixes $x_3$.
It therefore sends
$a x_1+b x_2+c x_3$ to $b x_1+a x_2+c x_3$.
This is an invertible linear map on coefficient vectors. Other
permutations give other linear maps, and doing two relabellings in
succession has the same effect as their composite.

In general, a \emph{representation} of $S_n$ is a vector space $M$
equipped with these compatible linear actions of the label
permutations. A \emph{module} is another name for such a space here;
later the same action is extended to linear combinations of permutations.
The vectors remain ordinary coefficient vectors; the representation
records how relabelling acts on them.

For a finite labelled set $\mathcal U_n$, the notation $K[\mathcal U_n]$
means the vector space with one independent basis vector $[u]$ for
each $u\in\mathcal U_n$. Its vectors are
$\sum_{u\in\mathcal U_n}a_u[u]$. If permutations relabel the objects
$u$, they permute these basis vectors. This is a \emph{permutation
module}. Monomials of degree at most $d$ give precisely such a basis
for our polynomial ambient space.

An \emph{invariant subspace}, also called a \emph{submodule}, is a
subspace carried to itself by every label permutation. In the displayed
example, $K(x_1+x_2+x_3)$ is invariant, whereas $Kx_1$ is not.
The consequence spaces are invariant because relabelling an entire
legal proof gives another legal proof from the uniform axioms.
An individual proof line need not be invariant.

\subsection{Maps that respect the relabelling}
A linear map $f:M\to M'$ is \emph{equivariant} when
\[
                         f(\pi v)=\pi f(v)
\]
for every label permutation $\pi$. The equation says that it makes
no difference whether we relabel first or apply the map first.
For example, sending an axiom index $[a]$ to its polynomial $f_{a,n}$
is equivariant by uniformity. Ordinary variable multiplication is
equivariant when the \emph{same} permutation acts on the variable
and the polynomial being multiplied.

Equivariant maps encode rules that are compatible with renaming labels.
Their symmetry lets one common expression in permutations interact
with all the maps at once. Chapter~\ref{tx:chap:deletion} gives the
elementary calculation behind this compatibility.

\subsection{Kernel, image and quotient: three different questions}
For a linear map $f:A\to B$, its \emph{kernel} is
$\ker f=\{a:f(a)=0\}$; its \emph{image} is
$\im f=\{f(a):a\in A\}$. They record respectively what the map kills
and what the map can produce. If $f$ is equivariant, both are invariant
subspaces. Indeed, $f(a)=0$ implies $f(\pi a)=\pi f(a)=0$, and
$\pi f(a)=f(\pi a)$ is again an image vector.

For a subspace $H\subseteq W$, the quotient $W/H$ identifies two
vectors whose difference belongs to $H$. The quotient map is
\[
                  \rho_H:W\to W/H,\qquad g\mapsto[g]=g+H.
\]
Its kernel is exactly $H$. Thus
\[
                         [g]=0\text{ in }W/H
                              \quad\Longleftrightarrow\quad g\in H.
\]
This will express the question of whether a target is derivable.
The zero \emph{class} does not mean that the original polynomial is
zero: a nonzero derived polynomial has zero class modulo the
consequence space.

For example, let $W=Kx\oplus Ky$ and $H=K(x+y)$. Then
$[x+y]=0$, while $[x]\ne0$ and $[y]\ne0$, with $[x]=-[y]$.
To see that $[x]\ne0$, an identity $x=a(x+y)$ would force both
$a=1$ and $a=0$ by comparing the two coefficients.
This calculation works in characteristic two as well.
If $H$ is invariant, relabelling the representatives defines an action
on the quotient: $\pi[g]=[\pi g]$.
A \emph{subquotient} of $W$ is a quotient $U/H$ with
$H\subseteq U\subseteq W$ invariant. It records the directions in
$U$ that are still missing from $H$.

\subsection{What the word exact will mean}
Suppose $A\xrightarrow{f}B\xrightarrow{g}C$ are linear maps.
The sequence is \emph{exact at $B$} if $\im f=\ker g$:
the vectors arriving from $A$ are exactly those killed by the map to
$C$. This is a statement about two specified subspaces of $B$.
For instance,
\[
        0\longrightarrow H\longrightarrow W
                      \xrightarrow{\rho_H}W/H\longrightarrow0
\]
says that inclusion is injective, its image is the kernel of the
quotient map, and the quotient map is surjective.

Later we will construct an operation $\Del$ taking suitable spaces and
maps at rank $N=n+q$ to spaces and maps at rank $n$.
To say that this operation is \emph{exact} means that it preserves
these image--kernel relationships, with their actual inclusions.
It must carry an injection to an injection, and an onto map to an
onto map, on the specified class. It must also identify the transported
kernel and image \emph{inside the transported ambient space}.
Matching their dimensions would not do this.

The one-dimensional space $K$, with every permutation acting trivially,
is called the \emph{trivial module}. It is \emph{simple}: it has no
proper nonzero subspace. We will use this particular observation for
invariant targets, rather than assume a general structure theorem
about representations.

\subsection{Background for the transport construction}
There are four layers of background here. The first is ordinary
linear algebra: spans, linear maps, kernels and quotient spaces.
The second is representation theory: we keep track of relabelling
while doing that linear algebra. Since the coefficient field has
positive characteristic, this is \emph{modular} representation theory.
One familiar shortcut is unavailable: averaging over all permutations
would divide by $n!$, which is zero in the field once $n\ge p$.
The later construction supplies a different operator rather than
using this averaging argument.

A representation is \emph{semisimple} if it is a direct sum of simple
representations. Here \emph{simple} means that there is no proper
nonzero invariant subspace. In positive characteristic, an invariant subspace
need not have an invariant complementary subspace. For example, the
swap of two coordinates over a characteristic-two field fixes
$K(x_1+x_2)$. Any invariant complementary line would be an eigenline
for the swap, but its only eigenvalue is one and its fixed space is
precisely that same line. Thus no invariant complement exists.
This elementary example explains why a decomposition into invariant
pieces cannot simply be assumed in the later construction.

The third layer is the language of categories and functors: a way to
say that one operation changes many spaces and their connecting
maps consistently. We spell out that requirement in
Section~\ref{tx:sec:grammar}. The fourth is the language of exactness,
usually taught in homological algebra. Here we use just its
linear-algebra content: preserving kernels, images and quotient
relations. No knowledge of chain complexes or homology groups is
needed. For optional background on this terminology, see
\cite{StacksExact}; the required calculation is developed in this book.

The tensor product introduced next belongs to ordinary multilinear
algebra. In our application it has a particularly concrete meaning:
one register for each possible variable multiplication.

\section{The data at one universe size}
\label{tx:sec:multiplication-registers}
We now put names on the ingredients of one ordinary PC step.
\emph{Nothing in this section changes the universe size.}
In particular, the map $m_n$ below is ordinary polynomial
multiplication at one fixed rank; it is not the lifting operator.
We first understand this familiar operation, then ask how to transport it.
The letter $z$ in the worked polynomial example below is just a
polynomial variable. In Chapter~\ref{tx:chap:deletion}, $z$ will instead
denote a linear operator made from block permutations; that operator
will be defined there.

\subsection{Three spaces with three different jobs}
Fix a universe size $n$ and a proof-degree budget $d\ge1$.
The variable names are the members $v$ of the finite set
$\mathcal V_n$; write $x_v$ for the corresponding polynomial variable.
There may be many more variables than universe labels.
For example, variables $x_{ij}$ give $n^2$ names.
The two numbers must not be identified.

The spaces we shall use are
\[
 \begin{array}{ll}
 P_d(n)=P_{n,\le d}
   &\text{all polynomials allowed as proof lines},\\
 L_d(n)=P_{n,\le d-1}
   &\text{all polynomials small enough to multiply once},\\
 X(n)=K[\mathcal V_n]
   &\text{one independent basis vector $[v]$ per variable name}.
 \end{array}
\]
$P_d(n)$ contains both derivable and underivable polynomials.
$L_d(n)$ also contains both kinds; it imposes only a degree restriction.
If $V_j(n)$ is the space reached after $j$ closure stages, the part
we may actually multiply is $V_j(n)\cap L_d(n)$.
The space $X(n)$ records which variable we use. We do \emph{not}
require that variable itself to have a proof.

For an explicit coordinate picture, take variables $x,y,z$ and $d=2$.
Then
\[
 \begin{split}
 P_2(n)&=\Span\{1,x,y,z,x^2,xy,xz,y^2,yz,z^2\},\\
 L_2(n)&=\Span\{1,x,y,z\},\\
 X(n)&=\Span\{[x],[y],[z]\}.
 \end{split}
\]
Here the brackets indicate variable-name basis vectors, not quotient
classes. This is a fixed-instance illustration; it does not specify
a uniform axiom family or compute a lower bound.

\subsection{Multiplication by one specified variable}
For each variable $x_v$, consider the map
\[
     M_v:L_d(n)\to P_d(n),\qquad M_v(f)=x_vf.
\]
It is linear: $x_v(af+bg)=a\,x_vf+b\,x_vg$.
It stays inside $P_d(n)$ because $\deg f\le d-1$.
For $d=2$, multiplication by $z$ sends the coefficient vector
$a+bx+cy+ez$ to $az+bxz+cyz+ez^2$.
This is an ordinary matrix calculation.

Now we want one map that packages \emph{all} the maps $M_v$.
Its input should allow a separate polynomial $f_v$ for each variable:
one slot for what is multiplied by $x$, one for what is multiplied by
$y$, and so on. On these slots the output should be
$\sum_v x_vf_v$. A direct sum of one copy of $L_d(n)$ for each
variable does exactly this.

\subsection{Variable-labelled copies and tensor products}\label{tx:sec:tensor-registers}
The space of variable-labelled inputs is $X(n)\otimes L_d(n)$.
Give it one basis vector
$[v]\otimes\mu$ for each pair consisting of a variable name $v$
and a monomial $\mu$ of degree at most $d-1$.
Extend by linear combinations. Grouping these basis vectors by $v$
identifies the space with those variable-labelled copies of $L_d(n)$:
\[
       X(n)\otimes L_d(n)\ \simeq\
                     \bigoplus_{v\in\mathcal V_n}L_d(n).
\]
Every vector has a unique expression $\sum_v[v]\otimes f_v$.
The notation $[v]\otimes f$ means: put the polynomial $f$ in the
slot labelled by $v$, with zero in every other slot.

For the three-variable example this space has dimension $3\cdot4=12$.
It remembers inputs such as $[x]\otimes y$ and $[y]\otimes x$
as distinct vectors, even though their products are the same
polynomial $xy$. It also allows sums such as
$[x]\otimes f+[z]\otimes g$, with two different multiplication requests.

The tensor symbol is different from an ordinary Cartesian product.
Polynomial multiplication on pairs
$(a,f)\in X(n)\times L_d(n)$, after identifying $a$ with its
linear polynomial, is \emph{bilinear}: it is linear in either
argument when the other is fixed. It is not a linear map on the
vector space of pairs. For example,
\[
 ( [x],0)+(0,y)=([x],y),
\]
whereas multiplication is zero on the two summands and is $xy$
on their sum. The tensor product makes all these bilinear requests
into one \emph{linear} map. In a basis it is just a matrix on the
pair coordinates described above.

In particular,
$([x]+[y])\otimes f=[x]\otimes f+[y]\otimes f$ and
$[x]\otimes(f+g)=[x]\otimes f+[x]\otimes g$.
These rules distribute sums in either input; the monomial-pair
description already implements them.
The variable-labelled-copy interpretation uses the specified basis
of $X(n)$. It is not a claim that a general tensor product comes
with a preferred decomposition into copies.

The same coordinate construction applies to arbitrary finite-dimensional
spaces $A$ and $B$. After choosing bases $a_i$ and $b_j$, their tensor
product has basis $a_i\otimes b_j$. A bilinear map $\beta:A\times B\to C$
extends to a linear map on these coordinates by sending
$a_i\otimes b_j$ to $\beta(a_i,b_j)$. For linear maps
$f:A\to A'$ and $g:B\to B'$, the tensor map is defined by
$(f\otimes g)(a\otimes b)=f(a)\otimes g(b)$ and extended linearly.
Changing either basis changes the coordinates by this same rule;
the construction does not depend on which bases were chosen.
If labels act on both spaces, their \emph{diagonal action} is
$\pi(a\otimes b)=(\pi a)\otimes(\pi b)$. Each summand uses the same
permutation in both positions. This is the tensor action used when
multiplication and later detection pairings are transported.

\subsection{The single multiplication map}
Combining the variable multiplications gives the linear map
\[
 \boxed{\quad
 m_n:X(n)\otimes L_d(n)\longrightarrow P_d(n),\qquad
 m_n\!\left(\sum_v[v]\otimes f_v\right)=\sum_v x_vf_v.
 \quad}
\]
The subscript $n$ only tells us at which universe size the variables
and polynomials are being used. On a basis pair it sends
$[v]\otimes\mu$ to $x_v\mu$.
This determines the whole map by linearity.
Some later formulas write $x_v$ itself for the basis vector of $X(n)$;
at its output, $x_v$ always means the actual polynomial variable.

This map need not be injective. For instance,
$[x]\otimes y-[y]\otimes x$ is a nonzero input vector with output zero.
Different multiplication requests can produce the same polynomial;
their difference belongs to the kernel.
The map is not being asked to recover the request from its output.
We use its \emph{image} to collect the polynomials obtained by legal
multiplications.

\subsection{A complete closure step with cancellation}
Return to the three axioms
\[
 f_1=xy+x,\qquad f_2=xy+1,\qquad f_3=xz
                         \quad\text{over }\mathbb F_2,
\]
and keep the degree budget $d=2$. We first take all their linear
combinations, $V_0=\Span\{f_1,f_2,f_3\}$. Comparing the coefficients
of $xy$ and $xz$, as in Chapter~\ref{tx:ch:closure}, gives
\[
                         V_0\cap L_2=\Span\{x+1\}.
\]
Both $f_1$ and $f_2$ were quadratic. Their sum $x+1$ is linear
and is now eligible for multiplication.
The input $[z]\otimes(x+1)$ lies in
$X\otimes(V_0\cap L_2)$, and
\[
                 m_n([z]\otimes(x+1))=xz+z.
\]
Adjoining the old vector $f_3=xz$ gives the target $z$.
The corresponding PC lines are completely explicit:
\[
 \begin{array}{c|l|l}
 \text{line}&\text{polynomial}&\text{reason}\\ \hline
 1&xy+x&\text{input }f_1\\
 2&xy+1&\text{input }f_2\\
 3&xz&\text{input }f_3\\
 4&x+1&\text{add lines 1 and 2}\\
 5&xz+z&\text{multiply line 4 by }z\\
 6&z&\text{add lines 5 and 3}.
 \end{array}
\]
Every line has degree at most two. Multiplying line 1 or line 2
directly by $z$ would instead give a cubic and exceed the budget.
This is why the operation must first identify the low-degree
\emph{combinations}, rather than select low-degree inputs.

For a general current stage the same calculation reads
\[
 \underbrace{V_{j+1}(n)}_{\text{next stage}}
 =\underbrace{V_j(n)}_{\text{old consequences}}
  +\underbrace{m_n\bigl(X(n)\otimes
                   (V_j(n)\cap L_d(n))\bigr)}_{\text{allowed new multiples}}.
\]
The plus sign is a sum of subspaces: take all sums of one vector
from each displayed space. The image term is itself a subspace
because $m_n$ is linear.

\subsection{The input, inclusion and target maps}
There are two other PC maps besides multiplication.
Let $\mathcal A_{n,\le d}$ be the axiom indices whose polynomials
fit the degree budget. A vector in $K[\mathcal A_{n,\le d}]$
stores one coefficient for each such index. The map
\[
        \alpha_n:K[\mathcal A_{n,\le d}]\to P_d(n),
                  \qquad [a]\mapsto f_{a,n}
\]
forms the resulting linear combination of usable axioms.
Thus $V_0(n)=\im\alpha_n$. Distinct input columns may be dependent;
the image, not the number of columns, is the initial consequence space.
The map
$\ell_n:L_d(n)\hookrightarrow P_d(n)$ is simply inclusion:
the same polynomial is now regarded as a possible degree-$d$ line.
Remembering this inclusion is necessary when forming an intersection.

For NS, instead put one coordinate on each permitted pair $(\mu,a)$,
where $\mu$ is a monomial and
$\deg\mu+\deg f_{a,n}\le d$. Calling this set $\mathcal C_d(n)$,
the certificate map is
\[
 \phi_n:K[\mathcal C_d(n)]\to P_d(n),\qquad[\mu,a]\mapsto\mu f_{a,n}.
\]
Its image $I_d(n)$ is one span, with no PC iteration.
The set $\mathcal C_d(n)$ here consists of certificate \emph{indices};
it is not the representation class $\mathcal C_r(N)$ introduced
in Chapter~\ref{tx:chap:deletion}.

Finally, the requested targets have a separate column map
\[
       J_n:K[\mathcal B_n]\to P_d(n),\qquad[b]\mapsto G_{b,n}.
\]
We use this map into $P_d(n)$ only when every full target degree
is at most $d$. A budget below a stabilised target degree cannot
derive that target; we do not truncate the target to make it fit.
The map records what we ask to derive, not extra input axioms.
Adding its columns to $\alpha_n$ would change the problem.

\begin{checkpoint}
$P_d$ contains possible proof lines; $V_j$ contains the lines already
reachable at stage $j$; $L_d$ is the degree test; $X$ labels the permitted
variables. The map $m_n$ assembles ordinary variable multiples.
Only after understanding these fixed-rank ingredients do we ask for
an operation relating them at two different ranks.
\end{checkpoint}

\section{Building a consequence space by linear operations}\label{tx:sec:grammar}
We build each next consequence space by four familiar tasks:
find an intersection, make variable-labelled copies, apply the
multiplication matrix, and add the old space. We now spell out that
recipe and why kernels and images describe it.

\subsection{A direct sum stores two vectors; a sum combines them}
If $U,V$ are subspaces of the same ambient $P$, then $U\oplus V$
is the space of \emph{pairs} $(u,v)$. In contrast, $U+V\subseteq P$
is the space of outputs $u+v$. The linear map
\[
              \sigma:U\oplus V\to P,\qquad (u,v)\mapsto u+v
\]
has image $U+V$.
These spaces need not have the same dimension.
For $U=V=Kx$, the input has two independent coordinates,
while every output is a multiple of $x$. The pair $(x,-x)$
belongs to the kernel. The direct sum remembers the two requests;
the sum forgets how an output was divided between them.

\subsection{An intersection is found by solving a difference equation}
Keep the same ambient space $P$ and define
\[
              \delta:U\oplus V\to P,\qquad(u,v)\mapsto u-v.
\]
A pair belongs to its kernel exactly when $u=v$. The common vector
then belongs to both $U$ and $V$. Thus the map
\[
                  U\cap V\longrightarrow\ker\delta,
                              \qquad h\mapsto(h,h)
\]
is an isomorphism, with inverse $(u,u)\mapsto u$.
Strictly, the kernel is a space of pairs, whereas the intersection
is a subspace of $P$. The inverse just given, followed by inclusion
in $P$, specifies how to identify them.
This small distinction matters when the whole calculation is transported.

In the cancellation example, $U=V_0$ and $V=L_2$.
A kernel vector records a linear combination of the three axioms
which equals a degree-one polynomial. Solving the difference equation
gives the common line $K(x+1)$. It finds a cancelled combination,
not merely a subset of the original input columns.

\subsection{Read the closure recurrence as a four-step algorithm}
Suppose the current embedded space $V_j(n)\subseteq P_d(n)$ is known.
\begin{enumerate}
\item Find $U_j(n)=V_j(n)\cap L_d(n)$ by the difference-map kernel
      above. These are the already derived polynomials eligible
      for one further multiplication.
\item Make one copy of $U_j(n)$ for each variable name.
      This is $X(n)\otimes U_j(n)$, sitting in
      $X(n)\otimes L_d(n)$ by the actual inclusion $U_j(n)\subseteq L_d(n)$.
\item Apply the restriction of $m_n$ to that input space, and take
      its image. Call the image $M_j(n)$. It consists of all sums
      $\sum_v x_vf_v$ with $f_v\in U_j(n)$.
\item Apply the sum map from $V_j(n)\oplus M_j(n)$ to $P_d(n)$.
      Its image is the next space $V_{j+1}(n)$.
\end{enumerate}
All these operations are finite linear algebra at one rank.
Finding kernels or images is Gaussian elimination; making the
variable-labelled copies is just organising coordinates.
Each vector thereby obtained has a legal proof, by the constructive
argument in Proposition~\ref{tx:prop:closure}.

An unchanged dimension would not tell us that the right space had
been obtained. For example, $Kx$ and $Ky$ are isomorphic
one-dimensional subspaces of $Kx\oplus Ky$, but only the first
contains the target $x$. We must keep the embeddings and the
specified maps, not just abstract isomorphism types.

\subsection{Diagram and functor: the language behind the recipe}
A \emph{diagram} in this book is a collection of specified vector
spaces and linear maps between them. Calling the data of Section 5.3
the \emph{degree-$d$ polynomial diagram} simply says that we remember
the ambient, low-degree, variable, input and target spaces together
with their maps. The inclusions are part of the data.

The vocabulary comes from category theory. A \emph{category} organises
objects and the maps allowed between them, with identity maps and
composition. Here the objects are spaces with label actions and the
allowed maps are equivariant linear maps. A \emph{functor} coherently
assigns a new object to each object and a new map to each allowed
map: it preserves identities and composition. Thus, for maps
$A\xrightarrow{f}B\xrightarrow{g}C$,
\[
                \Del(gf)=\Del(g)\Del(f).
\]
This equation ensures that a composite calculation can be
transported either all at once or step by step.
The operation used here is linear on the allowed maps and must
also preserve the finite direct sums needed in the recipe.

An identification between transported and smaller-rank spaces must
also respect their connecting maps. For fixed-data maps $f_N$ and
$f_n$, write these identifications as $\eta_A$ and $\eta_B$. The
required square is
\[
\begin{CD}
\Del(A_N) @>{\Del(f_N)}>> \Del(B_N)\\
@V{\eta_A}VV @VV{\eta_B}V\\
A_n @>{f_n}>> B_n .
\end{CD}
\]
It \emph{commutes} when the two paths have the same value:
$\eta_B\Del(f_N)=f_n\eta_A$. A \emph{natural} identification satisfies
this compatibility for the allowed maps. This is the content of
naturality used here; it is not a claim that auxiliary witness
coefficients have a unique choice.

We have \emph{not} constructed this operation yet.
Write $\Del$ for the operation we want, from a larger rank
$N=n+q$ to the retained rank $n$. We must first identify the transported
fixed data with the data at $n$, so that the transported multiplication
map really is $m_n$, rather than some unrelated matrix of the same size.
We must then show that it preserves the kernels, images and sums
used by the four-step recipe.

If those statements hold, the induction is straightforward.
Transport the initial image $\im\alpha_N$ to $\im\alpha_n$.
If the current space has been transported, preserve the difference-map
kernel, the variable-labelled copies with their inclusion, the
multiplication image, and finally the sum image.
The next stage is therefore transported too.
At a stage past both stabilisation points, this gives transport of
the entire degree-$d$ consequence space.
No individual proof length has to be bounded or followed through the
change of universe.

This is why exactness is useful here. It is the preservation of a
concrete linear-algebra recipe. The formal contract below records the
three things which will make this reasoning valid; its existence
is the subject of the later construction.

\begin{definition}[Transport contract for a polynomial diagram]\label{tx:def:polynomial-diagram}
At a fixed degree $d\ge1$, a transport from rank $N$ to rank $n$
has three obligations.
\begin{enumerate}
\item \textbf{Identify the data.} Identify the transported ambient,
lower-degree and basis spaces with their smaller-rank counterparts,
carrying $\alpha,\ell,m,\phi,J$ to the actual smaller-rank maps.
For multiplication, the identification
$\Del(X(N)\otimes Y)\simeq X(n)\otimes\Del(Y)$ must respect the
inclusions of the subspaces $Y\subseteq L_d(N)$ used in the recurrence.
\item \textbf{Preserve the calculation.} Preserve kernels, images,
finite direct sums and their defining maps throughout the closure and
target calculation. This transports intersections, sums and every
closure stage as embedded subspaces.
\item \textbf{Reflect target vanishing.} For $H_N=I_d(N)$ or $C_d(N)$,
put $Y_N=(H_N+\im J_N)/H_N$. Require
$\Del(Y_N)=0\Longrightarrow Y_N=0$. This last condition says that
transport cannot hide a requested target outside the consequence space.
\end{enumerate}
\end{definition}

The next two lemmas separate the work: the first two obligations
transport the closure, while the third turns that transport into a
two-way target-membership test.
Throughout this discussion, transport includes the coherent maps and
finite direct sums just described. The tensor identification is
compatible with the actual subspace inclusions, not an independently
chosen isomorphism for each space.

\begin{lemma}[Transporting a closure recurrence]\label{tx:lem:closure-transport}
Suppose an operation $\Del$ takes the larger-rank ambient, lower-degree,
variable and usable-axiom spaces and maps to their actual smaller-rank
counterparts. Suppose it preserves kernels and images on all the spaces
used by the recurrence, and has the natural permutation-first-factor
identification
\[
 \Del(X(N)\otimes Y)\simeq X(n)\otimes\Del(Y).
\]
Then, as embedded subspaces of the identified smaller ambient space,
\[
                         \Del(V_j(N))=V_j(n)
                         \quad\text{for every }j.
\]
It also transports the stable PC closure. Transport of the certificate
map gives $\Del(I_d(N))=I_d(n)$.
\end{lemma}
\begin{proof}
The initial equality is preservation of the actual axiom-map image.
Given the equality at stage $j$, express the intersection as the kernel
of $(v,l)\mapsto v-l$. Preserve that kernel and its inclusion. Apply
the natural tensor identification and the actual multiplication map.
Preserve its image and then its sum with the old stage. This gives the
next equality. Both finite-dimensional chains eventually stabilise;
choose a stage past both stabilisation points to obtain the closure
statement. NS is the single certificate-map image calculation.
\end{proof}

\begin{checkpoint}
The content is compatibility of the diagrams, not equality of dimensions.
An isomorphism chosen separately for every space need not carry the
intersection, multiplication map or target vector to the right one.
\end{checkpoint}

\section{Preserving consequences and detecting missing targets}\label{tx:sec:target-reflection}
We now reach a separate question. Suppose the calculation of the
consequence space has been transported correctly. Does that already
tell us, in both directions, whether the prescribed target is derivable?
The answer depends on whether transport can erase the obstruction
to deriving it. We start with the ordinary quotient test.

\subsection{A target leaves a remainder modulo the consequences}
Let $H_N$ be $I_d(N)$ for NS or $C_d(N)$ for PC, inside
$W_N=P_d(N)$. For a single target $g$, its class $[g]\in W_N/H_N$
is zero exactly when $g$ has a proof within the degree budget.
Think of the quotient as recording only directions not yet proved.

For several targets the composite map
\[
 K[\mathcal B_N]\xrightarrow{J_N}W_N
                         \xrightarrow{\rho_{H_N}}W_N/H_N
\]
sends each target index to that target's unproved class.
Every target is derived exactly when this map is zero.
Its image is naturally the space
\[
                       Y_N=(H_N+\im J_N)/H_N.
\]
To check the formula, map $H_N+\im J_N$ into $W_N/H_N$.
Its kernel is $H_N$, and its image is the target classes just described.
Thus $Y_N$ is a space of missing target directions, not a new
polynomial or a new axiom.
As before, the full target degrees must fit $d$ before we use
$J_N$ into $P_d(N)$; there is no degree truncation.

For the small example $W=Kx\oplus Ky$, $H=K(x+y)$, a request for
$x+y$ has zero obstruction. A request for both $x$ and $y$ has
obstruction $W/H\ne0$. Knowing that their sum is a consequence
does not derive the two summands individually.

\subsection{What exact transport establishes, and what is still missing}
If all the actual maps are identified correctly, exactness takes the
target image and its quotient to the corresponding calculation at $n$:
\[
                             \Del(Y_N)\simeq Y_n.
\]
If $Y_N=0$, then $Y_n=0$. A derived request therefore remains derived
after this downward transport.
But from $Y_n=0$ we know only $\Del(Y_N)=0$.
To conclude $Y_N=0$ we need to know that this operation cannot
send a \emph{nonzero} missing-target space to zero.
That additional implication is called \emph{reflection of vanishing}.
In ordinary language: a missing consequence must still be detectable
after transport.

There is no conflict with exactness. Even if an injection is
transported to an injection, both its transported source and its
transported destination might be zero.
The map $0\to0$ is injective and surjective.
Preserving relations among surviving spaces does not by itself
require every original space to survive.

\subsection{An exact operation which really loses an obstruction}
Here is a concrete illustration. Work with pairs of vector spaces
$(A,B)$, with maps given by pairs of linear maps
$(f_A,f_B)$. Maps act separately on the two components and are not
allowed to mix them. Define $F(A,B)=A$: retain the first component
and forget the second.

Kernels, images, quotients and direct sums are computed componentwise.
For example, $\ker(f_A,f_B)=(\ker f_A,\ker f_B)$, and forgetting the
second component leaves $\ker f_A$. The same check works for images.
Consequently $F$ is an exact operation on these paired spaces.
This is not a claim that ordinary projection preserves every
intersection of arbitrary subspaces of $K^2$; our objects and allowed
maps are the componentwise ones just specified.

Now take
\[
 W=(K,K),\qquad H=(K,0),\qquad Q=(0,K),\qquad
                         J:Q\hookrightarrow W.
\]
The target map includes the second component. Its missing direction
has obstruction
$Y=(H+\im J)/H=(0,K)\ne0$.
After transport,
\[
             F(W)=K,\qquad F(H)=K,\qquad F(Q)=0,\qquad F(Y)=0.
\]
All the image--kernel calculations have been preserved, but the
target obstruction was forgotten.
This example explains the logical issue only; it is not the
polynomial deletion construction or a counterexample to our final
theorem. The additional survival lemma is precisely what prevents
this phenomenon on the relevant polynomial subquotients.

\subsection{How the actual proof will keep missing directions visible}
If a polynomial $g$ is outside $H$, elementary linear algebra supplies
a functional $\lambda$ with $\lambda(H)=0$ and $\lambda(g)=1$:
extend a basis of $H$ by $g$ and prescribe these values.
For $H=K(x+y)$, $g=x$, the example is
$\lambda(ax+by)=a-b$.
Such a functional certifies nonmembership without requiring an
invariant representative of a proof.

The permutation basis of the polynomial space has a dot product,
so a functional can be represented by a coefficient vector.
All functionals annihilating $H$ form $H^\perp$.
Concretely, this pairing is
$\langle a,b\rangle=\sum_\mu a_\mu b_\mu$ on monomial coordinates.
It is nondegenerate because each coordinate can be tested separately.
There is no positivity requirement over the coefficient field.
In characteristic two, for example, $(1,1)$ pairs to zero with itself
while pairing nontrivially with $(1,0)$.
Accordingly $H^\perp$ means the \emph{annihilator} of $H$ under this
pairing, not a chosen complementary subspace.
Permutations preserve the pairing because they permute its coordinates
simultaneously. The induced action on a functional is
$(\pi\lambda)(v)=\lambda(\pi^{-1}v)$; this is the usual dual action.
Pairing them with the missing classes gives a map onto $K$ whenever
the obstruction is nonzero. The pairing is compatible with
relabelings as a whole; an individually chosen separating functional
need not be invariant.

The formal argument in Chapter~\ref{tx:chap:assemble-transfer} uses
exactness and a tensor compatibility to show that, if deletion erased
the obstruction, it would erase the source of this onto pairing.
It would then produce a map from zero onto the surviving line
$\Del(K)\simeq K$, which is impossible.
This is the reason for a doubled ambient support budget.
One monomial coordinate uses at most $kd$ labels; a pair of such
coordinates uses at most $2kd$. We are bounding the coordinates
of the pairing's ambient basis, not the union of labels in an
arbitrary polynomial or functional.

We have described the idea here, not assumed the survival theorem.
The next lemma records the implication once survival is proved.

\begin{lemma}[The target-reflection contract]\label{tx:lem:target-contract}
In addition to the closure transport, suppose the actual target map is
transported and every nonzero subquotient of $P_d(N)$ survives deletion.
Then
\[
        \im J_N\subseteq H_N\quad\Longleftrightarrow\quad
        \im J_n\subseteq H_n.
\]
\end{lemma}
\begin{proof}
The two membership conditions are $Y_N=0$ and $Y_n=0$.
Exactness identifies $\Del(Y_N)$ with $Y_n$. Since $Y_N$ is a
subquotient of the ambient polynomial space, survival of nonzero
subquotients reflects its vanishing. This proves both directions.
\end{proof}

The survival lemma will come from a nondegenerate permutation-basis
pairing, with doubled support. It is proved in
Chapter~\ref{tx:chap:assemble-transfer}; it is not an extra unproved
premise of the final theorem.

\section{The easier invariant target}
The simplest important request is refutation: derive the polynomial
$1$. More generally, let $G_N$ be fixed by every label permutation.
The input space of its target map has only one coordinate, recording
the scalar coefficient with which that target is requested.
For a fully invariant polynomial $G_N$, the source is just $K$:
\[
                   j_N:K\to P_d(N),\qquad a\mapsto aG_N.
\]
Its image in $P_d(N)/H_N$ is zero or a copy of the simple trivial
module. The deletion of Chapter~\ref{tx:chap:deletion} fixes the trivial
module $K$, so exactness already reflects
whether this composite is zero. This explains why the invariant-target
budget is smaller than the general indexed-target budget. The constant
one is only one example of an invariant target.

Here is the argument without the module terminology. A map from a
one-dimensional space has either zero image or a one-dimensional
image. If the missing class $[G_N]$ is nonzero, its line has trivial
label action because $G_N$ was invariant. Deletion preserves the
image map and keeps that trivial line nonzero. It therefore cannot
hide this particular obstruction. An indexed request can instead
have a higher-dimensional missing space, so the doubled-support
pairing argument is needed there.

\section{The budget table}
The fixed variable support is at most $k$, axiom-index support at most
$s$, and target-index support at most $t$. The spaces have the following
safe bounds.
\begin{center}
\begin{tabular}{@{}ll@{}}
\toprule
Basis or ambient space & Support bound\\
\midrule
Degree-$d$ monomials & $kd$\\
Multiplier--axiom pairs & $kd+s$\\
Variable with degree-$d-1$ monomial & $k+k(d-1)=kd$\\
Target index & $t$\\
Two degree-$d$ monomial coordinates & $2kd$\\
\bottomrule
\end{tabular}
\end{center}
One affine choice covering the whole diagram and the survival argument is
\[
                           R(d)=2kd+s+t.
\]
For an invariant target the pairing is unnecessary, so use $R(d)=kd+s$.
Nothing in the PC recurrence increases the bound because its intermediate
spaces are subspaces of those fixed ambient spaces. In particular, the
support budget does not grow with proof length.

\section{The connection with representation stability}\label{tx:sec:stability-context}
The contract belongs to a familiar mathematical pattern: hold the
complexity of a labelled object fixed and let the number of available
labels grow. Here the complexity parameter is proof degree $d$, and
one monomial uses at most $kd$ labels. The ambient spaces grow, but
the allowed equality patterns among the labels in a basis object
come from a finite list. Representation stability studies how such
families of symmetric-group representations are related across ranks.
Three references help locate our argument within that setting.
The terminology in this section concerns the surrounding literature.
A \emph{categorical equivalence} matches objects and their allowed
maps, with an inverse up to compatible isomorphisms. An
\emph{abelian category} is a setting in which kernels, images and
quotients have the formal properties familiar from modules.
Spaces with equivariant linear maps are an example. Our construction
below will establish the needed maps directly.

\subsection{Church--Ellenberg--Farb: the broad fixed-degree picture}
An FI-module is a functor from finite sets and injections to vector
spaces: every injection has an assigned linear map, and compositions
are respected. Church--Ellenberg--Farb~\cite{CEF} turn representation
stability into a finite-generation condition. Their characteristic-zero
results give eventual character polynomials and polynomial dimensions;
fixed-degree configuration-space cohomology is one motivating example.
This provides a useful comparison with fixing polynomial degree here.
A character is the trace of a permutation acting on a representation.
On a permutation basis it counts fixed basis objects. For example,
the single-label space has character equal to the number of fixed
labels; the space on ordered pairs of distinct labels has character
$a(a-1)$ when a
permutation fixes $a$ labels. A \emph{character polynomial} describes
such traces using counts of cycles of specified lengths.
The cited cohomology examples are background applications of that
framework, not prerequisites for this book's proof.

Our consequence spaces need not form an FI-submodule under ordinary
inclusion. For example, the uniform axiom $f_n=\sum_{i<n}x_i$ has
degree-one consequence space $Kf_n$. Under $[n]\hookrightarrow[n+1]$,
ordinary inclusion takes it to $K\sum_{i<n}x_i$, which is not contained
in $Kf_{n+1}$. The new coefficient of $x_n$ proves the failure directly.
Thus fixed-degree FI intuition does not supply the consequence-transfer
theorem automatically. We have to construct the maps which preserve
the actual uniform equations.

\subsection{Harman: the close modular-periodicity antecedent}
Harman~\cite[Theorem 2.5 and Corollary 2.11]{Harman} constructs
periodic equivalences for bounded-weight permutation-module categories
and extends them to the relevant abelian categories. The stable range
is linear: his permutation result compares $n,m>2r$. In positive
characteristic, the suitable residue of the rank matters, rather than
one characteristic-zero eventual-stability statement.
The structural principle of linear-range modular periodicity therefore
predates this proof-complexity application.
His Proposition~2.7 includes compatibility with natural operations on
permutation modules; Proposition~2.12 includes induction, restriction,
tensor products and duality on the abelian envelope. Compatibility with
these operations is therefore also an antecedent. Here we construct
the actual polynomial diagrams independently and prove that their
closure operations and prescribed-target obstructions survive deletion,
with the support and rank bounds needed for the degree application.
The bounded parameter in this comparison controls the small
non-growing part of a permutation-module index. In the tuple
models used here, the corresponding concrete parameter is the
number of labelled positions $b\le r$.

Keep the period convention precise. This book uses a prime power
\emph{strictly greater} than the support budget. Harman's printed
$p^{\lceil\log_p r\rceil}$ needs care when $r$ is itself a power of
$p$; the strict choice avoids that boundary. For example,
$\binom52=0$ and $\binom72=1$ modulo two, although $5$ and $7$ agree
modulo two. Counts of degree $r$ can detect a digit beyond that
non-strict modulus. The present argument proves its strict-period
statement directly.

\subsection{Hone and coauthors: an explicit functor viewpoint}
Hone, Klein, Pauwels, Sherman, Yacobi and Zhang~\cite{HKPSYZ}
construct explicit periodic functors through semisimplification,
motivated by the Duflo--Serganova functor in Lie superalgebra theory.
Their symmetric-group result in Theorem 8.5.3 gives a closely related
categorical setting. It helps explain why an operation which removes
a prime-power block can have more structure than a numerical
coincidence of dimensions. Its hypotheses and restricted categories
matter; this is not global exactness on all representations.
In that literature, \emph{semisimplification} quotients a category by
maps invisible to trace tests: $f:A\to B$ is negligible when
$\operatorname{tr}(gf)=0$ for every allowed $g:B\to A$.
This is an operation on maps, not merely replacing each representation
by a list of its simple factors.
The Duflo--Serganova analogy starts with a parity-changing square-zero
operator on a graded representation and takes its kernel modulo image.
For the present proof the concrete quotient below is sufficient;
none of the superalgebra machinery is assumed.

The concrete operation here is
$\Del(M)=\ker z_M/(\ker z_M\cap\im z_M)$. Here $z_M$ is a linear
operator on a representation $M$, formed from permutations of the
added block; Chapter~\ref{tx:chap:deletion} gives its definition.
Its powers mean repeated composition of that operator. In characteristic two,
$z^2=0$ makes this look like the homology of a square-zero operator.
Homology generally fails to preserve injections, and the two-point
example in Chapter~\ref{tx:chap:deletion} shows that failure directly.
For a square-zero map $d:M\to M$, the relation $d^2=0$ implies
$\im d\subseteq\ker d$. In homological language, kernel vectors are
\emph{cycles}, image vectors are \emph{boundaries}, and
$\ker d/\im d$ is the \emph{homology}: the cycles remaining after
boundaries are identified with zero.
This vocabulary describes a quotient of two concrete subspaces.
Our characteristic-two deletion has this form. In odd characteristic,
$z^p=0$ need not give $\im z\subseteq\ker z$, so the intersection
in the displayed quotient remains essential.
Our common group-algebra projector proves exactness on the required
bounded-support class, with the actual coefficient maps. In odd
characteristic the quotient is defined exactly as displayed, not
silently replaced by $\ker z/\im z^{p-1}$.

The proof which follows is self-contained: it constructs the witness,
proves exactness and transports the whole NS or ordinary-PC diagram.
The references explain the mathematical setting and the antecedents;
they are not unproved inputs to the argument. The application requires
prescribed-target reflection, the separate degree-zero case and the
onset/period estimates used in Chapter~\ref{tx:ch:profiles}. No claim
that modular periodicity itself originates here is needed.

\section{How the rest of the proof fulfils the contract}
We add $q=p^a>R(d)$ labels and retain at least $2R(d)-1$.
The prime-power block defines a nilpotent operator $z$. On permutation
bases a quotient of its kernel deletes exactly the block labels, but
that quotient is not automatically exact on subspaces. One common
group-algebra construction produces an idempotent $e$ which models the quotient
on every needed space at once: $e^2=e$, so $e$ acts as a projection.
That common expression supplies exactness
and map compatibility. The only hard quantitative question is how many
retained labels are needed to build it.

In characteristic two, the preliminary construction uses $1+2+\cdots+r$ labels for a support
budget $r$. The improved construction reuses one pool and requires only
$\max_{b\le r}(2b-1)$. Once these bounds are proved, this chapter's
contract becomes the actual transfer theorem rather than a conditional
plan. The next chapter first shows exactly how that theorem yields the
four types, so we know why the improvement from a sum to a maximum matters.

\chapter{From fixed-degree transfer to four global profiles}\label{tx:ch:profiles}
We now extract the global conclusions from the transfer theorem. This
chapter does not yet construct the label-deletion operator; the later
algebraic chapters do that. Its purpose is to show exactly what the
transfer statement buys us, which direction of transfer is used, and
where the linear lower bounds enter.

Fix a field of prime characteristic $p$, a uniform axiom template, and
a uniform target request. Let $h(n)$ be either its NS or ordinary-PC
profile. It may be an all-target profile or the profile of one fixed
target orbit. Recall the fixed bounds: $k$, $s$ and $t$ bound the
variable-label, axiom-index and target-index supports, respectively;
$D_G$ bounds target degree, and $n_0$ is the template's starting rank.
For general indexed targets put
\[
 \kappa=2k,\qquad \sigma=s+t,\qquad R(d)=\kappa d+\sigma,
 \qquad n_*=\max\{n_0,1,kD_G+t\}.
\]
For a single invariant target one can instead use
$\kappa=k$, $\sigma=s$, and $n_*=\max\{n_0,1,kD_G\}$. These are
fixed constants of the problem, while $d$ is a proposed proof degree.

\section{What lifting does, and what it does not do}
The transfer contract of Theorem~\ref{tx:thm:transfer} compares ranks
which differ by a $p$-power block $q>R(d)$. Once the retained rank is
large enough relative to $R(d)$ and the stable starting rank, it gives
\[
                         h(n)\le d\quad\Longleftrightarrow\quad
                         h(n+q)\le d.
\]
The degree $d$ is the same on both sides. We do not divide the degree
by the new universe size, lower it by one, or change the target to
its low-degree part. The operation transports the consequence space
and the prescribed target-membership test.

Suppose a forward lifting argument alone starts from a degree-$d$
proof at $n$. Iterating with the same admissible $q$ gives proofs of
degree at most $d$ at
\[
                            n,n+q,n+2q,\ldots.
\]
Their absolute degree bound remains $d$. The ratio $d/(n+\ell q)$
tends to zero, but that is not a statement that their absolute degrees
fall. Nor does this progression cover all sufficiently large ranks.
One cannot conclude that the entire profile is eventually bounded.

Forward lifting also need not preserve the exact minimum degree.
For a simple sequence, let $g(1)=2$ and $g(n)=1$ for $n\ge2$.
Every threshold predicate $g(n)\le d$ persists under the step $n\mapsto
n+1$, but the initial exact value two never reappears. Forward
persistence of a finite bound does force \emph{some} finite value to
recur, by the pigeonhole principle. The stronger conclusion that a
particular value repeats, and the exact periodic statements below,
come from two-way transfer at all the relevant thresholds.

\section{A period and a starting rank for each threshold}
A \emph{threshold predicate} answers a yes-or-no question at each
size: does the request have degree at most $d$? Its eventual
periodicity means that this answer repeats after the starting rank.
The period and the starting rank may both depend on $d$.
There need not be one period which works for the unbounded profile
as a whole.
Truncating $h$ at $d+1$ retains the exact values $0,\ldots,d$ and
records every larger value, including infinity, by the single symbol
$d+1$. This allows several threshold comparisons to be combined.

Use $N(d)$ for the periodicity onset, as in the paper. The argument
$(d)$ distinguishes this threshold from a larger universe size
$N=n+q$. Define
\begin{equation}\label{tx:eq:NQ}
 \begin{gathered}
 N(0)=n_*,\qquad Q(0)=1,\\
 N(d)=\max\{n_*,2(\kappa d+\sigma)+1\},\qquad
 Q(d)=\min\{p^a:p^a>\kappa d+\sigma\}\quad(d\ge1).
 \end{gathered}
\end{equation}
The conservative onset makes the retained-rank condition of the
transfer theorem automatic. The degree-zero predicate is handled
separately: beyond $n_*$ the constant-target and constant-axiom data
are stable, so that predicate is constant. In particular its period
is one, even if $\sigma$ is positive.

\begin{proposition}[Threshold and truncated periodicity]\label{tx:thm:periodicity}
For $n,m\ge N(d)$ with $n\equiv m\pmod{Q(d)}$, one has
\begin{equation}\label{tx:eq:truncated-periodicity}
                     \min\{h(n),d+1\}=\min\{h(m),d+1\}.
\end{equation}
The convention is $\min\{\infty,d+1\}=d+1$. For $d\ge1$,
\begin{equation}\label{tx:eq:period-bound}
                              Q(d)\le pR(d).
\end{equation}
\end{proposition}
\begin{proof}
Repeated one-block transfer gives equivalence of the threshold
$h\le d$ at ranks in the same residue class, since every intermediate
rank remains beyond the onset. The smaller thresholds $e\le d$
can also be compared there: $N(e)\le N(d)$, and $Q(e)$ divides $Q(d)$
because both are powers of $p$ and the budgets are nondecreasing.
The zero threshold is constant in this range.

If $h(n)=e\le d$, threshold $e$ holds at $m$, while threshold $e-1$
fails there when $e>0$. Hence $h(m)=e$; for $e=0$ the first condition
alone suffices. If $h(n)>d$, threshold $d$ fails at both ranks. This
proves the truncated equality.

Write $Q(d)=p^a$. Since $R(d)\ge1$, the preceding power exists and
minimality gives $p^{a-1}\le R(d)$. Multiplication by $p$ proves
\eqref{tx:eq:period-bound}.
\end{proof}

\begin{example}[A numerical budget]\label{tx:eg:numerical-period}
Take $p=2$, $\kappa=4$, $\sigma=3$, and $n_*=7$. For degree $d=3$,
\[
                       R(3)=15,\qquad Q(3)=16,\qquad N(3)=31.
\]
Thus the degree-at-most-three predicate agrees at $31,47,63,\ldots$.
If the exact minimum at 31 is three, the minimum is exactly three at
every size of that progression: the comparison of the lower thresholds
rules out a drop. This conclusion says nothing about a different
residue class modulo sixteen.

The inequality in the definition of $Q$ is strict. If a budget equals
$16$ in characteristic two, the next permitted power is $32$, not $16$.
\end{example}

\section{Which finite degrees come back infinitely often?}
\begin{definition}[Recurring finite degrees]\label{tx:def:recurring}
Let
\[
 E(h)=\{d\in\NN:\text{there exists }n\ge N(d)\text{ with }h(n)=d\}.
\]
\end{definition}

\begin{lemma}[An occurrence beyond its onset recurs]\label{tx:lem:recurrence}
The set $E(h)$ consists exactly of the finite values assumed at
arbitrarily large ranks. If $h(n)=d$ and $n\ge N(d)$, then
\[
                          h(n+\ell Q(d))=d\qquad(\ell\ge0).
\]
A nonrecurring finite value can occur only before its own onset.
\end{lemma}
\begin{proof}
Apply truncated periodicity at threshold $d$ to the displayed
progression. A truncated value equal to $d$ represents the exact value
$d$, not a value above the cutoff. Conversely, any value appearing at
arbitrarily large ranks appears beyond the fixed integer $N(d)$.
\end{proof}

The lemma is the central restriction on possible profiles. An arbitrary
integer sequence might take the value 100 for the last time at an
astronomically large rank. Here, if 100 will not recur, its final
appearance must precede the prescribed linear onset $N(100)$.
Inverting that restriction gives a linear degree lower bound.

\section{Inverting the onset}
Define the nonnegative integer
\begin{equation}\label{tx:eq:inverse-onset}
 \beta(n)=\left\lfloor\frac{(n-2\sigma-1)_+}{2\kappa}\right\rfloor,
                         \qquad x_+=\max\{x,0\}.
\end{equation}
For every $n\ge n_*$, one has $N(\beta(n))\le n$. If $\beta(n)=0$,
this is the special definition $N(0)=n_*$. Otherwise the floor inequality
gives $2\kappa\beta(n)\le n-2\sigma-1$, so the positive-degree formula
for $A$ applies. This separate zero case avoids subtracting natural
numbers as if negative values were allowed.

Consequently,
\begin{equation}\label{tx:eq:low-value-recurs}
 n\ge n_*,\quad h(n)\le\beta(n)
                   \quad\Longrightarrow\quad h(n)\in E(h).
\end{equation}
Indeed, the actual value $e=h(n)$ is finite and
$N(e)\le N(\beta(n))\le n$, so the recurrence lemma applies.

\section{The four alternatives}
We use limits in the extended order $\NN\cup\{\infty\}$.
The two tail bounds are
\[
 \begin{aligned}
 \liminf_{n\to\infty}h(n)&=\sup_{m\ge1}\inf_{n\ge m}h(n),\\
 \limsup_{n\to\infty}h(n)&=\inf_{m\ge1}\sup_{n\ge m}h(n).
 \end{aligned}
\]
The liminf measures the best degree level still visible arbitrarily
far out; the limsup measures the largest levels that keep returning.
For an integer-valued profile with a recurring finite minimum $e$,
the liminf is $e$ once the finitely many smaller values have stopped
occurring. A finite liminf therefore does not mean convergence.
In T4 below, bounded subsequences coexist with arbitrarily large
finite peaks.

\begin{theorem}[Four degree profiles]\label{tx:thm:profiles}\label{tx:thm:classification}
For each fixed uniform request, and separately for NS and ordinary PC,
exactly one of the following holds.
\begin{enumerate}[label=\textup{T\arabic*.}]
\item The profile is bounded by a finite constant at every sufficiently
large rank.
\item There is no recurring finite value. For every $n\ge n_*$,
\begin{equation}\label{tx:eq:linear-bound}
                                  h(n)>\beta(n).
\end{equation}
Infinity is allowed in this inequality; finite values satisfy a linear
lower bound.
\item The set $E(h)$ is finite and nonempty, but the profile is not
eventually bounded. Put $D=\max E(h)$. A nonempty proper set
$U\subseteq\{0,\ldots,Q(D)-1\}$ has the following properties from
rank $N(D)$ onward. On every class $c\in U$, the profile has one
constant finite value $v_c\le D$. Off $U$, inequality
\eqref{tx:eq:linear-bound} holds.
\item The set $E(h)$ is infinite. For every prescribed bound $B$, there
are integers $n,d\ge B$, with $d\in E(h)$ and $h(n)=d<\infty$, such that
\begin{equation}\label{tx:eq:finite-peaks}
 n\le(p+2)(\kappa d+\sigma),\qquad
 h(n)\ge\frac{n}{(p+2)\kappa}-\frac{\sigma}{\kappa}.
\end{equation}
Moreover $\liminf_{n\to\infty}h(n)=\min E(h)<\infty$.
\end{enumerate}
In T1 and T3 the liminf also equals the least recurring finite value;
in T2 it is infinity. Initial ranks below the specified onsets remain
unconstrained.
\end{theorem}
\begin{proof}
If $E(h)$ is empty, \eqref{tx:eq:low-value-recurs} excludes
$h(n)\le\beta(n)$ and proves T2. Since $\beta(n)$ grows linearly,
this is an $\Omega(n)$ lower bound wherever the minimum degree is
finite. It is stronger to state the inequality in the extended order:
an underivable target is not an exceptional counterexample.

Suppose next that $E(h)$ is finite and nonempty, with maximum $D$.
Truncated periodicity at $D$ divides the residue classes modulo $Q(D)$
into two kinds. On one kind the profile has a fixed finite value at
most $D$; on the other it is always greater than $D$. Let $U$ be the
classes of the first kind. A recurring degree appears beyond $N(D)$,
so $U$ is nonempty. If $U$ is the entire residue set, the profile is
eventually bounded and belongs to T1. Otherwise, at a size off $U$,
an inequality $h(n)\le\beta(n)$ would force a recurring value, hence
a value at most $D$, a contradiction. This proves T3. The easy values
$v_c$ need not be equal on different classes.

If $E(h)$ is infinite, choose recurring values $d$ arbitrarily large,
in particular so large that $N(d)=2R(d)+1$. An occurrence of $d$
beyond that onset determines a residue class modulo $Q(d)$. There is
a representative $n$ in the interval
\[
                           N(d)\le n<N(d)+Q(d).
\]
Two-way truncated periodicity ensures $h(n)=d$, even if the original
occurrence was much later. Integrality and \eqref{tx:eq:period-bound}
now give
\[
              n\le N(d)+Q(d)-1\le 2R(d)+pR(d)=(p+2)R(d).
\]
Rearrange for $d$ to obtain the second inequality in
\eqref{tx:eq:finite-peaks}. Taking $d$ larger also makes $n\ge N(d)$
larger, so both may exceed the prescribed bound $B$.

To finish the exhaustion, eventual boundedness forces $E(h)$ to be
finite and nonempty: only finitely many degrees can occur eventually,
and at least one occurs infinitely often. Thus a profile either is
eventually bounded, or falls into exactly one of the other three cases
according as $E(h)$ is empty, finite nonempty, or infinite.

If $e=\min E(h)$ exists, every value below $e$ occurs only finitely
often. There are finitely many such values, so eventually $h(n)\ge e$.
The value $e$ keeps recurring, proving the liminf assertion. When
$E(h)$ is empty, the same argument applies below each finite bound,
and $h(n)$ tends to infinity in the extended order.
\end{proof}

Notice the two different numerical inputs. The linear onset $N(d)$
bounds the last occurrence of a nonrecurring small degree; this yields
the T2 and hard-residue T3 bounds. The upper bound on $Q(d)$ additionally
puts a recurring large degree at a rank proportional to that degree;
this yields the \emph{finite} T4 peaks. Infinitely many infinite values
would not by themselves prove the latter assertion.

If the proof only gave $N(d)=O(d^2)$, inverting the onset would give a
square-root lower bound. If it gave no useful upper bound on $Q(d)$,
a recurring value might first appear very far along its progression,
and the finite linear peak argument would not follow. These observations
explain why the later construction spends effort on a linear retained
label budget rather than merely proving some eventual periodicity.

\section{Reading the four cases correctly}
In the examples below, $v_p(n)$ is the largest integer $a\ge0$ such
that $p^a$ divides the positive integer $n$. Thus $p^{v_p(n)}$ is the
largest power of $p$ dividing $n$.

The following sequences illustrate the shapes. They are abstract
profiles, not claims that these formulas are exact degrees of particular
encodings. For fixed prime $p$ and $n\ge1$ consider
\[
 \begin{array}{c|c|c}
 \text{shape}&h(n)&E(h)\\\hline
 \mathrm T_1&1&\{1\}\\
 \mathrm T_2&n&\varnothing\\
 \mathrm T_3&1\text{ if }p\nmid n,\quad n\text{ if }p\mid n&\{1\}\\
 \mathrm T_4&p^{v_p(n)}&\{1,p,p^2,\ldots\}.
 \end{array}
\]
They satisfy the abstract threshold contract with, for example,
$\kappa=1$, $\sigma=0$, and $n_*=1$. For the last row, the condition
$p^{v_p(n)}\le d$ says that $n$ is not divisible by the least $p$-power
strictly larger than $d$. This makes its threshold periodicity explicit.

T1 does not mean eventually constant; bounded periodic variation is
allowed. T2 can include infinite values, including the constant-infinity
profile. T3 is linear only on its hard residues. T4 has a finite liminf,
but that number is not an eventual value: infinitely many other finite
levels recur as well. None of these descriptions removes finite initial
exceptions. A sharp statement about every small rank must be proved
separately.

\begin{proposition}[T4 admits no fixed easy/hard periodic partition]\label{tx:prop:no-periodic-split}
In T4 there do not exist a period $q_0$, a set $U$ of its residues, a
finite bound $B$, and $c>0$ such that, at every sufficiently large rank,
$h(n)\le B$ on $U$ and $h(n)\ge cn$ off $U$.
\end{proposition}
\begin{proof}
Choose a recurring finite value $e>B$. At sufficiently late occurrences
of $e$, the rank cannot be on an easy class. But once $n>e/c$, it cannot
be on a hard class either. This contradicts recurrence.
\end{proof}

For a fully invariant target the same deductions use the smaller
parameters $\kappa=k$, $\sigma=s$. For the target $1$, $D_G=0$ and
$n_*=n_0$, giving the refutation corollary. The later comparison theorem
shows that NS and PC assign the same type to each fixed request. It
does not assert equal exact minimum degrees or equal recurring sets.

\section{The size consequences under full coefficient lists}
Recall that a full coefficient list for degree $e$ in $v$ variables
has $B_v(e)=\binom{v+e}{e}$ entries, including zeros. In this section
$D$ again denotes the fixed axiom-degree bound, rather than the largest
recurring degree used in the T3 statement.
For an indexed request define $S_S(n)$ as the maximum of the minimum
coefficient-list sizes of the individual target proofs. Use zero for
an empty request and infinity if any target is underivable. A bundle
of proofs for all targets differs from this maximum by a factor between
one and $|\mathcal B_n|$ when the request is nonempty. That factor is
polynomial in $n$ and does not change the cases below.

\begin{corollary}[Polynomial bounds and exponential lower bounds]\label{tx:cor:dense-types}
For full coefficient lists, the four degree types have these consequences.
\begin{enumerate}[label=\textup{T\arabic*.}]
\item The minimum sizes are polynomial in $n$ eventually.
\item They are at least $2^{cn}$ eventually, for some $c>0$, with
infinity interpreted in the extended order.
\item They have a polynomial upper bound on the easy residue classes
and an exponential lower bound on the hard ones, beyond fixed onsets.
\item They have a polynomial upper bound at infinitely many sizes and
finite exponential peaks at arbitrarily large sizes. In addition, for
every $e>0$ there is $b$ such that $n^e<S_S(n)\le n^b$ at infinitely
many sizes.
\end{enumerate}
There is no eventual fixed periodic partition into one uniformly
polynomial part and one exponentially hard part in T4. These are
exponential lower bounds, not universal matching upper bounds, and they
concern full coefficient lists rather than sparse monomial size.
\end{corollary}
\begin{proof}
Bounded supports give polynomially many variables and indices. If some
variable has support size $j>0$, its relabellings have every $j$-element
support, giving at least $\binom nj\ge n$ variables for all sufficiently
large $n$. A fixed degree therefore gives a polynomial coefficient-list
upper bound by Lemma~\ref{tx:lem:dense-size}. A finite linear degree
lower bound gives an exponential size lower bound, since both systems
have size at least $B_{v(n)}(d-D)$ for a target attaining a minimum
degree $d>D$, where $D$ bounds input degree. Such a target is nonzero.
The zero-target empty certificate causes no exception. Underivability
gives infinite size directly.

For the extra T4 assertion choose a recurring degree $d$ with
$\ell=d-D>e$. At its sufficiently late occurrences,
\[
                 S_S(n)\ge B_{v(n)}(\ell)
                         \ge n^\ell/\ell!>n^e,
\]
while the fixed degree $d$ gives $S_S(n)\le n^b$ for some $b$.
These intermediate polynomial subsequences cannot fit a proposed
partition with one fixed polynomial exponent and an exponential lower
bound on the other part.

If every variable has empty support, there are only finitely many
variables and relabelling fixes them. Each axiom or target index orbit
gives a fixed polynomial, and all finitely many orbits appear
simultaneously beyond a fixed rank. The systems and requests then
stabilise. Their degree is a finite constant, or infinity, producing
T1 or the infinite-size case of T2. This handles the degenerate case
without assuming a growing number of variables.
\end{proof}

\section{Exercises}
\begin{exercise}\label{tx:ex:period-arithmetic}
Let $p=2$, $\kappa=4$, $\sigma=4$, and $n_*=7$. Compute $N(0),Q(0)$,
$N(3),Q(3)$ and $\beta(62)$. Check $N(\beta(62))\le62$ and explain
why the strict inequality defining $Q$ matters.
\end{exercise}
\noindent{\small\hyperref[tx:ans:period-arithmetic]{Model answer (p.~\pageref*{tx:ans:period-arithmetic}).}}\par\medskip

\begin{exercise}\label{tx:ex:forward-only}
For $g(1)=2$ and $g(n)=1$ for $n\ge2$, verify forward persistence of
every threshold under $n\mapsto n+1$. Identify a reverse threshold
implication that fails. Explain why exact recurrence of the starting
value cannot be deduced from forward persistence alone.
\end{exercise}
\noindent{\small\hyperref[tx:ans:forward-only]{Model answer (p.~\pageref*{tx:ans:forward-only}).}}\par\medskip

\begin{exercise}\label{tx:ex:valuation-profile}
For $h(n)=p^{v_p(n)}$, determine every recurring value, its liminf,
and a sequence of finite linear peaks. Explain why no eventual linear
lower bound holds at all sizes.
\end{exercise}
\noindent{\small\hyperref[tx:ans:valuation-profile]{Model answer (p.~\pageref*{tx:ans:valuation-profile}).}}\par\medskip

\begin{exercise}\label{tx:ex:zero-T3}
Show that T3 cannot have $E(h)=\{0\}$. State exactly where stability of
the degree-zero predicate is used.
\end{exercise}
\noindent{\small\hyperref[tx:ans:zero-T3]{Model answer (p.~\pageref*{tx:ans:zero-T3}).}}\par\medskip


\chapter{What it means to delete a block of labels}
\label{tx:chap:deletion}\label{tx:sec:transfer}
We now describe the operation that the preceding contract requires.
It acts on spaces with label permutations, not by setting some polynomial
variables to zero. Its construction will justify the word ``deletion''
on bases, and then explain why extra work is needed on subspaces.

The construction has three ingredients: spaces on labelled tuples,
one nilpotent operator built from block translations, and a quotient
of its kernel. We develop these tools before stating the exactness
theorem. Its clauses identify the surviving basis, preserve the
actual coefficient maps, and transport the subspace calculation
of a proof. The common-operator construction in the later chapters
establishes these properties.

\section{Why substitution is the wrong operation}
Consider the totality axiom $\sum_jx_{ij}-1$. If the row label $i$ is
deleted and all its variables are set to zero, this polynomial becomes
$-1$. Such a substitution refutes the remaining system immediately.
It does not produce the intended smaller uniform instance.

Restriction of coefficient matrices to retained rows and columns is
correct on uniform data, but arbitrary subspaces do not restrict exactly.
A linear relation supported partly on deleted coordinates may change its
meaning. The proof therefore uses a quotient operation on modules with
label action, together with a theorem that this operation preserves the
required subspace calculations.

\section{Permutation modules and their subspaces}
Recall the label-action language of Section~\ref{tx:sec:action-primer}.
The monomial space is a permutation module, and the NS/PC consequence
spaces are invariant subspaces: relabelling a legal proof gives another
legal proof from the uniform axioms.

\subsection*{A vector space with named basis vectors}
For a finite set of labelled objects $\mathcal U_N$, the notation
$K[\mathcal U_N]$ means the vector space with one basis vector $e_u$
for each object $u\in\mathcal U_N$. A vector is a coefficient list
\[
                         v=\sum_{u\in\mathcal U_N}a_u e_u.
\]
The basis objects might be variables, monomials, axiom indices or
ordered tuples. A label permutation $g\in S_N$ acts by
$g e_u=e_{g u}$ and extends linearly. In matrix language it permutes
the basis coordinates. The phrase \emph{permutation module} means
exactly this vector space together with these permutation matrices.

For instance, the space on single labels has basis
$e_1,\ldots,e_N$, and the swap $(12)$ sends $e_1$ to $e_2$,
$e_2$ to $e_1$, and fixes the other basis vectors. More general
spaces below need not have a basis permuted by labels: an invariant
subspace or a quotient still carries linear permutation actions.
Such a space is called a representation, or module. No additional
algebraic operation on its vectors is intended by that word.

\subsection*{Subspaces, quotients and maps which respect relabelling}
A subspace $U\subseteq M$ is \emph{invariant} if every label
permutation maps it to itself. Its quotient $M/U$ consists of classes
$[v]$ with $v$ and $v'$ representing the same class exactly when
$v-v'\in U$. Relabelling acts by $g[v]=[gv]$, which is well defined
because $gU=U$.

As a small example, in the space on three labels the line
$U=K(e_1+e_2+e_3)$ is invariant. In the quotient the sum
$e_1+e_2+e_3$ becomes zero, while $[e_1]$ is still nonzero.
Thus a quotient does not discard a designated subset of coordinates:
it imposes the specified linear relations. A \emph{subquotient} is
a quotient of an invariant subspace, such as $U/V$ with
$V\subseteq U\subseteq M$ invariant. These are precisely the
objects which arise when we look at targets modulo consequences.

A linear map $f:M\to M'$ is \emph{equivariant} if
$f(gv)=g f(v)$ for every label permutation $g$. This just says that
the two ways of relabelling and applying $f$ give the same answer.
For example, on the space $E_N$ with basis $e_1,\ldots,e_N$,
the coordinate-sum map
\[
 \epsilon:E_N\longrightarrow K,
                 \qquad\epsilon\Bigl(\sum_i a_i e_i\Bigr)=\sum_i a_i
\]
is equivariant when labels act as the identity on the output $K$.
Forgetting one position of a tuple is another example. The axiom,
target and multiplication maps in the uniform proof calculation are
equivariant too; their actual coefficients specify which maps they are.

\subsection*{One expression in permutations, acting on many spaces}
The group algebra $K[S_N]$ consists of finite sums $\sum_g a_g g$.
Its multiplication is composition, extended bilinearly. It acts by the
corresponding linear combination of renamings. If a map
$f$ respects label permutations, then
\[
                           f\bigl((\sum_g a_gg)v\bigr)
                                      = (\sum_g a_gg)f(v).
\]
This elementary identity is the reason a common group-algebra expression
is so valuable: the same expression automatically commutes with every
map in the uniform diagram.

For example, $h=2g-\tau$ acts on $M$ as
$v\mapsto2g v-\tau v$. The same expression acts on $M'$ and
commutes with $f$ whenever $f$ is equivariant. We do not choose
unrelated matrices separately on the two spaces. Nor do we claim
that $h$ commutes with every label permutation: the statement is
that every equivariant map commutes with the common expression.
In positive characteristic, the coefficient $2$ in this example is
interpreted in $K$, so it is zero when the characteristic is two.

\subsection*{Why ordered tuples are the test spaces}
For $b\ge0$ let
\[
                 P_b(N)=K[\Inj([b],[N])]
\]
be the space on ordered tuples of distinct labels. Every orbit of a
support-$b$ index is a quotient of such a space. Choose an ordering of
its support; an injective tuple specifies where that support is sent.
Two extensions agreeing on the support act identically on the index.
Thus ordered tuples provide a finite universal list of spaces on which
to construct an operator valid for all support-bounded data.
Here $P_b(N)$ denotes a tuple permutation space; the bounded-degree
polynomial space is written $P_{N,\le d}$.

Here $[b]=\{1,\ldots,b\}$ names the \emph{positions}; an injection
assigns distinct universe labels to them. In particular:
\begin{itemize}
\item $P_0(N)$ is the one-dimensional space on the empty tuple.
Every permutation fixes its basis vector.
\item $P_1(N)$ has basis $e_1,\ldots,e_N$.
\item $P_2(N)$ has basis $e_{(i,j)}$ for $i\ne j$. It has
$N(N-1)$ basis vectors, and $g e_{(i,j)}=e_{(g(i),g(j))}$.
\end{itemize}
There is also an action permuting tuple \emph{positions}: swapping
the two positions sends $(i,j)$ to $(j,i)$. It is different from
swapping two particular universe labels. The two actions commute.
Throughout the construction, $b$ counts positions and $r$ bounds
their number; neither is a polynomial degree.

The universality claim in the preceding paragraph is concrete. If an
object $u$ has support $\{u_1,\ldots,u_b\}$, send
$e_{(t_1,\ldots,t_b)}$ to $e_{g u}$, where $g$ sends each $u_i$
to $t_i$. Any two such permutations agree on the support and hence
give the same object. The map is linear, equivariant and onto the
space on the orbit of $u$. Therefore any operator which vanishes on
$P_b(N)$ also vanishes on that orbit space.

\subsection*{The common-annihilator class}
Let $J_r(N)$ be the ideal of group-algebra elements acting as zero on
$P_b(N)$ for every $b\le r$. Write $\mathcal C_r(N)$ for the
finite-dimensional modules annihilated by $J_r(N)$. Every support-$r$
permutation basis, and every submodule or quotient of it, belongs to this
class. It is also closed under finite direct sums. We will not assume
that it is closed under every tensor product: the actual tensor products
used in the proof are placed inside known support-bounded permutation
ambients.

In formulas, the definition is
\[
 J_r(N)=\{h\in K[S_N]:h v=0\text{ for every }v\in P_b(N)
                                  \text{ and every }0\le b\le r\}.
\]
Calling it an \emph{ideal} records that sums and scalar multiples
remain in it, as do products on either side by group-algebra
elements. The latter fact follows because each $P_b(N)$ is preserved
by all those operators.

Membership of $M$ in $\mathcal C_r(N)$ means that \emph{every}
operator in this set acts as zero on $M$ too. Equivalently, all the
operator relations shared by the tuple spaces up to length $r$ hold
on $M$. This is a condition on the permutation action on a space,
not a claim that each of its vectors mentions only $r$ labels.
For example, $\sum_{i=1}^N a_i e_i$ can use every label, while its
single-label permutation space belongs to $\mathcal C_1(N)$.
Likewise the degree-$d$ monomial space has basis support at most
$kd$, although a polynomial in it can involve many different labels
across its monomials.
The symbol $J_r(N)$ denotes this set of operators, not the target
map $J_N$ from the previous chapter.

The closure assertions now have elementary proofs. If an operator
kills every vector of $M$, it kills every vector of an invariant
subspace and acts as zero on every quotient class. It also kills
a finite direct sum of such spaces. Hence the class includes the
subspaces and subquotients needed by a proof. It need not include
an arbitrary tensor product: combining two basis objects can add
their support sizes. We will check the relevant \emph{ambient}
support each time a tensor product is used.

\section{The regular prime-power block}
The parameters have different roles. We retain $n$ labels, add
$q=p^a$ block labels, and obtain $N=n+q$ universe labels in total.
Here $p$ is the characteristic, $a$ is a block exponent, and $r$
is the support budget for the spaces in the proof diagram. Choosing
$q>r$ does not mean choosing $q$ larger than $n$.

Write $N=n+q$ and identify the last $q=p^a$ labels with
$A=(\mathbb F_p)^a$. Translation by a coordinate vector gives a
permutation $\sigma_i$ of the block, fixing the retained labels. The
translations commute and have order $p$. Over a suitable coefficient
extension choose $c_1,\ldots,c_a$ independent over $\mathbb F_p$ and put
\[
                 z=\sum_{i=1}^a c_i(\sigma_i-1).
\]
Here $1$ is the identity permutation. On a representation $M$, this
group-algebra element acts as a linear operator $z_M$. A power of $z$
means repeated composition: $z^p=0$ will say that $p$ successive
applications send every vector to zero.
The characteristic-$p$ binomial formula gives $z^p=0$. The retained
permutations commute with $z$, so its kernel and the quotient below
retain an action of $S_n$.
In operator formulas, $I$ denotes the identity linear map, the action
of this identity permutation. Independence of the coefficients is
over the prime field $\mathbb F_p$: a relation
$\sum_i t_i c_i=0$ with $t_i\in\mathbb F_p$ must have all
$t_i=0$. It is not independence over the larger working field.

The word \emph{regular} means that the translations move any block
label to any other in exactly one way. For $p=2,a=2$, write the
four block labels as $00,10,01,11$. Translation in the first
coordinate swaps $00$ with $10$ and $01$ with $11$; translation
in the second swaps $00$ with $01$ and $10$ with $11$. They
leave all retained labels fixed. In general, the identity
\[
 z^p=\sum_i c_i^p(\sigma_i-I)^p
     =\sum_i c_i^p(\sigma_i^p-I)=0
\]
uses the commutation of these translations and the characteristic-$p$
binomial formula. It is a statement about matrices or composed
operators, not about a scalar whose power is zero.

The independence of the coefficients prevents a nontrivial translation
orbit from being accidentally invisible. It can require extending the
field. Membership of a target with original-field coefficients descends
by Lemma~\ref{tx:lem:field}; no proof over an arbitrary extension is
silently substituted for a proof over the original field.

\subsection*{A quotient made from a kernel and an image}
Define
\[
                 \Del(M)=\ker z_M/(\ker z_M\cap\im z_M).
\]
For $p=2$ this is $\ker z/\im z$. For odd $p$ that shorter formula
is not meaningful in general, because $z^2$ need not be zero.

The kernel consists of vectors $v$ with $z_Mv=0$. The image
consists of vectors $z_Mw$. We first keep the kernel vectors and
then identify two of them when their difference is also an image
vector. Thus an element of $\Del(M)$ is a class $[v]$ of a vector
already in the kernel. An arbitrary basis vector need not define
such a class.

For a fully explicit case, work in characteristic two, retain one
label $a$, and swap two block labels $b,c$. Choose
$z=(bc)-I$ on the space with basis $e_a,e_b,e_c$. Then
\[
 ze_a=0,\qquad ze_b=e_b+e_c,\qquad ze_c=e_b+e_c.
\]
Consequently
\[
 \ker z=K e_a\oplus K(e_b+e_c),\qquad
 \im z=K(e_b+e_c),\qquad \Del(M)\simeq K e_a.
\]
The class of $\alpha e_a+\beta(e_b+e_c)$ is represented by
$\alpha e_a$. The combination $e_b+e_c$ lies in the kernel,
so it qualifies as a representative; it also lies in the image,
so it represents the zero class.
The vector $e_a$ survives. The individual vector $e_b$ is not in
the kernel and hence is not a representative for this quotient.
This distinction is why deletion is not a coordinate restriction
applied to every vector.

To see what the quotient keeps, take a nilpotent chain with basis
$v_0,\ldots,v_{\ell-1}$ and
\[
                  zv_0=0,\qquad zv_j=v_{j-1}\quad(j>0).
\]
This is a nilpotent \emph{Jordan block}: applying $z$ moves each vector
one position down the chain and eventually sends it to zero.
The chain basis describes the linear operator $z$. It need not be a
permutation basis or an invariant decomposition for the full symmetric
group. This distinction will matter when we later build an operator
which works on many representations at once.

For $\ell=1$, the kernel is the whole line and the image is zero,
so the line survives. For $\ell\ge2$, the kernel is $Kv_0$ and
$v_0=zv_1$ already lies in the image, so the quotient is zero.
Thus deletion retains length-one chains and discards every longer chain.
The equation $z^p=0$ permits every length $1\le\ell\le p$.
The later odd-characteristic construction separately proves that the
relevant tuple generators have only lengths one and $p$; neither that
purity statement nor exactness follows from nilpotence alone.

\subsection*{How a linear map induces a map of the quotients}
Let $f:M\to M'$ commute with label permutations. The common
group-algebra definition gives $f z_M=z_{M'}f$. If $v\in\ker z_M$,
then $f(v)\in\ker z_{M'}$. If two kernel vectors differ by
$z_Mw$, their images differ by $z_{M'}f(w)$, which is again an
image vector and remains in the kernel. Therefore the rule
\[
                  \Del(f):\Del(M)\longrightarrow\Del(M'),
                              \qquad[v]\longmapsto[f(v)]
\]
is well defined. It is the actual induced map, not an independently
chosen map between spaces of the same dimension. Applying two maps
in succession before or after taking classes gives the same answer.

In the three-label example, the coordinate-sum map
$\epsilon:M\to K$ is equivariant. On a kernel vector it gives
\[
 \epsilon\bigl(\alpha e_a+\beta(e_b+e_c)\bigr)
                       =\alpha+\beta+\beta=\alpha.
\]
Thus its induced map sends the surviving class $[e_a]$ to $1$
and the zero class $[e_b+e_c]$ to $0$. On the trivial output
space $K$ every label permutation is the identity, so $z=0$ and
$\Del(K)=K$.

Finally, permutations of the retained labels commute with each
block translation. They preserve the kernel and image and act on
the quotient by $[v]\mapsto[g v]$. This is why the output carries
an action of the \emph{smaller} group $S_n$. The deleted block
does not become a set of new polynomial variables in that output.

\section{Why small supports make a basis object retained}
The support bound in the next argument belongs to an individual
\emph{basis object}, not to every vector in its span. A basis variable
$x_i$ uses label $i$; the monomial $x_i x_j$ uses labels $i,j$;
a sum of many such monomials may mention most of the universe.

Suppose $\mathcal U$ has support bound $u<q$. If an object is fixed by
all block translations, its support is a block-translation-invariant set.
The regular action on the block is transitive. Hence a support meeting
the block would contain all $q$ block labels, contrary to $u<q$.
The fixed objects are exactly those supported on retained labels.

This is the first structural use of the inequality $q>r$. It has nothing
to do with choosing $q$ larger than the universe itself. It must exceed
the support needed by the degree-$d$ calculation. A one-label change
usually has $q=1\le r$ and cannot satisfy this argument. This explains
why the natural shifts are prime powers rather than $n\mapsto n+1$.

The strict inequality is visible in an example. With a four-label
block and unary variables, the product of all four block variables
is a basis monomial fixed by every translation. It uses the whole
block, so it is excluded by a support bound $u<4$. A monomial
meeting the block in one, two or three labels cannot be fixed by
all translations.
The support argument identifies which basis objects are fixed;
the later orbit calculation must still prove that the moved part
has zero deletion quotient.

\subsection*{The theorem in plain language}
Choose a block larger than the supports we need to track, and keep
enough labels outside it. The theorem asserts that the block
coefficients can be chosen so that the quotient just defined does
three things at once. It retains the expected smaller basis; it
carries the actual uniform coefficient maps to the smaller maps;
and, on the class $\mathcal C_r$, it preserves the subspace
operations used by a proof. The last assertion is the substantial
one: neither the definition of the quotient nor $z^p=0$ proves it.

Here is the precise statement. The phrases ``exact'', ``natural''
and ``trivial module'' are read clause by clause immediately below.

\begin{theorem}[The exact deletion input]\label{tx:thm:deletion}
Work first over an algebraically closed field of characteristic $p$.
For $q=p^a>r$ and retained rank $n\ge\max\{0,2r-1\}$, the block
coefficients can be chosen so that $\Del$ is exact on $\mathcal C_r(n+q)$.
It fixes the trivial module. For every support-$u$ permutation basis with
$u<q$ there are natural identifications
\begin{align}
 \Del(K[\mathcal U_{n+q}])&\simeq K[\mathcal U_n],\label{tx:eq:permdelete}\\
 \Del(K[\mathcal U_{n+q}]\otimes Y)&\simeq
                  K[\mathcal U_n]\otimes\Del(Y),\label{tx:eq:tensordelete}
\end{align}
with a permutation module as the first tensor factor. Stable invariant
coefficient matrices are carried to their actual retained-rank matrices.
\end{theorem}

This finite-algebra theorem is proved by the chapters on
square-zero blocks, linear layer repair, rook realisation and top-power
projection. No representation-theoretic equivalence is
invoked as an external assumption.

\subsection*{Reading the clauses}
\paragraph{The working field.}
An algebraically closed field has the coefficients needed for the
independent block directions. This is where the construction is
carried out, not a change of the original input equations.
Lemma~\ref{tx:lem:field} brings membership of original-coefficient
targets back to the original field. The block size $q=p^a$ is a
number of universe labels, not a prescribed cardinality for that
field.

\paragraph{The support budget and the range.}
The integer $r$ bounds the tuple test spaces and the supports of the
proof-diagram ambients which must undergo exact transport.
The inequalities are $q>r$ and $n\ge\max\{0,2r-1\}$.
They concern block labels and retained labels, respectively. For
$r\ge1$, the latter is the sufficient retained-label budget
$2r-1$; for $r=0$, it is simply $n\ge0$. The proof will construct
one compatible witness for all lengths $0,\ldots,r$ on that one
retained pool. Template stabilisation and target-degree thresholds
are additional requirements when this finite-algebra input is
applied to polynomial families.

\paragraph{Exactness on the specified class.}
Take an equivariant map $f:M\to M'$ between spaces in
$\mathcal C_r(N)$. Its kernel and image are invariant subspaces.
Exactness says that their deleted spaces are the actual kernel and
image of $\Del(f)$, with the induced inclusions. In particular,
if $U\subseteq V$ are invariant spaces in this class, the induced
map $\Del(U)\to\Del(V)$ is injective, and
\[
                  \Del(V/U)\simeq\Del(V)/\Del(U).
\]
This is a claim about maps and their linear relations, not only
about dimensions. It is restricted to $\mathcal C_r(N)$;
arbitrary label representations are not being declared exact.

\paragraph{Fixing the trivial module.}
On the one-dimensional space $K$ with all label permutations acting
as the identity, every $\sigma_i-I$ acts as zero. Hence $z=0$,
the kernel is $K$, the image is zero and $\Del(K)=K$.
The constant line does not disappear. This clause will be needed
both for invariant targets and for the detection pairing.

\paragraph{The first basis identification.}
The first displayed isomorphism says that the surviving directions
in the space on $\mathcal U_{n+q}$ are precisely the basis
objects supported on retained labels. A retained object $u$ gives
the class of its basis vector $e_u$, identified with the same
object in $\mathcal U_n$. Moved basis vectors need not themselves
be kernel vectors; the assertion is that their entire moved-orbit
summand contributes zero to the quotient. For a single-label basis,
this turns the large label space into the space on the $n$
retained labels.

\paragraph{The tensor identification and its naturality.}
Here $Y$ is any space with the larger label action. The first
factor has a permutation basis. Regard a tensor
$\sum_u e_u\otimes y_u$ as a list of vectors from $Y$, one
packet for each basis object $u$. A label permutation acts on
\emph{both} the packet index and its contents:
$g(e_u\otimes y)=e_{gu}\otimes gy$.
The theorem keeps retained packet indices and applies deletion
to their contents. More precisely, for retained $u$ and
$y\in\ker z_Y$, it matches
\[
                         e_u\otimes[y]\quad\text{with}\quad[e_u\otimes y].
\]
It also proves that the moved packet summand disappears even
though translations act on its $Y$ contents. There is no separate
bounded-support assumption on $Y$ in this basis identification.
When exactness is used on a tensor source, we separately check
that source's ambient support budget.

``Natural'' means that applying any equivariant map $Y\to Y'$
to the packet contents before or after this identification gives
the same result. This is what lets the rule apply to a consequence
subspace $Y\subseteq L_d$ with its actual inclusion, rather than
to a separately chosen basis for $Y$. It is also why the tensor
clause specifically states which factor has a permutation basis.

\paragraph{The coefficient matrices are the actual maps.}
If $F_N$ sends a retained input basis vector $e_u$ to
$\sum_v a_N(v,u)e_v$, equivariance puts this image in the kernel
of $z$. Its retained coefficients are exactly $a_N(v,u)$ for
retained $v$. A stable rule gives $a_N(v,u)=a_n(v,u)$ on those
objects, and the moved kernel component vanishes in the quotient.
Thus the induced map is the prescribed $F_n$ entry by entry.
The statement does not allow us to replace it by an unrelated
matrix with the same rank. It applies in particular to the
uniform axiom, target and multiplication maps from Chapter~\ref{tx:chap:contract}.

\subsection*{How the later proof makes exactness possible}
The later chapters construct one group-algebra element $w$ with
\[
             z^{p-1}wz^{p-1}=z^{p-1}
                           \quad\text{on every }P_b(N),\ b\le r.
\]
This is a partial-inverse condition: on the image of $z^{p-1}$,
$w$ supplies a preimage. It does not invert the nilpotent operator
on the whole space. From $z$ and $w$ the proof constructs another
common operator $e$, using also the proved chain property of the
tuple generators. This $e$ is a projector with $e^2=e$, and identifies
$\Del(M)$ with its image $eM$ as a vector space. The letter $e$
here names an operator, unlike the basis-vector notation $e_u$.

Every actual equivariant map commutes with this same expression
in permutations. For example, if $f:V\to W$ is onto and
$y\in eW$, choose an ordinary preimage $v$ of $y$. Then
$f(ev)=e f(v)=ey=y$, so the restricted map $eV\to eW$
is still onto. The kernel and injection checks use the same
commutation. This is the source of exactness. Its hard part is
constructing one such expression which works on the entire
annihilator class; separate projectors on separate spaces would
not justify these equations. The quotient itself carries the
retained label action; the projector image is its vector-space
model.

\subsection*{An admissible numerical instance}
Take characteristic two, $r=2$, $q=4$, $n=3$, and $N=7$.
Work in an algebraic closure of $\mathbb F_2$, choose
$\omega^2+\omega+1=0$, and use coefficients $1,\omega$.
Retain labels $1,2,3$ and identify labels $4,5,6,7$ with
$00,10,01,11$, respectively. The two translations and their
common operator are
\[
 \sigma_1=(45)(67),\qquad\sigma_2=(46)(57),\qquad
 z=(\sigma_1-I)+\omega(\sigma_2-I).
\]
The numbers satisfy $q>r$ and $n=2r-1$.
The tuple test spaces are $P_0(7)$, $P_1(7)$ and $P_2(7)$,
of dimensions $1,7,42$. The moved-orbit calculation, which uses
the independence of $1,\omega$ over $\mathbb F_2$, identifies
their deletion quotients with retained tuple spaces of dimensions
$1,3,6$. This illustrates the basis calculation for the displayed
operator.

The all-class exactness assertion is separate. For these same
parameters, Theorem~\ref{tx:thm:deletion} provides a suitable choice
of block coefficients for which deletion is exact on
$\mathcal C_2(7)$. The present numerical illustration does not
certify that its specific choice $1,\omega$ has that additional
property or supply a common witness for it. The later construction
establishes the exactness input with its required coefficient choice.

For comparison with polynomial notation, the ordinary degree-at-most-two
monomial space in seven unary variables has $36$ basis monomials:
one constant, seven linear monomials, seven squares and twenty-one
products of distinct variables. Each monomial uses at most two
labels. Its deletion is the corresponding space in three unary
variables, with $10$ monomials. This is a consequence of the basis
identification, not a computation of any minimum NS or PC degree.
Squares are separate coordinates because this is the ordinary
polynomial ambient, not the space of functions on Boolean assignments.
An application to a proof must additionally cover the actual axiom
and target maps, multiplication domains and, when needed, the
doubled-support detector. Those data may require a larger $r$.

The instance illustrates the theorem's parameters and conclusions;
its concrete operator illustrates basis deletion, and the theorem's
existential coefficient choice supplies exactness. It does not
independently construct a common witness or prove general lifting
by finite experiment.

\section{Why deletion on bases is not enough}
In characteristic two, let two block labels be swapped. Their permutation
space has basis $e_1,e_2$ and $z e_1=z e_2=e_1+e_2$. Its kernel equals
its image, so its deletion is zero. The invariant line
$T=K(e_1+e_2)$ has $z=0$ and survives. Therefore deletion of
$T\hookrightarrow Ke_1\oplus Ke_2$ is the zero map on a nonzero space.
An injection was not preserved.

Here there are no retained labels: $n=0$, $q=2$ and the single-label
basis has support budget $r=1$. The retained-rank hypothesis would
require $n\ge1$. The example therefore explains the need for a
range restriction rather than contradicting Theorem~\ref{tx:thm:deletion}.
The vector $e_1+e_2$ represents a nonzero quotient class when
the operation is applied to $T$ alone, but represents the zero
class after inclusion into the larger two-point space. Inducing
maps of quotients is always possible; preserving injectivity is
an additional theorem.

The extra retained labels in Theorem~\ref{tx:thm:deletion} rule out
this problem in the relevant class by providing a common projector.
The basis orbit calculation and the exactness construction have distinct
jobs. Forgetting the second job would leave the proof unable to handle
intersections and hence ordinary-PC degree falls.

\section{What exactness will buy}
In the contract of Definition~\ref{tx:def:polynomial-diagram}, matrix
naturality identifies the actual axiom, multiplication and target maps.
Exactness preserves the kernels, images, intersections and sums of the
closure calculation. The doubled-support pairing supplies the remaining
target-vanishing reflection for general indexed targets.
The roles should remain distinct. Exactness carries a consequence
space and its actual inclusion to the smaller rank. The target-map
compatibility then identifies the deleted request. To reflect
underivability, we still need to know that a nonzero target class
in the relevant quotient cannot disappear. That last step is the
nonzero-subquotient detector, not a consequence of the word
``exact'' alone.

The next chapter gives the related upward lift and explains how it
interacts with multiplication without pretending to be a ring map.

\chapter{How to lift a proof without increasing its degree}
\label{tx:chap:lift}
The deletion operation goes from a larger universe to a smaller one.
The lifting operation goes in the other direction: take a low-degree
argument, enlarge the universe by a suitable block, and construct
a larger-universe consequence. These are related viewpoints.
Deletion is an operation on entire invariant spaces and their maps.
The lift is an ordinary linear map on polynomial vectors.

We first specify and verify the lift from the identities of a common
projector. The following technical chapters construct that projector.
Thus the present chapter gives the formula and its mathematical role,
while the later construction proves that an operator with the required
identities exists.

\section{An elementary reason for prime-power steps}
Before the operator construction, consider the simplest counting
phenomenon behind it. Adding an arbitrary number of labels changes
counting expressions. Adding a sufficiently large prime-power block
leaves small binomial counts unchanged in the field.

The basic tool is the \emph{Frobenius map}, $u\mapsto u^p$, in a
commutative ring of characteristic $p$.
Since the intermediate binomial coefficients $\binom pi$ vanish,
$(u+v)^p=u^p+v^p$; multiplication is preserved as well.
Repeating this identity gives the analogous formula for $q=p^a$.
It applies in the polynomial ring $K[T]$ used below, and to the
commuting block-translation operators used later.
Commutation is part of the latter application; no such formula is
being asserted for arbitrary noncommuting matrices.
Integer counts are read in $K$ through the map $m\mapsto m\,1_K$.

\begin{lemma}[Counts below a prime-power block]\label{tx:lem:prime-counts}
Let $K$ have characteristic $p$, let $q=p^a$, and let $n\ge0$.
For every $0\le j<q$, the integer binomial coefficients have the same
image in $K$:
\[
                         \binom{n+q}{j}=\binom nj.
\]
\end{lemma}
\begin{proof}
In $K[T]$, repeated application of the characteristic-$p$ binomial
identity gives $(1+T)^q=1+T^q$. Consequently
\[
                  (1+T)^{n+q}=(1+T)^n(1+T^q).
\]
Multiplication by $T^q$ contributes no coefficient in degrees below
$q$. Comparing the coefficient of $T^j$ proves the assertion.
\end{proof}

For example, in characteristic two, adding four labels preserves
$\binom nj$ for $j=0,1,2,3$. At $n=7$ and $n+4=11$ all four values
are one in the field. Degree four can see the change:
$\binom74=35$ is odd, whereas $\binom{11}4=330$ is even.
This explains why the inequality between block size and complexity
budget is strict.

The count has a polynomial interpretation. For unary variables put
\[
 E_{j,n}=\sum_{1\le i_1<\cdots<i_j\le n}
                           x_{i_1}\cdots x_{i_j}.
\]
Use $E_{0,n}=1$ and $E_{j,n}=0$ for $j>n$.
Its value at the all-one assignment is $\binom nj$.
More strongly, setting the $q$ added variables to one gives, for $j<q$,
\[
 E_{j,n+q}(x_1,\ldots,x_n,1,\ldots,1)
   =\sum_{h=0}^j\binom qh E_{j-h,n}=E_{j,n}.
\]
The intermediate coefficients $\binom qh$, $0<h<q$, vanish in $K$.

This is a complete elementary calculation, but not a lifting theorem
for arbitrary axioms. Setting new variables to one need not preserve
their equations. Even stable coefficient counts do not yet preserve
the intersections which allow PC to multiply a polynomial after
cancellation. The full lift must preserve that subspace calculation.
The lemma explains the choice of step; the common projector below
provides the extra structure needed for proofs.

\section{Low degree, a large starting rank, and repeated lifting}
Write $N(d)$ for a sufficient starting rank for degree-$d$ transfer.
If $n\ge N(d)$, one proof with that budget gives proofs at every
size $n+\ell q$, $\ell\ge0$, with one fixed degree bound $d$.
This is already enough to contradict a separate proof that the
minimum degree tends to infinity along that progression. It is the
argument used in Chapter~\ref{tx:ch:running-order}.

The earlier LICS refutation result~\cite[p.~276]{Riis} summarised
estimates $A_{\NS}(d)\le2^{cd}$ and
$A_{\PC}(d)\le2^{2^{cd}}$, extracted there from earlier work of
Kraj\'i\v cek. A degree bounded by a sufficiently small multiple of
$\log n$ or $\log\log n$, respectively, therefore lies in that
historical transfer range. Inverting the starting-rank estimates gives
the earlier logarithmic and double-logarithmic gaps.
These estimates concern the historical first-order refutation setting;
we do not use them as inputs to the general target theorem here.

A small budget is easier to fit inside a given universe, but does not
by itself supply a simple construction of the lift. The first complete
construction in this book is the characteristic-two argument of
Chapter~\ref{tx:ch:quadratic}. Its starting rank is quadratic in $d$.
The later constructions improve this to a linear starting rank and
cover every positive characteristic. Each stage answers the same
question: how many retained labels are sufficient to transport the
entire degree-$d$ proof calculation?

\section{What we are allowed to keep fixed}
We now specify the operator which makes those repeated lifts possible.

Recall the three support bounds: one variable uses at most $k$ labels,
one axiom index at most $s$, and one target index at most $t$.
A degree-$d$ monomial therefore uses at most $kd$ labels. The number
$r$ below covers all coordinates of the required diagram, rather than
just one monomial; this is why indexed-target detection has a larger
budget than closure transport alone.

Fix a proof degree $d$ and a support budget $r$. The consequence-space
calculation needs $r\ge kd+s$. To combine it with general indexed target
reflection we will use $r=2kd+s+t$. Fix $q=p^a>r$, retain enough labels
for exactness, and set $N=n+q$. Every later iteration may use the same
$d,r,q$. Only the universe size and the particular witness coefficients
change.

\subsection{A numerical trace of the retained-label test}
Take characteristic two, variable support $k=2$, axiom-index support
$s=1$, input-degree bound $D=1$, input stability rank $n_0=3$,
target-degree bound $D_G=2$,
and proposed proof degree $d=2$. Compare a fully invariant target with
an indexed request whose target-index support is $t=1$.
The final transfer theorem in Chapter~\ref{tx:chap:assemble-transfer}
gives the following sufficient ranges:
\begin{center}
\begin{tabular}{@{}lccc@{}}
\toprule
Request & Budget $r$ & Least $2$-power $q>r$ & Retained rank $n$\\
\midrule
Invariant & $kd+s=5$ & $8$ & $n\ge\max\{3,1,4,9\}=9$\\
Indexed & $2kd+s+t=10$ & $16$ & $n\ge\max\{3,1,5,19\}=19$\\
\bottomrule
\end{tabular}
\end{center}
For the invariant request, a degree-two proof at size $9$ lifts to
sizes $17,25,\ldots$. For the indexed request, the corresponding
progression starts at $19$ and runs through $35,51,\ldots$.
These are implications from an existing proof, not assertions that
either example request has such a proof. At each step, $n$ is the
\emph{smaller} rank and $N=n+q$ the larger one. A starting rank below
the stated range supplies no conclusion from this test.

The count $2r-1$ is one retained pool sufficient for every tuple
length $b\le r$: the individual requirements $2b-1$ are bounded by
$2r-1$. They are not added together as the witness is repaired.
The block size $q$ is additional to that pool. If $r\ge1$, the least
$p$-power strictly above $r$ satisfies $r<q\le pr$; at a $p$-power
budget the upper equality can occur. Once $d$ is fixed, all three
numbers $r,q,2r-1$ stay fixed through the progression.

The lift should not be thought of as replacing every variable by a new
polynomial and then substituting into every product. Its coefficients
are shared across all relevant spaces through a single group-algebra
expression. This gives linear and diagrammatic compatibility, not a
multiplicative substitution law.

\section{The proof can be reconstructed from the closure}
We spell out reconstruction for PC, whose consequence space is the
iterated closure $C_d$. For NS the same construction transports the
single certificate-map image $I_d$ described in
Chapter~\ref{tx:sec:pc-space}.
Suppose an original PC proof produces $f$. The exact consequence-space
characterisation says $f\in C_d(n)$. If a linear lift sends that vector
into $C_d(N)$, there is a legal PC proof at degree $d$ in the larger
universe. One obtains it by storing derivations for bases of closure
stages, expressing the required lifted vectors in those bases, and
concatenating the permitted inference steps.

This reconstruction is finite, but no efficient algorithm or bound on
proof length is asserted. Degree and length are different resources.
The closure lets us explain why a proof exists without forcing every
larger proof to have the same list of intermediate lines as the original.

\section{A lifting operator for bounded-degree consequences}
Fix a proposed proof degree $d\ge1$. Degree zero uses the separate constant-space calculation from Chapter~\ref{tx:sec:pc-space}. Write $W_n=P_{n,\le d}$,
$H_n=P_{n,\le d-1}$ and $X_n=K[\mathcal V_n]$.
Let $V_j(n)$ be the closure stages from \eqref{tx:eq:closure} and write
$C_d(n)$ for their stable union. A proof ending in a polynomial $f$ says exactly that
$f\in C_d(n)$. For the uniform-consequence theorem we must additionally
compare the prescribed targets at different ranks.

Choose a support budget $r\ge kd+s$ and add a block of $q=p^a>r$ labels and write $N=n+q$.
The required properties of the linear operator are:
\begin{equation}\label{tx:eq:newcontract}
 \mathsf L:W_n\longrightarrow W_N,\qquad
 \deg\mathsf L(f)\le\deg f,\qquad \mathsf L(1)=1,\qquad
 \mathsf L(V_j(n))\subseteq V_j(N)\quad\hbox{for every }j.
\end{equation}
The degree assertion for zero is interpreted using its harmless zero
image. This contract carries every degree-$d$ consequence to its specified
projector lift and, since $\mathsf L(1)=1$, carries a refutation
upwards. A legal proof is recovered from the larger closure. Iteration reaches $n+\ell q$ for every $\ell\ge0$.
The maps at different steps need not have the same coefficients.

For refutations, full periodicity needs a little more than the forward
contract. It needs
\begin{equation}\label{tx:eq:newiff}
 1\in C_d(n)\quad\Longleftrightarrow\quad1\in C_d(n+q),
\end{equation}
and it needs this at each lower degree budget too. The exact deletion
construction will supply both directions. A forward lift by itself
would not justify all the arguments below about exact recurring values and representatives of T4 peaks. General targets require the same two-way conclusion for their actual coefficient map, obtained by the image-subquotient detection argument in Chapter~\ref{tx:chap:assemble-transfer}.

Let $\iota:W_n\to W_N$ include each retained monomial, and let
$\rho:W_N\to W_n$ discard all other monomials. We also seek
$\rho\mathsf L=\mathrm{id}$. Plain inclusion has this property, but
usually does not carry smaller axioms to consequences of larger axioms.
The sum-axiom example below will show the failure
explicitly. The remedy is to compose inclusion with a common projector.

A map $\mathsf L$ with $\rho\mathsf L=\mathrm{id}$ is a
\emph{section}, or right inverse, of $\rho$. It chooses one larger
representative for each smaller vector, and is injective because
$\mathsf L(f)=0$ implies $f=\rho\mathsf L(f)=0$.
The lifting problem asks for a section which also respects the
consequence spaces. The term will recur later for sections of
coordinate-deletion maps.

The following verification uses Theorem~\ref{tx:thm:deletion} and the common projector constructed in Chapters~\ref{tx:sec:block-algebra}--\ref{tx:sec:odd}. It describes the actual operator and its compatibility with the PC inference rules.

\section{The lift from a common projector}\label{tx:lift-sec:formula}
We now give the map, leaving the construction of its certificate to the
next sections. Temporarily allow a coefficient extension $E/K$ and use
the same notation for the extended polynomial spaces. Section~\ref{tx:lift-sec:descent}
returns the lifting statement to $K$.

Identify the added block $B$ with $A=\Fp^a$, acting on it by translations
and fixing $[n]$. Let $\sigma_1,\ldots,\sigma_a$ be the coordinate
translations. Choose $c_1,\ldots,c_a\in E$ linearly independent over
$\Fp$ and set
\[
 z=\sum_{i=1}^a c_i(\sigma_i-1)\in E[S_N],\qquad z^p=0.
\]
The group algebra $E[S_N]$ consists of linear combinations of label
permutations; it acts on polynomials by the corresponding linear
combinations of variable renamings. The permutations commute on the
block, and $\sigma_i^p=1$, which proves $z^p=0$.

The finite algebraic target is one element $w\in E[S_N]$ for which
\begin{equation}\label{tx:lift-eq:witness}
 u w u=u\quad\text{on every }P_b(N)=E[\Inj([b],[N])],\quad b\le r,
 \qquad u=z^{p-1}.
\end{equation}
Thus the same expression must work on every tuple length up to the
support budget. The identity makes $w$ a partial inverse for the
top power $u=z^{p-1}$ on these spaces; it does not invert the nilpotent
operator $z$. Here $P_b(N)$ is a tuple permutation space; the
polynomial ambient space remains $W_N=P_{N,\le d}$.
For $y=uv$ in $\im u$, the witness identity says $uwy=uwuv=uv=y$.
Thus $wy$ is one preimage of $y$ under $u$; the word \emph{partial}
specifies that this inversion is required on $\im u$ only.
The top power sends vectors to the bottom of length-$p$ nilpotent
chains. Its common witness will let the construction remove those
chains while retaining the length-one directions.

For characteristic two put
\[
                         e=(1-zw)(1-wz).
\]
For general $p$, put $m=p-1$ and
\begin{equation}\label{tx:eq:projector}
 S_0=1,\qquad S_{j+1}=(1-z^jwz^{m-j})S_j\ (0\le j<m),
 \qquad e=(1-z^mw)S_m.
\end{equation}
The two formulas agree at $p=2$. For odd characteristic, the witness
identity is used together with a property of these particular tuple
spaces: their nilpotent chains have length one or $p$. The projector
lemma in Chapter~\ref{tx:ch:prime} proves the formula under that
property; nilpotence $z^p=0$ alone is not sufficient. Once these inputs
are established, the formulas are explicit finite algebra expressions.
The lifting operator is
\begin{equation}\label{tx:lift-eq:actual-lift}
 \boxed{\quad\Lift_{n,N,d}(f)=e\,\iota(f).\quad}
\end{equation}
Equivalently, writing $e=\sum_{g\in S_N}a_g g$,
\[
                    \Lift(f)=\sum_{g\in S_N}a_g\,g\bigl(\iota(f)\bigr).
\]
This explains why degree does not increase: inclusion and every
renaming preserve degree. The difficult point is to choose the
\emph{common} coefficients $a_g$ so that all consequence spaces are
preserved. An arbitrary sum of renamings does not have that property.

\subsection{The two facts that make the formula work}
Let $J_r$ be the ideal of group-algebra elements acting as zero on all
$P_b(N)$ for $b\le r$, and let $\mathcal M_r$ consist of the finite
dimensional representations annihilated by $J_r$. This class contains
all the bounded-support permutation spaces needed by PC, and their
invariant subspaces and quotients. It is the class $\mathcal C_r(N)$
of Chapter~\ref{tx:chap:deletion}, now over the auxiliary field $E$.
Define
\[
               \Del(M)=\ker z_M/(\ker z_M\cap\im z_M).
\]
The construction establishes two facts.
\begin{enumerate}
\item On permutation bases, $\Del$ retains exactly the objects supported
on $[n]$. On polynomial spaces its identification is induced by $\rho$:
\[
 \ker z_{W_N}/(\ker z_{W_N}\cap\im z_{W_N})\simeq W_n.
\]
Stable axiom matrices and the actual multiplication maps are transported
to their smaller-universe counterparts. With a permutation space as
first tensor factor, the identification is natural on the other factor.
\item The common $e$ is an idempotent with $ze=ez=0$ on
$\mathcal M_r$, and for $v\in\ker z$,
\begin{equation}\label{tx:lift-eq:onker}
                             ev=v-z^{p-1}wv.
\end{equation}
It identifies $\Del(M)$ naturally, as a vector space, with $eM$.
Consequently $\Del$ preserves kernels, images and exact sequences in
this class.
\end{enumerate}
These facts are proved in
Chapters~\ref{tx:ch:square-zero}--\ref{tx:ch:prime} from the common
witness and the block action. They are stronger than a count of dimensions. They refer to
the actual maps used in the consequence calculation.
The idempotent property means $e^2=e$: applying the common projection
twice has the same effect as applying it once.

\section{Compatibility with the entire PC closure}\label{tx:lift-sec:verify}
Apply the two facts just stated to the recurrence
\eqref{tx:eq:closure}. The usable axioms form the image of a stable
coefficient matrix, so deletion takes $V_0(N)$ to $V_0(n)$. Suppose
it has already taken $V_j(N)$, with its inclusion in $W_N$, to $V_j(n)$.
The intersection $V_j(N)\cap H_N$ is the kernel of
\[
              V_j(N)\oplus H_N\longrightarrow W_N,
              \qquad (f,h)\longmapsto f-h.
\]
Exactness preserves it. The tensor identification and the actual
multiplication map then preserve the multiplication image, and exactness
preserves its sum with $V_j(N)$. Induction gives
\begin{equation}\label{tx:lift-eq:stage-transport}
       \rho\bigl(V_j(N)\cap\ker z_{W_N}\bigr)=V_j(n)
                         \qquad\text{for every }j.
\end{equation}
The tensor domain is a subspace of $X_N\otimes H_N$, whose support bound
is $kd$. No closure of $\mathcal M_r$ under arbitrary tensor products
is assumed.

\begin{proposition}[The lifting formula preserves and reflects consequences]\label{tx:lift-prop:lift}
Assume the two facts in Section~\ref{tx:lift-sec:formula}. Then
$\Lift=e\iota$ preserves degree, $\rho\Lift=\mathrm{id}$,
$\Lift(1)=1$, and
\[
                  \Lift(f)\in V_j(N)\quad\Longleftrightarrow\quad
                  f\in V_j(n).
\]
The same equivalence holds for the stable spaces $C_d$.
\end{proposition}
\begin{proof}
Every retained monomial is fixed by the block action, so $z\iota(f)=0$.
Equation~\eqref{tx:lift-eq:onker} gives
\[
                  \Lift(f)=\iota(f)-z^{p-1}w\iota(f).
\]
The image of $z$ has no retained component: $z$ kills retained basis
vectors and preserves the span of moved ones. Hence $\rho\Lift(f)=f$.
On constants, every permutation acts trivially, so the correction term
vanishes and $\Lift(1)=1$.

If $f\in V_j(n)$, equation~\eqref{tx:lift-eq:stage-transport} supplies
$v\in V_j(N)\cap\ker z$ with $\rho v=f$. The permutation-space
identification says that $v-\iota(f)\in\ker z\cap\im z$. Since
$e\im z=0$, we have $\Lift(f)=ev$. The space $V_j(N)$ is invariant
under all label permutations, hence under $e$, and $ev\in V_j(N)$.
Conversely, $\Lift(f)\in V_j(N)$ is killed by $z$, so
\eqref{tx:lift-eq:stage-transport} and $\rho\Lift(f)=f$ give $f\in V_j(n)$.
Degree preservation follows from the renaming formula. Pass to the
unions for the stable statement.
\end{proof}

This proves both the forward lifting implication and
\eqref{tx:eq:newiff}. The image of the lift need not be all of $C_d(N)$:
the larger universe has many additional consequences. Every smaller consequence has this projector lift. For $n\ge n_*$ and a fully invariant target whose stabilised degree
is at most $d$, naturality of its map from $K$ gives
$\Lift(G_n)=G_N$: the target map commutes with the common projector,
which acts as the identity on $K$. More concretely, invariance gives
$zG_N=0$ and $eG_N=G_N$, while stable coefficients give
$\rho G_N=G_n$. The retained-vector identification and the argument
in Proposition~\ref{tx:lift-prop:lift} therefore yield $e\iota(G_n)=G_N$. Hence invariant targets, including
$1$, are exactly preserved. For indexed targets this equality need
not hold. The image-subquotient argument in
Chapter~\ref{tx:chap:assemble-transfer} proves transfer of
prescribed indexed targets without assuming it.

\subsection{Axioms and joint tensor multiplication}
For an axiom index $a$ supported in $[n]$, let $F_n$ and $F_N$ be the
usable-axiom coefficient maps. Write $\iota_A$ for inclusion of retained
axiom indices and $e_A$ for the action of the same group-algebra element
on their source space. Naturality gives
\begin{equation}\label{tx:lift-eq:axiomlift}
             \Lift(f_{a,n})=F_N(e_A\iota_A(a)).
\end{equation}
The right side is a linear combination of actual usable axioms of
$\Gamma_N$. This is the precise sense in which axioms lift.

For the order example, a retained predecessor row has axiom
$\sum_{j<n}F_{ij}-1$. Plain inclusion omits the new summands with
$j$ in the added block, so it is not the larger row axiom.
Equation~\eqref{tx:lift-eq:axiomlift} repairs precisely this mismatch:
the lifted smaller row is a linear combination of the actual larger
axioms. The order, Boolean, functionality and predecessor-guard
equations are handled by the same usable-axiom map. We need no
injectivity of the predecessor function. The common projector fixes
the last line $1$, and joint tensor compatibility handles every
multiplication step, including one made possible by cancellation.

Now suppose $f\in V_j(n)\cap H_n$, so the smaller proof may derive
$x_vf$. Apply the same $e$ to the \emph{joint} tensor input and expand
by its variable coordinate:
\begin{equation}\label{tx:eq:joint-lift}
 e_{X_N\otimes H_N}\bigl(\iota(x_v)\otimes\iota(f)\bigr)
                  =\sum_{y\in\mathcal V_N}x_y\otimes g_y.
\end{equation}
Exactness and the natural tensor identification show that this vector
lies in $X_N\otimes(V_j(N)\cap H_N)$. Thus each $g_y$ belongs to
$V_j(N)$ and has degree at most $d-1$. Multiplication commutes with
the common group-algebra action, so
\begin{equation}\label{tx:lift-eq:liftmul}
                         \Lift(x_vf)=\sum_y x_yg_y.
\end{equation}
Each term is a permitted multiplication of a derivable polynomial.

There is even an explicit expression for these coefficients:
\[
            g_y=\sum_{g\in S_N:\,g(v)=y}a_g\,g\bigl(\iota(f)\bigr).
\]
Individual summands need not be known consequences; their grouped sum
$g_y$ is a consequence by the tensor argument. One must not replace the
joint operator in \eqref{tx:eq:joint-lift} by $e_X\otimes e_H$. That would
discard the correlations supplied by the common permutation $g$.

To produce a literal proof sequence, compute bases of the closure
stages while storing derivations for those bases. Express each $g_y$ in
the appropriate basis by linear algebra, append the multiplication
steps in \eqref{tx:lift-eq:liftmul}, and add. This is a finite procedure, with
no efficiency or proof-length claim. It explains exactly how
consequence-space lifting supplies a legal PC derivation.

\subsection{A multiplication step with all its proof lines}
The distinction between a joint action and two independent actions is
visible in a very small calculation. Work over any field, with three
labels $0,1,2$, variables $x_0,x_1,x_2$, and the uniform axioms
$\Gamma_3=\{x_0,x_1,x_2\}$. These mean $x_i=0$. Set
$d=2$, $X=\Span\{x_0,x_1,x_2\}$ and $H=P_{3,\le1}$.
Take $f=x_0+x_1$, let $\tau=(0\ 2)$ rename the labels, and put
$T=1-\tau$ in the group algebra. Applying $T$ to the joint input gives
\begin{align*}
 T_{X\otimes H}(x_0\otimes f)
 &=x_0\otimes(x_0+x_1)-x_2\otimes(x_2+x_1)\\
 &=x_0\otimes g_0+x_2\otimes g_2,
 \qquad g_0=x_0+x_1,\quad g_2=-x_2-x_1.
\end{align*}
Both grouped coefficients are derivable linear combinations of the
axioms. Multiplying the two tensor coordinates produces a literal
degree-two PC derivation:
\begin{center}
\begin{tabular}{@{}cll@{}}
\toprule
Line & Polynomial & Reason\\
\midrule
1 & $x_0$ & axiom\\
2 & $x_1$ & axiom\\
3 & $x_2$ & axiom\\
4 & $x_0+x_1$ & lines 1 and 2\\
5 & $-x_2-x_1$ & lines 3 and 2\\
6 & $x_0^2+x_0x_1$ & multiply line 4 by $x_0$\\
7 & $-x_2^2-x_2x_1$ & multiply line 5 by $x_2$\\
8 & $x_0^2+x_0x_1-x_2^2-x_2x_1$ & add lines 6 and 7\\
\bottomrule
\end{tabular}
\end{center}
Lines 4 and 5 have degree one, so their multiplication is legal within
the degree-two budget. The last line is $T(x_0f)$. There has been no
reduction modulo Boolean equations and no multiplication by a
polynomial of degree greater than one.

By contrast, acting independently on the two factors gives
\[
 (T_X\otimes T_H)(x_0\otimes f)
       =(x_0-x_2)\otimes(x_0-x_2),
\]
whose product is $(x_0-x_2)^2$. This is a different polynomial, also in
characteristic two: the joint result still has the terms involving
$x_1$. The common permutation in each summand is therefore essential.

This calculation illustrates the inference mechanism at one rank.
The element $T$ is not a deletion projector and does not fix $1$.
For the actual larger-universe lift, replace $T$ by the common $e$
and use \eqref{tx:eq:joint-lift}. The grouped $g_y$ are then certified
consequences by exactness and the tensor identification, rather than
by the two-line calculation above. Once that certificate is available,
the final multiplication and addition steps are exactly the same.


\section{A one-label lifting example}\label{tx:lift-sec:small}
Work in characteristic two. Retain the label $a$ and add $b,c$.
Let $z=1+(bc)$ and $w=(ab)$ act on the linear variables
$x_a,x_b,x_c$. On their three-dimensional span,
\[
 z=\begin{pmatrix}0&0&0\\0&1&1\\0&1&1\end{pmatrix},\qquad
 w=\begin{pmatrix}0&1&0\\1&0&0\\0&0&1\end{pmatrix}.
\]
Direct multiplication gives $zwz=z$ and
\[
                       e=(1-zw)(1-wz)
                         =\begin{pmatrix}1&1&1\\1&1&1\\1&1&1\end{pmatrix}.
\]
Consequently
\begin{equation}\label{tx:lift-eq:smalllift}
                  \Lift(\alpha+\beta x_a)
                      =\alpha+\beta(x_a+x_b+x_c).
\end{equation}
Its retained coordinate is the original polynomial, and $\Lift(1)=1$.

Take the uniform degree-one axiom $\sum_{i<\ell}x_i=0$ at every size $\ell$.
The all-zero assignment satisfies this system at every size.
At size one the axiom is $x_a$; at size three it is $x_a+x_b+x_c$.
Formula~\eqref{tx:lift-eq:smalllift} sends the former exactly to the latter and
lifts every degree-one consequence. Plain inclusion would leave $x_a$,
which is not a consequence of the larger system: the assignment
$(x_a,x_b,x_c)=(1,1,0)$ satisfies its sum but not $x_a=0$.

Here the invariant target $G_\ell=\sum_{i<\ell}x_i$ is transported to
the prescribed larger target. In contrast, the indexed expression
$G_{a,\ell}=x_a$ has projector lift $x_a+x_b+x_c$, not $x_a$. The map
is linear but is not a ring homomorphism and does not preserve every
indexed expression literally. This calculation illustrates the distinction but lies outside the
indexed transfer range: for $k=d=t=1$ and $s=0$, where $t=1$ bounds
the support of the target index, that theorem uses
$R(d)=3$, whereas here $q=2$ and $n=1$. It is therefore no
counterexample to indexed periodicity.

The failure of a literal indexed lift requirement also has a general
explanation. At an admissible rank, naturality for the invariant target
$\sum_{i<n}x_i$ gives
$\Lift(\sum_{i<n}x_i)=\sum_{i<N}x_i$. If every retained indexed
target $x_i$ had lift exactly $x_i$, linearity would instead give
$\sum_{i<n}x_i$, a different polynomial. Thus at least one indexed
target must change under the projector lift. This is why the proof
uses actual target-map images and nonzero-subquotient detection.
For the displayed invariant example,
$r=1$ and $q=2>r$, so a two-label block suffices. Larger
support budgets require larger blocks. The sufficient witness threshold
$n\ge2r-1$ is attained in this example; the general template theorem
uses a slightly more conservative onset.

\section{Descent to the coefficient field}\label{tx:lift-sec:descent}
Here \emph{descent} means returning a consequence established with
auxiliary coefficients to the original coefficient field.
The auxiliary choices have to meet the conditions of the block and
common-witness constructions; they are not probabilistic assumptions.
Their existence is supplied by those constructions.

The generic block coefficients may require an extension field. This
does not change the derivation degree of any target with coefficients in $K$: at fixed degree, scalar extension
commutes with the finite-dimensional spans, intersections and images
in \eqref{tx:eq:closure}. In particular,
\[
                       V_j(n;E)=E\otimes_K V_j(n;K),
\]
with the analogous equality for $C_d$. Membership of every $K$-coefficient target is therefore unchanged,
so each target-membership equivalence proved over $E$ holds over $K$.

There is also a concrete $K$-linear lifting map satisfying the forward
requirements \eqref{tx:eq:newcontract}. All coefficients used in the
construction lie in a finite extension $E/K$ inside the algebraic
closure. Choose a $K$-linear functional $\lambda:E\to K$ with
$\lambda(1)=1$, for example the coefficient of $1$ in a vector-space
basis containing $1$. Apply it coefficientwise to polynomials and put
\begin{equation}\label{tx:lift-eq:fieldlift}
                 \widehat\Lift(f)=\lambda_*\bigl(\Lift_E(f)\bigr)
                      \qquad(f\in W_n(K)).
\end{equation}
The finite-extension assertion uses only finitely many coefficients
of the block operator and group-algebra witness. Each is algebraic
over $K$ when chosen in its algebraic closure; adjoining this finite
list gives a finite-dimensional extension.

The functional $\lambda$ is a linear projection, not a field map.
For example, over $\mathbb F_2\subseteq\mathbb F_4$ with
$\omega^2+\omega+1=0$, take $\lambda(a+b\omega)=a$.
Then $\lambda(\omega)=0$, but $\lambda(\omega^2)=1$.
This explains why descent works through linear consequence spaces:
it is not substitution by a multiplicative map of coefficients.

This map preserves degree, satisfies $\rho\widehat\Lift=\mathrm{id}$
and fixes $1$. If $f\in V_j(n;K)$, then $\Lift_E(f)$ is an
$E$-linear combination of a $K$-basis of $V_j(N;K)$. Applying
$\lambda_*$ leaves a $K$-linear combination of that same basis.
Thus $\widehat\Lift(f)\in V_j(N;K)$, as required.

This is descent of consequence-space membership, not coefficientwise
projection of an arbitrary proof with extension-field scalars. Such a
projection need not preserve each scalar inference. A proof over $K$
is recovered from the $K$-closure bases. We use only the forward lifting
properties for $\widehat\Lift$; the stronger reflection of individual
consequences in Proposition~\ref{tx:lift-prop:lift} refers to the projector lift
over $E$. The two-way prescribed-target statement over $K$ follows separately
from target-map transport, nonzero-subquotient detection and field
invariance, as above.


\section{A summary of the map's three different uses}
For any smaller consequence $f$, the projector lift gives a larger
consequence $\mathsf L(f)$ of no greater degree. For a full stable
invariant target fitting the budget, this vector is exactly the
prescribed larger target $G_N$. For general indexed targets, literal
pointwise equality is not available: the actual target-map image and
the quotient $(C_d(N)+\im J_N)/C_d(N)$ supply the transfer instead.
For NS, use the analogous quotient by $I_d(N)$. These quotients test
whether a target lies outside the relevant consequence space.

These uses share one construction but have different logical inputs.
Keeping them separate is the simplest way to avoid proving only a
statement about an unspecified lifted vector while claiming a theorem
about the specified target.

\begin{checkpoint}
A single common projector makes the entire closure calculation compatible.
It does not replace ordinary-PC multiplication by a new proof rule.
The larger consequence remains derivable by the original legal rules.
\end{checkpoint}

\part{Constructing a lift: from square-root to linear bounds}

\chapter{Nilpotent blocks and a common projector}\label{tx:sec:block-algebra}
\label{tx:ch:square-zero}
\section*{The property we need}
The lifting contract asks for one operation which works on an entire
proof diagram. It must transport the input-axiom map, the maps used for
multiplication, the inclusions of consequence spaces and the actual
target map. The first problem is to find an operation which has the
right effect on bases. The second, more difficult problem is to make
that operation exact on all the subspaces arising from a proof.

This chapter separates those problems. In characteristic two, a
square-zero operator gives a very concrete quotient: vectors killed
by the operator, modulo vectors in its image. On a permutation basis,
this quotient deletes the block labels. That calculation supplies
actual basis and tensor maps. It does \emph{not} prove exactness.
For exactness we will derive a projector from one group-algebra
witness, common to every tuple length up to the chosen support budget.
The next three chapters construct that witness rather than leaving
its existence as a hypothesis.

Keep the five roles separate:
\begin{center}
\begin{tabular}{@{}ll@{}}
$n$ & number of retained universe labels,\\
$q=2^a$ & number of block labels to be deleted,\\
$N=n+q$ & size of the larger universe,\\
$r$ & support budget for all spaces in the proof diagram,\\
$b\le r$ & current tuple length in the witness construction.
\end{tabular}
\end{center}
Neither $b$ nor $r$ is the maximum degree of a polynomial proof.
For an application to degree-$d$ consequences, that degree determines
$r$: $r=kd+s$ for invariant targets and $r=2kd+s+t$ for the general
indexed request. We first solve the finite algebra problem with $r$
as an independent integer.

The construction proves Exact Deletion
(Theorem~\ref{tx:thm:deletion}) independently of the equations and targets.
We begin with the square-zero operator in characteristic two,
then build a witness with a quadratic retained-label cost, replace that
cost by the linear threshold, and finally handle odd characteristic.

\section{A block of labels and a square-zero operator}\label{tx:sec:block2}
For this part $\operatorname{char}K=2$. Choose $q=2^a>r$ and split the labels into retained labels $R=[n]$ and a block $B$ with $q$ labels. Identify $B$ with the vector space $A=\Ftwo^a$. Translation by the $i$th unit vector gives a permutation $\sigma_i$ of $B$, extended by the identity on $R$. These permutations commute and square to the identity.

Extend the coefficient field if necessary so that it contains $a$
elements linearly independent over $\Ftwo$. In this finite-algebra
construction, $K$ denotes that working field; Lemma~\ref{tx:lem:field}
later returns degree-bounded derivability to the original field.
Choose $c_1,\ldots,c_a$ linearly independent over $\Ftwo$ in $K$, and put
\begin{equation}\label{tx:eq:z2}z=\sum_{i=1}^ac_i(\sigma_i-1).\end{equation}
This notation means an operator on every vector space on which label permutations act. Products of operators mean composition, so $z^2=0$ means that applying $z$ twice sends every vector to zero; it does not mean that $z$ itself is zero. Since the summands commute, cross terms cancel in characteristic two, and
\[z^2=\sum_i c_i^2(\sigma_i-1)^2=0.\]
A linear operator is called \emph{nilpotent} when some positive
power is zero. Here that power is already two, so every vector
in $\im z$ lies in $\ker z$. The nonzero two-point example below
shows that nilpotence need not make the operator itself zero.
For $q=2$, $z$ can be the swap of the two block labels plus the identity. For $q=4$, work over $\mathbb F_4=\Ftwo[\omega]/(\omega^2+\omega+1)$ and use $(c_1,c_2)=(1,\omega)$.

The independence is essential. On a two-point orbit where $\sigma_1$ and $\sigma_2$ induce the same swap, choosing $c_1=c_2=1$ makes $z=0$. The orbit is moved, yet would survive deletion. Independent coefficients ensure that every nonzero binary combination of the $c_i$ remains nonzero.

For a space $M$ with label action define
\begin{equation}\label{tx:eq:D2}\D(M)=\ker z/\im z.\end{equation}
This quotient is legal because $z^2=0$ implies $\im z\subseteq\ker z$. Two vectors killed by $z$ represent the same class if their difference is $zw$ for some $w$. A map commuting with label permutations commutes with $z$ and induces a map between these quotients. Permutations of the retained labels commute with $z$, so they still act on the quotient.

\subsection{Two actions which must not be confused}
The coordinates of a tuple have \emph{positions}, whereas its entries
are \emph{universe labels}. For example, the tuple $(a,b)$ has two
positions and uses two labels. A label permutation $g$ sends it to
$(g(a),g(b))$. A position swap sends it to $(b,a)$. These actions
commute, but they play different roles. Group-algebra elements in
$K[S_N]$ combine label permutations. The auxiliary inverses used later
must commute with position permutations so that their matrix entries
can be expressed as functions of an unordered matching between input
and output labels.

For a concrete check, take labels $1,2,3$ and the tuple $(1,3)$.
The label swap $(12)$ sends it to $(2,3)$; the position swap sends
it to $(3,1)$. Doing both, in either order, gives $(3,2)$. It is
this commuting pair of actions which allows the same correction to
be discussed both as an operator on ordered tuples and as a function
on rook placements.

The block subgroup $A$ acts on labels, never on tuple positions. The
operator $z$ therefore commutes with every position permutation.
An auxiliary operator chosen by the underlying \emph{set} of labels,
instead of by a preferred ordering of a tuple, has the same property.
That observation will matter when constructing an inverse without
averaging over $S_b$.

The usual averaging recipe would replace an operator $h$ by
$(b!)^{-1}\sum_{\tau\in S_b}\tau h\tau^{-1}$, thereby forcing
it to commute with positions. In characteristic two, $b!$ is zero
in the field when $b\ge2$, so this recipe is unavailable. Instead
we construct $h$ from the unordered label set of a tuple. That
choice is already unchanged by every reordering of its positions.

\subsection{The two-point orbit calculation}
For the illustrative two-label choice $z=\sigma-I$, on basis vectors $e_\beta,e_\gamma$ swapped by $\sigma$,
\[ze_\beta=ze_\gamma=e_\beta+e_\gamma.\]
The kernel and image are both the one-dimensional span of $e_\beta+e_\gamma$, so the quotient is zero. On a fixed basis vector $e_i$, $z=0$, and its line survives. This is the prototype of deletion: moved directions disappear, fixed directions remain.

\paragraph*{A useful failure to commute.}
On this two-point orbit let $S$ be the swap and let $D$ multiply
the first basis vector by $0$ and the second by $1$. Their matrices
are
\[
 S=\begin{pmatrix}0&1\\1&0\end{pmatrix},\qquad
 D=\begin{pmatrix}0&0\\0&1\end{pmatrix},\qquad DS-SD=S
 \quad(\operatorname{char}K=2).
\]
The difference $DS-SD$ is called a \emph{commutator}; it measures
the failure of the two operators to commute. Here that difference
is invertible. The next calculation constructs the same useful
failure to commute on every moved orbit, and uses it to write any
kernel vector as an image vector.

\subsection{The calculation for every moved orbit}
An $A$-set is a set of basis objects on which the translations in
$A=\mathbb F_2^a$ act. The orbit of $x_0$ consists of all $v x_0$,
$v\in A$; its stabiliser $H$ consists of translations fixing $x_0$.
Since the orbit is moved, $H$ is a proper subspace. A basis of $H$
can be extended to a basis of $A$, so a nonzero linear functional
$\ell$ can be chosen to vanish on $H$. It assigns a binary coordinate
to the orbit, independent of the translation used to reach an object.
The coordinate-reading operator $D$ below is diagonal, just as in
the two-point illustration.

Let $X$ be a finite $A$-set with no point fixed by all of $A$. On each orbit choose $x_0$. Its stabiliser $H$ is a proper binary subspace. Choose a nonzero linear functional $\ell:A\to\Ftwo$ vanishing on $H$. The function
\[\pi(vx_0)=\ell(v)\]
is well defined: two choices of $v$ differ by an element of $H$. Let $D$ multiply $e_x$ by $\pi(x)$, and put $e_i=\ell(\varepsilon_i)$. These are constants on that orbit. If $S_i$ denotes the action of $\sigma_i$, direct evaluation at a basis vector gives
\[DS_i-S_iD=e_iS_i,\qquad Dz-zD=W:=\sum_i c_ie_iS_i.\]
Moreover $Wz=zW$ and
\[W^2=\Bigl(\sum_i c_ie_i\Bigr)^2I=\alpha^2I,\qquad\alpha\ne0.\]
In particular $W^{-1}=\alpha^{-2}W$. Thus the inverse used here
is explicit; no choice of a complement or inverse matrix is hidden
in the notation. The coordinate vector $\varepsilon_i$ is the
$i$th unit vector in $A$, and $e_i=\ell(\varepsilon_i)$ is a
scalar, distinct from a basis vector $e_x$.
Thus $\Theta=W^{-1}D$ satisfies $\Theta z-z\Theta=I$. For $zm=0$ this gives $m=-z\Theta m$, hence $\ker z\subseteq\im z$. The reverse inclusion was automatic, so the quotient vanishes.

The calculation has three ingredients, each serving a different
purpose. The orbit coordinate $\pi$ produces a diagonal operator whose
commutator with translation is computable. Independence of the $c_i$
makes the commutator invertible. The commutator identity then turns a
kernel vector into an image vector. We do not choose a complement of
$\ker z$ by hand and hope it will respect later maps.

For a two-point orbit of $A=\mathbb F_2^2$ with stabiliser
$H=\{(0,0),(1,1)\}$, both coordinate translations induce the same
swap $S$. Choose $\ell(v_1,v_2)=v_1+v_2$. With
$(c_1,c_2)=(1,\omega)$, the operator is
$z=(1+\omega)(S-I)$, and $\alpha=1+\omega\ne0$.
Thus the one-dimensional kernel is exactly the image. If both
coefficients had instead been $1$, $z$ would be zero on this moved
orbit. This explicit orbit is why using one coefficient for every
translation is not a harmless simplification.

This calculation also works on $K[X]\otimes Y$ for \emph{any} space $Y$ carrying the label action. Let $D$ inspect only the $X$ coordinate and let $S_i$ act on both factors. We still have $S_i^2=I$ and the same commutator calculation. This stronger statement will justify multiplication, whose domain contains a variable factor tensored with a consequence subspace.

\section{Natural deletion on permutation bases}
Let $\mathcal U$ be a basis system whose objects have support less than $q$. An object fixed by all of $A$ cannot have support meeting $B$. Otherwise invariance of its support under the regular action would force its support to contain all $q$ block labels. Thus its fixed objects are precisely $\mathcal U_n$.

Split the basis into fixed and moved objects. The preceding calculation kills the moved summand and fixes the retained one. Hence
\begin{align}
\D(K[\mathcal U_{n+q}])&\cong K[\mathcal U_n],\label{tx:eq:basis}\\
\D(K[\mathcal U_{n+q}]\otimes Y)&\cong K[\mathcal U_n]\otimes\D(Y).\label{tx:eq:tensor}
\end{align}
The second identification sends $e_u\otimes[y]$ to $[e_u\otimes y]$ for retained $u$. It is natural in $Y$: applying a map to $y$ before or after taking its class gives the same result. For a product of two permutation bases, it simply keeps pairs of retained objects. There is no choice of unrelated isomorphisms here.

Suppose a matrix from one permutation basis to another has stable relabelling-invariant coefficients. On a retained input vector, its image is killed by $z$. Its retained component has exactly the smaller-rank coefficients; its moved component is in $\ker z=\im z$ on the moved summand, so vanishes in the quotient. Thus the \emph{actual matrix}, not merely its rank, is transported correctly.

\subsection{What naturality says about actual matrix entries}
Let $F_N:K[\mathcal U_N]\to K[\mathcal V_N]$ be an equivariant
coefficient map. On a retained basis vector $e_u$, its image is
\[
                       F_N(e_u)=\sum_v a_N(v,u)e_v.
\]
It is killed by $z$ because $e_u$ is killed by $z$ and $F_Nz=zF_N$.
Split this image into its retained and moved parts. Both parts are
killed separately because $z$ preserves those two summands. The moved
part is in $\im z$ by the orbit calculation, and so contributes zero
to the quotient. The remaining coefficients are $a_N(v,u)$ for
retained $u,v$. A stable coefficient rule identifies each of these
with $a_n(v,u)$. The smaller matrix is therefore obtained entry by
entry; we have not merely proved equality of two ranks or dimensions.

For a tensor map, apply the same argument to the permutation first
factor. A retained basis vector in that factor is fixed by the block,
so the quotient sees precisely the class of the second factor. A moved
first-factor orbit has zero quotient even when the block acts
nontrivially on the second factor. This is why the strengthened orbit
calculation, with a diagonal action on $K[X]\otimes Y$, was needed.
The tensor rule subsequently applies to a variable factor tensored
with a consequence subspace.

\subsection{A small failure of exactness}
Take the two-point permutation space $E=Ke_\beta\oplus Ke_\gamma$ and its invariant line $T=K(e_\beta+e_\gamma)$. Inclusion $T\hookrightarrow E$ is injective. But $\D(T)=T$ whereas $\D(E)=0$. Deletion takes an injection to the zero map on a nonzero space. Therefore the basis calculation alone does not prove that deletion preserves kernels, images or intersections.

This failure is precisely why a proof cannot stop at ``the monomials delete correctly.'' We must show that, with enough retained labels, every subspace produced by the PC calculation lies in a class where the quotient operation is exact.

\section{A common group-algebra witness}\label{tx:sec:common}
The group algebra $K[S_N]$ consists of finite formal sums $\sum_ga_gg$ of label permutations. It acts on every space with an $S_N$ action by the corresponding linear combination of operators. Its multiplication is composition. This provides a common language for certificates valid on many spaces at once.

For $b\ge0$ let
\[P_b(N)=K[\Inj([b],[N])],\]
the space on ordered injective $b$-tuples. In these construction chapters $P_b(N)$ denotes this tuple space, not the polynomial ambient space $P_{N,\le d}$ from the consequence calculation. The integer $b$ counts tuple positions, whereas $d$ bounds polynomial degree. Label permutations act on entries; a second group $S_b$ permutes tuple positions. These actions commute. Define
\[J_r(N)=\{a\in K[S_N]:aP_b(N)=0\text{ for every }b\le r\}.\]
It is a two-sided ideal. Write $\mathcal C_r(N)$ for the finite-dimensional spaces with label action annihilated by this ideal: every element of $J_r(N)$ acts as zero on each such space. This is the same annihilator class called $\mathcal M_r$ in Chapter~\ref{tx:chap:lift}. Subspaces invariant under labels, quotients and finite direct sums remain in this class.

Every bounded-support permutation basis of support at most $r$ belongs to the class. For an orbit of an object, order its support of size $b\le r$. Sending an injective tuple to the corresponding relabelling of that object gives a surjective equivariant map from $P_b(N)$ onto the orbit space. The support condition makes this well defined. Elements of $J_r$ kill that quotient as well.

The required witness is a \emph{partial inverse}, not an inverse of
the nilpotent operator $z$. If $zwz=z$ and $v=zt$ is an image
vector, then $z(wv)=zwzt=zt=v$: applying $w$ supplies a preimage
under $z$. The witness condition below requires this one expression
to work on all tuple lengths simultaneously. Its error need only
lie in $J_r$, since such an error acts as zero on the entire class.

\begin{definition}
A support-$r$ witness is one element $w\in K[S_N]$ satisfying
\begin{equation}\label{tx:eq:witness2}zwz-z\in J_r(N).
\end{equation}
\end{definition}
Such a $w$ gives $zwz=z$ on \emph{every} object of $\mathcal C_r(N)$. Choosing a separate matrix inverse on each space would not give this compatibility. A module map here is a linear map commuting with all label permutations. It therefore satisfies $f(wv)=w f(v)$ for a common group-algebra element, which is exactly what is needed.

\subsection{The projector calculation, without an exactness assumption}
A projector is an operator $e$ with $e^2=e$. It fixes every vector
in its image $eM$ and sends every vector in $\ker e$ to zero;
it need not be a coordinate projection.
We seek one which selects representatives
of the quotient $\ker z/\im z$ in a way compatible with every
equivariant map. The factors in the formula below act from right
to left. They are not assumed to commute.

Suppose $z^2=0$ and $zwz=z$. Set
\[e=(I-zw)(I-wz).\]
Then
\[ze=z(I-wz)=z-zwz=0,\qquad ez=(I-zw)z=0.\]
If $zv=0$, then $ev=v-zwv$. Also
\[(zw)e=(zw-zwzw)(I-wz)=0,\]
since $zwzw=zw$. Applying the formula on $\ker z$ to $ev$ gives $e^2v=ev$. Thus $e$ is a projector.

The map $[v]\mapsto ev$ is an isomorphism $\ker z/\im z\to eM$: it is well defined because $ez=0$, its image lies in $\ker z$, and $ev-v=-zwv$ is in $\im z$ when $zv=0$. Its inverse sends a vector of $eM$ to its quotient class.

Now let $0\to U\to V\to W\to0$ be an exact sequence in $\mathcal C_r(N)$. All maps commute with the same $e$. Restricting to $eU,eV,eW$ preserves injectivity. A vector of $eW$ has a lift $v\in V$; $ev$ is a lift in $eV$, proving surjectivity. A vector $ev$ mapping to zero lies in $U$ and equals its own projection, so lies in $eU$. This proves exactness explicitly. By the natural isomorphism, $\D$ is exact.

The projector need not commute with permutations of the retained labels. We use its image as a vector-space model to prove exactness, and keep the original quotient for the residual label action. Nothing here asserts an unwarranted equivariant direct-sum decomposition for that residual group.

\begin{example}[One retained label]
Let the retained label be $a$ and the two block labels be $b,c$. Put $z=I+(bc)$ and $w=(ab)$. On $P_1(3)$,
\[zwz(e_b)=zw(e_b+e_c)=z(e_a+e_c)=e_b+e_c=ze_b,\]
and similarly for $e_c$; both sides vanish on $e_a$. The projector is the all-ones $3\times3$ matrix. Its image is the constant line, naturally representing the retained one-dimensional quotient. With no retained label, every element of $K[S_2]$ satisfies $zwz=0$, so there is no such witness. The spare label makes a real difference.
\end{example}


\section{Reading the projector as an exactness certificate}
In a short exact sequence the first map $i$ is injective, the last
map $f$ is onto, and $\ker f=\im i$. Thus $U$ is a specified
subspace of $V$, and $W$ records the quotient by it. A
\emph{diagram chase} checks these assertions after transport by
following an individual vector and its preimages through the maps.
Here the common projector makes those vector checks possible.

The exactness proof is worth restating as a diagram chase, because
it is the point at which a basis calculation becomes a tool for
ordinary PC. Suppose
\[
 0\longrightarrow U\xrightarrow{i}V\xrightarrow{f}W\longrightarrow0
\]
is exact and all three spaces are in the common annihilator class.
The equations $ei=ie$ and $ef=fe$ hold because $i,f$ commute with
every label permutation and hence with their common linear combination
$e$. A separate projector selected independently on each space would
not justify these equations.

If $u\in eU$ maps to zero, its original image is zero and injectivity
of $i$ gives $u=0$. If $w\in eW$, take any lift $v\in V$ with
$f(v)=w$. Then $f(ev)=e f(v)=ew=w$. Finally, if $v\in eV$ maps
to zero, exactness of the original sequence gives $v=i(u)$ for some
$u\in U$. Applying $e$ gives $v=ev=i(eu)$. These are respectively
injectivity, surjectivity and equality of image with kernel for the
projected sequence. The natural quotient-to-projector isomorphism
transfers all three conclusions to deletion.

This argument proves more than the existence of some complement to
$\im z$. It proves that every actual map in the consequence diagram
uses compatible complements. Intersections, sums and restricted
multiplication images can therefore be transported. The joint tensor
operation in \eqref{tx:eq:joint-lift} uses the same group-algebra
expression on both factors together; it is not the tensor product
of two independently selected projectors.

\subsection{The three-label matrix in full}
With retained label $a$ and block labels $b,c$, use the ordered basis
$(e_a,e_b,e_c)$. Then
\[
 z=\begin{pmatrix}0&0&0\\0&1&1\\0&1&1\end{pmatrix},\qquad
 w=\begin{pmatrix}0&1&0\\1&0&0\\0&0&1\end{pmatrix}.
\]
Matrix multiplication gives $zwz=z$ and
\[
 e=(I-zw)(I-wz)=
       \begin{pmatrix}1&1&1\\1&1&1\\1&1&1\end{pmatrix}.
\]
Its square is itself: every entry of the square is $1+1+1=1$ in
characteristic two. The quotient can be represented by the retained
coordinate, while the projector realises that coordinate as
$e_a+e_b+e_c$. Thus the lift of a retained vector can involve all
labels. This is already the mechanism that turns a smaller totality
sum into its larger counterpart.

The example also shows why deletion is not substitution of zero for
block variables. The projector lift of $e_a$ is a sum of three basis
vectors. Its block terms are required to preserve the relations of
the larger system. A support budget does not assert that lifted proof
lines mention only retained labels; it limits the labels in each
\emph{basis object} used by the transport construction.

\section*{What has been proved, and what remains}
We have proved natural deletion of permutation bases and its
permutation-first-factor tensor rule. We have also proved that the
single identity $zwz=z$ on all tuple lengths up to $r$ gives a common
projector and therefore exact deletion on the full annihilator class.
The identity has not been assumed as an external theorem: its
construction is the subject of the next chapters. The preliminary
construction spends fresh labels at each length; the improved
construction realises a smaller, specially shaped error on one reused
pool.

\chapter{Lifting in characteristic two: a quadratic label budget}\label{tx:sec:quadratic-route}
\label{tx:ch:quadratic}
\section*{The property we need}
We remain in characteristic two. The block has $q=2^a$ labels,
and $z=\sum_i c_i(\sigma_i-I)$ is its square-zero translation
operator. A common witness is a group-algebra element $w$ satisfying
$zwz=z$ on every injective-tuple space $P_b$ with $b\le r$.
The block stays fixed throughout the following induction; the
retained labels are the resource we count.

We now construct the witness used in the previous chapter. The first
construction is intentionally generous with labels: at tuple length
$b$, it allocates $b$ labels which have not been used before. It
therefore proves exact deletion with a quadratic retained-label
budget. This gives a complete route to a square-root degree lower
bound and isolates exactly where that route loses a factor of $r$.

Keep the progression picture in view. Given a degree-$d$ proof at
one size eligible for this construction, we obtain a degree-$d$ proof
at every later size separated by the same block $q$. Thus $d$ is constant along
that progression. The quadratic cost below governs which starting
proofs are eligible; its inverse is the square-root lower-bound scale.

The induction maintains one group-algebra element, not a list of
independently chosen inverses. At the start of stage $b$, this element
satisfies $zwz=z$ on every shorter tuple module. We must repair length
$b$ while preserving all those previous identities. Three steps do
that: normalise the old element, localise its residual to tuples
without the fresh labels, and simulate a position-compatible inverse
by permutations involving those fresh labels.

Keep the order of those steps. A correct partial inverse of $z$ on one
vector space does not automatically give a partial inverse for the old residual, and an
operator which works after projecting back to old labels need not
work on all larger tuples. Normalisation settles the first issue;
localisation settles the second.

\section{Fresh-label construction}\label{tx:sec:quadratic}
We now prove that a support-$r$ witness exists in characteristic two when
\[n\ge T(r):=1+2+\cdots+r=\frac{r(r+1)}2,\qquad r<q.\]
This proof is independent of the rook theorem used later. Its ingredients are finite algebra, a position-compatible auxiliary inverse, and an explicit way to simulate that inverse using fresh labels.

The phrase ``generic coefficients'' below refers to the independent
block coefficients from Chapter~\ref{tx:ch:square-zero}. We choose
them in a finite field so that the subsequent ring calculation takes
place in a finite set. If $\Omega$ is the current finite label set,
$F[S_\Omega]$ has only $|F|^{|S_\Omega|}$ elements: each permutation
has one coefficient in $F$. The finite field is a construction device,
not a restriction on the field of the original polynomial request.

Choose the generic coefficients in a finite field $F$ of characteristic two, large enough to contain $a$ independent elements. We may work inside the algebraic closure of the original field; the degree question descends by the field-extension argument. The finiteness of $F$ will be used only to normalise an element of a finite group algebra.

\subsection{Normalisation of the residual}
Suppose $w_0$ works on shorter tuples but not yet on length $b$. Put $y=z-zw_0z$. An auxiliary linear map $h$ may satisfy $zhz=z$ on the larger tuple space. It is tempting to conclude $yhy=y$. In general the expansion contains an extra term $zw_0zw_0z$, and this need not collapse. Before correcting the error we arrange the additional identity $w_0zw_0=w_0$ in the actual group algebra. This normalisation was a genuine mathematical step in the earlier construction.

Here \emph{normalisation} is an algebraic replacement: it modifies
a ring element to satisfy the
additional algebraic identity needed for the residual calculation.
A unital ring has addition, associative multiplication and an identity
$1$; its multiplication need not commute. A ring representation applies
its elements as operators, respecting addition, composition and the
identity. Thus an identity proved in the ring holds in all its
representations at once. In the lemma, $v$ is a ring element and
the letters $i,\ell,h$ in the proof are integers recording powers;
this local $h$ is not the auxiliary linear operator used elsewhere.

\begin{lemma}[Normalisation in a finite ring]\label{tx:lem:normal}
Let $v,z$ belong to a finite unital ring. There is an element
$v^\natural$ in the same ring with
$v^\natural z v^\natural=v^\natural$, such that every ring
representation in which $zvz=z$ also satisfies
$zv^\natural z=z$.
\end{lemma}
\begin{proof}
Put $x=vz$. Finiteness gives $x^{i+\ell}=x^i$ for some $i\ge0$
and $\ell>0$. Multiplication by powers of $x$ gives
$x^{j+\ell}=x^j$ whenever $j\ge i$. Choose $h\ge i$ such that
$h+1$ is divisible by $\ell$. Then $2h+1-h=h+1$ is a multiple
of $\ell$, so $x^{2h+1}=x^h$. Define $v^\natural=x^hv$.
Associativity, without commutation of factors, gives
\[
 v^\natural z v^\natural
      =x^h(vz)x^hv=x^{2h+1}v=x^hv=v^\natural.
\]
In a representation with $zvz=z$, induction gives
$z(vz)^j=z$ for every $j\ge0$. Hence
\[
                     zv^\natural z=z(vz)^{h+1}=z.
\]
The exponent was selected inside the original finite ring. The
normalisation therefore preserves the witness relation simultaneously
in every representation where it was already true.
\end{proof}
In the construction below we denote the normalised old witness by $b_0$. It is one ring element shared by every shorter tuple representation.
The lemma preserves $zvz=z$ wherever it was already valid. It does
not assert that a formerly unsuccessful tuple length has been repaired.
The new identity $b_0zb_0=b_0$, obtained in the ring itself, is what
makes the following sandwich of the residual work.

\begin{lemma}[Normalised residual]\label{tx:lem:residual}
If $b_0zb_0=b_0$ and $zhz=z$, then, for $y=z-zb_0z$,
\[yhy=y.\]
\end{lemma}
\begin{proof}
Expand without commuting any factors:
\begin{align*}
yhy&=zhz-zhzb_0z-zb_0zhz+zb_0zhzb_0z\\
&=z-2zb_0z+zb_0zb_0z\\
&=z-zb_0z=y.
\end{align*}
The last equality uses the normalisation. This identity works in any ring; characteristic two is not needed for this particular calculation.
\end{proof}

\subsection{A residual calculation over \texorpdfstring{$\mathbb F_4$}{F4}}
The notation $E_{ij}$ denotes a matrix unit: it sends basis vector
$e_j$ to $e_i$ and sends the other basis vectors to zero. Thus
$E_{12}(s,t)=(t,0)$ and $E_{21}(s,t)=(0,s)$ on a two-dimensional
coordinate space. Applying $E_{12},E_{21},E_{12}$ successively
returns $E_{12}(s,t)$, which makes the partial-inverse identity
in the calculation below transparent.

The normalisation is not optional algebraic tidying. Let
$\omega^2+\omega+1=0$, let $z=E_{12}$ on $\mathbb F_4^2$, and
let $h=E_{21}$. Then $zhz=z$. If we take the unnormalised old
operator $w_0=\omega E_{21}$, its residual is
\[
                   y=z-zw_0z=(1-\omega)z=\omega^2z.
\]
But $yhy=\omega^4z=\omega z\ne\omega^2z=y$. Thus a partial inverse
for $z$ does not automatically serve as an inverse for $y$. The
finite-ring lemma repairs exactly the missing identity used in the
expansion of $yhy$.

This example is a two-dimensional operator calculation explaining
one algebraic step. It is not being used as a label-permutation
construction or as an asymptotic counterexample. The group-algebra
construction obtains its normalisation in the finite ring
$F[S_\Omega]$ itself, which is why the corrected identity holds in
all the required modules at once.

\subsection{An auxiliary inverse compatible with tuple positions}
Let $\Omega$ be the current label set, consisting of the block $B$ and the retained labels already allocated. On $Q=P_b(\Omega)$, partition the tuples according to the $A$-orbit of their \emph{underlying label set}. A set meeting $B$ has proper stabiliser because it has size $b<q$. Choose the functional $\ell$ and coordinate $\pi$ as in Section~\ref{tx:sec:block2}, now on these sets. Let $D$ read $\pi$ from the underlying set, and use the same $W$ and nonzero $\alpha$ as before. The choices and the scalar $\alpha$ are made separately on each orbit; the resulting operators are combined on their direct sum. On each moved orbit,
\[Dz-zD=W,\quad Wz=zW,\quad W^2=\alpha^2I,\quad z^2=0.\]
Multiplying the commutator by $z$ shows $zDz=zW$. Thus
\[h=\alpha^{-2}WD\quad\text{satisfies}\quad zhz=z.\]
On tuples wholly outside $B$, put $h=0$, since $z=0$ there. Because $D$ depends on the underlying set and translations act on labels, $h$ commutes with every permutation of tuple positions. We have not claimed that $h$ is a group-algebra element. It is just an explicitly constructed linear operator.
Here and in this section, ``inverse'' is shorthand for the partial
identity $zhz=z$, not for $zh=hz=I$.

\subsection{Fresh labels simulate any position-compatible operator}
Write $C=\{\kappa_1,\ldots,\kappa_b\}$ for the fresh labels.
The old tuple space sits inside the larger one by including its basis
vectors; projection back discards the basis vectors using a fresh label.
The composite of that inclusion and projection on the old space is
the identity. We want only the old-to-old part of a larger permutation
combination, which is the meaning of the compressed operator below.

Adjoin $b$ new labels $\kappa_1,\ldots,\kappa_b$ to $\Omega$. Let $i:P_b(\Omega)\to P_b(\Omega\cup C)$ include the old tuples and let $p$ discard every tuple using a fresh label. We prove that any operator $h$ commuting with position permutations has the form
\begin{equation}\label{tx:eq:compression}h=pci\end{equation}
for some $c\in F[S_{\Omega\cup C}]$.
The order of operations in $pci$ is: include the old input, apply
$c$, then project the result to old outputs. The symbol $p$ here
names this projection, not the prime characteristic. The formula
does not say that $c$ preserves the old space or agrees with $h$
on every larger tuple; it identifies its old-to-old matrix corner.

Fix a $b$-element set $U=\{u_1,\ldots,u_b\}\subseteq\Omega$, with one chosen ordering. Form
\[E_U=\prod_{j=1}^b\bigl(1-(u_j\ \kappa_j)\bigr).\]
The transpositions are disjoint. On an old tuple missing some $u_j$, the corresponding factor kills the tuple, and the other factors cannot introduce $u_j$. On a tuple using all the $u_j$, the identity term leaves the tuple old, while every nonidentity term introduces a fresh label. Therefore $pE_Ui$ is precisely the projector onto the orderings of $U$.

The $b!$ orderings of $U$ form a \emph{free orbit} for position
permutations: exactly one reordering carries a chosen ordering to
any other. Consequently a position-compatible map is determined
on that orbit by its value at one ordering. All other values are
obtained by the same position reordering of the output. This is
why one can prescribe the map without averaging or dividing by $b!$.

For an ordered output tuple $v$, choose a permutation $g\in S_\Omega$ taking $(u_1,\ldots,u_b)$ to $v$. Extend $g$ by fixing fresh labels. Then $pgE_Ui$ takes that input ordering to $v$, takes its position permutations to the corresponding permutations of $v$, and kills all other underlying sets. The orderings of each $U$ form a free orbit for $S_b$. Consequently an operator commuting with $S_b$ is completely determined by its value on one ordering for each $U$. Expanding those values as linear combinations of output tuples and using the operators $pgE_Ui$ proves \eqref{tx:eq:compression}. There is no division by $b!$.

\begin{example}[Two registers]
For $U=\{u_1,u_2\}$,
\[E_U=1-(u_1\kappa_1)-(u_2\kappa_2)+(u_1\kappa_1)(u_2\kappa_2).
\]
On $(u_1,u_2)$, projection $p$ kills the last three terms. On $(u_1,v)$ with $v\ne u_2$, the identity contribution is cancelled by the second swap factor, which fixes that tuple. The detector recognises the set $U$, while leaving the two possible orderings available for a position-compatible output rule.
\end{example}

\subsection{Why compression is enough only after localisation}
Equation~\eqref{tx:eq:compression} asserts $h=pci$ on the old tuple
space. It does not assert that $c$ acts like $h$ on every tuple in
the larger universe. Terms of $c$ may introduce fresh labels. This
would normally be a serious problem: a proof must preserve all basis
objects, not just an arbitrarily selected corner of a matrix.

The residual has a special support in the larger tuple space. It is
zero on every summand with a fresh label, because such a summand uses
fewer than $b$ old labels and the old witness is already valid at
those shorter lengths. Therefore the two residual factors in
$ycy$ project to the old corner before and after the middle operator.
The identity $h=pci$ is exactly the information the sandwich needs.
The next calculation verifies this point as an equality of maps,
rather than assuming that compression itself is a global realisation.

In matrix terms, order the basis with old tuples first and tuples
using fresh labels second. Under the induction hypothesis the
calculation below gives
\[
                       y=\begin{pmatrix}y_Q&0\\0&0\end{pmatrix}.
\]
In $ycy$ the rightmost $y$ removes fresh input directions, and the
leftmost $y$ removes fresh output directions. Only the old-to-old
corner $pci=h$ of the middle operator matters. The zeros here come
from validity of the old witness at shorter lengths, not from a
claim that an arbitrary permutation combination has this block form.

\subsection{The induction, including localisation of the error}
For support zero take $w_0=0$: $z$ is zero on constants. Suppose a witness exists on the block and $T(b-1)$ retained labels, collectively $\Omega$. Normalise it in the finite ring $F[S_\Omega]$ to $b_0$. It still works on all tuple lengths below $b$ and satisfies $b_0zb_0=b_0$ in the ring itself. Set $y=z-zb_0z$.

Adjoin $b$ fresh labels. Decompose $P_b(\Omega\cup C)$ according to the positions occupied by fresh labels and their actual values. The group $S_\Omega$ fixes each fresh label, so every summand with at least one fresh label is a copy of $P_j(\Omega)$ for $j<b$. On it $y=0$ by the old witness relation. The summand with no fresh labels is $Q=P_b(\Omega)$. Thus, as actual maps,
\[y=i y_Q p.\]
This localisation is what allows a compressed inverse to suffice. Construct the auxiliary $h$ on $Q$ above. Lemma~\ref{tx:lem:residual} gives $y_Qhy_Q=y_Q$. Choose $c$ with $pci=h$. Then
\[ycy=i y_Q(pci)y_Qp=i y_Qp=y.\]
Define the new group-algebra element
\[w_b=b_0+(1-b_0z)c(1-zb_0).\]
Since $z(1-b_0z)=(1-zb_0)z=y$,
\[zw_bz=zb_0z+ycy=zb_0z+y=z\]
on length $b$. On every shorter tuple space the same fresh-label decomposition makes $y=0$, so the relation remains valid there too. The new cost is exactly $b$ retained labels, giving
\[T(b)=T(b-1)+b.\]
If still more retained labels are present, extend the constructed permutations by fixing them. Decomposing according to those outside labels reduces every tuple of length at most $b$ to a shorter tuple on the constructed labels. Thus the same witness works for all $n\ge T(b)$.

\begin{theorem}[Quadratic retained-label budget]
For every $r<q=2^a$ and $n\ge r(r+1)/2$, there is a common witness $w$ with $zwz=z$ on all tuple modules of length at most $r$. Therefore deletion is exact on the support-$r$ class. It transports the
entire degree-$d$ PC consequence space with $r=kd+s$, beyond the fixed
template threshold. For prescribed indexed targets use the budget
$r=2kd+s+t$ to include the target map and the doubled-support detector.
\end{theorem}
Every construction needed for this theorem has now been supplied. In particular, the proof is not an inference from witnesses found for small tuple lengths.


\begin{remark}[The square-root scale]
For either prescribed-target budget, the construction gives an onset
$O(d^2)$ for degree-$d$ transfer.
Inverting that onset gives an $\Omega(\sqrt n)$ lower bound when no
finite degree recurs, and on the hard residues of the periodic case.
The loss is precisely the cumulative allocation $T(r)=r(r+1)/2$.
The following construction replaces this sum by the maximum $2r-1$,
using one retained pool at every layer.
\end{remark}

\section{The first three induction stages}
At length zero there is one basis vector, the empty tuple. Every
permutation fixes it, so $z=0$ and $w_0=0$ is a valid witness. No
retained label is needed.

At length one, the auxiliary operator recognises the moved label
orbits. One fresh retained label realises its compressed inverse.
The witness can therefore be built on the block plus one retained
label. For a block of two labels, the previous chapter exhibited
this witness explicitly as a transposition between one block label
and the retained label.

At length two, normalise that same witness first. Allocate two
additional retained labels $\kappa_1,\kappa_2$. A tuple containing
either fresh label has at most one old entry; the residual is zero
on it. Only the tuples with both entries old require repair. The
detector
\[
                 (1-(u_1\ \kappa_1))(1-(u_2\ \kappa_2))
\]
selects an old two-element set. Combining its translates realises
any position-compatible inverse on the old corner. The total number
of retained labels is now $1+2=3$.

At length three, the same argument spends three further labels.
Tuples with one, two or three of those fresh labels reduce to old
tuple lengths two, one or zero, respectively. All three pieces
already have zero residual. The only unsolved piece consists of
old three-tuples, and a detector with three disjoint swaps compresses
its auxiliary inverse. The retained total is $1+2+3=6$.

The block size must exceed the largest tuple length being treated.
For example, one may choose a block of size four when explaining
stages one through three. A two-label block illustrates stage one,
but cannot be reused as a valid block for stage three. Retained
labels and block labels have different jobs in this count.

\begin{center}
\begin{tabular}{@{}c|rrrr@{}}
Tuple length $r$ &1&2&3&4\\
Fresh-label total $T(r)$ &1&3&6&10\\
Linear total proved later $2r-1$ &1&3&5&7
\end{tabular}
\end{center}
The first strict improvement is at length three. The linear method
will not allocate the third batch separately: it will exploit the
residual's zero marginals and reuse a common retained pool.

\section{From quadratic label cost to square-root degree hardness}
Return now from tuple spaces to one fixed polynomial request and one
proof system, NS or PC. Its profile $h(n)$ is the least degree for
one fixed target orbit, or the maximum of the least degrees for an
all-target request. It is infinity if any requested target is
underivable; an empty request has value zero. A finite value recurs if that exact value is
attained at infinitely many universe sizes. The block construction
will force recurrence whenever such a value occurs beyond its
eligible starting rank.

Let $r(d)=\alpha d+\beta$, where
$(\alpha,\beta)=(k,s)$ for an invariant target or
$(\alpha,\beta)=(2k,s+t)$ for an indexed request. The fixed starting
rank $n_*$ includes degree stabilisation of the targets. A convenient
safe onset for this preliminary construction is
\[
 N_{\rm quad}(d)=
  \max\{n_*,\,2r(d)+1,\,T(r(d))\}\qquad(d\ge1).
\]
The extra linear term is harmless and makes this onset comparable
with the eventual linear statement. Exactness of the quadratic
construction, together with the target-map and detection arguments,
gives degree-threshold periodicity beyond this onset. Take the least
power of two $Q(d)>r(d)$; then $Q(d)\le2r(d)$.

Here $\alpha,\beta$ are integer coefficients of the support budget,
unrelated to the nonzero field scalar $\alpha$ in the orbit
calculation. The onset estimates concern positive degree budgets.
Degree zero uses its previously established fixed predicate beyond
$n_*$; it therefore cannot be a nonrecurring value on a hard tail.
The hard-value inversion below concerns positive finite values,
and the unbounded recurring values chosen for peaks are positive too.

Suppose a profile has no recurring finite degree. If $h(n)=e<\infty$
and $n\ge N_{\rm quad}(e)$, threshold periodicity at all budgets
through $e$ repeats that exact value on an infinite arithmetic
progression. This is a contradiction. Thus every finite value must
satisfy $n<N_{\rm quad}(e)$. Since
$N_{\rm quad}(e)\le A(e+1)^2$ for a constant $A$ depending only
on the fixed request, it follows that
\[
                          h(n)>\sqrt{n/A}-1
\]
whenever $n$ is beyond the fixed initial range. The assertion also
holds when $h(n)=\infty$. The same argument works on hard residues
in a fixed periodic easy/hard split: a value attained beyond its
onset would recur and would therefore be one of the bounded easy
values.

For a fluctuating profile, choose recurring values $e$ tending to
infinity. A representative of its recurring residue can be selected
in
\[
       N_{\rm quad}(e)\le n<N_{\rm quad}(e)+Q(e).
\]
Both the onset and the period are bounded by a constant times
$(e+1)^2$, so these representatives are finite square-root-scale
peaks, $e\ge c\sqrt n-O(1)$, for a fixed $c>0$. They are not an
all-size lower bound; low recurring values remain at arbitrarily
large sizes.

Nothing in this argument makes the degree of an individual proof
shrink. At fixed $d$ we compare derivability at two universe sizes.
The lower bound follows by comparing that \emph{fixed budget} with
its allowed universe-size onset. The preliminary route just proved
is in characteristic two. The later prime-characteristic chapter
will give the full linear result directly over every fixed positive
characteristic.

\chapter{Towards linear lifting: reuse one pool of labels}\label{tx:sec:linear-construction}
\label{tx:ch:linear-layers}
\section*{The property we need}
Keep the characteristic-two block from the previous chapters:
$N=n+q$, with $n$ retained labels, $q=2^a>r$ block labels and
the square-zero operator $z$. We seek one $w$ with $zwz=z$ on
all ordered-tuple spaces of lengths at most $r$. The final retained
label requirement will be $n\ge2r-1$; the block and the budget $r$
stay fixed while the tuple length $b$ increases.

The quadratic construction repaired an arbitrary position-compatible
operator on the whole old $b$-tuple space. That is more freedom than
we need. The old witness error already annihilates every direction
coming from shorter tuples and has zero marginal in every output
position. We will use that shape to construct a correction on the
same retained label pool.

There are two separate contracts. The present chapter constructs a
position-compatible operator $\varphi$ with $z\varphi z=-\varepsilon$,
whose image has zero marginals and which kills the incoming shorter
layers. The next chapter proves that every operator with exactly
this shape is the action of a group-algebra element annihilating the
shorter tuple modules. It uses no additional labels. Combining the
two contracts yields one common witness at retained rank $2r-1$.

This is the improvement which changes the degree gap. With fresh labels,
the retained-label cost was a sum over all tuple lengths. Reusing one
pool makes it a maximum. A proof with degree $d$ can therefore be lifted
once the universe is only linearly large in $d$. If that degree cannot
recur on an infinite progression, every occurrence must have been
before this linear starting rank.

Semisimplicity would supply invariant complements to every invariant
subspace (Section~\ref{tx:sec:action-primer}); here we construct the
particular compatible complements needed by the lift.
The construction builds sections using label-avoiding decompositions,
then makes those sections commute with the block operator. In
positive characteristic, a request for an equivariant section cannot
be justified by averaging. Every such section below is explicit.

A complement to a subspace $U\subseteq V$ is a subspace $C$ with
$V=U\oplus C$. A basis extension always supplies a vector-space
complement, but it need not respect a permutation action or the
operator $z$. A compatible complement therefore needs a construction.
For a surjection $f:E\to Q$, a \emph{section} is a linear map
$\sigma:Q\to E$ with $f\sigma=I$: it chooses a preimage of each
vector, consistently with linear combinations. Such a section gives
$E=\ker f\oplus\im\sigma$. Requiring it to commute with an action
is an additional condition, not a consequence of this decomposition.

\section{The zero-marginal layer}\label{tx:sec:layers}
We still work in characteristic two. Fix a tuple length $b\ge1$, distinct from the proof-degree budget $d$. Let $V=P_b(N)$. For each position $i$ there is a deletion map
\[\partial_i:P_b(N)\longrightarrow P_{[b]\setminus\{i\}}(N)
\]
which forgets the entry in position $i$. On a vector of coefficients it sums over that entry, so we call it a marginal. The target is a copy of $P_{b-1}(N)$ with the surviving positions named. Give tuple spaces the dot product making their bases orthonormal, and put
\begin{equation}\label{tx:eq:anatomy}
V'=\sum_i\im\partial_i^{\mathsf T},\qquad
\mathcal Z_b(N)=\bigcap_i\ker\partial_i=(V')^\perp,\qquad T_b(N)=V/V'.
\end{equation}
Here $\mathcal Z_b$ is the zero-marginal space, and $T_b$ is the top quotient. This notation distinguishes the tuple zero-marginal space from the
polynomial consequence space $C_d(n)$. All statements below use the
displayed definitions.

These are marginal sums, rather than projections onto some of the
original coordinates. At length two, for example,
\[
 \partial_1\Bigl(\sum_{i\ne j}a_{ij}e_{ij}\Bigr)
                  =\sum_j\Bigl(\sum_{i\ne j}a_{ij}\Bigr)e_j.
\]
All entries with the same surviving label contribute to one output
coefficient. A zero marginal says that these contributions cancel.

The dot product here is the bilinear coefficient pairing
$\langle v,w\rangle=\sum_t v_t w_t$, without conjugating coefficients.
It is nondegenerate because the tuple basis is orthonormal. It is not
a positive inner product: over $\mathbb F_2$, the nonzero vector
$e_1+e_2$ pairs to zero with itself. Thus $(V')^\perp$ is the space
of vectors pairing to zero with every vector in $V'$, and need not
be a complementary subspace to $V'$. The transpose is characterised
by $\langle Tv,w\rangle=\langle v,T^{\mathsf T}w\rangle$.

The name \emph{top quotient} records what remains after quotienting
by the incoming directions from shorter tuple spaces. It does not
assert that this quotient is semisimple.

The transpose map has a direct combinatorial meaning. For a
$(b-1)$-tuple $u$ with position $i$ missing,
\[
 \partial_i^{\mathsf T}e_u
       =\sum_{x\notin\im u} e_{u[i\mapsto x]}.
\]
It inserts each admissible label once. There is no factor $1/N$ and
no division by the number of available labels. This is the incoming
subspace $V'$, and it is a sum of images, not an arbitrarily selected
copy of $P_{b-1}$.

The incoming space $V'$ contains directions lifted from shorter tuples. The zero-marginal space contains directions invisible after deleting any one position. The quotient $T_b$ forgets the incoming directions. Label permutations commute with all marginal maps, and position permutations permute the marginal maps. Thus these spaces carry the two appropriate actions.

\begin{example}[Lengths one and two]
For $b=1$, $\partial$ sums coordinates, $V'$ is the constant line and $\mathcal Z_1(N)$ is the zero-sum space. For $b=2$, a coefficient vector is an off-diagonal array $f(i,j)$, $i\ne j$. Its marginals are its row and column sums. Thus $\mathcal Z_2(N)$ consists of off-diagonal arrays with every row and column sum zero. On three labels it is spanned by
\[e_{12}+e_{23}+e_{31}-e_{21}-e_{32}-e_{13}.
\]
In characteristic two the six coefficients are all one, but every line still sums to zero because it contains two ones.
\end{example}

\subsection{The old error already has a special shape}
The annihilator ideal $J_{b-1}(N)$ consists of combinations of label
permutations that act as zero on every tuple space of length below $b$.
Suppose $w_0$ is a witness on all lengths below $b$, and set
\[\varepsilon=zw_0z-z\quad\text{on }V.\]
This is the action of an element of $J_{b-1}(N)$. For any $x\in J_{b-1}(N)$,
\[\partial_i x=x\partial_i=0,\qquad x\partial_i^{\mathsf T}=\partial_i^{\mathsf T}x=0.
\]
In these equations the same group-algebra element $x$ acts on both
tuple spaces. The marginal and its transpose commute with label
permutations, while $x$ acts as zero on the shorter $(b-1)$-tuple
space by the definition of $J_{b-1}(N)$.
Therefore $x$ maps $V$ into $\mathcal Z_b(N)$ and kills $V'$. In particular $\varepsilon$ factors through
\[V\longrightarrow T_b(N)\longrightarrow \mathcal Z_b(N)\longrightarrow V.
\]
To \emph{factor through} the quotient means that
$\bar\varepsilon([v])=\varepsilon v$ is well defined: replacing $v$
by $v+v'$ with $v'\in V'$ does not change the answer. Its values
belong to $\mathcal Z_b$, and the final arrow is that space's actual
inclusion into $V$. This turns the error into a map from the quotient
by incoming shorter-tuple directions to the zero-marginal output space.
It also commutes with position permutations. This is much more special than an arbitrary operator on $V$. The improved route exploits exactly this special form.

There are two tasks. First, construct a position-compatible correction of this shape with $z\varphi z=-\varepsilon$. Second, realise that correction by an element of the actual label group algebra which kills all shorter tuples. The first task needs retained labels; the second will be solved by a rook-placement theorem without adding labels.

\subsection{A length-three zero-marginal vector}
For three positions and four labels, write
$\epsilon_{ijkl}$ for the sign of the permutation $(i,j,k,l)$
of $(1,2,3,4)$, and let $l$ be the unique missing label. Then
\[
           Y=\sum_{i,j,k\text{ distinct in }[4]}
                        \epsilon_{ijkl}e_{ijk}
\]
has all three marginals zero. Fix, for example, the second and third
entries $2,3$. The only contributions are $e_{123}$ and $e_{423}$,
with opposite coefficients, so their sum under deletion of the first
position is zero. The same argument for any surviving ordered pair
uses the transposition between the missing label and the deleted
entry. It changes the sign and pairs the two terms. In characteristic
two the signs become equal, and the two coefficients still sum to
zero. Thus the zero-marginal condition is concrete at lengths one,
two and three; it is not a name for an unspecified representation.

The linear witness proof does not require this particular vector to
span $\mathcal Z_3$. It illustrates the cancellations which the
partition and cone maps must respect for \emph{every} vector in
that space.

\section{Partitioning a zero-marginal vector so that one label is avoided}\label{tx:sec:partition}
For a finite position set $S$ and label set $L$, let $P_S(L)=K[\Inj(S,L)]$ and $\mathcal Z_S(L)=\bigcap_{i\in S}\ker\partial_i$. If $L'\subseteq L$, regard $\mathcal Z_S(L')$ as the vectors of $\mathcal Z_S(L)$ supported on tuples using only $L'$. Write $\mathcal Z_c$ when there are $c$ positions.

Suppose a block label $\beta$ occupies a new position and $y\in \mathcal Z_c(L)$ occupies the other positions. The vector $e_\beta\otimes y$ has zero marginals in those other positions, but its new-position marginal is $y$. We want a retained correction with the same marginal. If
\[y=\sum_{x\in L}\pi_x(y),\qquad \pi_x(y)\in \mathcal Z_c(L\setminus\{x\}),
\]
then
\begin{equation}\label{tx:eq:standin}
e_\beta\otimes y-\sum_{x\in L}e_x\otimes\pi_x(y)
\end{equation}
has every marginal zero. Avoiding $x$ is essential: it makes each displayed tuple injective. The maps must also commute with permutations of the old positions, so they can be used consistently in all layers.
Here $e_\beta\otimes y$ means inserting $\beta$ into the designated
new position of every tuple in $y$; the other positions retain their
names. It is tuple-coordinate notation, not multiplication of
polynomials.

\begin{example}[One old position]
Take $L=\{1,2,3\}$ and $y=a e_1+b e_2-(a+b)e_3$. Define
\[\pi_1(y)=b(e_2-e_3),\quad \pi_2(y)=a(e_1-e_3),\quad\pi_3(y)=0.
\]
Each summand avoids its subscript and has coordinate sum zero; their sum is $y$. For $y=e_1-e_2$, formula \eqref{tx:eq:standin} is the six-term vector
\[e_{\beta1}-e_{\beta2}+e_{12}-e_{13}-e_{21}+e_{23}.
\]
Check its two marginals by grouping terms with the same surviving label. With only labels $1,2$, every $\mathcal Z_1(L\setminus\{x\})$ is zero, while $e_1-e_2\ne0$. The third label is what permits the construction.
\end{example}

The phrase \emph{partition of unity} below refers to the identity
$\sum_x\pi_x=I$ between linear maps. Each map distributes part of
the input into a space avoiding its label. The outputs may overlap;
the maps are not required to be positive or idempotent projections.
Their equivariance concerns permutations of the tuple positions.
They can choose distinguished labels and need not commute with all
permutations of the label set.

\begin{lemma}[Partition of unity]\label{tx:lem:unity}
If $|L|\ge2c+1$, there are position-equivariant linear maps
\[\pi_x:\mathcal Z_c(L)\longrightarrow \mathcal Z_c(L\setminus\{x\}),\qquad
\sum_{x\in L}\pi_x=I.
\]
This statement and its proof work over any field.
\end{lemma}
\begin{proof}
We induct on $c$. At $c=0$ the space is $K$ on the empty tuple; choose one label $x$, make $\pi_x=I$ and all others zero.

Let $c\ge1$. Choose $\xi\in L$ and put $L'=L\setminus\{\xi\}$. For $y\in \mathcal Z_c(L)$ and position $i$, let $y^{(i)}$ be the coefficients of those tuples in which $\xi$ appears in position $i$, with that position then removed. This lies in $\mathcal Z_{[c]\setminus\{i\}}(L')$. To see this, fix $\xi$ in position $i$ and sum another position: it is a marginal of $y$, hence zero. There cannot be a second occurrence of $\xi$ in an injective tuple.

By induction, partition $y^{(i)}$ on $L'$ as $\sum_{x\in L'}\pi'_x(y^{(i)})$, each term avoiding $x$. For a vector $v\in \mathcal Z_{[c]\setminus\{i\}}(L'\setminus\{x\})$, partition once more on $L'\setminus\{x\}$:
\[v=\sum_{x'\in L'\setminus\{x\}}\pi''_{x'}(v).
\]
Both applications are justified: $|L'|\ge2c$ and $|L'\setminus\{x\}|\ge2c-1=2(c-1)+1$.

Construct $F_{i,x}(v)$ by inserting $\xi$ into position $i$ of $v$, then subtracting the sum obtained by inserting $x'$ into position $i$ of $\pi''_{x'}(v)$. Every term avoids $x$. All tuples are injective because the lower-dimensional vector avoids the inserted label. A marginal in a position other than $i$ vanishes term by term. The marginal in position $i$ is $v-\sum_{x'}\pi''_{x'}v=0$. Thus
\[F_{i,x}(v)\in \mathcal Z_c(L\setminus\{x\}).
\]
Its $\xi$-part in position $i$ is $v$, and its $\xi$-parts in other positions are zero.

For $x\in L'$ set
\[\pi_x(y)=\sum_{i=1}^cF_{i,x}\bigl(\pi'_x(y^{(i)})\bigr),\qquad
\pi_\xi(y)=y-\sum_{x\in L'}\pi_x(y).
\]
The $\xi$-part of the last expression in every position is $y^{(i)}-\sum_x\pi'_x(y^{(i)})=0$. It therefore uses no label $\xi$, and belongs to $\mathcal Z_c(L')$. The sum identity is built into the definition.

For equivariance, use the same inductively constructed maps on every position set of size $c-1$, transported along a bijection. The transport is independent of the chosen bijection because the inductive maps commute with position permutations. Permuting the $c$ positions simply permutes the terms indexed by $i$ in the displayed sum. No averaging over a symmetric group is used.
\end{proof}

The count $2c+1$ is now visible in an actual induction: reducing one position removes two labels, the label $\xi$ being cleared and the label $x$ being avoided. This is a sufficient bound proved by the construction, not an appeal to a dimension formula.

\subsection{Why the two recursive partitions have different jobs}
In the proof just given, the first partition decomposes the slice
where the distinguished label occurs in position $i$. It chooses
which label $x$ the completed vector will avoid. The second partition
then chooses a label $x'$ to insert in the correction replacing the
distinguished label. Avoiding $x'$ is necessary for injectivity;
avoiding $x$ is necessary for the final destination of $\pi_x$.
One partition cannot simply be reused as the other, since the two
labels are excluded for different reasons.

The residual component $\pi_\xi(y)$ is not justified only by the sum
identity. Its slice containing $\xi$ must vanish at every position.
The inserted terms were designed so their $\xi$-slice at position $i$
is exactly the first partition of $y^{(i)}$, and their $\xi$-slices at
all other positions are zero. This proves that the residual really
avoids $\xi$. Subtracting zero-marginal vectors then keeps all its
marginals zero. These two checks establish the stated codomain of
$\pi_\xi$, as well as the identity $\sum_x\pi_x=I$.

The distinguished label called $\xi$ in that induction is local
bookkeeping; it is unrelated to the proof-degree budget. All the
requirements in the induction count retained labels. To avoid any
confusion, the subsequent applications describe this chosen label
as the distinguished label rather than as a degree parameter.

\section{Adding positions coherently}\label{tx:sec:cone}
Writing $\partial_J$ means forgetting every position in $J$, by
successive marginal maps. Their order does not matter: they are sums
over distinct coordinates. Inserting an extra label and then forgetting
it resembles the elementary combinatorial operation of adding a vertex
to a configuration, often called coning. Here the inserted label must
be avoided by the old tuple, so the partition just proved supplies a
linear sum of admissible insertions. The construction is a section of
a specified marginal map; it does not assume a topological contraction
theorem.

Let $R_0$ and $J$ be disjoint position sets of sizes $c$ and $\ell$. Here $\ell$ counts positions being added; the target-index support parameter $t$ is not involved in this layer construction. We want to add the positions in $J$ while keeping the marginals in $R_0$ zero. If $|L|\ge2c+\ell$, the map deleting all positions in $J$ has a position-$R_0$-equivariant section
\begin{equation}\label{tx:eq:cone}
\mathcal Z_{R_0}(L)\longrightarrow
\{\Phi\in P_{R_0\cup J}(L):\partial_i\Phi=0\ (i\in R_0)\}.
\end{equation}
Being a section means that deleting all the new positions $J$
after this insertion returns the original vector. The equivariance
required here concerns the old positions $R_0$, not arbitrary
permutations of the label set.
Here is the construction. Suppose some positions $J'\subsetneq J$ have already been added. Split the vector according to their injective labelling $\gamma$. On that component the remaining zero-marginal vector uses the label set $L\setminus\im\gamma$. Its size is at least
\[2c+\ell-|J'|\ge2c+1.
\]
Partition that vector by Lemma~\ref{tx:lem:unity}. For a new position $j$, insert label $x$ in the summand $\pi_x$. Summing the new position returns the original vector, and the old zero marginals stay zero. Repeat until all $\ell$ positions have been added. Every step is linear and commutes with the $R_0$ position action, proving \eqref{tx:eq:cone}.

\subsection{Sorting a zero-marginal vector by its block positions}
Split the universe as retained labels $[n]$ and block labels $B$. Every vector $f\in P_b(n+q)$ has components $f_D$, for subsets $D\subseteq[b]$, consisting of tuples whose block positions are exactly $D$:
\[f_D\in K[\Inj(D,B)]\otimes P_{[b]\setminus D}(n).
\]
Write $\epsilon_i$ for summing a block entry and $\partial_i^{\rm ret}$ for summing a retained entry. The condition of zero marginal in position $i\notin D$ is
\begin{equation}\label{tx:eq:cycleeq}
\partial_i^{\rm ret}f_D+\epsilon_i f_{D\cup\{i\}}=0.
\end{equation}
It simply separates the possible value of the deleted entry into a retained label or a block label.

Let $\mathcal Z^{\le \ell}$ be the zero-marginal vectors with no components using more than $\ell$ block positions. The first space is $\mathcal Z^{\le0}=\mathcal Z_b(n)$. The top components of a vector in $\mathcal Z^{\le \ell}$ have zero retained marginals by \eqref{tx:eq:cycleeq}. We obtain a map onto the prospective layer
\[G_\ell=\bigoplus_{|D|=\ell}K[\Inj(D,B)]\otimes \mathcal Z_{[b]\setminus D}(n).
\]
Its kernel is $\mathcal Z^{\le \ell-1}$. We next prove surjectivity with a position-compatible section, so this is an actual quotient description rather than a count of dimensions.

A \emph{filtration} is simply this increasing sequence of subspaces.
Its layer at $\ell$ records what appears after allowing $\ell$ block
positions, modulo everything already allowed at $\ell-1$. A quotient
layer is not yet a chosen subspace of the larger space. The section
below will choose representatives, and the subsequent operator argument
will make that choice compatible with $z$. Here $G_\ell$ names a
space of tuple components, not a target polynomial $G_n$.

In the displayed tensor product, the first factor records which block
labels occupy the positions in $D$, while the second carries the
retained-tuple coefficients. An elementary tensor $e_\beta\otimes y$
places the retained vector $y$ at the block pattern $\beta$. General
vectors are sums of these tensors. The direct sum keeps different
sets $D$ of block positions in separate components.

\subsection{An explicit lift of one top component}
Take $J=[\ell]$, an injective block tuple $\beta:J\to B$, and $y\in \mathcal Z_{R_0}(n)$ with $R_0=[b]\setminus J$. If
\[n\ge2(b-\ell)+\ell=2b-\ell,
\]
the section \eqref{tx:eq:cone} gives $\Phi$ on the $b$ retained positions such that $\partial_J\Phi=y$ and $\partial_i\Phi=0$ for $i\in R_0$. For $J'\subseteq J$, set
\[\phi_{J'}=(-1)^{|J'|}\partial_{J\setminus J'}\Phi.
\]
For $D\subseteq J$ define
\[f_D=e_{\beta|D}\otimes\phi_{J\setminus D},
\]
and put the other components zero. At $D=J$ this is $e_\beta\otimes y$, the required top component. For $i\in R_0$, the retained marginal vanishes. For $i\in J\setminus D$,
\[\partial_i\phi_{J\setminus D}=-\phi_{J\setminus(D\cup\{i\})},
\]
which cancels the block marginal of the adjacent component in \eqref{tx:eq:cycleeq}. Thus the whole lift belongs to $\mathcal Z^{\le \ell}$.

To lift all components equivariantly, choose an ordering of each $\ell$-element subset of block labels. Use it in the representative positions $[\ell]$. Its stabiliser among position permutations consists exactly of permutations of the remaining $b-\ell$ positions. The constructed lift commutes with this stabiliser, by the cone construction. Transport it along position permutations to every other placement of the same ordered block labels. If two permutations give the same placement, their difference is in the stabiliser, so the transported lift is the same. Extend linearly over all block subsets and retained vectors. This supplies the promised $S_b$-equivariant section without division by any group order.

For $1\le \ell\le b$, the largest requirement $2b-\ell$ occurs at $\ell=1$. Therefore the single hypothesis
\begin{equation}\label{tx:eq:linearlabel}n\ge2b-1\end{equation}
supplies all these sections.

\subsection{Two worked lifts at tuple length three}
The signs in the lift formula can be checked without abstract
notation. First take one block position, a block label $\beta$, and
\[
 y=e_{12}+e_{23}+e_{31}-e_{21}-e_{32}-e_{13}
\]
on the other two positions. The vector $y$ uses only retained labels
$1,2,3$, so the particular vector
\[
                    e_\beta\otimes y-e_4\otimes y
\]
is injective and has every marginal zero. The first marginal is
$y-y=0$; the two other marginals are zero because those of $y$ are
zero. This is a single-vector instance of the one-block-position
construction. Constructing a linear map for \emph{all} $y$ is the
harder task solved by the partition of unity, and it is what needs
$n\ge5$ at $b=3$.

For two block positions let the retained position be first and use
$y=e_1-e_2$. With retained labels $1,2,3,4$, the cone can use
\[
                         \Phi=e_{134}-e_{234}.
\]
It has zero first marginal, and deleting the last two positions
returns $y$. For distinct block labels $\beta,\gamma$, the signed
lift is
\begin{align*}
 f={}&e_{1\beta\gamma}-e_{2\beta\gamma}
       -e_{13\gamma}+e_{23\gamma}
       -e_{1\beta4}+e_{2\beta4}
       +e_{134}-e_{234}.
\end{align*}
Deleting position one cancels each adjacent pair. Deleting position
two cancels the terms with $\beta$ against those with $3$, and
those with both retained values against the remaining terms.
Deleting position three gives the analogous cancellations between
$\gamma$ and $4$. Its top component is exactly
$e_{\beta\gamma}\otimes(e_1-e_2)$. The count for this layer is
$2b-\ell=4$; the one-block layer is the more expensive layer, requiring
five retained labels. This is why a maximum over the layers, rather
than a sum, controls the completed construction.

\section{Making the splitting compatible with the block operator}\label{tx:sec:blocksplitting}
The sections just constructed commute with position permutations, but need not commute with $z$. We now fix that, without adding labels.

The grading used here sorts basis patterns into two classes. It is not
polynomial degree. Its diagonal operator $D$ reads the chosen orbit
coordinate, whereas $z$ moves patterns by translations. Their
\emph{commutator} $[D,z]=Dz-zD$ measures the difference between doing
these two operations in opposite orders. The following calculation
will turn that difference into an invertible operator $W$.

A block pattern is an injective partial tuple $\beta$ with a specified nonempty position domain of size $\ell$, and values in $B$. Its image is a nonempty subset of $B$ of size $\ell<q$. Such a subset cannot be fixed by all block translations: a nonempty translation-invariant subset of the regular block is the whole block. Choose the orbit coordinate $\pi$ from Section~\ref{tx:sec:block2} on these \emph{image subsets}, and define
\[
                 D e_\beta=\pi(\operatorname{im}\beta)e_\beta.
\]
Let $E$ be the vector space on all the current block patterns, and let $E_0,E_1$ be the spans of the patterns with coordinate $0,1$, respectively. Thus $E=E_0\oplus E_1$ before any use of $z$. They are stable under position permutations and split by the position domain, since neither operation changes the image subset. On each subset orbit, let $\lambda:A\to\Ftwo$ be the nonzero linear functional vanishing on its stabiliser used to define $\pi$. The previous calculation gives
\[
 [D,z]=W=\sum_i c_i\lambda(\varepsilon_i)S_i,
 \qquad W^2=\Bigl(\sum_i c_i\lambda(\varepsilon_i)\Bigr)^2I.
\]
The scalar is nonzero by independence of the block coefficients, so $W$ is invertible on the whole pattern space.

For either $e=0,1$, suppose $v_0+zv_1=0$ with $v_0,v_1\in E_e$. Apply $D-eI$. The first term vanishes and the second gives $Wv_1$, so $v_1=0$, and then $v_0=0$. Thus the map $(v_0,v_1)\mapsto v_0+zv_1$ from $E_e\oplus E_e$ is injective, giving $2\dim E_e\le\dim E$. Since $E=E_0\oplus E_1$, both inequalities must be equalities. In particular
\[E=E_0\oplus zE_0.
\]
Tensor each domain component with its retained zero-marginal factor. Because $z$ fixes retained labels, the layer has the analogous decomposition
\[G_\ell=Q_\ell\oplus zQ_\ell
\]
with $Q_\ell$ position-stable and $z:Q_\ell\to zQ_\ell$ injective.

\subsection{Turning a section into a commuting section}
A section of $f:E\to Q$ may choose good representatives but fail to
respect $z$. The two-level decomposition supplies a way to repair it:
choose the representative of $q_0\in Q_0$, and require the
representative of $zq_0$ to be its image under $z$. Injectivity of
$z$ on $Q_0$ makes the two inputs uniquely identifiable. The next
lemma checks that this prescription is linear and compatible with
positions as well. Once a section $\tau$ is available, the splitting
is explicit:
\[
             v=\bigl(v-\tau f(v)\bigr)+\tau f(v),
       \qquad v-\tau f(v)\in\ker f.
\]
When $f$ and $\tau$ commute with $z$, both subspaces are stable under $z$.

\begin{lemma}[A section commuting with a square-zero operator]\label{tx:lem:commuting-section}
Let $E,Q$ carry square-zero operators $z$ commuting with position permutations, and let $f:E\to Q$ be a surjection commuting with both actions. Suppose $Q_0$ is position-stable,
\[
                Q=Q_0\oplus zQ_0,
                \qquad z:Q_0\longrightarrow zQ_0\text{ is injective}.
\]
If $\sigma$ is a position-equivariant section of $f$, then a section commuting with both actions is given by
\[\tau(q_0+zq_1)=\sigma(q_0)+z\sigma(q_1).
\]
\end{lemma}
\begin{proof}
The direct sum and injectivity make both $q_0,q_1$ unique. We have $f\tau=I$. Moreover
\[z\tau(q_0+zq_1)=z\sigma(q_0)=\tau(zq_0),
\]
using $z^2=0$. Stability of $Q_0$ and commutation of $z,\sigma$ with positions show that $\tau$ also commutes with positions.
\end{proof}
The construction lifts the first member of each pair $q_0,zq_0$ freely and obtains the second by applying $z$. Apply the lemma to each quotient
$\mathcal Z^{\le \ell}\to G_\ell$. Its kernel is
$\mathcal Z^{\le \ell-1}$, and the commuting section makes the next
space the direct sum of that kernel and the image of the section.
That image is $\tau(Q_\ell)\oplus z\tau(Q_\ell)$, since the quotient map
is injective on the section image. Inducting on $\ell$ gives the direct
sum of the retained layer and these two-level summands. Let $C$ be
the direct sum of all $\tau(Q_\ell)$ for $\ell\ge1$. Every summand is
position-stable and the map $z:C\to zC$ is injective. This proves
the splitting of the filtration as actual subspaces, not merely
an equality of their dimensions. We have proved
\begin{equation}\label{tx:eq:socsplit2}
\mathcal Z_b(n+q)=\mathcal Z_b(n)\oplus C\oplus zC,
\end{equation}
where $C$ is position-stable, $z:C\to zC$ is an isomorphism, and $z$ kills the retained summand. This is a concrete direct-sum decomposition, valid under $b<q$ and $n\ge2b-1$.

\subsection{The four consequences used for correction}
Write $V=P_b(N)$. The next four properties are precisely the interface used in the
correction. They do not assert that every invariant subspace of $V$
is free for $z$.
Recall $V'=\sum_i\im\partial_i^{\mathsf T}$,
$\mathcal Z_b=(V')^\perp$ and $T_b=V/V'$. Their dual pairing is
$\langle[v],\zeta\rangle=\langle v,\zeta\rangle$ for
$\zeta\in\mathcal Z_b$. It is independent of the representative
$v$ because every vector in $V'$ is orthogonal to $\mathcal Z_b$.
This is the pairing used in the transpose argument below.
It is a perfect pairing: it identifies $T_b$ with the space of
linear functionals on $\mathcal Z_b$. Nondegeneracy follows from
$((V')^\perp)^\perp=V'$ in the finite-dimensional tuple pairing.
The notation $z^{-1}(U)$ in the proof means the \emph{inverse image}
$\{v:zv\in U\}$, not an inverse for the nilpotent operator.
For subspaces $U_1,U_2$ of the tuple ambient, the orthogonal-complement
calculations use the elementary rules
\[
 (U_1+U_2)^\perp=U_1^\perp\cap U_2^\perp,\qquad
 (U_1\cap U_2)^\perp=U_1^\perp+U_2^\perp,\qquad
 (zU_1)^\perp=(z^{\mathsf T})^{-1}(U_1^\perp).
\]
The last identity is exactly the transpose pairing formula. These
rules apply over our field without positivity or semisimplicity.

\begin{lemma}[The layer correction interface]\label{tx:lem:layer-interface}
Under $b<q$ and $n\ge2b-1$, the following hold.
\begin{enumerate}
\item $zV\cap \mathcal Z_b(N)\subseteq z\mathcal Z_b(N)$.
\item $zV\cap V'\subseteq zV'$.
\item There is a position-equivariant $w_C$ on $\mathcal Z_b(N)$ with $zw_Cz=z$ there.
\item There is a position-equivariant $w_T$ on $T_b(N)$ with $zw_Tz=z$ there.
\end{enumerate}
\end{lemma}
\begin{proof}
For the third assertion, let $w_C$ send $zc$ back to $c$ and kill $C$ and the retained summand. For the first, the kernel of $z$ in \eqref{tx:eq:socsplit2} is $\mathcal Z_b(n)\oplus zC$. An element of $zV\cap \mathcal Z_b(N)$ is therefore $s+zc$. But $zV$ consists of vectors on tuples meeting the block; so does $zc$. The remaining $s$ is supported entirely on retained tuples, forcing $s=0$.

For the other two assertions use transpose. In characteristic two each $\sigma_i$ is an involution, so $z^{\mathsf T}=z$. The pairing identifies $T_b(N)$ with $\mathcal Z_b(N)^*$, giving $w_T$ by duality. To see the second inclusion explicitly, take orthogonal complements. The first inclusion implies
\[z^{-1}(\mathcal Z_b)=\ker z+\mathcal Z_b.
\]
Indeed, if $zv\in \mathcal Z_b$, it equals $zc$ for some $c\in \mathcal Z_b$, so $v-c\in\ker z$. But
\[(zV')^\perp=z^{-1}(\mathcal Z_b),\qquad
(zV\cap V')^\perp=\ker z+\mathcal Z_b.
\]
The complements agree, hence so do the spaces. This proves the required inclusion, with no appeal to semisimplicity.
\end{proof}

\section{Correcting the error on one layer}\label{tx:sec:layerstep}
The correction has a reusable algebraic form. Write $\pi:V\to T_b$ for the quotient and $\iota_{\mathcal Z}:\mathcal Z_b\to V$ for inclusion. Since both spaces are stable under labels, $z$ induces operators $z_T,z_{\mathcal Z}$ satisfying
\[
       \pi z=z_T\pi,\qquad z\iota_{\mathcal Z}=\iota_{\mathcal Z}z_{\mathcal Z}.
\]
The distinction between these three domains keeps the construction independent of any choice of representatives.
The maps $w_C$ and $w_T$ below are auxiliary inverses on these two
spaces, not yet a common group-algebra witness. They construct one
correction matrix; realisation by permutations is the separate task
of the next chapter.

The auxiliary maps reverse only the relevant arrows in nilpotent
chains. On a pair $c\mapsto zc\mapsto0$, for example, $w_C$ sends
$zc$ back to $c$; it is not an inverse of $z$ on the whole space.
The correction will control its output with $w_C$ on $\mathcal Z_b$
and its input with $w_T$ on $T_b$. Composition is read from right
to left through these two different spaces. The two inverse
conditions in the next lemma are what make the defect cancel;
we do not invert the residual $\bar\varepsilon$.

\begin{lemma}[The two-sided correction contract]\label{tx:lem:two-sided-correction}
Let $1\le b\le N$, and suppose $w_0\in K[S_N]$ satisfies $zw_0z=z$ on every tuple length below $b$. On $V=P_b(N)$ write its residual as
\[
 \varepsilon=zw_0z-z=\iota_{\mathcal Z}\bar\varepsilon\pi,
 \qquad \bar\varepsilon:T_b\longrightarrow\mathcal Z_b.
\]
Suppose position-equivariant maps $w_C:\mathcal Z_b\to\mathcal Z_b$ and $w_T:T_b\to T_b$ satisfy the two inverse conditions
\begin{equation}\label{tx:eq:two-sided-inverses}
 z_{\mathcal Z}w_C\bar\varepsilon=\bar\varepsilon,
 \qquad \bar\varepsilon w_Tz_T=\bar\varepsilon.
\end{equation}
Then the position-equivariant correction
\[
 \bar\varphi=-w_C\bar\varepsilon w_T:T_b\to\mathcal Z_b,
 \qquad \varphi=\iota_{\mathcal Z}\bar\varphi\pi:V\to V
\]
has image in $\mathcal Z_b$, kills $V'$, and satisfies $z\varphi z=-\varepsilon$. If an element $j\in J_{b-1}(N)$ realises $\varphi$, then $w_0+j$ is a witness on every length at most $b$.
\end{lemma}
\begin{proof}
The factorisation through quotient and inclusion gives the two shape conditions. The residual $\bar\varepsilon$ commutes with positions because the label group-algebra action does, and hence so does $\varphi$. The inverse conditions give the fully typed calculation
\[
 z\varphi z
 =-\iota_{\mathcal Z}z_{\mathcal Z}w_C\bar\varepsilon w_Tz_T\pi
 =-\iota_{\mathcal Z}\bar\varepsilon\pi
 =-\varepsilon.
\]
Thus $z(w_0+j)z-z=\varepsilon+zjz=0$ on $V$. On each shorter tuple module $j$ is zero, so its addition preserves the old witness identity.
\end{proof}

\subsection{Verifying the two inverse conditions}
Return to the residual constructed at the start of the chapter. It kills $V'$ and has image in $\mathcal Z_b$, so its factorisation through $\bar\varepsilon$ is available. Its image also lies in $zV$. The first assertion of Lemma~\ref{tx:lem:layer-interface} therefore gives $\im\varepsilon\subseteq z\mathcal Z_b$, and the third gives
\[
                       z_{\mathcal Z}w_C\bar\varepsilon=\bar\varepsilon.
\]
Here we use the surjectivity of $\pi$ and injectivity of
$\iota_{\mathcal Z}$ to express the identity on the indicated domains.
This is the first condition of \eqref{tx:eq:two-sided-inverses}.

For the other side, if $zv\in V'$, the second interface assertion gives $zv=zv'$ for some $v'\in V'$. Since $\varepsilon=(zw_0-I)z$,
\[
       \varepsilon v=(zw_0-I)zv'=\varepsilon v'=0.
\]
Now $z_Tw_Tz_T=z_T$ by the fourth assertion. For any representative $v_T\in V$ of $w_Tz_T\pi(v)$, it follows that $z(v_T-v)\in V'$. The implication just proved gives $\varepsilon v_T=\varepsilon v$, or
\[
               \bar\varepsilon w_Tz_T\pi=\bar\varepsilon\pi.
\]
Surjectivity of $\pi$ gives the second inverse condition. This calculation needs no section $T_b\to V$: the value of $\varepsilon$ is independent of representatives because it kills $V'$.

The contract now supplies the correction. Theorem~\ref{tx:thm:realization}, proved in the next chapter, realises its shape by an actual $j\in J_{b-1}(N)$ using the existing labels. This is the step that replaces the fresh-label compression of the quadratic construction.

\begin{theorem}[The characteristic-two linear witness]\label{tx:thm:witness2}
Let $q=2^a>r$, and let $n\ge\max\{0,2r-1\}$. With the independent
block coefficients specified in Chapter~\ref{tx:ch:square-zero}, one
common element $w\in K[S_{n+q}]$ satisfies $zwz=z$ on every
$P_b(n+q)$, $0\le b\le r$. Deletion is therefore exact on
$\mathcal C_r(n+q)$, with the natural basis and tensor maps proved
there.
\end{theorem}
\begin{proof}
At length zero take $w_0=0$. Suppose one element $w_{b-1}$ works
on lengths $0,\ldots,b-1$. Its residual on length $b$ has the
shape proved at the beginning of this chapter. The layer interface
verifies the two inverse conditions, and Lemma~\ref{tx:lem:two-sided-correction}
constructs $\varphi$ with $z\varphi z=-\varepsilon$. Theorem~\ref{tx:thm:realization}, whose
complete proof is in the next chapter, realises it by
$j\in J_{b-1}(n+q)$. Set $w_b=w_{b-1}+j$. Then
\[
 z w_b z-z=\varepsilon+zjz=0
\]
on length $b$. On every shorter length $j$ is zero, so the old
witness identity is unchanged. Every step uses only the same labels;
$n\ge2r-1$ implies $n\ge2b-1$ for all $b\le r$. Induction gives
$w=w_r$. The common-projector proof of Chapter~\ref{tx:ch:square-zero} gives
exactness, while its orbit calculation gives the natural maps.
\end{proof}

\section*{Why the cost is a maximum}
The filtration sections at length $b$ require $2b-\ell$ retained labels
for its layer with $\ell$ block positions. All such layers use the same
retained set, so their total requirement is $2b-1$. The group-algebra
realisation then uses that same board without adding labels. At the
next tuple length, the new residual is again special and the same
pool can be used. The final requirement is therefore
\[
                 \max_{1\le b\le r}(2b-1)=2r-1,
\]
not $\sum_b b$. The entire improvement is a change in how labels
are allocated to correction maps. It changes the onset for a fixed
degree budget from quadratic to linear. The later consequence
classification inverts that onset to obtain the degree lower bound.

\chapter{Realising the lift's corrections by permutations}\label{tx:sec:rook-realization}
\label{tx:ch:rooks}
\section*{The property we need}
The previous chapter constructed a correction with zero marginals in
its outputs and zero incoming shorter layers. To repair a common
witness, that correction must be an actual linear combination of
label permutations. A different matrix correction for each tuple
length would not be sufficient. We now prove the required realisation
statement from finite combinatorics, in this main chapter.

Recall the spaces on which that correction acts:
$V=P_b(N)$ is the ordered injective-$b$-tuple space,
$V'=\sum_i\im\partial_i^{\mathsf T}$ contains the incoming
shorter-tuple directions, and $\mathcal Z_b=\bigcap_i\ker\partial_i$
is the zero-marginal space. We want a single element of
$J_{b-1}(N)$, the ideal of permutation combinations acting as zero
on every tuple length below $b$, whose action on $V$ is the required
correction.

The proof has two parts. First, values on a carefully selected set
of rook placements determine any zero-line-sum function. Second,
those selected values can be prescribed by triangular elimination
in permutation coefficients. The triangular system has diagonal
coefficient one, so it can be solved by subtraction over any
commutative ring. It needs no averaging, semisimplicity or division
by a factorial.

The order of the board is a bookkeeping tool for this construction.
It is not a built-in order in the uniform polynomial input. We are
free to order the existing finite labels while constructing a
group-algebra element; the final element acts on the original
uninterpreted labels. No label is added.

\section{Rook coordinates and realisation}\label{tx:sec:rooks}
The representation-theoretic term for maps commuting with every
permutation of tuple positions is their \emph{centraliser algebra}.
Addition and composition preserve that condition. Label permutations
commute with position permutations, so their linear combinations
belong to this algebra. The realisation problem asks whether the
particular correction matrix can be obtained from such a combination,
with its required zero action on shorter tuples as well.
Commutation alone does not supply the coefficients; the following
coordinate construction does.

The correction contract in Lemma~\ref{tx:lem:two-sided-correction} leaves one finite spanning problem: realise the correction's whole matrix by permutations. Rook coordinates express its two shape conditions as line sums. The argument is independent of the block and works over any commutative ring.

\subsection{From a map to a board of non-attacking rooks}
A matrix entry of a position-equivariant map $\varphi:P_b(N)\to P_b(N)$ is indexed by an output tuple $s$ and input tuple $t$. Permuting their positions together leaves the entry unchanged. Thus it depends only on
\[\rho=\{(s_i,t_i):i=1,\ldots,b\},
\]
a set of $b$ cells with distinct rows and distinct columns on an $N\times N$ board. Call this a $b$-rook. Board rows and columns are \emph{labels}; they are different from the positions of the original tuples.

For instance, output $(3,4)$ and input $(1,2)$ give the cells
$\{(3,1),(4,2)\}$. Reordering both tuples to $(4,3)$ and $(2,1)$
gives the same cells, so position equivariance gives the same matrix
entry. Reordering only one tuple would give different cells. Passing
to rook coordinates forgets the simultaneous ordering, while retaining
which input label is paired with which output label.

The matrix therefore defines a function $F$ on $b$-rooks. For a $(b-1)$-rook $\sigma$, choose an unused column $y$. A column line through $\sigma$ is the family of rooks $\sigma\cup\{(x,y)\}$ as $x$ ranges over unused rows. Define row lines analogously. The map has image in $\mathcal Z_b$ precisely when its column lines sum to zero. It kills $V'$ precisely when its row lines sum to zero. Write $\mathcal R_b^{0}$ for this zero-line-sum space of rook functions.

A label permutation $g$ has graph $\{(g(y),y)\}$. Its matrix entry is one exactly when $\rho$ is contained in that graph. Set
\[\chi_g(\rho)=[\rho\subseteq\operatorname{graph}g],\qquad
L_b\Bigl(\sum_ga_gg\Bigr)=\sum_ga_g\chi_g
\]
on $b$-rooks. For a fixed unused row or column, exactly one extension of a rook $\sigma$ lies in the graph of $g$ if $\sigma$ itself does, and none otherwise. Therefore
\begin{equation}\label{tx:eq:lines}
\text{each line sum of }L_b(x)\text{ through }\sigma=L_{b-1}(x)(\sigma).
\end{equation}
If $F\in \mathcal R_b^{0}$ is realised as $L_b(x)$, this identity forces $L_{b-1}(x)=0$, then every smaller $L_j(x)=0$. Thus $x\in J_{b-1}(N)$ automatically. It remains to prove that every $F\in \mathcal R_b^{0}$ is realisable.

Here an indicator $[\rho\subseteq\operatorname{graph}g]$ is the
scalar $1$ or $0$ according to whether the displayed containment
holds. The map $L_b$ reads the matrix of a permutation combination
in these coordinates; it is unrelated to the lower-degree polynomial
space $L_d(n)$. A line sum varies one cell while the other $b-1$
cells are held fixed. It is the rook-coordinate version of a marginal.

\subsection{A first example: matrices with zero row and column sums}
At $b=1$, $\mathcal R_1^{0}$ is the space of square matrices with zero row and column sums. If row and column $N$ are designated as last, such a matrix is determined by its upper-left $(N-1)\times(N-1)$ entries and is
\[\sum_{i,j<N}F_{ij}(E_{ij}-E_{iN}-E_{Nj}+E_{NN}).
\]
Each four-entry matrix is a difference of two permutation matrices: choose a permutation taking column $j$ to row $i$ and column $N$ to row $N$, and subtract the permutation with these two outputs swapped. The other entries cancel. This proves realisation at length one over every ring. The general rook argument is a higher-length replacement for this elementary construction.

\subsection{Length two: a four-term correction}
Take four labels and disjoint swaps $\tau_1=(13)$ and $\tau_2=(24)$.
The group-algebra element
\[
                  j=(1-\tau_1)(1-\tau_2)
                    =1-\tau_1-\tau_2+\tau_1\tau_2
\]
annihilates every one-tuple. A one-tuple omits at least one whole
pair; the corresponding swap fixes it, and the disjoint factors
commute. On the two-tuple $(1,2)$ it gives
\[
                       e_{12}-e_{14}-e_{32}+e_{34}.
\]
Deleting either position makes this vector zero. More generally,
its matrix has zero output and input line sums because $j$ kills
one-tuples and commutes with the marginal maps. This is a concrete
correction in $J_1(4)$.

\subsection{Length three: an eight-term correction}
On six labels take $\tau_1=(14)$, $\tau_2=(25)$ and $\tau_3=(36)$.
Then $j=(1-\tau_1)(1-\tau_2)(1-\tau_3)$ annihilates every tuple
of length at most two: such a tuple leaves one whole pair unused.
Its action on $(1,2,3)$ is
\begin{align*}
 j e_{123}={}&e_{123}-e_{423}-e_{153}-e_{126}
                  +e_{453}+e_{426}+e_{156}-e_{456}.
\end{align*}
Each marginal pairs terms differing only in the deleted coordinate,
with opposite coefficients. In characteristic two the signs coincide
and the paired terms still cancel. Thus $j\in J_2(6)$ provides a
literal length-three example of the shape required by the correction
contract.

These products illustrate some realisable corrections. They do not
prove that \emph{every} correction is realisable, and their use of
$2b$ displayed labels is not a new bound on the retained-label pool.
The forcing theorem below proves the spanning statement on the
existing $N\times N$ board, for arbitrary zero-line-sum data.

\section{A forcing set of coordinates}\label{tx:sec:forcing}
Work on any two equally large finite sets of integer labels $R,C$,
of size $M$, with their inherited orders. In this combinatorial proof
$M$ is the board size and $c$ is the number of rooks; the application
later sets $M=N$ and $c=b$. Numerical thresholds such as
$y+1$ denote the cut strictly above the integer label $y$. A permutation here is a bijection $g:C\to R$. An increasing chain is a subset of its graph increasing in both coordinates. Let $\LIS(g)$ be the largest chain length. Among the column sets of longest chains choose the lexicographically first one, $I_g$: at the first differing column, the set containing that column is earlier. Define
\[\rho_g=\{(g(y),y):y\notin I_g\}.
\]
If $\LIS(g)=M-c$, this is a $c$-rook. Let $\Psi_c$ consist of these rooks.

The longest increasing subsequence can be read directly from the
list $g(1),\ldots,g(M)$ when the labels are $1,\ldots,M$: choose
increasing column numbers whose listed values also increase.
For the values $(2,1,3)$ there are two length-two choices,
with columns $\{1,3\}$ and $\{2,3\}$. The first is lexicographically
earlier because it contains the first differing column, $1$.
Removing it leaves the cell $(1,2)$ as the canonical complement.
This rule fixes a coordinate even when the longest chain is not unique.

To call $\Psi_c$ a \emph{forcing set} means that a zero-line-sum
function is determined by its values there. If two such functions
agree on $\Psi_c$, their difference vanishes there, and the theorem
below will show that it vanishes everywhere. This is uniqueness under
the line equations, rather than a statement about arbitrary functions
on the board. Existence of permutation coefficients is a separate
task after forcing has been proved.

\begin{theorem}[Forcing]\label{tx:thm:force}
A function in $\mathcal R_c^{0}$ that vanishes on $\Psi_c$ vanishes on every $c$-rook.
\end{theorem}


\begin{example}[The $3\times3$ board]
We write permutations in one-line notation: $312$ means
$g(1)=3$, $g(2)=1$, $g(3)=2$.
The permutations with longest increasing subsequence of length two are $312,231,213,132$. Their canonical complement rooks are respectively
\[(3,1),\ (1,3),\ (1,2),\ (2,3).
\]
If a zero-line-sum matrix is zero on these four cells, row one first forces $(1,1)$ to be zero. Column one then forces $(2,1)$; row two forces $(2,2)$; column two forces $(3,2)$; and row three forces $(3,3)$. This is the small model of the forcing argument.
\end{example}

\subsection{Why a longest chain determines a useful coordinate}
A permutation is a complete placement of $M$ rooks. Removing a longest
increasing chain leaves its canonical complement $\rho_g$. The chain
can be reconstructed once that complement is fixed: the remaining
columns and rows must be paired in increasing order. The forcing
coordinates therefore encode a rigid increasing completion, while
still having the same number of cells as the corrections we want to
realise.

The lexicographic choice controls ties when a corner is added. The two contractions in Example~\ref{tx:eg:corner-contractions} show why a minimum-row corner can be allowed a tie, whereas a minimum-column corner must be excluded from every longest chain. To keep that information through induction, we pair a chain bound above a cut with a vanishing condition based on unused labels below it.

\subsection{Strengthened induction statement}
To remove the bottom row and the first column, we must know which canonical longest chains survive contraction. Bounds on longest chains in upper parts of the board keep track of this. These are bookkeeping conditions for the induction, not new mathematical assumptions on $F$.

Let $L=M-c$. A condition is $(s,\theta,\eta)$, where $s$ is ``row'' or ``column'', $\theta$ is a label threshold, and $\eta\ge1$. A bijection satisfies it if the longest increasing chain using only cells with that coordinate at least $\theta$ has length at most $L-\eta$. Bounds are interpreted as integer inequalities; a negative upper bound is impossible for a nonnegative chain length.

For a collection $\mathcal E$ of conditions let $\Psi(\mathcal E)$ be the canonical complement rooks of permutations with $\LIS=L$ satisfying every condition. Let $Q(\mathcal E)$ be the set of $c$-rooks for which, for at least one condition, at most $\eta-1$ \emph{unused} labels of the specified side lie below $\theta$.

The two descriptions are paired deliberately. A bound on chains above a threshold will, through a line-sum identity, force a value at a rook with too few unused labels below it.

An induction invariant is the collection of facts carried to each
smaller board. Here $M$ is its size; for a $c$-rook, $c$ is the number
of occupied cells, and $L=M-c$ is the number of unused rows and also of unused
columns. A condition with deficit $\eta$ permits at most $L-\eta$
members of an increasing chain above its cut. The set $\Psi(\mathcal E)$
records canonical coordinates obeying those chain restrictions;
$Q(\mathcal E)$ records additional zeros established through line
equations. These are different kinds of information.
For example, with row labels $1,2,3,4$, the row condition with cut
$3$ and $\eta=2$ puts a one-rook in $Q$ when at most one of rows $1,2$
is unused, that is, when its occupied row is $1$ or $2$.
The extra conditions below are introduced within the induction;
the desired forcing theorem starts with no such conditions.

\begin{lemma}[Strengthened forcing]\label{tx:lem:GP}
If $F\in \mathcal R_c^{0}$ vanishes on $\Psi(\mathcal E)\cup Q(\mathcal E)$, then $F=0$.
\end{lemma}
Theorem~\ref{tx:thm:force} is its case of no conditions.

The \emph{induction invariant} is the following data at the current board:
\[
 L=M-c,\qquad F\in\mathcal R_c^0,\qquad
 F|_{\Psi(\mathcal E)}=0,\qquad F|_{Q(\mathcal E)}=0.
\]
Each use of induction must preserve the three assertions about $F$ with the updated conditions. The first two sweeps replace $(M,c)$ by $(M-1,c-1)$, keeping $L$ fixed; the final restriction replaces it by $(M-1,c)$, decreasing $L$ by one. The conditions record exactly what these contractions change.

\subsection{Three facts about adding a corner cell}
Let $r_0=\min R$ and $y_0=\min C$. The following statements describe the only possible new increasing chains.
\begin{enumerate}
\item Remove row $r_0$ and column $y$, and let $\widetilde g$ be a bijection on the smaller board with longest chain length $L$. Add the cell $(r_0,y)$. Every new chain using it continues with a chain in columns strictly above $y$. If that region had longest chain at most $L-1$, the new permutation has longest chain $L$. Its lexicographically first longest chain is still the old one: the old chain must use some column below $y$, whereas every new chain through the corner uses columns at least $y$. Therefore $\rho_g=\rho_{\widetilde g}\cup\{(r_0,y)\}$. For a column threshold $\theta\le y$, the new regional maximum is
\[\max\{\LIS(\widetilde g|_{\mathrm{col}\ge\theta}),\ 1+\LIS(\widetilde g|_{\mathrm{col}>y})\}.
\]
\item Remove row $x$ and column $y_0$, and add $(x,y_0)$. If the old region of rows strictly above $x$ has longest chain at most $L-2$, every chain using the new corner has length at most $L-1$. Thus all longest chains remain old, the canonical one remains unchanged, and $\rho_g=\rho_{\widetilde g}\cup\{(x,y_0)\}$. The analogous maximum formula holds for row thresholds at most $x$. The stronger deficit two eliminates a tie involving the smallest column.
\item Remove both minimal labels and add $(r_0,y_0)$. This corner can be placed before every increasing chain of the smaller board. The longest length increases by one, every new longest chain contains the corner, and $I_g=\{y_0\}\cup I_{\widetilde g}$. Thus $\rho_g=\rho_{\widetilde g}$.
\end{enumerate}
These facts follow just by splitting chains according to whether they use the added cell. They do not rely on a representation-theoretic result.

\begin{example}[One board, two contractions]\label{tx:eg:corner-contractions}
Take $M=4,c=1$, so $L=3$, and the permutation $2134$. Rows increase upwards and columns to the right. The vertical rule marks the cut at column $3$; the starred cell is the corner $(1,2)$ to be contracted:
\[
\begin{array}{c|cc|cc}
 &1&2&3&4\\\hline
4&\cdot&\cdot&\cdot&\bullet\\
3&\cdot&\cdot&\bullet&\cdot\\
2&\bullet&\cdot&\cdot&\cdot\\\hline
1&\cdot&\star&\cdot&\cdot
\end{array}
\quad\xrightarrow{\text{remove row }1\text{ and column }2}\quad
\begin{array}{c|c|cc}
 &1&3&4\\\hline
4&\cdot&\cdot&\bullet\\
3&\cdot&\bullet&\cdot\\
2&\bullet&\cdot&\cdot
\end{array}
\]
The contracted permutation has the increasing chain with columns $\{1,3,4\}$. Above the cut, its chain length is $2=L-1$. Adding the corner creates another length-three chain with columns $\{2,3,4\}$, but $\{1,3,4\}$ is lexicographically earlier. Thus deficit one suffices to preserve the canonical chain in corner fact~1. The removed column $2$ is absent from the contracted label set; the labels below the cut are now just $\{1\}$. For a rook on that board, the corresponding $Q$ condition counts whether this remaining label is unused.

Instead contract $(2,1)$ from the same full board. The horizontal rule now marks the row cut at $3$:
\[
\begin{array}{c|ccc}
 &2&3&4\\\hline
4&\cdot&\cdot&\bullet\\
3&\cdot&\bullet&\cdot\\\hline
1&\bullet&\cdot&\cdot
\end{array}
\]
The old longest chain uses columns $\{2,3,4\}$, and the region of rows above $2$ again has length $L-1=2$. Adding $(2,1)$ creates the competing chain $\{1,3,4\}$, which is now earlier. This explains the stronger requirement in corner fact~2: imposing regional length at most $L-2$ makes every chain through the added corner have length at most $L-1$, so there can be no tie. The example exhibits why the two sweeps require different deficits.
\end{example}

\subsection{Reading the cuts as integer inequalities}
Row labels and column labels may have gaps; they need not be
$1,\ldots,M$ after a contraction. The notation $y+1$ means the
integer cut strictly above the particular label $y$. On the contracted
column set, which excludes $y$, being below $y+1$ is exactly being
strictly below $y$. This observation justifies the unused-label
counts in both sweeps.

The inequalities are ordinary integer inequalities. In particular,
$\LIS\le L-\eta$ is impossible if $L-\eta<0$; it does not silently
mean a truncated natural-number bound of zero.

\subsection{Base cases and degenerate conditions}
The proof uses two levels of induction. The outer induction reduces
the board size $M$. Within a fixed board, a decreasing sweep remembers
which corner choices have already been eliminated. Contracting a
specified corner means fixing that cell and deleting its row and
column: $G(\sigma)=F(\sigma\cup\{\text{corner}\})$.
Thus $(M,c)$ becomes $(M-1,c-1)$ and $L$ stays unchanged. The final
restriction has no fixed cell and leaves $c$ unchanged, so $L$ drops
by one. Keeping the integer row and column names, including their
gaps, makes the threshold conditions unambiguous.

Prove Lemma~\ref{tx:lem:GP} by induction on $M$. If $c>M$ there are no rooks. If $c=M>0$, each line through an $(M-1)$-rook contains exactly one $M$-rook, so the zero-line-sum equation forces its value to zero.

A condition is degenerate if the entire board has at most $\eta-1$ labels below its threshold. Then every rook belongs to $Q$, and the conclusion is immediate. Assume no condition is degenerate. Every column threshold is therefore above $y_0$ and every row threshold above $r_0$.

If $c=0$, consider the increasing bijection. It has longest length $M$. For a column condition, at least $\eta$ columns lie below $\theta$, so the part above $\theta$ has length at most $M-\eta$; the same holds for rows. Thus its empty complement rook belongs to $\Psi(\mathcal E)$, and the sole value of $F$ is zero. We may now assume $1\le c<M$.

\subsection{First sweep: eliminate every rook meeting the bottom row}
For $y\in C$ in decreasing order, contract the cell $(r_0,y)$ and define
\[G(\sigma)=F(\sigma\cup\{(r_0,y)\})
\]
on $(c-1)$-rooks of the board with that row and column removed. Its line sums are line sums of $F$, so $G\in \mathcal R_{c-1}^{0}$. The smaller board has size $M-1$ and complement length $(M-1)-(c-1)=L$.

Retain the old conditions and add
\begin{align*}
&(\mathrm{column},y+1,1),\\
&(\mathrm{column},y+1,\eta+1)
\quad\text{for each old }(\mathrm{column},\theta,\eta)\text{ with }\theta\le y.
\end{align*}
The first new condition bounds the right region by $L-1$, giving corner fact 1. The second bounds it by $L-(\eta+1)$ whenever needed. An old column condition with $\theta\le y$ is preserved upon adding the corner because
\[\max\{L-\eta,1+L-(\eta+1)\}=L-\eta.
\]
Old column conditions above $y$ and old row conditions are unaffected by the corner. Thus a canonical rook satisfying the new conditions extends to one in $\Psi(\mathcal E)$, and $G$ vanishes on its new $\Psi$ set.

We must also verify every case of its new $Q$ set. An old condition passes directly to the enlarged rook: the deleted row and column are occupied by the added cell, so the count of unused labels below an old threshold is unchanged.

For a new condition $(\mathrm{column},y+1,e)$, let
\[U=\{y'\in C:y'<y,\ y'\text{ unused by }\sigma\}.
\]
On the contracted board this is exactly the unused set below $y+1$, so $|U|\le e-1$. If $e=1$, there are no such columns. If the condition came from an old $(\mathrm{column},\theta,\eta)$, then $e=\eta+1$ and $|U|\le \eta$. For each $y'\in U$, the enlarged rook $\sigma\cup\{(r_0,y')\}$ has at most $\eta-1$ unused columns below $\theta$. Indeed, if $y'<\theta$, its occupation removes one member from a subset of $U$; if $y'\ge\theta$, $y'$ itself lies outside that subset, again leaving at most $|U|-1$. Thus every such rook lies in the old $Q$ set.

Now use the row line through $\sigma$ in row $r_0$. Its terms at columns greater than $y$ are zero by the decreasing inner induction. Its terms below $y$ are either absent ($e=1$) or zero by the old $Q$ condition. The only remaining term is $G(\sigma)$, so it is zero. This proves vanishing on the new $Q$ set in all cases. The outer induction on $M$ now gives $G=0$.

Completing the decreasing sweep shows that $F$ vanishes on every rook meeting row $r_0$.

\subsection{A sweep term with its unused labels displayed}
Here is one instance of the new-$Q$ verification. Use row and column
labels $1,\ldots,7$, take $c=3$ and hence $L=4$, and suppose an old
condition is $(\mathrm{column},3,1)$. In the bottom-row sweep at
$y=4$, contract row $1$ and column $4$. The smaller board has
$(M,c,L)=(6,2,4)$; its column labels are $1,2,3,5,6,7$.
The new condition $(\mathrm{column},5,2)$ comes from that old one.

Choose the smaller rook $\sigma=\{(2,1),(3,3)\}$. Below the new cut
$5$, the contracted board has columns $1,2,3$, of which only $2$
is unused. Thus $\sigma$ is in the new $Q$ set since $1\le2-1$.
The row-line equation on the original board is exactly
\[
 0=F(\sigma\cup\{(1,2)\})+F(\sigma\cup\{(1,4)\})
       +\sum_{y'\in\{5,6,7\}}F(\sigma\cup\{(1,y')\}).
\]
The last three terms were eliminated earlier in the decreasing sweep.
In the first term, both columns below the \emph{old} cut $3$, namely
$1,2$, are occupied. Its unused count is zero, so it lies in the old
$Q$ set and is also zero. The remaining term is the contracted value
$G(\sigma)$, which is therefore zero. This trace shows why unused
labels are counted after contraction, why the deficit increases by one,
and how the new vanishing condition is proved by a line equation.

\subsection{Second sweep: eliminate every rook meeting the first column}
Adjoin the condition $(\mathrm{row},r_0+1,1)$ to obtain $\mathcal E^+$. Its new $Q$ rooks are exactly those using row $r_0$: there must be no unused row below $r_0+1$, and that row is the only candidate. They have just been proved zero. Its $\Psi$ set is a subset of the old one. Thus the vanishing assumptions still hold for $\mathcal E^+$.

Sweep $x\in R\setminus\{r_0\}$ in decreasing order. Contract $(x,y_0)$ and keep every old condition in $\mathcal E^+$. For each row condition $(\mathrm{row},\theta,\eta)$ with $\theta\le x$, add $(\mathrm{row},x+1,\eta+1)$. In particular the recently added condition produces $(\mathrm{row},x+1,2)$, because $x>r_0$. Therefore the region above $x$ has maximum at most $L-2$, exactly the hypothesis of corner fact 2.

The verification for $\Psi$ is now explicit: row thresholds at most $x$ use the same maximum calculation as in the first sweep; row thresholds above $x$ are unchanged; every column threshold exceeds $y_0$, so excludes the new corner. For $Q$, old conditions pass unchanged. A new condition from deficit $\eta$ says that there are at most $\eta$ unused rows below $x$. When any such row $x'$ is occupied instead of $x$, the number unused below its source threshold is at most $\eta-1$, by the same two cases $x'<\theta$ and $x'\ge\theta$. These terms lie in $Q(\mathcal E^+)$. The column-line equation at $y_0$ then forces the remaining term to zero, since larger rows were handled earlier. The induction on $M$ proves the contracted function zero.

The case $x=r_0$ was already covered by the first sweep. Thus $F$ now vanishes on every rook meeting either row $r_0$ or column $y_0$.

\subsection{The remaining smaller board}
Restrict $F$ to $(R\setminus\{r_0\})\times(C\setminus\{y_0\})$ without changing $c$. A line in this board is an old line minus one term using a removed label; that term is zero, so the restriction has zero line sums.

For every condition in $\mathcal E^+$ with deficit $\eta\ge2$, keep its side and threshold and replace its deficit by $\eta-1$. Discard the conditions with deficit one. The smaller complement length is $L-1$.

A permutation on the smaller board satisfying these conditions and having longest length $L-1$ extends by the minimal corner. Corner fact 3 gives longest length $L$ and the same complement rook. Every old threshold excludes that corner. For $\eta\ge2$ the regional bound is
\[(L-1)-(\eta-1)=L-\eta;
\]
for $\eta=1$ it follows from the whole smaller board having longest length $L-1$. So the restricted function vanishes on the new $\Psi$ set.

If a smaller rook belongs to the new $Q$ set, its unused count below a threshold is at most $\eta-2$. In the original board the removed minimum label is also unused by this rook and lies below that threshold, increasing the count by exactly one. The resulting count is at most $\eta-1$, so the rook lies in the old $Q(\mathcal E^+)$. Thus the new $Q$ vanishes too. The induction on $M$ finishes the proof of Lemma~\ref{tx:lem:GP}, and hence of the forcing theorem.


\subsection{What each part of the induction accomplished}
The chain conditions ensure that canonical coordinates survive adding a corner. The unused-label conditions ensure, through a line equation, that the other terms of a sweep have already vanished. Once both sweeps are finished, restriction loses only zero terms. The final deficit change preserves the same induction invariant with $L$ decreased by one. These two records let a local contraction prove a global forcing theorem.

\section{Triangular realisation and the common correction}\label{tx:sec:realization}
The forcing theorem says that the selected values on $\Psi_b$
determine a zero-line-sum function. We now prescribe those values by
permutation coefficients while simultaneously making every smaller
tuple action zero. Recall that $L_k(x)$ records the matrix entries of
a permutation combination $x$ on $k$-rooks.

This second ingredient is simpler. If $\rho_g$ lies on the graph of $g'\ne g$, then $g'$ agrees with $g$ outside $I_g$ and permutes the remaining values $g(I_g)$. On $I_g$ the values of $g$ are increasing. At the first column where $g'$ differs, it must use a larger remaining value. Thus $g'$ is lexicographically larger than $g$.

This is triangular elimination in the ordinary linear-algebra sense.
Equations of the form
\[
                   a_i+\sum_{j>i}t_{ij}a_j=b_i
\]
are solved from the last unknown backwards. For example,
$a_1+a_2=b_1$, $a_2=b_2$ give $a_2=b_2$ and $a_1=b_1-b_2$.
No division is needed because the coefficient of each new unknown
is $1$. In the present system, $t_{ij}$ is the indicator that the
later permutation contains the earlier one's canonical rook. The
prescribed right-hand side is $F$ at length $b$ and zero at every
shorter length. Solving those lower equations is what preserves
the witness identities on all shorter tuple spaces.

Fix $F\in \mathcal R_b^{0}(N)$ and let $\mathcal T=\{g:\LIS(g)\ge N-b\}$. Its canonical complement rooks have sizes at most $b$. Define coefficients $a_g$ in decreasing lexicographic order by
\begin{equation}\label{tx:eq:triangular}
a_g=\begin{cases}F(\rho_g)&|\rho_g|=b,\\0&|\rho_g|<b\end{cases}
-\sum_{\substack{g'\in\mathcal T,\ g'\ne g\\\rho_g\subseteq\operatorname{graph}g'}}a_{g'}.
\end{equation}
Every coefficient on the right was already defined by the preceding observation. Put $x=\sum_{g\in\mathcal T}a_gg$, and let $H_k=L_k(x)$ for each $k\le b$. The construction makes $H_k$ zero on $\Psi_k$ for $k<b$ and makes $H_b=F$ on $\Psi_b$.

At level zero the only coordinate is the empty rook, and it is prescribed zero when $b>0$. Inductively suppose $H_{k-1}=0$. Equation~\eqref{tx:eq:lines} shows that $H_k$ has zero line sums. For $k<b$, it vanishes on $\Psi_k$, so forcing gives $H_k=0$. At $k=b$, $H_b-F$ has zero line sums and vanishes on $\Psi_b$, so forcing gives $H_b=F$. The case $b=0$ is immediate by a scalar multiple of any permutation.

We have proved that every zero-line-sum rook function is realised by a group-algebra element, and the realising element annihilates all smaller tuple modules. Translating back to matrices proves:
\begin{theorem}[Realisation of the layer correction]\label{tx:thm:realization}
For $1\le b\le N$, every position-equivariant map $\varphi:P_b(N)\to P_b(N)$ with image in $\mathcal Z_b(N)$ and killing $V'$ is the action of an element $j\in J_{b-1}(N)$. No additional labels are required.
\end{theorem}

\begin{example}[Actual permutation coefficients]
On a $3\times3$ board take
\[F=\begin{pmatrix}0&1&-1\\1&-1&0\\-1&0&1\end{pmatrix}.
\]
The recursion gives coefficients $-1,-1,2,1,-1$ for $312,231,213,132,123$, respectively. Therefore
\[F=-P_{312}-P_{231}+2P_{213}+P_{132}-P_{123}.
\]
Here $P_g$ has its one in row $g(y)$ of column $y$. The coefficients sum to zero, so the action on the empty tuple is zero, as required for $J_0$. Reducing this integer identity modulo two or any odd prime gives the corresponding realisation in that characteristic.
\end{example}

\begin{example}[The six-term cycle from a two-rook function]\label{tx:eg:two-rook-cycle}
Take $b=2,N=3$. Every two-rook $\rho$ has a unique completion to a permutation $g_\rho$: its one unused column must be paired with its one unused row. Define
\[
                       F(\rho)=\operatorname{sign}(g_\rho).
\]
Each row or column line has two rooks, and their completions differ by a transposition. Their signs add to zero, so $F\in\mathcal R_2^0$. The only permutation with longest chain length $N-b=1$ is $321$. Its chosen chain uses column $1$, and hence
\[
            \Psi_2=\bigl\{\{(2,2),(1,3)\}\bigr\},
            \qquad F(\rho_{321})=-1.
\]
This is the single forcing coordinate. In the triangular construction, every smaller canonical rook is prescribed zero. The complete decreasing recursion is:
\begin{center}
\begin{tabular}{@{}cll@{}}
\toprule
$g$ & $\rho_g$ & $a_g$\\
\midrule
$321$ & $\{(2,2),(1,3)\}$ & $-1$\\
$312$ & $\{(3,1)\}$ & $0-a_{321}=1$\\
$231$ & $\{(1,3)\}$ & $0-a_{321}=1$\\
$213$ & $\{(1,2)\}$ & $0-a_{312}=-1$\\
$132$ & $\{(2,3)\}$ & $0-a_{312}=-1$\\
$123$ & $\varnothing$ & $0-(a_{321}+a_{312}+a_{231}+a_{213}+a_{132})=1$\\
\bottomrule
\end{tabular}
\end{center}
Thus the realising element is
\[
 j=123-132-213+231+312-321
   =\sum_{g\in S_3}\operatorname{sign}(g)g.
\]
Here the three-digit strings denote permutations in the group algebra. On a one-tuple, the two permutations taking its label to any fixed output label have opposite signs; on the empty tuple, all six coefficients sum to zero. Consequently $j\in J_1(3)$. Its action on $e_{12}$ is
\[
 j e_{12}=e_{12}+e_{23}+e_{31}-e_{21}-e_{32}-e_{13},
\]
the zero-marginal cycle of the preceding chapter. Also $L_2(j)=F$ entry by entry, because only $g_\rho$ contributes at $\rho$. This example connects the tuple vector, its forcing coordinate and the actual permutation coefficients. Reducing the integer calculation in characteristic two turns every sign into one while preserving all cancellations.
\end{example}

\section{Why the triangular system has diagonal one}
The graph of $g$ contains $\rho_g$, so its contribution at that coordinate is exactly one. The lexicographic argument above places every other contributing permutation later. Subtraction in decreasing order therefore solves the system without dividing by a dimension, orbit size or group order, even over a ring with zero divisors.

The lower-level prescriptions are essential: they make $H_0=0$, then make each $H_k$ a zero-line-sum function by \eqref{tx:eq:lines}. Forcing successively gives $H_k=0$ for $k<b$ and $H_b=F$. Thus the same triangular construction proves both realisation and annihilation of shorter tuples, as the table in Example~\ref{tx:eg:two-rook-cycle} shows.

The stated realisation theorem is for $1\le b\le N$. If $b>N$,
the tuple module is zero and its only endomorphism is realised by
the zero group-algebra element. In the witness construction we
always have $b\le r<q\le N$, so the nonempty-board case is the
one being used.

\section*{What this theorem gives the lifting proof}
Return to $\bar\varphi:T_b\to\mathcal Z_b$ in the two-sided correction contract and its composite
\[
             \varphi=\iota_{\mathcal Z}\bar\varphi\pi:P_b\to P_b.
\]
Position equivariance turns this composite into a rook function. Its image in
$\mathcal Z_b$ gives zero column lines, and its vanishing on $V'$
gives zero row lines. The realisation theorem supplies an actual
$j\in K[S_N]$ with this whole matrix, and with $jP_k=0$ for every
$k<b$. Lemma~\ref{tx:lem:two-sided-correction} therefore applies: adding that single $j$ to the old witness repairs length $b$ and leaves every shorter length unchanged.

No fresh labels were used. The construction only ordered the
existing rows and columns, performed finite forcing, and subtracted
already known coefficients. It proves the realisation contract in
every characteristic, and that is why the same correction mechanism
will survive the transition from characteristic two to arbitrary
prime characteristic in the next chapter.

\chapter{Lifting in every positive characteristic}\label{tx:sec:odd}
\label{tx:ch:prime}
\section*{The property we need}
Recall the \emph{block operator} used to construct the lift:
\[
                       z=\sum_i c_i(\sigma_i-I).
\]
Here $\sigma_i$ translates the added block in its $i$th coordinate,
$I$ is the identity operator, and the $c_i$ are the chosen field
coefficients. Thus $z$ is a linear operator on our labelled vector
spaces, not a scalar. The equation $z^2=0$ says that applying $z$
\emph{twice} gives zero: $z(z(v))=0$ for every vector $v$. It does
not require $z(v)=0$ on the first application.

For a concrete example over $\mathbb F_2$, subtract the identity
from the operator which swaps two coordinates. The resulting matrix is
\[
 z=\begin{pmatrix}1&1\\1&1\end{pmatrix}\ne0,
 \qquad z^2=\begin{pmatrix}0&0\\0&0\end{pmatrix}.
\]
Each entry in its square is $1+1=0$ in the field. Such an operator is
called \emph{square-zero}; it sends its own image to zero.

The characteristic-two proof used two special properties of the
general block operator: it was square-zero, and its matrix on the
permutation basis equalled its transpose. The latter holds because
each coordinate translation is an involution, so its permutation
matrix is symmetric. In odd characteristic these properties require
different replacements. The combinatorial core, however, used signed
sums and triangular elimination without division by a group order,
and remains available.

This chapter proves every replacement needed for the same linear
retained-label budget. The block has $q=p^a$ labels and $z^p=0$.
The deletion quotient must discard every nilpotent chain longer than
one, not just the full length-$p$ chains. A common witness for the
top power $u=z^{p-1}$ yields a multi-factor projector for that actual
quotient. We first prove the requisite purity on the tuple generators,
then obtain identities in their common annihilator algebra, and only
then apply those identities to arbitrary subspaces and quotients.

The layer splitting now has $p$ levels. Duality uses the block
relabelling $v\mapsto-v$ to handle transpose. Finally, the same rook
realisation repairs the top-power residual without new labels.
These steps supply the exact-deletion theorem in every prime
characteristic.

The conclusion has the same progression form as before: fix $d$,
choose $q=p^a$ larger than its support budget, and lift with degree $d$
at every step. Odd characteristic changes the algebra which builds the
lift, not the meaning of repeated lifting or the linear label cost.

\section{The block operator in prime characteristic}
A linear operator is \emph{nilpotent} when some number of successive
applications gives zero. A nilpotent chain has independent members
$c,zc,\ldots,z^{j-1}c$, with $z^jc=0$; applying $z$ moves one step
down this chain. These are the elementary blocks often called Jordan
chains. The equation $z^p=0$ bounds their length by $p$, but does not
by itself exclude lengths between $2$ and $p-1$. That exclusion is
a property to prove for the actual block action.

Let now $\operatorname{char}K=p$ for any prime. Keep $N=n+q$ for the larger universe, with $n$ retained labels and $q$ added block labels. The support budget $r$ bounds the tuple lengths needed by the proof diagram; it is not the block exponent $a$ or the proof degree $d$.
As in the characteristic-two construction, we temporarily extend the
coefficient field if needed and write $K$ for the working field below.
The descent in Lemma~\ref{tx:lem:field} preserves the original degree
question. Take a regular block $A=\Fp^a$ of size $q=p^a>r$, choose
$c_1,\ldots,c_a$ independent over $\Fp$ in $K$, and put
\[z=\sum_i c_i(\sigma_i-1).
\]
The generators commute and have order $p$, so the binomial formula in characteristic $p$ gives
\[z^p=\sum_i c_i^p(\sigma_i^p-1)=0.
\]
For $p>2$ we cannot replace this by $z^2=0$. For example, for a three-cycle $\sigma$ in characteristic three,
\[z=\sigma-I,\qquad z^2=\sigma^2+\sigma+I\ne0
\]
on its three-point permutation space.

The binomial step uses the elementary characteristic-$p$ identity
$(U+V)^p=U^p+V^p$ for \emph{commuting} operators: each intermediate
binomial coefficient is divisible by $p$. It is valid for the coordinate
translations because they commute. It is not a rule for arbitrary
matrix sums, nor a statement that powers of unrelated operators commute.

\subsection{The deletion quotient}
Since $\im z$ need not be contained in $\ker z$, the characteristic-two expression $\ker z/\im z$ is no longer defined. Use
\begin{equation}\label{tx:eq:Dp}
\D(M)=\ker z/(\ker z\cap\im z).
\end{equation}
A nilpotent chain has basis $v_0,\ldots,v_{j-1}$ with $zv_0=0$ and $zv_i=v_{i-1}$ for $i\ge1$. Its kernel is $Kv_0$. If $j=1$, its image is zero and the chain survives in \eqref{tx:eq:Dp}. If $j\ge2$, $v_0$ is in the image and the chain disappears. Thus deletion keeps the one-dimensional chains and discards every longer chain.

One might instead try $\ker z/\im z^{p-1}$. It is a defined quotient, but it is the wrong one: a chain of length two in characteristic three would survive because $z^2=0$ there, even though it should be discarded by \eqref{tx:eq:Dp}. Keeping the quotient precise is important.

\section{The commutator argument in characteristic \texorpdfstring{$p$}{p}}\label{tx:sec:commutatorp}
Work on one moved orbit of a permutation basis. Choose a basis
object $x_0$, and let $H=\{v\in A:vx_0=x_0\}$ be its stabiliser.
Because the orbit is moved, $H$ is a proper subspace of $A$.
Choose a nonzero linear functional $\ell:A\to\Fp$ vanishing on
$H$, and define $\pi(vx_0)=\ell(v)$. This is well defined because
two translations taking $x_0$ to the same object differ by an
element of $H$. Let $D$ multiply each basis vector by its orbit
coordinate $\pi$, and let $S_i$ act by translation in coordinate
$i$. With $\varepsilon_i$ the $i$th unit vector, put
$e_i=\ell(\varepsilon_i)$ and $\alpha=\sum_i c_ie_i\ne0$.
The choices are made orbit by orbit. Exactly as before,
\[Dz-zD=W=\sum_i c_ie_iS_i,\quad Wz=zW,\quad W^p=\alpha^pI.
\]
The commutator $Dz-zD$ compares the order of two operations.
Here $W$ names that commutator operator, distinct from a polynomial
ambient space $W_N$ or a consequence space $W_d(n)$.
Its value $W$ is invertible for a concrete reason:
$W^{-1}=\alpha^{-p}W^{p-1}$. Independence of the $c_i$ makes the
scalar $\alpha$ nonzero. Thus it is $W$ whose inverse is used here,
not the nilpotent operator $z$.
Consequently the auxiliary linear operator $B=W^{-1}D$ satisfies
\begin{equation}\label{tx:eq:commutator}[B,z]=Bz-zB=I.\end{equation}
Here $B$ names an operator, not the set of block labels used in the
earlier chapters; brackets denote the displayed commutator of
composed operators.
This proof works also on a moved permutation basis tensored with an arbitrary representation: $D$ reads only the basis coordinate, while the $S_i$ act diagonally and still have order $p$.

There is an elementary model for this commutator. On the space with
basis $1,T,\ldots,T^{p-1}$, work modulo $T^p$, let $z_0$ multiply
by $T$, and let $B_0$ differentiate. Differentiation is well defined
modulo $T^p$ because $(T^p)'=pT^{p-1}=0$. For $j<p-1$,
$(B_0z_0-z_0B_0)T^j=((j+1)-j)T^j=T^j$.
At $j=p-1$ the first term is zero and the second has coefficient
$-(p-1)=1$ in the field. Hence $[B_0,z_0]=I$ on the whole space.
The actual orbit construction above produces its commutator directly;
this model explains why a finite-dimensional commutator can equal the
identity in characteristic $p$.

If $zm=0$, commuting $z$ past powers of $B$ yields
\[zB^jm=-jB^{j-1}m,
\qquad z^jB^jm=(-1)^jj!\,m.
\]
The first identity follows by induction from $zB=Bz-I$; applying it repeatedly gives the second. At $j=p-1$ the scalar $(p-1)!$ is nonzero in characteristic $p$, because every factor $1,\ldots,p-1$ is nonzero. Therefore
\begin{equation}\label{tx:eq:freep}\ker z=\im z^{p-1}
\end{equation}
on the moved summand. The reverse inclusion uses $z^p=0$.

The displayed formula starts with a kernel vector $m$ and constructs
a preimage at the full height $p-1$ by applying $B^{p-1}$ and dividing
by $(-1)^{p-1}(p-1)!$. Thus every vector killed by $z$ on a moved
orbit already comes from the top power. A full chain is a copy of
the truncated polynomial model just described; a direct sum of such
chains is what the earlier phrase ``free for $z$'' means here.
This is a statement about the block operator, not freeness for the
full symmetric-group action.
Equivalently, let the variable $T$ in the ring $K[T]/(T^p)$ act
as $z$. A full chain generated by $c$ has basis
$c,zc,\ldots,z^{p-1}c$ and is one free copy of that ring.
A direct sum of full chains is therefore a free module for this
truncated polynomial ring. This is the precise module meaning of
freeness used in the calculation.

This is the full replacement for the square-zero freeness calculation. It does not divide by $p$, by $|A|$, or by a symmetric-group order. It divides only by a nonzero orbit coefficient and by $(p-1)!$.

On a bounded-support permutation basis the retained summand has $z=0$ and the moved summand satisfies \eqref{tx:eq:freep}. Therefore the identifications with the retained basis and the tensor identifications from Chapter~\ref{tx:sec:block-algebra} still hold, with quotient \eqref{tx:eq:Dp}. Stable matrices are transported by exactly the same retained-component argument.

\subsection{A three-cycle computation}
In characteristic three, let $\sigma$ cycle $e_0,e_1,e_2$.
The operator $z=\sigma-I$ has
\[
 z^2=I+\sigma+\sigma^2,\qquad
 z^2 e_0=e_0+e_1+e_2,\qquad z^3=0.
\]
Its kernel is the line spanned by $e_0+e_1+e_2$, and this line is
exactly $\im z^2$. It lies in $\im z$, so the deletion quotient
is zero on this moved orbit. A separate fixed basis vector has
zero image and survives. This is the prime-characteristic counterpart
of the moved two-point calculation, with a full chain of three
levels instead of two.

The commutator argument proves the same conclusion for every moved
orbit, including those with nontrivial stabilisers and those tensored
with an arbitrary representation. The factorial is $(p-1)!$, not
$p!$: each factor being divided by is nonzero. It is this explicit
calculation that excludes shorter nontrivial nilpotent chains in
the original moved tuple modules.

\section{The top-power witness}\label{tx:sec:topwitness}
Set $m=p-1$ and $u=z^m$. Here $m$ records this fixed power
exponent, rather than a tuple length or support size. Like $z$, $u$ is a linear operator; $u^2=0$ because
$2(p-1)\ge p$ and $z^p=0$. Recall that $J_r(N)$ consists of the
group-algebra elements acting as zero on all ordered-tuple spaces
$P_b(N)$ with $b\le r$. The common certificate becomes
\begin{equation}\label{tx:eq:topwitness}uwu-u\in J_r(N),\qquad\text{or }uwu=u\text{ on every }P_b(N),\ b\le r.
\end{equation}
This says that $w$ takes the bottom of a full length-$p$ chain back to a preimage at its top. It is enough to build a projector for the actual quotient, but more steps are needed than in characteristic two.

\subsection{Fixed directions, full chains, and purity}
The condition called \emph{purity} below is
$\ker z\cap\im z\subseteq\im z^{p-1}$. It says that a kernel
vector which already has a preimage under $z$ also has a preimage
under its top power. On a chain of length $j\ge2$, the bottom vector
is in $\ker z\cap\im z$. It is in $\im z^{p-1}$ only when the
chain has the full length $p$. On a fixed line the intersection is
zero. This is why purity distinguishes fixed directions and full
chains from shorter nontrivial chains.

\begin{definition}[Fixed and full-chain decomposition]\label{tx:def:full-chains}
For an operator $z$ with $z^p=0$, a fixed and full-chain decomposition is
\[
             V=F\oplus\bigoplus_{j=0}^{p-1}z^jC,
             \qquad zF=0,
\]
where the sum is direct and $z^{p-1}:C\to z^{p-1}C$ is injective.
Thus each basis vector $c$ of $C$ supplies the full chain
$z^{p-1}c,\ldots,zc,c$; $F$ consists of fixed directions.
\end{definition}
``Fixed'' here means killed by $z$. It need not mean fixed by every
label permutation. The direct-sum decomposition is a statement about
this nilpotent operator, not an equivariant decomposition for the
full symmetric group.

This definition will describe the layer splitting constructed below.
For the original tuple modules, purity follows directly from the orbit
calculation, without choosing chain bases:
\begin{equation}\label{tx:eq:purity}\ker z\cap\im z\subseteq\im z^{p-1}.
\end{equation}
The image has no retained component, and a moved kernel vector is in
the top-power image by \eqref{tx:eq:freep}. The algebraic contract is
therefore: prove purity on these generators; use one top-power witness
to obtain projector identities there; transfer the identities through
$J_r(N)$ to prove exactness on the whole annihilator class. A chosen
chain complement need not be compatible with submodules or quotients.

\begin{example}[Why the distinction matters]
In characteristic three, a lone length-two chain has $u=z^2=0$. The relation $uwu=u$ is then vacuous for any $w$, but $\ker z\cap\im z$ is nonzero and $\im u=0$. Such a chain does not satisfy \eqref{tx:eq:purity}. The argument must establish its projector identities on the actual tuple modules first; the top-power relation alone on an arbitrary vector space is insufficient.
\end{example}

\subsection{The ordered projector}
For $p>2$, the square-zero operator $u=z^{p-1}$ can have a larger
kernel than $z$. In characteristic three, on the chain $zv_2=v_1$,
$zv_1=v_0$, $zv_0=0$, the top power kills $v_1$, while $z$ does not. If $w$
sends $v_0$ to $v_2$ and kills the other two vectors, the
characteristic-two expression $(I-uw)(I-wu)$ would leave $v_1$
unchanged. It therefore cannot model the deletion quotient for $z$.
The ordered product below uses intermediate powers of $z$ to remove
every level, not merely the top and bottom ones.

For a vector satisfying $z^{h+1}v=0$, call $h$ an upper bound on
its chain height. This height is independent of polynomial degree.
The proof controls the height of the image after each successive
factor. The permitted height decreases; the degree budget of a PC
proof is not being decreased. Only the completed product is claimed
to be the required projector; its factors are not assumed to commute.

\begin{lemma}[Projector from purity and a top-power witness]\label{tx:lem:prime-projector}
Let $z,w$ be endomorphisms of a vector space over a field of
characteristic $p$, and put $m=p-1$. Suppose $z^{m+1}=0$,
$z^mwz^m=z^m$, and \eqref{tx:eq:purity} holds. Define
\begin{equation}\label{tx:eq:prime-projector}
S_0=I,\qquad S_{k+1}=(I-z^kwz^{m-k})S_k\quad(0\le k<m),
\qquad e=(I-z^mw)S_m.
\end{equation}
Then $ze=ez=0$, $e^2=e$, and $ev=v-z^mwv$ whenever $zv=0$.
\end{lemma}
\begin{proof}
Operator products act from right to left. Their order is part of the
construction; no factor is commuted past $w$.

\paragraph{Eliminating high levels.}
At step $k$ put $h=m-k$. Inductively $z^{h+1}S_k=0$; this is
true at $k=0$ because $z^{m+1}=0$. For $k<m$ we have $h\ge1$, so
$y=z^hS_kv$ belongs to $\ker z\cap\im z$, hence to $\im z^m$.
The witness gives $z^mwy=y$, and therefore
\[
 z^hS_{k+1}v
       =z^h(I-z^{m-h}wz^h)S_kv=y-z^mwy=0.
\]
The permitted height drops from $h$ to $h-1$ at each step. Thus
$z^{m+1-k}S_k=0$ for every $k$, and finally $zS_m=0$.
As $z^{m+1}=0$, the last factor also gives $ze=0$.

\paragraph{Action on a kernel vector.}
If $zv=0$, every factor in $S_m$ fixes $v$, because its rightmost power of $z$ has exponent at least one. Therefore
\begin{equation}\label{tx:eq:ekernel}ev=v-z^mwv\quad(zv=0).
\end{equation}

\paragraph{Killing the image.}
The first factor $I-wz^m$ fixes $\im z$, since $z^{m+1}=0$. All later factors preserve $\im z$ because their correction terms have a left factor $z^k$ with $k\ge1$. Hence $S_mzv\in\ker z\cap\im z\subseteq\im z^m$. The final factor $I-z^mw$ kills this vector by the witness relation. Thus $ez=0$.

\paragraph{Idempotence.}
The witness gives
\[z^mwe=(z^mw-z^mwz^mw)S_m=0.
\]
Since $ev\in\ker z$, formula \eqref{tx:eq:ekernel} now gives $e^2v=ev$. We have proved $ze=ez=0$ and $e^2=e$.
\end{proof}

The final factor kills the bottom of each full chain as well.
Formula \eqref{tx:eq:ekernel} is especially useful for lifting:
only the final correction acts on a kernel vector. It is the identity needed for
$\Lift=e\iota$ and the retained-coordinate right inverse.

\subsection{Passing from tuples to every relevant space}
Recall that $\mathcal C_r(N)$ contains every finite-dimensional
label representation killed by $J_r(N)$, including the invariant
subspaces and quotients used by the consequence calculation. The
point of a common witness is that identities verified on all the
tuple generators then hold on every space in this class.

Concretely, $K[S_N]/J_r(N)$ is the algebra in which two permutation
combinations are identified when their difference kills every tuple
test space of length at most $r$. The action on each $M\in\mathcal C_r(N)$ factors
through this quotient. Consequently an identity such as $e^2=e$
proved on all the tuple spaces is an identity of this algebra and
acts on every such $M$. A chosen chain complement need not restrict
to a submodule, but an identity of common operators does. The argument
therefore transfers identities rather than separately chosen bases.

Naturality here concerns maps $f$ of label representations: they
intertwine every common group-algebra operator, giving
$f(ev)=e f(v)$. This does not assert that $e$ commutes with every
label permutation. Its image is used as a vector-space model for
the deletion quotient, which carries the appropriate residual action.

For the actual block action, purity on every tuple generator has
already been proved in \eqref{tx:eq:purity}. The proposition below
isolates the implication from that property and a common witness to
exactness. Its purity premise records the input to this algebraic
step; it is not an outstanding assumption of the block construction.

\begin{proposition}[Common identities give exact deletion]\label{tx:prop:prime-exactness}
For the block operator $z$ above, $z^p=0$. Suppose purity holds on
every $P_b(N)$, $b\le r$, and one
$w\in K[S_N]$ satisfies \eqref{tx:eq:topwitness}. The group-algebra
expression $e$ in \eqref{tx:eq:prime-projector} gives a natural
vector-space isomorphism $\D(M)\simeq eM$ for every
$M\in\mathcal C_r(N)$, and $\D$ is exact on this class.
\end{proposition}
\begin{proof}
The lemma proves the identities on the tuple generators, so $ze$, $ez$
and $e^2-e$ belong to $J_r(N)$. They act as zero on every object of
$\mathcal C_r(N)$, whether or not that object has a fixed and full-chain
decomposition. Formula \eqref{tx:eq:ekernel} follows there directly from
the ordered product, since all factors in $S_m$ fix a kernel vector.

The quotient-to-image map sends $[v]$ to $ev$. It is well defined
because $ez=0$, and $eM\subseteq\ker z$ because $ze=0$.
For $zv=0$, the difference $v-ev=z^{p-1}wv$ belongs to both
$\im z$ and $\ker z$. Taking the quotient class of $ev$ therefore
returns $[v]$. Conversely, idempotence makes every element of $eM$
fixed by $e$. These maps are inverse for the actual quotient
\eqref{tx:eq:Dp}.

Every module map commutes with the same group-algebra expression $e$.
For an exact sequence $0\to U\xrightarrow{i}V\xrightarrow{f}W\to0$,
restriction to $eU\to eV$ preserves injectivity. If $y\in eW$, any
lift $v\in V$ gives the projected lift $ev$, since $f(ev)=ef(v)=y$.
If $v\in eV$ maps to zero, write $v=i(u)$ and then
$v=ev=i(eu)$. Thus the projected sequence is exact. The natural
quotient-to-image isomorphisms prove exactness of deletion, without a
separate purity assumption on subobjects. As in
Chapter~\ref{tx:ch:square-zero}, $eM$ is a vector-space model;
the quotient retains the residual label action.
\end{proof}

\begin{example}[A full chain in characteristic three]
Let $zv_2=v_1$, $zv_1=v_0$, $zv_0=0$, and let $wv_0=v_2$, with $w$ zero on the other chain vectors. Then $z^2wz^2=z^2$. The projector is
\[e=(I-z^2w)(I-zwz)(I-wz^2).
\]
The rightmost factor removes $v_2$, the middle factor removes $v_1$, and the leftmost removes $v_0$. If a separate fixed vector $v_{\rm fix}$ is added and $wv_{\rm fix}=0$, every factor fixes $v_{\rm fix}$. Thus the projector removes exactly the full chain and retains the fixed direction. For $p=2$ the same formula has only the two factors $(I-zw)(I-wz)$ from Chapter~\ref{tx:sec:block-algebra}.
\end{example}

\section{The linear construction in odd characteristic}\label{tx:sec:oddlinear}
We now adapt the characteristic-two construction from
Chapter~\ref{tx:ch:linear-layers}, using characteristic-independent
combinatorial ingredients and new operator arguments. The combinatorial
ingredients are the partition of unity in Section~\ref{tx:sec:partition}, the cone
and signed component lifts in Section~\ref{tx:sec:cone}, and the
rook forcing and permutation realisation in
Chapter~\ref{tx:ch:rooks}. Their insertion and cancellation formulas
use $1$ and $-1$ without a characteristic-two cancellation, so these
combinatorial proofs remain valid in every field. They supply
the same retained-label count $2b-1$.

The operator arguments do require replacements. The two-level
splitting must become a full $p$-level splitting, and the
characteristic-two identity $z^{\mathsf T}=z$ must be replaced by
conjugacy with the transpose. The following subsections establish
these replacements before using the common correction contract
with $u=z^{p-1}$.

Fix the current tuple length $b\le r$, and recall
$V=P_b(N)$, $V'=\sum_i\im\partial_i^{\mathsf T}$,
$\mathcal Z_b=\bigcap_i\ker\partial_i$ and $T_b=V/V'$.
Thus $V'$ is the incoming shorter-tuple space, $\mathcal Z_b$ is
the zero-marginal space, and $T_b$ is the quotient by the incoming
directions. The filtration $\mathcal Z^{\le\ell}$ consists of
zero-marginal vectors supported on tuples with at most $\ell$ block
positions. Its prospective quotient layer
is
\[
G_\ell=\bigoplus_{|J|=\ell}
       K[\Inj(J,A)]\otimes\mathcal Z_{[b]\setminus J}(n),
\]
where $J\subseteq[b]$ is a set of positions and $A$ is the block-label set.
These are the same layer spaces used in the characteristic-two
construction; only the block action and its chain structure change.

\subsection{Constructing a fixed and full-chain decomposition}
A nonempty block pattern has an image set of size $\ell<q$.
Give it a grade $e\in\Fp$ using the orbit coordinate of that image
set, as in Chapter~\ref{tx:ch:linear-layers}. Let $E$ be the space on
all these block patterns, and let $E_e$ span the
patterns of grade $e$, and let $D$ read the grade. The same commutator
$W=[D,z]$ is invertible and commutes with $z$.
In this subsection $e$ is a scalar grade; the projector denoted by
$e$ in the preceding sections is an operator and plays no role in
the grade decomposition.

Fix a grade $e$ and suppose $\sum_{j=0}^{p-1}z^jv_j=0$ with
$v_j\in E_e$. The commutator identity gives
\[
 (D-eI)z^jv_j=jWz^{j-1}v_j\quad(j\ge1).
\]
Applying $D-eI$ to the relation and cancelling $W$ leaves
$\sum_{j=1}^{p-1}jz^{j-1}v_j=0$. Repeat $p-1$ times to obtain
$(p-1)!v_{p-1}=0$, hence $v_{p-1}=0$. Remove that term and repeat
with one fewer power. All $v_j$ vanish. This proves both directness
of $\sum_jz^jE_e$ and injectivity of each $z^j$ on $E_e$.
It does not require $W$ to preserve the grade: after each cancellation
the coefficients are still the original grade-$e$ vectors.

Consequently $p\dim E_e\le\dim E$ for each of the $p$ grades.
As they partition the basis, $\sum_e\dim E_e=\dim E$. Summing the
inequalities forces equality in every one. In particular
\[
                   E=\bigoplus_{j=0}^{p-1}z^jE_0.
\]
The grade depends only on the image set, so it is stable under
position permutations and splits by position domain. It ignores the
retained tuple. Tensoring each domain component with its retained
zero-marginal factor therefore gives
$G_\ell=\bigoplus_{j=0}^{p-1}z^jQ_\ell$, with $Q_\ell$
position-stable and $z^{p-1}$ injective on $Q_\ell$.

This is the $p$-level version of the two-level section construction.
The grade $e\in\mathbb F_p$ is a coordinate on an image-set orbit;
$\ell$ instead counts block positions. To make a section commute
with $z$, choose preimages of the chain starts $q\in Q_\ell$,
then prescribe the preimage of $z^jq$ to be $z^j\sigma(q)$.
Directness and injectivity make all the chain coordinates unique.
The last step goes to zero because $z^p=0$, so the prescription
is compatible with $z$ at every level without an averaging operation.

For a layer quotient map $f$ and the position-equivariant section
$\sigma$ supplied by the signed lift, define
\[
 \tau\Bigl(\sum_{j=0}^{p-1}z^jq_j\Bigr)
                  =\sum_{j=0}^{p-1}z^j\sigma(q_j).
\]
Uniqueness of the decomposition makes the formula well defined;
$f\tau=I$ follows from $fz=zf$. Commuting with $z$ is checked by
shifting each power up by one; the term at power $p$ is zero.
Its image splits that filtration quotient as an actual subspace.
Position stability of $Q_\ell$ and position equivariance of $z$ and
$\sigma$ also make $\tau$ position-equivariant.
Successively splitting the quotients gives
\begin{equation}\label{tx:eq:socsplitp}
\mathcal Z_b(n+q)=\mathcal Z_b(n)\oplus\bigoplus_{j=0}^{p-1}z^jC,
\end{equation}
where $C$ is position-stable and $z^{p-1}$ is injective on $C$.
This is a fixed and full-chain decomposition in the sense of
Definition~\ref{tx:def:full-chains}. The hypotheses remain
$n\ge2b-1$ and $b<q$: making the sections commute with $z$ adds no
labels to the retained-label cost of the signed lift.

\subsection{Top and zero-marginal spaces}
Write $V=P_b(N)$. We verify the four properties used in
Lemma~\ref{tx:lem:layer-interface}, now with $u=z^{p-1}$ in place of
$z$. In \eqref{tx:eq:socsplitp}, every vector has a unique expression
$s+\sum_{j=0}^{p-1}z^jc_j$, with $s\in\mathcal Z_b(n)$ and
$c_j\in C$. Since $zs=0$,
\[
 z\Bigl(s+\sum_{j=0}^{p-1}z^jc_j\Bigr)
                         =\sum_{j=0}^{p-2}z^{j+1}c_j.
\]
Directness and injectivity of the powers of $z$ on $C$ force
$c_0=\cdots=c_{p-2}=0$ in a kernel vector. Hence
\[
 \ker(z|_{\mathcal Z_b})=\mathcal Z_b(n)\oplus uC,
                         \qquad u\mathcal Z_b=uC.
\]
If a vector belongs to $uV\cap\mathcal Z_b$, it is therefore $s+uc$.
Both that vector and $uc$ have zero retained coordinates: $u$ kills
retained tuples and preserves the moved-tuple summand. Thus $s=0$, and
\[
                         uV\cap\mathcal Z_b=u\mathcal Z_b.
\]
Define $w_C$ on $\mathcal Z_b$ by sending $uc$ to $c$ and killing
the retained summand and the levels $C,zC,\ldots,z^{p-2}C$.
Injectivity of $u$ on $C$ makes this well defined; position stability
makes it position-equivariant. It satisfies $uw_Cu=u$ there.

\subsection{Duality with the actual transpose operator}
The transpose changes which space carries the inverse used by
duality. In the tuple pairing, $\langle uv,w\rangle
=\langle v,u^{\mathsf T}w\rangle$. Thus the operator induced by
$u$ on the quotient $T_b$ is paired with $u^{\mathsf T}$ on
$\mathcal Z_b$, rather than automatically with $u$ itself.
This is the part of the characteristic-two argument that must change.

All transposes here are ordinary bilinear matrix transposes. They do
not conjugate the coefficients $c_i$. A permutation matrix transposes
to the matrix of its inverse, giving
\[
 z^{\mathsf T}=\sum_i c_i(\sigma_i^{-1}-1).
\]
The map
$\iota(v)=-v$ on the block is a genuine label permutation, fixes
retained labels and commutes with position permutations. It therefore
conjugates the entire tuple-layer construction for $z$ to that for
$z^{\mathsf T}$: $\iota\sigma_i\iota^{-1}=\sigma_i^{-1}$ and
$\iota u\iota^{-1}=u^{\mathsf T}$. It preserves $V'$ and
$\mathcal Z_b$, since these are defined by label-equivariant marginals.
The symbol $\iota$ here denotes this block inversion and its action
on tuple spaces; it is different from the inclusion
$\iota_{\mathcal Z}:\mathcal Z_b\hookrightarrow V$ used in the
correction formulas.

For a three-element block labelled by $\mathbb F_3$, inversion fixes
$0$ and swaps $1,2$. Conjugating translation by $1$ with this swap
gives translation by $-1$. This is the elementary meaning of
$\iota z\iota^{-1}=z^{\mathsf T}$: the change is an actual
relabelling of the block, so the marginal constructions remain valid.
The inverse-image notation $(u^{\mathsf T})^{-1}(U)$ below means
$\{v:u^{\mathsf T}v\in U\}$, not an inverse matrix for $u^{\mathsf T}$.

In particular,
$u^{\mathsf T}V\cap\mathcal Z_b\subseteq
u^{\mathsf T}\mathcal Z_b$. This gives
\[
 (u^{\mathsf T})^{-1}(\mathcal Z_b)
                   =\ker u^{\mathsf T}+\mathcal Z_b.
\]
Taking orthogonal complements in the nondegenerate tuple pairing,
\[
 (uV')^\perp=(u^{\mathsf T})^{-1}(\mathcal Z_b),\qquad
 (uV\cap V')^\perp=\ker u^{\mathsf T}+\mathcal Z_b.
\]
Their equality proves $uV\cap V'=uV'$.

The perfect pairing $T_b\times\mathcal Z_b\to K$ identifies a
quotient vector with a linear functional on $\mathcal Z_b$.
The dual of a map $h$ sends such a functional $\lambda$ to
$\lambda\circ h$. This reverses composition order, as with ordinary
matrix transposes. In the following construction, this dual map is
$w_T$; the displayed pairing calculation verifies its inverse
condition on the correct space.

To specify the quotient witness, put
\[
 h=\iota w_C\iota^{-1}:\mathcal Z_b\longrightarrow\mathcal Z_b,
                  \qquad u^{\mathsf T}hu^{\mathsf T}=u^{\mathsf T}.
\]
The tuple dot product induces the perfect pairing
$\langle\ ,\ \rangle:T_b\times\mathcal Z_b\to K$.
Define $w_T:T_b\to T_b$ by
$\langle w_Tt,\zeta\rangle=\langle t,h\zeta\rangle$.
Then
\[
 \langle uw_Tut,\zeta\rangle
 =\langle t,u^{\mathsf T}hu^{\mathsf T}\zeta\rangle
 =\langle ut,\zeta\rangle,
\]
so $uw_Tu=u$. The pairing is invariant under position permutations,
and $h$ commutes with them; equivalently, their matrices transpose to
their inverses. Thus $w_T$ is position-equivariant. This supplies all
four correction properties with the actual transpose operator.

\subsection{A trace from transpose to correction in characteristic three}
The following small linear-algebra model displays the maps in the
interface. Let $V=Kf\oplus Ke_0\oplus Ke_1\oplus Ke_2$ over
$K=\mathbb F_3$, let $\sigma$ fix $f$ and cycle the three $e_i$, and
put $z=\sigma-I$, $u=z^2$. Choose the incoming space $V'=Kf$.
For the ordinary dot product its orthogonal complement is
$\mathcal Z=\Span\{e_0,e_1,e_2\}$, and $T=V/V'$ is paired perfectly
with $\mathcal Z$. These are model spaces for the interface, rather
than an assertion about a particular tuple length.

Write $v_{\Sigma}=e_0+e_1+e_2$. On $\mathcal Z$ the top power satisfies
$ue_i=v_{\Sigma}$, so its matrix is the all-one $3\times3$ matrix. The
chain-start space can be chosen as $C=Ke_0$: the vectors
$e_0,ze_0,v_{\Sigma}$ are independent. Define $w_C$ by
\[
 w_Ce_2=e_0,\qquad w_Ce_0=w_Ce_1=0.
\]
It sends $v_{\Sigma}$ back to $e_0$ and kills $C$ and $zC$, as the construction
requires. Block inversion fixes $e_0$ and swaps $e_1,e_2$; hence
$h=\iota w_C\iota^{-1}$ sends $e_1$ to $e_0$ and kills the other
two basis vectors. Taking the dual in the displayed pairing gives
\[
 w_T[e_0]=[e_1],\qquad w_T[e_1]=w_T[e_2]=0.
\]
For example, $uw_Tu[e_2]=uw_T[v_{\Sigma}]=u[e_1]=[v_{\Sigma}]=u[e_2]$.
The same calculation applies to every basis vector. This shows the
transpose reversal explicitly: $w_T$ is the dual of $h$, with its
specified domain $T$, rather than a reuse of $w_C$ on quotient vectors.

Now start this model's repair with $w_0=0$, so $\varepsilon=-u$.
The typed correction
$\varphi=\iota_{\mathcal Z}(-w_C\bar\varepsilon w_T)\pi$
sends $e_0$ to $e_0$ and kills $f,e_1,e_2$: indeed $w_T[e_0]=[e_1]$,
$\bar\varepsilon[e_1]=-v_{\Sigma}$, and $-w_C(-v_{\Sigma})=e_0$.
Therefore $u\varphi ue_i=v_{\Sigma}=-\varepsilon e_i$, while both sides
vanish on $f$. Thus the residual is cancelled on every vector.

Using $w=\varphi$ in the ordered projector makes the three stages
equally visible. The factor $I-wu$ kills $e_0$ and fixes $ze_0$ and
$v_{\Sigma}$; $I-zwz$ then kills $ze_0$ and fixes $v_{\Sigma}$; $I-uw$ kills $v_{\Sigma}$.
Every factor fixes $f$. The full chain is removed and the fixed line
survives. On actual tuple spaces the preceding construction supplies
position equivariance and the incoming/zero-marginal conditions;
the rook theorem supplies the common permutation combination.

\subsection{The same correction, with the correct power}
The remaining correction is the algebraic contract from
Chapter~\ref{tx:ch:linear-layers}, with $u$ in place of $z$.
Its typed composition uses the quotient $T_b$ at the input and the
zero-marginal space at the output. It uses neither an arbitrary choice
of quotient representatives nor a characteristic-two cancellation.
The preceding chain and transpose arguments have supplied its two
inverse conditions; the rook theorem then realises its matrix on the
same existing label set.

Suppose $w_0$ works on shorter tuples for the top-power relation. Its error is
\[\varepsilon=uw_0u-u.
\]
It maps into $\mathcal Z_b$ and kills $V'$, giving the factorisation
$\varepsilon=\iota_{\mathcal Z}\bar\varepsilon\pi$, where
$\bar\varepsilon:T_b\to\mathcal Z_b$. Write $u_T$ and
$u_{\mathcal Z}$ for the induced operators. We now verify both inverse
relations explicitly. The error's image lies in
$uV\cap\mathcal Z_b=u\mathcal Z_b$, so the witness on that
space gives $u_{\mathcal Z}w_C\bar\varepsilon=\bar\varepsilon$.
If $uv\in V'$, the
inclusion $uV\cap V'\subseteq uV'$ gives $uv=uv'$ for some
$v'\in V'$. As $\varepsilon=(uw_0-I)u$ and $\varepsilon V'=0$,
\[
                  \varepsilon v=(uw_0-I)uv'
                                  =\varepsilon v'=0.
\]
The quotient witness $u_Tw_Tu_T=u_T$ says that a representative
$v_T$ of $w_Tu_T\pi(v)$ satisfies $u(v_T-v)\in V'$. The preceding
implication therefore gives $\varepsilon v_T=\varepsilon v$.
Injectivity of $\iota_{\mathcal Z}$ gives
$\bar\varepsilon w_Tu_T\pi=\bar\varepsilon\pi$; surjectivity of
$\pi$ then gives $\bar\varepsilon w_Tu_T=\bar\varepsilon$.
No compatible choice of representatives is assumed.

With $\pi:V\to T_b$ and $\iota_{\mathcal Z}:\mathcal Z_b\to V$,
define the actual endomorphism
\[
       \varphi=\iota_{\mathcal Z}
                      (-w_C\bar\varepsilon w_T)\pi.
\]
Lemma~\ref{tx:lem:two-sided-correction}, applied with $u$ in place
of $z$, now gives $u\varphi u=-\varepsilon$. The map
is position-compatible, has image in $\mathcal Z_b$ and kills $V'$.
The characteristic-independent realisation theorem supplies
$j\in J_{b-1}(N)$ with this endomorphism. Hence
\[
        u(w_0+j)u-u=\varepsilon+uju=0
\]
on length $b$. On every shorter tuple module, $j=0$, so its old
witness relation is unchanged. The same label pool is used throughout.

Inducting on $b$ proves \eqref{tx:eq:topwitness} on the same label pool whenever $n\ge2r-1$. The longer projector gives exactness, the moved-orbit calculation gives natural basis and tensor deletion, and the original PC closure transport gives the same period statement.


\begin{proof}[Completion of the proof of Theorem~\ref{tx:thm:deletion}]
Begin at tuple length zero with $w_0=0$, since $z$ acts as zero on
constants. At length $b$, the preceding correction and
Theorem~\ref{tx:thm:realization} replace $w_{b-1}$ by $w_b$ satisfying
$z^{p-1}w_bz^{p-1}=z^{p-1}$ on every length at most $b$. The same
retained labels are used at every step. The largest requirement is
$\max_{1\le b\le r}(2b-1)=2r-1$; length zero needs no label.
The projector identities then imply exactness on $\mathcal C_r(N)$.
The orbit commutator calculation supplies the natural retained-basis
and tensor identifications, including the actual stable coefficient
maps. These arguments prove every assertion of Theorem~\ref{tx:thm:deletion}.
\end{proof}

\section*{Returning to the prescribed-target theorem}
The construction has proved the common top-power witness on every
tuple length up to an abstract budget $r$, at retained rank
$n\ge\max\{0,2r-1\}$ and block size $q=p^a>r$. The tuple purity
calculation and the explicit projector prove exact deletion on the
full support class. The moved-orbit commutator proves natural basis
and permutation-first-factor tensor deletion. Stable coefficient
matrices therefore represent the actual smaller axiom and target
maps, and exactness transports the full PC closure.

For a fully invariant target the application uses $r=kd+s$. For
the indexed request it uses $r=2kd+s+t$, covering the certificate
source and the doubled polynomial pairing used to reflect target
membership. These choices are made by the target-transfer proof;
the finite algebra just proved is independent of the equations.
The target remains the prescribed full polynomial, including when
the degree budget is below its actual degree. No target is added
to the axioms and no semantic-entailment inference is introduced.

Finally, extend scalars only while constructing the generic block
operator and its common witness. Lemma~\ref{tx:lem:field} returns
actual derivability of each original-coefficient target to the original
field. Thus the construction applies to every coefficient field of
fixed positive characteristic. It has not asserted that a coefficientwise
projection of an arbitrary extension-field proof preserves each proof
line: the descent argument uses the finite-dimensional consequence
spaces and recovers a legal proof from their bases.

Both the onset and the period are linear in the fixed proof-degree
budget. Chapter~\ref{tx:chap:assemble-transfer} now assembles the actual
target-transfer theorem; the classification argument in
Chapter~\ref{tx:ch:profiles} inverts its two quantitative estimates,
with the exact easy-residue and finite-peak qualifications. The proof
engine is now complete in every prime characteristic.

\part{Completing lifting and comparing NS with PC}

\chapter{Lifting the actual axioms, proof rules and targets}
\label{tx:chap:assemble-transfer}
All the finite-algebra obligations have now been supplied. We return
to the original question $\Gamma_n\vdash G_{b,n}$: can every requested
target be derived from the given axioms within degree $d$? We verify each part of the contract
in Chapter~\ref{tx:chap:contract}. The point of doing the assembly slowly
is to make clear that a prescribed target, rather than an arbitrary
lifted vector, is transported.

\section{The theorem we can now state without hidden assumptions}
Let $\Gamma_n$ be a fixed uniform axiom template and let $G_{b,n}$ be
a fixed uniform target family. Put
\[
 R(d)=2kd+s+t,\qquad n_*=\max\{n_0,1,kD_G+t\}.
\]
Here $k,s,t$ bound the label supports of a variable, an axiom index and
a target index; $D_G$ bounds target degree. The rank $n_0$ is the axiom
stability threshold. These are fixed data; only $d$ is a proposed proof degree.

\begin{theorem}[Actual target transfer]\label{tx:thm:transfer}
Over any field of fixed prime characteristic $p$, suppose
\[
 q=p^a>R(d),\qquad n\ge\max\{n_*,2R(d)-1\},\qquad N=n+q.
\]
For either NS or ordinary PC, every prescribed rank-$N$ target has a
derivation of degree at most $d$ if and only if every prescribed rank-$n$
target does. The assertion holds for every $d\ge0$, with the separate
zero-degree interpretation given below. It also holds for any
representative of one fixed global target orbit at each eligible rank.
For a fully invariant target use the sharper budget $kd+s$ and start
$\max\{n_0,1,kD_G\}$.
\end{theorem}

The theorem concerns original-field polynomials and actual derivations.
It does not require invariant certificates, does not assume that the
axioms are inconsistent, and does not add the target to them. If a
stabilised target degree exceeds $d$, the threshold predicate is simply
false. The empty target request is true at every degree.

\section{Transporting the entire consequence space}
Fix a positive degree budget $d$ and write $W_N=P_{N,\le d}$
for the ordinary polynomial space of degree at most $d$. Recall
\[
 I_d(N)=\Span_K\{\mu f_{a,N}:\deg\mu+\deg f_{a,N}\le d\},
 \qquad C_d(N)=V_\infty(N).
\]
The first space records NS identities of total-product degree at most
$d$, with $\mu$ ranging over monomials; the second is the closure of the usable axioms under the legal
ordinary-PC rules (Proposition~\ref{tx:prop:closure}). Boolean equations,
when present, remain inputs in this polynomial space.
We first transport the spaces $I_d(N)$ and $C_d(N)$, without choosing
any target. This separates the proof-system calculation from the
question of whether a particular uniform expression belongs to it.
The operation $\Del(M)=\ker z_M/(\ker z_M\cap\im z_M)$ deletes the
block through the label action; the common projector represents its
classes, and $\Lift=e\iota$ lifts smaller-rank vectors. Exactness means
that kernels, images and quotients are preserved by this operation.

Concretely, a deletion class is represented by a vector $v$ with
$z_Mv=0$. Two such vectors represent the same class if their difference
lies in $\ker z_M\cap\im z_M$. The projector provides representatives
which can be compared with smaller-rank vectors. We use the actual
inclusions and maps between these classes, not just equal dimensions
of the resulting spaces (Chapter~\ref{tx:chap:deletion}).

For NS the certificate map has source basis the pairs $(\mu,a)$ with
\[
 \deg\mu+\deg f_{a,N}\le d,
 \qquad (\mu,a)\longmapsto\mu f_{a,N}\in W_N.
\]
Its source has support at most $kd+s$, and $W_N$ has support at most
$kd$. The matrix entry at a target monomial $\nu$ is zero unless
$\mu\mid\nu$, and otherwise is $c(a,\nu/\mu)$. These are the actual
rank-independent invariant entries of the certificate map. The
stability threshold $n_0$ makes each axiom degree independent of rank.
Exactness and matrix naturality therefore identify the embedded image
$\Del(I_d(N))$ with $I_d(n)$.

For PC let $X_N=K[\mathcal V_N]$ and $L_N=P_{N,\le d-1}$.
The tensor product records a variable and a lower-degree line together:
a pure tensor $x_v\otimes f$ is sent to their polynomial product,
and sums of pure tensors record sums of such products. The degree
restriction is imposed on $f$ before multiplication. This is the
variable-labelled register of Chapter~\ref{tx:chap:contract}.
The multiplication map
\[
       m_N:X_N\otimes L_N\longrightarrow W_N,\qquad
                 x_v\otimes f\longmapsto x_vf
\]
has stable entries. Its source is supported on at most
$k+k(d-1)=kd$ labels. The usable axioms give the image $V_0(N)$
of their coefficient map. The recurrence is
\[
 V_{j+1}(N)=V_j(N)+\im\bigl(m_N|_{X_N\otimes(V_j(N)\cap L_N)}\bigr).
\]
Every stage is a submodule of $W_N$, and every multiplication domain
is a submodule of the fixed support-bounded space $X_N\otimes L_N$.
Suppose stage $j$ has been identified at both ranks. Put
$U_N=V_j(N)\cap L_N$ and $U_n=V_j(n)\cap L_n$.
The intersection is the kernel of $(v,l)\mapsto v-l$ on
$V_j(N)\oplus L_N$, so exactness identifies $\Del(U_N)$ with
$U_n$, including their inclusions $i_N:U_N\hookrightarrow L_N$
and $i_n:U_n\hookrightarrow L_n$.

Write $m'_N=m_N\circ(1\otimes i_N)$ and
$m'_n=m_n\circ(1\otimes i_n)$ for the restricted multiplication
maps. The required compatibility is the commutative square
\begin{equation}\label{tx:eq:restricted-multiplication}
\begin{CD}
 \Del(X_N\otimes U_N) @>{\Del(m'_N)}>> \Del(W_N)\\
 @V{\simeq}VV @V{\simeq}VV\\
 X_n\otimes U_n @>{m'_n}>> W_n.
\end{CD}
\end{equation}
The left identification is \eqref{tx:eq:tensordelete}, natural in
the second factor: it sends $\Del(1\otimes i_N)$ to
$1\otimes i_n$. The ambient multiplication map $m_N$ is sent to
$m_n$ by stable matrix naturality. Composing these two statements
proves the square for the restriction. Exactness now transports its
image and the sum with $V_j$, giving stage $j+1$. The initial stage
is the actual usable-axiom image, so induction identifies the embedded
$\Del(V_j(N))$ with $V_j(n)$ for every $j$.

Both ascending chains stabilise because their ambient spaces are finite
dimensional. Choose a stage beyond both stabilisation points. We obtain
\begin{equation}\label{tx:eq:consequence-transport}
 \Del(I_d(N))=I_d(n),\qquad \Del(C_d(N))=C_d(n),
\end{equation}
where equality means equality after the natural identification
$\Del(W_N)\simeq W_n$. This transports every cancellation in a PC
proof, including a low-degree result subsequently multiplied by a
variable. No bound on proof length and no symmetry of individual proof
lines was used.

\section{Why deletion also reflects target membership}
Preserving an image is not by itself enough: an exact operation can
annihilate a nonzero quotient. The next lemma supplies the reflection
needed for a general target source.

The notation $S'/S$ means vectors in $S'$ considered equal when their
difference lies in $S$; it is a \emph{subquotient} of the ambient space.
In the lemma, $\mathcal C_{2u}(N)$ is the common-annihilator class of
Chapter~\ref{tx:chap:deletion}, not a PC consequence space.
Here $u$ is an integer support bound, unrelated to the top-power
operator $u=z^{p-1}$ used in the witness construction.
The doubled budget has a concrete purpose: a basis pair in
$P\otimes P$ can use the union of two support-$u$ label sets.
The annihilator class is not being assumed closed under arbitrary
tensor products.

\begin{lemma}[Nonzero subquotients survive]\label{tx:lem:detect}
Let $P=K[\mathcal U_N]$ have support at most $u$. Suppose deletion is
exact on $\mathcal C_{2u}(N)$, its block size exceeds $2u$, and the
natural identifications \eqref{tx:eq:permdelete}--\eqref{tx:eq:tensordelete}
hold. If $S\subsetneq S'\subseteq P$ are invariant subspaces, then
$\Del(S'/S)\ne0$.
\end{lemma}
\begin{proof}
The permutation basis gives a nondegenerate dot product on $P$,
equivariant under $S_N$. It induces a pairing
\[
                         S^\perp\otimes(S'/S)\longrightarrow K.
\]
This pairing is surjective: for $v\in S'\setminus S$, a linear
functional vanishing on $S$ can take value one at $v$; the dot product
represents that functional by an element of $S^\perp$.
Let $Y=S'/S$. Tensoring over a field is exact, so the displayed
source lies in $P\otimes Y$, itself a subquotient of $P\otimes P$.
The last permutation module has support at most $2u$, justifying
every use of exactness. If $\Del(Y)=0$, permutation-first-factor tensor
compatibility gives $\Del(P\otimes Y)=0$. Deleting the inclusion
$S^\perp\otimes Y\hookrightarrow P\otimes Y$ gives an injection
\[
                  \Del(S^\perp\otimes Y)\hookrightarrow0,
\]
so the pairing source also deletes to zero. Exactness applied to the
surjective pairing then gives $0\twoheadrightarrow\Del(K)\simeq K$,
which is impossible because $K\ne0$.

We use this surjective pairing, not a trace or an averaged inner
product. In characteristic $p$, composing a pairing with coevaluation
can yield a vanishing dimension scalar and would not prove the lemma.
\end{proof}

In the detector proof, $S^\perp$ consists of vectors whose dot product
with every vector of $S$ is zero. Thus pairing $h\in S^\perp$ with
$v+S\in S'/S$ is well-defined: replacing $v$ by $v+s$ changes the
value by $\langle h,s\rangle=0$. The tensor product is the linear
source of this bilinear pairing. Surjectivity says that some such
pair has value one, rather than that a dimension scalar is nonzero.

\begin{lemma}[An invariant target needs no doubled budget]\label{tx:lem:constant}
Let $j:K\to W$ and $H\subseteq W$ be equivariant and belong to the
exactness class. Then $j(1)\in H$ if and only if its deleted image
belongs to $\Del(H)\subseteq\Del(W)$.
\end{lemma}
\begin{proof}
Membership is the vanishing of $K\xrightarrow{j}W\to W/H$.
Its image is either zero or isomorphic to the simple trivial module
$K$. Exactness preserves this image, and $\Del(K)\simeq K\ne0$.
Thus the composite vanishes exactly when its deleted composite does.
The vector $j(1)$ may be any invariant target, not just the constant
polynomial one.
\end{proof}

\section{The actual uniform target map}\label{tx:sec:target-transfer}
For indexed targets, let $\mathcal B$ be the fixed index basis and
let $\gamma(b,\mu)$ be its stable coefficient rule. Write $\delta$ for the largest stabilised target degree, taking
$\delta=0$ for an empty request. At a rank where
all target degrees are at most $d$, define
\[
 J_N:K[\mathcal B_N]\longrightarrow W_N,\qquad
                [b]\longmapsto G_{b,N}.
\]
Here $J_N$ is a linear target map. It is distinct from the group-algebra
annihilator ideal $J_r(N)$ used to construct the common projector.
Its entry in row $\mu$, column $b$ is precisely $\gamma(b,\mu)$.
Equivariance and rank independence are the hypotheses of the natural
matrix identification in Theorem~\ref{tx:thm:deletion}. Consequently
$\Del(J_N)$ is the actual smaller target map $J_n$.
This statement concerns the prescribed full polynomials $G_{b,n}$;
we do not replace a target by its degree-$d$ truncation.

Let $H_N$ be either $I_d(N)$ or $C_d(N)$. Every target belongs to
$H_N$ exactly when $\im J_N\subseteq H_N$. Put
\[
          S_N=H_N+\im J_N,\qquad Y_N=S_N/H_N.
\]
The condition is $Y_N=0$. Exactness, target-map naturality and
\eqref{tx:eq:consequence-transport} identify its deletion with
\[
          Y_n=(H_n+\im J_n)/H_n.
\]
The polynomial permutation basis of $W_N$ has support at most $kd$.
Lemma~\ref{tx:lem:detect}, applied with the original monomial support bound,
reflects vanishing of $Y_N$ once the budget covers $2kd$.
The affine budget
\[
                            R(d)=2kd+s+t
\]
also covers the certificate source ($kd+s$), the target source ($t$)
and every multiplication domain. It therefore proves membership in
both directions. The target source need not be simple; reflection
comes from a subquotient of the ambient polynomial target space.
Doubling $R(d)$ once again would be unnecessary: the pairing uses
$W_N\otimes W_N$, whose support is $2kd$.

\begin{proof}[Proof of Theorem~\ref{tx:thm:transfer}]
First extend scalars to an algebraic closure, as permitted by
Lemma~\ref{tx:lem:field}. The degree of every indexed target stabilises
beyond $n_*=\max\{n_0,1,kD_G+t\}$. If $d$ is smaller than the
largest stable target degree, the all-target predicate is false at both
ranks. Otherwise the full map $J$ just described applies.
Take $N=n+q$, $q=p^a>R(d)$ and retained rank
$n\ge\max\{n_*,2R(d)-1\}$. Exact deletion transports $H_N$, $J_N$
and their sum and quotient. The detection lemma yields
\[
 \bigl(\forall b\in\mathcal B_N,\ G_{b,N}\in H_N\bigr)
 \quad\Longleftrightarrow\quad
 \bigl(\forall b\in\mathcal B_n,\ G_{b,n}\in H_n\bigr).
\]
Field descent holds for each actual target separately, so the equivalence
holds over the original field. Repeated steps connect any eligible
ranks in a common residue class and prove the stated periodicity.
For one fixed global target orbit, replace $\mathcal B$ by that orbit;
at eligible ranks, equivariance makes every representative equivalent.

At degree zero the stabilised targets are either all zero, or are
scalars with at least one nonzero value. In the latter case every target
is derivable in degree zero exactly when the original axioms include a
nonzero constant. If a target has positive stable degree, degree-zero
all-target membership fails. These are fixed predicates beyond $n_*$
and give the separate zero-level assertion.

For a fully invariant target use $j_N:K\to W_N$, $j_N(1)=G_N$,
in place of $J_N$. Lemma~\ref{tx:lem:constant} reflects membership without
the tensor pairing. The budget reduces to $kd+s$, with stable onset
$\max\{n_0,1,kD_G\}$. Taking $G_N=1$ gives the sharper refutation
special case. No target has been added to the axiom system.
\end{proof}


\section{From one deletion step to a residue class}
Write $h_S(n)$ for the maximum least degree of the requested targets
in system $S$. Recall the safe onset $N(d)$ and chosen $p$-power period
$Q(d)$ from \eqref{tx:eq:NQ}; their separate zero-degree values are
$N(0)=n_*$ and $Q(0)=1$.
Fix $d$ and an eligible power $q$. If two sizes at or above the chosen
onset have the same residue modulo $q$, repeatedly subtract $q$ from
the larger until reaching the smaller. Every retained size in those
steps remains eligible. Applying Theorem~\ref{tx:thm:transfer} in both
directions proves threshold periodicity.

To preserve an exact minimum degree rather than just the assertion
$h(n)\le d$, repeat the argument at every lower threshold. The periods
are powers of the same prime and divide the chosen degree-$d$ period;
zero degree has a stable separate predicate. The vector of all those
threshold answers determines the truncated value
$\min\{h(n),d+1\}$. This completes the input used in
Theorem~\ref{tx:thm:profiles} for the four profiles.

\section{The complete lower-bound argument in a paragraph}
Fix a minimum finite degree $d$ occurring at a size beyond its own onset.
The two-way threshold theorem makes that exact degree recur along a
prime-power progression. Therefore a degree which is not allowed to recur
must occur before its onset. Since the onset is $O(d)$, those isolated
occurrences must have degree $\Omega(n)$. On T3 hard residues the same
contradiction excludes every recurring easy level. On T4 we instead use
unbounded recurring levels and the short period to bring a representative
down to size $O(d)$, producing finite linear peaks. No iteration ever
reduces the absolute degree: it is the availability of recurrence and
its onset which produce the gap.

\chapter{An independent uniform NS--PC comparison}\label{tx:sec:comparison}
The linear PC lower bounds were proved by transporting the PC closure
itself. We now establish a different result: every bounded-degree PC
consequence has an NS identity of degree bounded independently of
universe size. The comparison applies to arbitrary polynomials, not
only invariant targets or refutations. Its qualitative consequence is
that NS and PC assign every fixed uniform request the same type.
Throughout, $\Gamma_n$ is a fixed uniform template over a field $K$
of fixed prime characteristic $p$. Recall that $I_d(n)$ is the span of
axiom products of total-product degree at most $d$, while $C_d(n)$ is
the ordinary-PC consequence closure inside $P_{n,\le d}$.

\section{What comparison can and cannot prove}
Suppose you already know that NS degree is large. It is tempting to use
$h_{\PC}\le h_{\NS}$ as a PC lower bound. The inequality points in the
wrong direction. Even an upper comparison $h_{\NS}\le F(h_{\PC})$
need not preserve a linear scale if $F$ is large. Our linear PC bound
therefore came from its own closure transport, not from inverting $F$.

The purpose of comparison is instead to say that bounded PC degree can
never coexist with unbounded NS degree on the same sequence of sizes.
This is its role in the running example of
Chapter~\ref{tx:ch:running-order}: the elementary induction-chain
restriction supplies a growing NS obstruction, and comparison rules
out recurring bounded PC degree for the original uniform order
family. The linear PC constant then comes from PC's own lifting
range, without inverting the possibly large function $F$.
It also controls recurring finite levels well enough to identify the
same one of the four types. The proof will make this statement for
\emph{every} polynomial vector in the PC consequence space, before
specialising to any uniform target.

\section{Why one plateau makes a space closed}
The ideal generated by $\Gamma_n$ consists of finite sums $\sum_f a_ff$
with arbitrary polynomial multipliers. The spaces $I_e(n)$ impose a
product-degree bound on those sums, so they need not be ideals:
multiplication may leave the allowed degree window. The argument below
finds enough closure inside a fixed polynomial ambient without asserting
that a truncated NS space is closed under arbitrary multiplication.

Imagine a chain of NS consequence spaces with increasing allowed degree,
but look only at their intersection with the fixed degree-$d$ ambient.
Thus $I_{d+j}(n)\cap P_{n,\le d}$ allows larger certificates while
retaining only consequences of polynomial degree at most $d$.
An invariant-chain bound forces two successive intersections to be equal.
At that equality, multiplication of a degree-$d-1$ vector cannot create
a new degree-$d$ vector: it has only used one extra unit of NS certificate
degree, precisely the step covered by the plateau. Thus the plateau is
closed under the ordinary-PC rules.

This is the core idea. The first lemma supplies the uniform chain bound;
the second argument turns a plateau into whole-consequence containment.
Neither step assumes symmetry of the proof lines or the target.

The length being bounded is the number of strict inclusions in an
ascending chain. For example, $A_0\subsetneq A_1=A_2\subsetneq A_3$ has two strict
inclusions. Invariance means that relabelling sends each subspace to
itself. Exact deletion transports the inclusions, and the
detector keeps a strict inclusion strict by preserving its nonzero
quotient. This is why a dimension bound at small rank can control
chains at every rank.

\begin{lemma}[Uniform chain bound]\label{tx:lem:chains}
For any bounded basis system $\mathcal U$, strict $S_n$-invariant
subspace chains in $K[\mathcal U_n]$ have a length bound independent of $n$.
\end{lemma}
\begin{proof}
Let $r$ bound the basis support. Over an algebraic closure, choose a
fixed $q=p^a>2r$. At $N=n+q$ with $n\ge4r+1$, deletion is exact at the
doubled budget $2r$. Lemma~\ref{tx:lem:detect} shows that each nonzero
subquotient $B/A$ of $K[\mathcal U_N]$ survives. Thus
$A\subsetneq B$ implies
$\mathcal D(A)\subsetneq\mathcal D(B)$: the whole strict chain becomes
a strict chain of the same length in $K[\mathcal U_n]$.

Repeatedly subtract $q$ while the rank exceeds $M=4r+q$. Every retained
rank in such a step exceeds $4r$, so all steps are eligible. The final
chain has at most $\dim K[\mathcal U_M]$ strict inclusions, because a
smaller-rank basis injects into the rank-$M$ basis. The same bound covers
the initial ranks at most $M$. Each step may choose its own generic
coefficients: the detector applies to arbitrary invariant subspaces,
without a naturality assumption on the chain across ranks. Finally,
scalar extension preserves strict inclusions, giving the result over
the original coefficient field.
\end{proof}

Only a finite exactness threshold at each fixed support is needed for
this argument. In characteristic two the earlier quadratic construction also suffices.
For an arbitrary positive characteristic the completed linear construction
already provides the needed threshold. The
bound may be large; its role is uniformity in $n$, not a sharp comparison
of the two minimum degrees.

\begin{theorem}[Uniform comparison for all consequences]\label{tx:thm:comparison}
For every uniform axiom template over a field of prime characteristic
$p$, there is a nondecreasing $F:\NN\to\NN$, independent of $n$ and of
the target, with $F(d)\ge d$ and
\begin{equation}\label{tx:eq:whole-comparison}
                              C_d(n)\subseteq I_{F(d)}(n)
                       \qquad(n,d\ge0).
\end{equation}
Consequently, every polynomial with an ordinary-PC derivation of degree
at most $d$ from $\Gamma_n$ has a total-product NS identity of degree at
most $F(d)$ from the same axioms. For the aggregate profiles of any
stable uniform target family, and for the profiles of any one fixed
target orbit,
\begin{equation}\label{tx:eq:target-comparison}
                      h_{\PC}(n)\le h_{\NS}(n)\le F(h_{\PC}(n)),
                      \qquad F(\infty)=\infty.
\end{equation}
Their NS and PC types in Theorem~\ref{tx:thm:classification} therefore agree.
\end{theorem}
\begin{proof}
Fix $d$. Monomials of degree at most $d$ form a permutation basis of
support at most $kd$. Let $L_d$ bound strict invariant-subspace chains in
$P_{n,\le d}$, uniformly in $n$, by Lemma~\ref{tx:lem:chains}.
Consider the invariant spaces
\[
              W_j=I_{d+j}(n)\cap P_{n,\le d},
                    \qquad 0\le j\le L_d+1.
\]
They form an ascending chain. Some $W_j=W_{j+1}$ with $j\le L_d$,
otherwise there would be more than $L_d$ strict inclusions.

This plateau contains every axiom usable within degree $d$ and is closed
under linear combinations. If $f\in W_j$ and $\deg f\le d-1$, then
multiplication of an NS identity for $f$ by a variable gives
$x_vf\in I_{d+j+1}(n)$. Its degree is at most $d$, hence
$x_vf\in W_{j+1}=W_j$. For $d=0$ there is no permitted multiplication
of a nonzero line; the zero line already lies in $W_j$.
Induction on \emph{any} degree-$d$ PC derivation now places every one of
its lines in $W_j$. In particular,
\[
                          C_d(n)\subseteq I_{d+L_d}(n).
\]
This induction retains cancellations: a lower-degree linear combination
can be multiplied because it lies in the plateau, even if its separate
summands have degree $d$.

Choose
\[
                          F(d)=d+\sum_{u=0}^d L_u.
\]
It is nondecreasing, satisfies $F(d)\ge d$, and gives
\eqref{tx:eq:whole-comparison}. No symmetry of a target was used.
Conversely, expansion of NS identities into monomial--axiom products
gives a same-degree PC derivation, so the first inequality in
\eqref{tx:eq:target-comparison} holds. If the aggregate PC degree is finite,
every requested target belongs to $C_d(n)$ at that budget; the inclusion
puts every target in $I_{F(d)}(n)$. Empty requests have degree zero.
Infinite PC degree forces infinite NS degree by the first inequality.
These observations prove the complete extended-order comparison.

For type agreement write $g=h_{\PC}$ and $h=h_{\NS}$.
The inequalities $g\le h\le F(g)$ transfer eventual boundedness in both
directions. If $g$ takes one finite value $d$ infinitely often, then on
those sizes $h$ lies in the finite interval $[d,F(d)]$; some finite value
of $h$ recurs. A recurring finite value of $h$ likewise forces a recurring
finite value of $g$. Thus the recurring sets are empty together.
If recurring values of $g$ are unbounded, the same finite-interval
argument yields unbounded recurring values of $h$. Conversely, choose a
recurring value of $h$ exceeding $F(B)$. On its infinite level set,
$g>B$ and $g\le h$, so some finite value above $B$ recurs for $g$.
Thus infinitude of the recurring sets agrees as well. Boundedness,
emptiness and infinitude determine exactly the four alternatives.
\end{proof}

\begin{corollary}[Common finite peaks]\label{tx:cor:common-peaks}
If a fixed uniform request has common type T4, there is a nondecreasing
$F$ such that, beyond every bound $B$, some size $n\ge B$ has
\[
 h_{\PC}(n)=d\ge B,\qquad d\in E(h_{\PC}),\qquad
 d\le h_{\NS}(n)\le F(d)<\infty,
\]
with $n\le(p+2)(\kappa d+\sigma)$. Both degrees satisfy the rational
linear lower bound in \eqref{tx:eq:finite-peaks} at these same sizes.
Here $E(h)$ is the set of recurring finite degrees, and
$R(d)=\kappa d+\sigma$ is the applicable transfer budget: $(\kappa,\sigma)
=(2k,s+t)$ for indexed targets and $(k,s)$ for an invariant target.
\end{corollary}
\begin{proof}
Take the finite PC peaks from Theorem~\ref{tx:thm:classification}.
The pointwise lower comparison gives the NS lower bound, and the uniform
upper comparison gives its finiteness at those sizes.
\end{proof}

The comparison does not identify exact degrees, easy-class values or
recurring degree sets. Its potentially large function $F$ cannot be
inverted to obtain a linear PC lower bound from a linear NS lower bound.
That quantitative result came directly from the PC transfer theorem.

\section{Reading the result as three tests of a profile}
A profile has a recurring finite level if it takes one finite value at
infinitely many sizes. The comparison sends that infinite set into a
finite interval of degrees in the other system, so some value of the
other profile repeats infinitely often. It also preserves whether
recurring levels are unbounded: choose levels above $F(B)$ and use
pigeonhole inside a finite interval. Finally, it preserves eventual
boundedness directly. These three tests distinguish T1, T2, T3 and T4.

This is why one does not need equality of exact degrees, equality of
recurring levels, or equal easy-class values to obtain same-type
agreement. For common finite T4 peaks, take a finite PC peak; pointwise
PC\,$\le$\,NS supplies its lower bound, and the upper comparison supplies
finiteness of the NS degree at that same size.

\part{Examples, witnesses and logical consequences}

\chapter{Examples with prescribed targets}\label{tx:ch:examples}
The classification allows both finite proof degrees and infinity. An
example should make clear which kind of hardness it illustrates. We
already have a finite T2 family from a Boolean first-order sentence:
the descending-predecessor example of
Chapter~\ref{tx:ch:running-order}, with its full lifting proof. Here we
begin with a consistent family whose target is derivable on one residue
class and underivable on the others. We then obtain finite hard proofs
of a nonconstant target. Finally we describe the four first-order
refutation families used by the undecidability argument.
Recall the test in Theorem~\ref{tx:thm:classification}: T1 is eventually
bounded; T2 has no recurring finite degree; T3 is not eventually bounded and has a finite
nonempty set of recurring finite degrees; T4 has infinitely many. For a request,
$h_S$ is its least-degree profile; $d_S(\Gamma)$ denotes the special
least degree of a refutation, whose target is $1$.

\section{A consistent family with an exact residue split}
Fix characteristic $p$ and, for $n\ge1$, take
\[
                  \Gamma_n=\{x_i-1:i<n\},\qquad
                  G_n=\sum_{i<n}x_i.
\]
The axioms have the solution $x_i=1$ for every $i$, so they are
consistent at every size. Both the axioms and the invariant target
have stable coefficient descriptions.

\begin{proposition}\label{tx:prop:consistent-T3}
For both NS and PC, the target profile is exactly
\[
 h_S(n)=\begin{cases}1,&p\mid n,\\\infty,&p\nmid n.\end{cases}
\]
In particular, it is T3. With the same axioms but the indexed request
$G_{i,n}=x_i$ for every $i<n$, the profile is constantly infinity and
has type T2.
\end{proposition}
\begin{proof}
Summing the axioms gives $G_n-n$. If $p\mid n$, this is $G_n$, so it
has an NS certificate, and hence a PC proof, of degree one. The target
is nonzero of degree one, proving equality. If $p\nmid n$, the root
with all coordinates one evaluates $G_n$ to the nonzero field element
$n$. Soundness excludes any derivation. The same root evaluates each
indexed target $x_i$ to one, proving the final assertion.
\end{proof}

This example also distinguishes an invariant sum from its individual
summands. The statement that the sum is a consequence does not say
that each summand is a consequence. The all-target profile is a maximum
of separate proof degrees, not the degree of a proof of their sum.

The easy residue class here is zero modulo $p$; its degree is exactly
one. The other classes have infinite degrees. This is a genuine
consistent-axiom example of the classification, but it does not yet
provide finite proofs requiring degree $\Omega(n)$. The next
construction addresses that different question.

\section{A fresh variable as a nonconstant target}\label{tx:sec:target-examples}
The degree of a target need not reflect the degree needed to derive it.
A polynomial of degree one may require arbitrarily high proof degree.
We first give an exact construction, then apply it to a uniform family
whose hard values are finite.

A simple invariant target illustrates the distinction between finite
hardness and failure of derivability. For $G_n=\sum_{i<n}x_i$, the input
$\{1\}$ gives degree exactly one in both systems for $n\ge1$. With no
inputs the same target is underivable and has degree $\infty$. Both
fit the uniform definition, but neither supplies a finite hard example.
The next construction does.

Here \emph{fresh} means a new indeterminate in the polynomial ring,
independent of every old variable and absent from every old axiom.
Any coefficient polynomial has a unique expansion
$a(x,y)=\sum_j a_j(x)y^j$; the notation $[y^1]a$ selects $a_1(x)$.
This linear operation is different from substituting a value for $y$.
The proof uses substitution for PC and coefficient extraction for NS,
because they control the two degree conventions differently.

\begin{proposition}[Adjoining a fresh target]\label{tx:prop:fresh-target}
Let $K$ be any field, let $\Gamma\subseteq K[x_v:v\in V]$ be any axiom
set, and adjoin a fresh variable $y$. Each certificate uses only
finitely many axioms. Let $\Gamma^\uparrow$ be the same axioms viewed
in $K[x_v:v\in V,y]$, with no additional axiom involving $y$. Then
\begin{equation}\label{tx:eq:fresh-identities}
 \begin{aligned}
 h_{\PC}(\Gamma^\uparrow\vdash y)
       &=\max\{1,d_{\PC}(\Gamma)\},\\
 h_{\NS}(\Gamma^\uparrow\vdash y)
       &=d_{\NS}(\Gamma)+1.
 \end{aligned}
\end{equation}
These equalities include degree zero and infinity, with
$\infty+1=\infty$.
\end{proposition}
\begin{proof}
For PC, embed a refutation of $\Gamma$ in the larger ring and multiply
its last line $1$ by $y$. If the refutation had degree $d$, the resulting
proof of $y$ has degree $\max\{1,d\}$: only the new last line is added,
and its degree is one.

Conversely, substitute $y=1$ in every line of a PC proof of $y$.
The old axioms are unchanged and linear combinations commute with the
substitution. Multiplication by an old variable remains a multiplication
step; multiplication by $y$ becomes a repeated line, which the linear
combination rule admits. No degree increases, and the last line becomes
$1$. Every proof of $y$ also has degree at least one. These two lower
bounds prove the first equality, including the case when no refutation
exists.

For NS, multiplying the coefficients of an identity
$1=\sum_fa_ff$ by $y$ gives an identity for $y$ with degree one larger.
For the reverse inequality, suppose
\[
                    y=\sum_f a_f(x,y)f(x)
\]
has total-product degree at most $e$. Take the coefficient of $y^1$.
Since every $f$ is independent of $y$, this gives the exact identity
\[
                    1=\sum_f [y^1]a_f(x,y)\,f(x).
\]
Every nonzero product on the right has degree at most $e-1$: it is the
coefficient of $y^1$ in a product of total degree at most $e$. Zero
products may be discarded. Hence $d_{\NS}(\Gamma)\le e-1$.
In particular a finite proof of $y$ exists exactly when a finite
refutation exists. If the original refutation degree is zero, both
new degrees are one, as asserted.
\end{proof}

\begin{example}[Why the two exact identities differ]
Take $\Gamma=\{x,1-x\}$. Both refutation degrees are one. PC derives
$1=x+(1-x)$ and then multiplies that last constant by $y$; every line
still has degree at most one. An NS identity for the fresh target is
\[
                             y=yx+y(1-x).
\]
Its products have degree two. No degree-one NS identity can produce
$y$, since its multipliers of these degree-one axioms would have to
be constants. Thus the new PC degree is one and the NS degree is two.
PC can finish the refutation before multiplying its last line; an NS
identity carries the extra variable in its coefficient polynomials.
\end{example}

The absence of a $y$-axiom is part of the construction; in particular,
we do not add $y^2-y=0$. For a uniform template, give the new variable
empty support and keep all old axiom indices and coefficients. Set
all coefficients of axiom monomials involving $y$ to zero. The result
is again a uniform template with the same bounds $k,s,D,n_0$, and
$G_n=y$ is a fully invariant target of degree one at every rank.
Thus \eqref{tx:eq:fresh-identities} holds rank by rank.

This operation preserves each of the four types, separately for NS
and PC. The NS recurring finite degrees are shifted by one. For PC,
the map $d\mapsto\max\{1,d\}$ merges only degrees zero and one;
all larger finite degrees and infinity are unchanged. Eventual
boundedness, absence or unboundedness of recurring finite levels, and
the mixed periodic case are consequently preserved.

The notation $v_p(n)$ is the exponent of $p$ in the factorisation of
$n$: writing $n=p^a u$ with $p\nmid u$ gives $v_p(n)=a$.
For instance, $v_3(12)=1$ and its associated power is $q=3$.
A window such as $[q,q+4]$ bounds a finite proof degree; the condition
$n\ge2q$ must still be checked before using its lower endpoint.

\begin{corollary}[A finite hard nonconstant consequence family]\label{tx:cor:fresh-peaks}
Let $K$ have prime characteristic $p$, and let $\widehat B_n$ be the
normalised weak-bijection system defined in
\eqref{tx:eq:bhat-witness}, with all of its function and coverage witnesses
retained. Adjoin a fresh variable $y$ as above, and request the target
$G_n=y$. Both target-degree profiles are of type T4.

More precisely, write $n=p^a u$, where $p\nmid u$, and put $q=p^a$.
Whenever $n\ge2q$,
\[
 \begin{gathered}
 q\le h_{\PC}(n)\le q+4,\qquad
 q+1\le h_{\NS}(n)\le q+5,\\
                         h_{\PC}(n)\le h_{\NS}(n).
 \end{gathered}
\]
In particular, at every size $n_a=(p+1)p^a$,
\begin{equation}\label{tx:eq:fresh-peaks}
 \frac{n_a}{p+1}\le h_{\PC}(n_a)\le h_{\NS}(n_a)
                              \le\frac{n_a}{p+1}+5<\infty.
\end{equation}
\end{corollary}
\begin{proof}
The refutation bounds \eqref{tx:eq:bhat-window}, proved below together with the rectangular obstruction in the next
chapter, give
$q\le d_{\PC}(\widehat B_n)\le d_{\NS}(\widehat B_n)\le q+4$.
Apply \eqref{tx:eq:fresh-identities}; since $q\ge1$, the maximum in the
PC identity changes nothing here. At $n_a=(p+1)p^a$ we have
$v_p(n_a)=a$ and $n_a\ge2p^a$, yielding \eqref{tx:eq:fresh-peaks}.

For each fixed $a$, infinitely many sizes
$n=p^a(p(j+1)+1)$ have the same valuation and satisfy the displayed
bounds. Some finite degree in each of the bounded windows must
therefore recur infinitely often. The windows tend to infinity with
$a$, proving T4 in each system. On the other hand, the old upper
bound gives $h_{\PC}(n)\le5$ and $h_{\NS}(n)\le6$ whenever
$p\nmid n$. The linear peaks thus coexist with infinitely many easy
sizes, rather than giving an eventual linear lower bound at all sizes.
\end{proof}

This example has finite hard proofs of a nonconstant target, but its
axioms are inconsistent: the difficulty is inherited from refuting
$\widehat B_n$. It does not establish a family of satisfiable axioms
with the same hard consequence profile. The small consistent example
in Chapter~\ref{tx:ch:consequences} illustrates ordinary-PC consequences in
a different setting; no growing lower bound is asserted for it.

\section{The refutation special case: four first-order examples}
\label{tx:sec:examples}

We now take the constant target $G_n=1$. This is the refutation special
case of the uniform-consequence theorem; the sharper invariant-target
budget applies. Write $\psi_1,\psi_2,\psi_3,\psi_4$ for the four sentences described in order below.
The four cases all occur among fixed first-order sentences with no finite
nonempty models. We give examples over a field of fixed prime characteristic
$p$. Each sentence is translated by the graph-and-sum convention: relation
and function-graph variables are Boolean, and every function has a totality
sum and pairwise functionality equations, as in
Chapter~\ref{tx:sec:systems}. These are ordinary polynomial inputs,
including the Boolean equations. The parameter $n\geq1$ is the
universe size. For the matrix notation in this section, we relabel the
universe as $\{1,\ldots,n\}$; this changes no equation or degree.

\paragraph{T1: a local contradiction.}
The sentence $\forall x\,(R(x)\wedge\neg R(x))$ contributes $r_i$ and
$1-r_i$. Their sum refutes the system in degree one, independently of $n$.
The contradiction is already visible at one element.

\paragraph{T2: an injection that misses a point.}
Let $\eta$ say that a unary function $g$ is injective and that a constant
$c$ is outside its image. Write $x_{ij}$ for $g(i)=j$ and $c_j$ for $c=j$.
In addition to the Boolean equations, its system $\Gamma_n^\eta$ consists of
\begin{equation}\label{tx:eq:G-witness}
\begin{gathered}
 \sum_jx_{ij}-1=0,\qquad
 x_{ij}x_{ik}=0\ (j\ne k),\qquad
 x_{ij}x_{hj}=0\ (i\ne h),\\
 \sum_jc_j-1=0,\qquad
 c_jc_k=0\ (j\ne k),\qquad x_{ij}c_j=0.
\end{gathered}
\end{equation}
A solution would inject $n$ elements into the $n-1$ elements other than $c$,
so every instance is unsatisfiable. For either proof system, every
degree-$d$ refutation satisfies
\[
 n<3pd\qquad(n\geq3).
\]
Here is how the internal matching bound gives this estimate. Choose the
largest $p$-power $q\leq n/3$ and put $r=n-q$. Set $c_n=1$, all other
$c_j=0$, and restrict the image of $g$ to $[r]$. After transposition, every
remaining axiom is an axiom of the bijection system $B_{r,n}$.
Here $B_{r,m}$ is the Boolean $r$-by-$m$ matrix system with every row and
column sum equal to one and both row and column collisions equal to zero.
If $d<q$, then $2d\leq r$ and $q\mid n-r$, contradicting the rectangle
obstruction of Lemma~\ref{tx:lem:rectangle-obstruction}. The substitutions do
not increase ordinary-PC or NS degree. Thus $d\geq q$, while maximality of
$q$ gives $n<3pq\leq3pd$. This proves T2 without an external pigeonhole
lower-bound theorem.

\paragraph{The injection certificate and the linear scale.}
The paper's matching upper bound is useful to inspect directly. Put
\[
 R_i=\sum_{j=1}^n x_{ij},\qquad C=\sum_{j=1}^n c_j,
                    \qquad P=\prod_{i=1}^n R_i.
\]
Every monomial in $CP$ has the form $c_j\prod_i x_{i,\lambda(i)}$.
If $\lambda$ is not injective, two factors give a column-collision
axiom. If it is injective, it is a permutation of $[n]$, and its
unique row $i$ with $\lambda(i)=j$ gives the omission axiom
$c_jx_{ij}$. In either case the monomial is an axiom multiple of
degree $n+1$. Summing these identities expresses $CP$ in the raw NS
span at that degree.

The second certificate is the exact telescoping identity
\[
 CP-1=(C-1)+C\sum_{i=1}^n(R_i-1)\prod_{h<i}R_h.
\]
Its summands use only constant and row totality axioms, with each
complete product of degree at most $n+1$. Subtracting this certificate
from the one for $CP$ derives $1$ at degree at most $n+1$.
For a two-element universe the telescoping list is
$(C-1)+C(R_1-1)+CR_1(R_2-1)$, with degrees $1,2,3$.
A term $c_1x_{11}x_{21}$ in $CP$ is a column-collision multiple;
$c_1x_{11}x_{22}$ is an omission-axiom multiple. These are the two
possibilities in the general expansion.

Thus, over any field and for every $n\ge1$,
\[
 d_{\PC}(\Gamma_n^\eta)\le d_{\NS}(\Gamma_n^\eta)\le n+1.
\]
The PC inequality uses the degree-preserving simulation of an NS
certificate by repeated variable multiplication. No Boolean reduction
is used. With the preceding $n<3pd$ lower bound for $n\ge3$, both
degrees are $\Theta(n)$ for each fixed positive characteristic.
This family proves sharpness of the general linear scale; it does
not supply an upper bound for every uniform family. This certificate
and the sharpness deduction are written mathematics, as in the paper;
no additional Lean endpoint is claimed for them.

\paragraph{T3: hardness on one residue class.}
Let $\chi_p$ assert
\[
 \forall x\qquad f^p(x)=x\quad\hbox{and}\quad f(x)\ne x.
\]
Its finite models are disjoint unions of directed $p$-cycles. In particular,
it has a $0/1$ graph assignment exactly when $p\mid n$. An explicit
normalisation uses the function equations for $f$, the equations $f_{ii}=0$,
and
\begin{equation}\label{tx:eq:cycle-path}
 f_{i_0i_1}f_{i_1i_2}\cdots f_{i_{p-1}i_p}=0
 \qquad(i_p\ne i_0).
\end{equation}
If $p\nmid n$, summing all length-$p$ walks gives an NS refutation of degree
at most $p$. Indeed, totality derives the sum of their monomials minus $n$
within degree $p$. Nonclosed walks vanish by \eqref{tx:eq:cycle-path}.
A closed walk with a repeated vertex either contains a functionality
collision or follows a shorter deterministic cycle; since $p$ is prime,
the latter is a forbidden fixed point. Each remaining cycle contributes its
monomial $p$ times, according to its starting vertex, and hence contributes
zero. The resulting nonzero scalar $n$ yields the refutation.

Now take $\psi_3=\eta\wedge\chi_p$ in disjoint nonlogical symbols.
For $p\mid n$, substitute a $p$-cycle assignment into the second block.
This preserves proof degree and shows that the conjunction has exactly the
degree of $\eta$, separately for NS and PC. For $p\nmid n$, the second
block refutes in degree at most $p$. Consequently this all-size
unsatisfiable family is bounded off the multiples of $p$ and has linear
degree on the multiples: it has type T3. No exact degree is needed on the
bounded residue classes.

\paragraph{T4: the valuation of the universe size controls the difficulty.}
The weak-bijection example is more revealing. Let unary functions $f,g$
be injective with disjoint images, and require their images together to
cover the universe. Use a unary witness $h$ for coverage:
\[
 \forall z\qquad f(h(z))=z\ \vee\ g(h(z))=z.
\]
Let $a_{ij},b_{ij},w_{ji}$ represent $f(i)=j$, $g(i)=j$, and $h(j)=i$.
The combined $2n$-by-$n$ matrix $(a;b)$ has row sums equal to one and both
row and column collisions equal to zero. The witness block consists of
\begin{equation}\label{tx:eq:bhat-witness}
 \sum_iw_{ji}-1=0,\qquad
 w_{ji}w_{ji'}=0\ (i\ne i'),\qquad
 w_{ji}(1-a_{ij})(1-b_{ij})=0,
\end{equation}
together with the Boolean equations. Call this complete system
$\widehat B_n$. Its witnesses are retained: it is not identified with the
simpler matrix system $B_{n,2n}$.

Put $q=p^{v_p(n)}$. The proved bounds give
\begin{equation}\label{tx:eq:bhat-window}
 q\leq d_{\mathrm{PC}}(\widehat B_n)
 \leq d_{\mathrm{NS}}(\widehat B_n)\leq q+4
 \qquad\hbox{whenever } n\geq2q.
\end{equation}
The upper bound holds at every positive size, without the displayed size
condition. We give both degree-preserving reduction and certificate
explicitly.

A substitution is used here as a map of proofs, not merely as a map
of solutions. An affine substitution replaces a variable by a linear
polynomial plus a constant, so it does not increase polynomial degree.
One must also check that every substituted input is derivable from the
new inputs within the permitted degree. The following identities supply
that check for the coverage and Boolean witness equations.

For the lower bound transpose the combined matrix and substitute
$w_{ji}=a_{ij}+b_{ij}$ into every witness variable. Its totality becomes
the corresponding totality of $B_{n,2n}$. The identities
\begin{align*}
(a+b)^2-(a+b)&=(a^2-a)+(b^2-b)+2ab,\\
(a+b)(1-a)(1-b)&=-(1-b)(a^2-a)-(1-a)(b^2-b)
\end{align*}
put the substituted Boolean and coverage equations in the raw NS spans
at degrees two and three respectively. A substituted witness collision
expands into four column collisions of the raw combined matrix.
Thus every substituted input has a derivation within its own degree.
Substitution is affine: a multiplication by an old variable becomes a
sum of permitted multiplications by new variables, with the preceding
line still of degree at most $d-1$. A degree-$d$ ordinary-PC refutation
would therefore refute $B_{n,2n}$ within the same degree. The rectangle
obstruction excludes this when $2d\le n$ and $d<q$, yielding the lower
bound in \eqref{tx:eq:bhat-window}.

For the upper bound fix a column $j$ and abbreviate
\[
 u_i=a_{ij}+b_{ij},\qquad T=\sum_i u_i,\qquad
 Z=\sum_i w_{ji}u_i.
\]
The coverage axiom and disjointness collision give the exact identity
\[
 Z-1=(\sum_iw_{ji}-1)
 -\sum_iw_{ji}(1-a_{ij})(1-b_{ij})
 +\sum_iw_{ji}a_{ij}b_{ij}.
\]
Consequently $Z-1$ lies in the degree-three NS span. If $x$ is one entry
of this column, then
\[
 Tx-x=(x^2-x)+\sum_{y\ne x}yx
\]
lies in the degree-two span, by Booleanity and column collisions.
Adding this for $a_{ij}$ and $b_{ij}$ gives $Tu_i-u_i$ in that span.
Finally
\[
 T-1=\sum_iw_{ji}(Tu_i-u_i)-(T-1)(Z-1)
\]
puts the missing column totality in the degree-four span. This retains
the coverage witnesses throughout the derivation.

For completeness, consider any matrix with row sums $R_i$ and column
sums $C_j$, with their differences from one in the NS span at degree
$e$, and with both collision families in degree two. For a fixed
$k$-element row set $S$, telescoping the product of its row sums gives
$\prod_{i\in S}R_i-1$ in degree at most $e+k$. Summing over $S$ gives
\[
 \sum_{|S|=k}\prod_{i\in S}R_i-\binom{\#\mathrm{rows}}k
\]
in that same span. Expand the products as partial functions from $k$
rows to columns. A noninjective partial function contains a column
collision, so its monomial is already in the degree-$k$ span.
The remaining monomials are exactly those of $k$-edge matchings.
The column calculation gives the same matching sum, since inversion
pairs partial injections with identical monomials. Subtracting derives
\[
 \binom{\#\mathrm{rows}}k-\binom{\#\mathrm{columns}}k
\]
within degree $e+k$. Applied to the $2n$-by-$n$ matrix with $e=4$, this
is the desired counting certificate.

Write $n=qu$ with $q=p^{v_p(n)}$ and $p\nmid u$. In characteristic
$p$, $(1+t)^n=(1+t^q)^u$. The coefficient of $t^q$ is therefore $u$;
for $(1+t)^{2n}$ it is $2u$. Thus
$\binom{2n}q-\binom nq=u\ne0$ in $K$. Rescaling the counting
certificate gives an NS refutation of degree at most $q+4$.

For each fixed $a$, the sizes $n=p^a(p(k+1)+1)$, $k\geq0$, have
$v_p(n)=a$ and satisfy \eqref{tx:eq:bhat-window}. Hence some finite degree in
$[p^a,p^a+4]$ recurs infinitely often. These windows are unbounded as $a$
varies, proving T4 for both proof systems. On the other hand, $p\nmid n$
implies degree at most five. There are explicit simultaneous linear peaks:
at $n_a=(p+1)p^a$,
\[
 \frac{n_a}{p+1}\leq d_{\mathrm{PC}}(\widehat B_{n_a})
 \leq d_{\mathrm{NS}}(\widehat B_{n_a})
 \leq\frac{n_a}{p+1}+4.
\]
For example, in characteristic two every odd $n$ has degree at most
five, whereas at $n=3\cdot2^a$ both degrees lie between $2^a$ and
$2^a+4$. The peaks are linear in their universe size even though the
profile returns to bounded values infinitely often.
Thus linear peaks coexist with infinitely many easy sizes. This is the
reason the T4 conclusion concerns peaks rather than an eventual linear
lower bound at every size.

\section{Exercises}
\begin{exercise}\label{tx:ex:sum-versus-individual}
For the consistent family of Proposition~\ref{tx:prop:consistent-T3},
compute the two profiles at sizes $1,2,3,4,5,6$ in characteristic three.
Compare the invariant sum request with the indexed request for every
individual variable. Which statements have finite proof degrees?
\end{exercise}
\noindent{\small\hyperref[tx:ans:sum-versus-individual]{Model answer (p.~\pageref*{tx:ans:sum-versus-individual}).}}\par\medskip

\begin{exercise}\label{tx:ex:fresh-small}
For $\Gamma=\{x,1-x\}$ over an arbitrary field, write a degree-one PC
proof of the fresh target $y$ and prove that its least NS degree is two.
Do not insert an axiom involving $y$.
\end{exercise}
\noindent{\small\hyperref[tx:ans:fresh-small]{Model answer (p.~\pageref*{tx:ans:fresh-small}).}}\par\medskip

\begin{exercise}\label{tx:ex:fresh-window}
In characteristic three, apply the fresh-$y$ bounds to the normalised
weak-bijection family at $n=6$ and $n=36$. What can the stated bounds
say at $n=9$? Distinguish the lower-bound size condition from the
all-size upper certificate.
\end{exercise}
\noindent{\small\hyperref[tx:ans:fresh-window]{Model answer (p.~\pageref*{tx:ans:fresh-window}).}}\par\medskip

\begin{exercise}\label{tx:ex:fresh-soundness}
Suppose $\Gamma$ has a root over some extension field of $K$. Explain
why the fresh target $y$ is underivable from $\Gamma^\uparrow$.
If one instead adds $y^2-y$ as an axiom, identify the extra term which
appears when extracting the coefficient of $y^1$ from an NS identity.
Why does the original coefficient-extraction proof no longer apply
unchanged?
\end{exercise}
\noindent{\small\hyperref[tx:ans:fresh-soundness]{Model answer (p.~\pageref*{tx:ans:fresh-soundness}).}}\par\medskip


\chapter{Matching designs and the concrete lower bounds}
\label{tx:app:matching}

The examples use a lower bound for a concrete rectangular system.  Its
proof illustrates why a functional excluding a Nullstellensatz certificate
does not by itself exclude a Polynomial Calculus derivation.  To handle
degree falls we construct a whole space of functionals, together with an
extension operator between successive degrees.
The use of designs as linear obstructions to Nullstellensatz proofs
has a direct antecedent in Buss~\cite{BussDesign}. The argument here
also proves the extension property needed to handle ordinary-PC lines.
Here the target is $1$. The capital $D$ below is the proposed proof-degree
budget, rather than the input-degree bound of a general template.
The integers $N,M$ describe row and column counts in the auxiliary boards.

\section{Why this separate example proof is informative}
The general linear theorem used an exact operation on whole consequence
spaces. The concrete rectangular family gives a second way to see why
ordinary-PC degree falls must be handled honestly. Here we will use
linear functionals instead of a group-algebra projector.

A single functional annihilating the permitted axiom products excludes
an NS certificate. It need not remain valid when a cancelled low-degree
polynomial is multiplied in PC. The remedy is not to announce that NS
and PC are the same. We construct a whole space of designs and prove
that every lower-degree design extends. The PC invariant will quantify
over \emph{all} designs. This universal quantifier is what survives a
multiplication after cancellation.

The proof has four steps: choose a terminating signed reduction of
matchings, prove independence of its integral coordinates in a square,
compress complete histories to obtain a binomial identity in rectangles,
and use that identity to extend designs. The last step applies the
functionals to every legal PC line. We give the complete construction
in this chapter, rather than referring to a pigeonhole lower-bound paper.

For positive integers $r\leq m$, let $B_{r,m}$ consist of the equations
\[
 \sum_{j=1}^{m}x_{ij}-1=0\quad(i\in[r]),\qquad
 \sum_{i=1}^{r}x_{ij}-1=0\quad(j\in[m]),
\]
the Boolean equations $x_{ij}^{2}-x_{ij}=0$, and the row and column
collision equations
\[
 x_{ij}x_{i\ell}=0\quad(j\ne\ell),\qquad
 x_{ij}x_{hj}=0\quad(i\ne h).
\]
Here $[r]=\{1,\ldots,r\}$.  A Boolean solution would describe a bijection
between the two sides, so a proper rectangle has no solution.  Nevertheless,
in positive characteristic the two linear totality sums can agree because
$m-r$ vanishes in the field.  The following statement identifies a range
in which this obstruction persists through ordinary PC.

\begin{lemma}[Rectangle obstruction]\label{tx:lem:rectangle-obstruction}
Let $K$ be a field of prime characteristic $p$, and let $a,D$ be
nonnegative integers.  If
\[
 1\leq r\leq m,\qquad 2D\leq r,\qquad D<p^a,
 \qquad p^a\mid m-r,
\]
then $B_{r,m}$ has no ordinary-PC refutation of degree at most $D$.
Consequently it has no Nullstellensatz refutation of total-product degree
at most $D$.
\end{lemma}

We give the full construction over the integers before passing to $K$.
An \emph{edge set} is a subset of $[N]\times[M]$.  It is a matching if
no two distinct edges share a row or a column.  Write $\operatorname{rows}P$
and $\operatorname{cols}P$ for its two projections.  A degree-$d$ design
may take values in a coefficient ring: it is a function $\psi$ on edge
sets, zero on nonmatchings, for which
\begin{equation}\label{tx:eq:matching-design}
 \sum_{j=1}^{M}\psi(P\cup\{(i,j)\})=\psi(P),\qquad
 \sum_{i=1}^{N}\psi(P\cup\{(i,j)\})=\psi(P)
 \quad(|P|<d).
\end{equation}
The indices $i\in[N]$ and $j\in[M]$ are arbitrary. For $K$-valued
designs, denote the resulting vector space by $E_d(N,M)$. The signed
coordinate reductions below are constructed integrally and then
interpreted in $K$.  Normalisation, $\psi(\varnothing)=1$, is a separate
condition.  We will construct a normalised member of $E_D(r,m)$ and prove
that, for $D\ge1$, \emph{every} member of $E_{D-1}(r,m)$ extends to $E_D(r,m)$ while
preserving its values through degree $D-1$.

To interpret the table, associate an edge set $P$ with the squarefree
monomial $x_P=\prod_{e\in P}x_e$. Its proposed value is $\psi(P)$.
For an actual matching assignment this value would be its containment
indicator; a design instead imposes only the stated linear consistency
conditions in the tested degree window. A normalised design can be read
as a \emph{pseudoexpectation} in the following algebraic sense: a linear
assignment of values to monomials satisfying those tests and giving the
constant monomial value one, without a probability distribution or a
positivity assumption. Its extension to ordinary monomials is spelled
out in the last section. Unnormalised designs are also needed for the
extension argument.

\section{A degree-one design which is not a solution}
Over $\mathbb F_2$, consider four rows and six columns. Give the six
columns respectively the selected row labels $1,1,1,2,3,4$, and let
$\psi(\{(i,j)\})$ be one on those six selected edges and zero elsewhere.
Let $\psi(\varnothing)=1$. Every column sum is one. The four row
sums are $3,1,1,1$, all one in the field. Thus these values satisfy the
degree-one totality conditions, although the selected edges do not form
an actual bijection: three columns use the same row.

This is an honest linear witness, not a Boolean model. Degree one has
not yet tested the quadratic collision axioms. Higher-degree designs
must also assign compatible values to matching edge sets. The extension
construction below supplies these values exactly in the range of the
rectangle obstruction, without division by a factorial.

\begin{checkpoint}
An actual solution evaluates every polynomial multiplicatively. A design
is a consistent linear response only to the tested degree window. It
can exist for an inconsistent rectangle and still rule out every proof
within that window.
\end{checkpoint}

\section{A terminating signed reduction}
The reduction will replace one matching by a formal integer linear
combination of others. The word \emph{signed} allows negative
coefficients; it does not count matchings with positive weights.
The required identity is that every design gives the same value before
and after a replacement. Keeping integer coefficients first lets us
use the same identities after reducing modulo any characteristic.


A chain in a matching is a subset whose row coordinates and column
coordinates both increase.  Let $h(P)$ be the maximum chain length, and
write $P_{\leq b}$ for the edges in rows at most $b$.  A matching $P$ is
\emph{terminal} if, for every positive integer $b$,
\[
 h(P_{\leq b})\leq b-|\operatorname{rows}(P_{\leq b})|,
\]
and the corresponding inequality holds for column prefixes.  The right
side counts unused labels in the prefix.  This definition depends only
on the labelled matching, not on the ambient rectangle.  Prefixes beyond
its last occupied row or column cannot introduce a new failure: their
left side stays fixed and their right side increases.

For a nonterminal matching, use the least failed row prefix if one exists;
otherwise use the least failed column prefix.  In the selected prefix
choose a longest chain, breaking ties by the lexicographically greatest
list of its edges, with the edges themselves in row-first lexicographic
order.  In a column step use this rule after transposition.  Thus all
choices are deterministic.

Consider a row step with failed prefix $b$ and chosen chain $A$.  Put
\[
 P_0=P\setminus A,\qquad S=\operatorname{cols}A,\qquad
 B=\{i\leq b:i\notin\operatorname{rows}P\},
\]
\[
 U=\operatorname{rows}A\ \cup\
       \{b<i\leq N:i\notin\operatorname{rows}P\}.
\]
The sets $B$ and $U$ are disjoint and partition all rows unused by $P_0$.
Moreover, $|B|<|A|$, because the selected prefix fails.  A longest chain
in the first failed prefix ends in row $b$.  Indeed, if it ended earlier,
the same chain would belong to the preceding prefix, whereas the number
of unused labels cannot decrease when one adds the next row.  That
preceding prefix would already fail.

For sets of columns $S$ and rows $V$, a full injection from $S$ into $V$
means a matching $J$ with $\operatorname{cols}J=S$ and
$\operatorname{rows}J\subseteq V$.  A partial injection has only
$\operatorname{cols}J\subseteq S$.  Replace $P$ by
\[
 \sum_{\substack{J\text{ partial injection}\\S\to B}}
       (-1)^{|J|}[P_0\cup J]
 \ -\!
 \sum_{\substack{J\text{ full injection}\\S\to U\\J\ne A}}
       [P_0\cup J],
\]
omitting colliding terms.  Square brackets here denote formal basis
elements, not polynomial equivalence classes.  The transposed rule defines
a column step.

To verify the identity for a degree-$d$ design with $|P|\leq d$, repeatedly
apply its column equation to the columns of $S$.  At each application
the partial matching has fewer than $|P|$ edges, so the degree restriction
is satisfied.  An occupied row contributes zero by collision vanishing.
For a still unassigned column the sum over available rows in $U$ is
therefore the previous value minus the sum over rows in $B$.  Expanding
these successive differences gives
\[
 \sum_{J:S\to U}\psi(P_0\cup J)
  =\sum_{A'\subseteq S}(-1)^{|A'|}
       \sum_{J:A'\to B}\psi(P_0\cup J).
\]
Every surviving term is an injection, because collisions have value zero.
The right side contains each partial injection exactly once, indexed by
its column set.  Isolating the term $J=A$ on the left proves the stated
reduction.  The column step follows by interchanging rows and columns in
the design equations.  No assertion that the deterministic reduction
itself commutes with transposition is needed.

We next prove termination, including the equal-degree case.  For the
fixed grid $G=[N]\times[M]$, set $H=N+M$, $c=NM+1$, and define
\[
 w(i,j)=c^{H-i-j},\qquad W(P)=\sum_{(i,j)\in P}w(i,j),\qquad
 \rho(P)=|P|(W(G)+1)+W(P).
\]
Every lower term has fewer edges, since an injection into $B$ has at most
$|B|<|A|$ edges.  Thus its $\rho$ is smaller irrespective of its weight.
For a full injection $J\ne A$ into $U$, let $e$ be the first edge of $A$
not belonging to $J$, in increasing row order.  Every earlier chain edge
still belongs to $J$.  If $x\in J\setminus A$ had row smaller than that
of $e$, its row would be an earlier chain row already occupied by its
retained edge; this contradicts the matching condition.  A row from the
other part of $U$ is greater than $b$ and hence cannot be smaller.
Likewise, a column of $x$ smaller than the column of $e$ belongs to an
earlier retained chain edge, again giving a collision.  Consequently
\[
 e_1\leq x_1,\qquad e_2\leq x_2,\qquad e_1+e_2<x_1+x_2.
\]
At least one inequality is strict because $x\ne e$.  Each newly added
edge has weight at most $w(e)/c$, and there are at most $NM<c$ such
edges.  Their total weight is strictly smaller than $w(e)$.  Cancelling
the common edges of $A$ and $J$ shows $W(P_0\cup J)<W(P)$.
The identical argument applies to a column step, since $i+j$ is
unchanged by transposition.  Every child therefore has strictly smaller
nonnegative integer rank $\rho$.  The tree is finitely branching, and
every path has length at most its initial rank.  Hence the complete
signed reduction tree is finite.

For terminal $T$, let $C_{N,M}(P,T)\in\mathbb Z$ be the sum, over paths
from $P$ to $T$, of the products of their step signs; set it to zero for
colliding $P$.  The empty path gives
\[
 C_{N,M}(S,T)=\delta_{S,T}\quad(S\text{ terminal}),\qquad
 C_{N,M}(P,T)=0\quad(|T|>|P|).
\]
Induction on $\rho(P)$ using the reduction identity gives, for every
degree-$d$ design,
\begin{equation}\label{tx:eq:terminal-expansion}
 \psi(P)=\sum_{\substack{T\text{ terminal}\\|T|\leq d}}
              C_{N,M}(P,T)\psi(T)\qquad(|P|\leq d).
\end{equation}
Thus the terminal values determine every design in the relevant range.
It remains to show that these coordinates themselves satisfy the required
totality equations and can be restricted from a square to a rectangle.

For a small illustration of the preceding reduction, in a $2$-by-$2$
square the matching $\{(1,2)\}$ fails the first row prefix. Its chosen
chain is that one edge, $B$ is empty, and $U=\{1,2\}$. The rule gives
\[
 [\{(1,2)\}]\longmapsto[\varnothing]-[\{(2,2)\}].
\]
Here $\{(2,2)\}$ is terminal. For a design of degree at least one,
the corresponding identity is simply the column sum $\psi(\{(1,2)\})+\psi(\{(2,2)\})=\psi(\varnothing)$.
Thus a terminal coordinate can be negative; it is not a probability.

\section{Integral independence in a square}
An indicator here is a function on actual permutations of the square:
it takes value one when the permutation contains the given matching,
and zero otherwise. The argument is \emph{integral} because it uses
integer coefficients throughout. Here it proves independence after
passing to any coefficient ring in the theorem, as well as over
the integers. The signed moments below isolate a coefficient
with multiplier $1$ or $-1$, both invertible in every ring. This is the
reason for proving a unit diagonal rather than a nonzero factorial.


Let $T\subseteq[N]\times[N]$ be terminal.  For an edge $e\in T$, let
its height be the maximum chain length in the southwest downset
$\{f\in T:f_1\leq e_1,\ f_2\leq e_2\}$.  A longest such chain ends
at $e$: otherwise $e$ can be appended, because distinct edges of a
matching have distinct row and column coordinates.  If $e$ is southwest
of $f$ with both coordinates strictly smaller, appending $f$ raises the
height.  It follows that every equal-height level $T_t$ is decreasing:
as its row coordinate increases, its column coordinate decreases.
The levels partition $T$.

Let $h=h(T)$.  The terminal inequality at prefix $N$ implies that there
are at least $h$ unused rows and $h$ unused columns.  Let $a_t$ and $b_t$
be the $t$th smallest unused row and column, respectively, for
$1\leq t\leq h$.  These anchors precede every coordinate in $T_t$.
For example, if $e\in T_t$, a chain of length $t$ ending at $e$ belongs
to its row prefix.  Terminality requires at least $t$ unused rows in
that prefix; row $e_1$ is occupied, so the $t$th unused row is strictly
less than $e_1$.  The column argument is identical.  Form the paired blocks
\[
 A_t=\operatorname{rows}T_t\cup\{a_t\},\qquad
 B_t=\operatorname{cols}T_t\cup\{b_t\}.
\]
They are disjoint on each side, their paired sizes agree, and their
minima are $a_t,b_t$.  Their complements also have equal sizes.  Fix the
increasing bijection on the complements, and an increasing base bijection
$A_t\to B_t$ in each block.  Vary each base bijection by a permutation
$\sigma_t$ of $A_t$.  The union of all these bijections is an actual
permutation $\pi_\sigma$ of $[N]$.

For any edge set $P$, define the integer signed moment
\[
 M_T(P)=\sum_{\sigma}
     \left(\prod_{t=1}^{h}\operatorname{sgn}\sigma_t\right)
     \boldsymbol{1}_{P\subseteq\pi_\sigma}.
\]
If a block has two rows absent from $\operatorname{rows}P$, interchange
their images.  This is a fixed-point-free involution on the choices,
changes the sign, and preserves the containment condition.  Its terms
cancel over $\mathbb Z$.  In particular, $M_T(P)=0$ when $|P|<|T|$:
otherwise every block would use at least $|A_t|-1=|T_t|$ rows, forcing
$|P|\geq\sum_t|T_t|=|T|$.

The diagonal is a unit.  In each block the edges of $T_t$ prescribe the
bijection on every row except $a_t$; its remaining value is necessarily
$b_t$.  Thus exactly one choice $\sigma$ contains $T$, and
\[
 M_T(T)\in\{1,-1\}.
\]
This is a single signed contribution; no factorial or positive integer
multiplicity must be inverted.

For equal-degree triangularity, suppose $|P|=|T|$ and $M_T(P)\ne0$.
There is a choice containing $P$, so $P$ is a matching.  Cancellation
shows that at most one row in each block is unused.  Counting over the
disjoint blocks forces exactly $|T_t|$ edges in block $A_t\times B_t$,
one unused row in every block, and no edge in the complement.  Write
$P_t$ for this block matching.

Among all matchings of size $|T_t|$ in $A_t\times B_t$, the decreasing
matching $T_t$ has uniquely least weight $W$.  Here is a direct proof.
Choose a weight minimiser $Q$ from this finite set.  If it uses the
minimum row $a_t$, move that edge to the unique unused row.  This is a
legal matching, has the same cardinality, and strictly smaller weight,
because its row coordinate has increased.  Thus $a_t$ is unused.
The same argument shows that the minimum column $b_t$ is unused.
If $Q$ contains an increasing pair $(i,j),(i',j')$ with
$i<i'$ and $j<j'$, replace it by the crossed pair
$(i,j'),(i',j)$.  Both added edges have coordinate sum strictly larger
than that of $(i,j)$, and their combined weight is smaller than its
weight, since $c=N^2+1>2$ when such a pair exists.  The unchanged context
is disjoint from the crossed pair.  This again contradicts minimality.
Therefore $Q$ is decreasing, uses precisely $A_t\setminus\{a_t\}$
and $B_t\setminus\{b_t\}$, and must equal $T_t$: on two fixed ordered
sets there is exactly one order-reversing bijection.  Any distinct
matching with the same minimum weight would also be a minimiser and
would give the same conclusion.  Hence
\[
 W(P_t)\geq W(T_t),\quad\text{with equality only if }P_t=T_t.
\]
Summing over the disjoint blocks proves that
\[
 |P|=|T|,\quad P\ne T,\quad M_T(P)\ne0
 \quad\Longrightarrow\quad W(P)>W(T).
\]
Together with the lower-degree cancellation, this says that $M_T(P)=0$
for every $P\ne T$ with $\rho(P)\leq\rho(T)$.

Now suppose a finite linear combination of terminal matching indicators
vanishes at every permutation of $[N]$, with coefficients in any
commutative ring $R$.  Among its nonzero coefficients choose a terminal
$T$ of greatest rank $\rho$.  Apply the integer signed moment $M_T$ to
this relation, then map the integers to $R$.  All other nonzero terms
vanish by triangularity; the result is the coefficient of $T$ multiplied
by $1$ or $-1$.  This contradicts its being nonzero.  The terminal
indicators are therefore independent over every commutative ring,
including fields of characteristic two and rings with zero divisors.

Each square permutation satisfies the design equations in every degree.
Equation~\eqref{tx:eq:terminal-expansion}, applied to its containment
indicator, expresses every matching indicator in terminal coordinates.
For a fixed row $i$, the permutation identity
\[
 \sum_{j=1}^{N}\boldsymbol{1}_{P\cup\{(i,j)\}\subseteq\pi}
       =\boldsymbol{1}_{P\subseteq\pi}
\]
holds because $\pi$ has exactly one edge in that row.  Expand all terms
in terminal coordinates and use the integral independence just proved.
For each terminal $T$ this yields
\[
 \sum_{j=1}^{N}C_{N,N}(P\cup\{(i,j)\},T)=C_{N,N}(P,T).
\]
The column identity follows in the same way from the unique edge in a
column.  Colliding sets have zero coordinates and zero permutation
indicators, so they are included.  Thus every square coordinate is an
unnormalised design, in every degree, over the integers.

\section{Whole-history compression and the marked binomial identity}
The last matching of a path need not remember every row visited by
that path. We must group paths using their \emph{entire} histories.
Once the total signed weight is the same for every chosen $k$-element set
of appended rows, summing over those sets produces $\binom{L}{k}$.
The mark ensures that $k$ is positive, even if its initial edge is
later deleted. This will remove the unwanted appended-row contribution
when the binomial coefficients vanish in $K$.


Fix $P,T\subseteq[r]\times[m]$, with $P$ a matching and $T$ terminal, and a column
$j\notin\operatorname{cols}P$.  Append $L$ rows, and consider all pairs
consisting of a marked initial row $i\in\{r+1,\ldots,r+L\}$ and a signed
reduction path from $P\cup\{(i,j)\}$ to $T$.  The weight of a pair is the
product of its path signs.  Its \emph{history support} is the set of
appended rows appearing anywhere on the path, including the initial
state and edges that are subsequently deleted.

If $|P|+1\leq D$ and $2D\leq r$, every such support is nonempty and has
size at most $|P|+1$.  To prove the latter assertion, give one token to
each initial edge.  In a row rewrite, transport a token by its column
to the unique child edge in that column, or delete it if that column
is absent.  The child uses only columns of the parent, so this describes
every child edge exactly once.  In a column rewrite transport by row,
with the analogous deletion rule.  No deleted token reappears, and no
two surviving tokens represent the same edge.

For any matching $Q$ of at most $D$ edges, failure at row prefix $b$
implies
\[
 b< h(Q_{\leq b})+
       |\operatorname{rows}(Q_{\leq b})|\leq 2|Q|\leq2D\leq r.
\]
Thus every chosen row chain lies in old rows.  An edge already in an
appended row is part of the unchanged context of a row rewrite and is
retained verbatim.  A column rewrite preserves its token's row or
deletes the token.  It follows that a token can encounter at most one
appended row throughout its lifetime.  There are $|P|+1$ initial tokens,
so the full history has at most that many appended rows.  The initial
mark belongs to it even when its token later disappears.

We next verify that histories of the same size carry the same total
integer weight.  Let $S$ be a prescribed $k$-element subset of the $L$
appended rows, and let $f$ be the increasing injection fixing $[r]$ and
sending $r+1,\ldots,r+k$ to the elements of $S$ in order.  For $t\geq1$, set $f(r+k+t)=r+L+t$.
This extends $f$ increasingly and makes the inverse image of the large
row interval $[r+L]$ exactly $[r+k]$.  Rename the row coordinate of every
edge by $f$ and retain its column and step sign.

First, this renaming preserves the deterministic selector.  Every failed
row prefix lies below $r$, as just proved; those prefixes and their
unused-row counts are unchanged.  Their chains are unchanged as labelled
edge sets.  For a column prefix, the number of occupied columns and all
chain lengths are preserved by the increasing row renaming.  The first
failed column prefix is therefore unchanged.  The lexicographic tie
choice is also preserved: a strictly increasing map preserves the
lexicographic order of sorted edge lists, including the transposed lists
used in a column step.  Orientation, first failure and chosen chain all
agree under renaming, and terminality is preserved.

Second, renaming gives a bijection on the branches whose child edges
remain in $[r]\cup S$.  In a row step, the inverse images of the sets
$B$ and $U$ are exactly the corresponding sets in the compressed
rectangle: order comparisons with the old endpoint and occupied-row
membership are preserved.  An image-supported replacement therefore
has a unique inverse injection.  Its column set, cardinality, exclusion
of the original chain, collision status and sign are all preserved.
For a column step the permissible column sets are unchanged; the same
argument, using row membership, applies.  A child cannot encode two
different replacement injections, since its replacement edges are
recovered by removing the fixed context: their columns are disjoint
from the context columns in a row step, and their rows are disjoint
in a column step.  The lower and upper branch families are disjoint
because their cardinalities differ.  Thus no multiplicity is lost.
There may be other branches using rows outside $S$, but these belong to
different history supports and are not part of this correspondence.

Applying the branch bijection successively gives a two-way, sign-preserving
bijection between marked paths with history exactly
$\{r+1,\ldots,r+k\}$ in the compressed rectangle and marked paths with
history exactly $S$ in the large rectangle.  Its inverse simply renames
every state and the initial mark back.  The endpoint $T$ is fixed because
it uses only old rows.  All intermediate states have at most $D$ edges,
so the selector argument applies at every step.

Define $c_k(P,T,j)$ to be the total weight of the marked paths in the
rectangle $[r+k]\times[m]$ whose full history is all $k$ appended rows.
Partitioning the marked paths by their history support, and then by its
cardinality, proves the integer identity
\begin{equation}\label{tx:eq:marked-coordinate}
 \sum_{i=r+1}^{r+L}
       C_{r+L,m}(P\cup\{(i,j)\},T)
    =\sum_{k=1}^{D}\binom{L}{k}c_k(P,T,j).
\end{equation}
The coefficients are independent of $L$.  Its constant term is absent
because each path retains its initial mark in the \emph{history}, not
necessarily in its final state.  This distinction is essential.

\section{Rectangular designs and their extensions}

We now pass from integer reduction coordinates to the fixed field $K$.
The two goals are a design with $\psi(\varnothing)=1$ and an extension
of every degree-$(D-1)$ design; the first will detect $1$, and the second
will make the obstruction survive ordinary-PC multiplication.
Assume now the hypotheses of Lemma~\ref{tx:lem:rectangle-obstruction}, with
$D\geq1$, and put $q=p^a$ and $L=m-r$.  For $1\leq k\leq D<q$ the
integer $\binom{L}{k}$ is zero in $K$.  Indeed, $q\mid L$ and in
characteristic $p$,
\[
 (1+z)^L=(1+z^q)^{L/q};
\]
the right side has no coefficient in these degrees.  Take the square
of size $m$ and an old terminal matching $T\subseteq[r]\times[m]$.
The function
\[
 \theta_T(P)=C_{m,m}(P,T)\quad\text{in }K,
       \qquad P\subseteq[r]\times[m],
\]
is a degree-$D$ rectangular design.  Its row equations are the square
row equations, with the same column set.  Its column equations follow
from the square column equation after deleting the appended-row sum,
which vanishes by~\eqref{tx:eq:marked-coordinate} for $|P|<D$ and an
unused column.  If the column is occupied, every appended insertion
collides, so the omitted sum is already zero.  If $P$ collides, every
insertion still collides and both sides are zero.  This covers all
design equations.

On old terminal inputs, the square reduction is still terminal and
\[
 \theta_T(S)=\delta_{S,T}.
\]
In particular, $\theta_{\varnothing}$ is normalised.  For an arbitrary
$\psi\in E_{D-1}(r,m)$ define explicitly
\[
 \widetilde\psi(P)=
    \sum_{\substack{T\text{ terminal in }[r]\times[m]\\|T|\leq D-1}}
                     \psi(T)\theta_T(P).
\]
This is a member of $E_D(r,m)$ and agrees with $\psi$ on all terminal
inputs of size at most $D-1$.  Apply~\eqref{tx:eq:terminal-expansion} in the
old rectangle to both designs, after restricting the new design to
degree $D-1$.  They agree on every edge set of size at most $D-1$.
Thus the extension exists for every unnormalised design, with no
choice of an unspecified lift.

\section{Why extension is precisely the multiplication rule}
The polynomial spaces $P_{\le0}\subseteq P_{\le1}\subseteq\cdots$
form a \emph{filtration}: nested spaces indexed by the degree allowed.
The corresponding design problem asks for compatible values on each
window. Extending a design means keeping every tested value on edge
sets of size at most $D-1$ and supplying the extra degree-$D$ values.
It is stronger than constructing one normalised design at degree $D$;
the multiplication argument needs an extension for whichever
lower-degree table a given edge produces.

Fix an edge $e$. To test a new line $x_e f$ of degree at most $D$,
start with any degree-$D$ design $\psi$. Force the edge into the design
by defining $s_e\psi(P)=\psi(P\cup\{e\})$. This is a degree-$(D-1)$
design. Extend it back to degree $D$. The extension evaluates every
monomial of $f$ exactly as the original design evaluates that monomial
multiplied by $x_e$, because $f$ has degree at most $D-1$.

The induction hypothesis says that \emph{every} degree-$D$ design kills
earlier lines. It therefore kills $f$ under this chosen extension,
which proves that the original $\psi$ kills $x_e f$. This is the exact
point at which a space of designs succeeds where one isolated witness
would be insufficient.

\section{The invariant through ordinary-PC degree falls}

For a design $\psi$ define a linear functional on the original polynomial
ring by
\[
 \Lambda_\psi\!\left(\sum_\alpha b_\alpha x^\alpha\right)
      =\sum_\alpha b_\alpha\psi(\operatorname{supp}\alpha),
\]
where $\operatorname{supp}\alpha$ is the set of edges with positive
exponent.  This formula is a functional on ordinary polynomials; it
does not replace a PC proof by a proof in a Boolean quotient.
Every degree-$D$ design annihilates each axiom usable in a degree-$D$
proof.  Totalities use~\eqref{tx:eq:matching-design} at $P=\varnothing$;
collisions use vanishing on nonmatchings; and the Boolean equations
use the equality of the supports of $x_e^2$ and $x_e$.

For an edge $e$, put $(s_e\psi)(P)=\psi(P\cup\{e\})$.
If $\psi\in E_D(r,m)$, then $s_e\psi\in E_{D-1}(r,m)$.
Collision vanishing is immediate.  In either totality equation at
$|P|<D-1$, apply the degree-$D$ equation to $P\cup\{e\}$, whose
cardinality is at most $D-1$.  The insertion operations commute,
including repeated insertion of the same edge.  Also, directly from
the monomial formula,
\[
 \Lambda_\psi(x_ef)=\Lambda_{s_e\psi}(f).
\]

Induct through a degree-$D$ ordinary-PC derivation, asserting that every
line is annihilated by \emph{every} member of $E_D(r,m)$.  Input axioms
were handled above, and linear combinations preserve the assertion
even when their highest-degree terms cancel.  Consider multiplication
of an earlier line $f$ by $x_e$.  If $f=0$ there is nothing to prove.
Otherwise ordinary polynomial degree gives
$\deg(x_ef)=\deg f+1$, so $\deg f\leq D-1$.
For a given $\psi\in E_D(r,m)$, extend $s_e\psi$ to a design
$\nu\in E_D(r,m)$ using the preceding construction.  The earlier line
is annihilated by $\nu$ by the inductive assertion.  Since the two
designs agree on all edge sets of size at most $D-1$, and every monomial
of $f$ has support of that size or smaller,
\[
 \Lambda_\psi(x_ef)=\Lambda_{s_e\psi}(f)=\Lambda_\nu(f)=0.
\]
This proves the invariant for multiplication as well.  Finally the
normalised design $\theta_{\varnothing}$ sends the constant $1$ to $1$,
so that constant cannot be the last line of the derivation.

For $D=0$, all displayed axioms have positive degree because both sides
of the rectangle are nonempty.  No nonzero constant axiom is available,
and multiplication of a nonzero polynomial cannot produce a degree-zero
line.  Linear combinations of the available zero lines remain zero.
Thus a degree-zero refutation is impossible directly.  This completes
the ordinary-PC proof in every case.  The NS conclusion follows from
the degree-preserving simulation of a total-product NS certificate by
ordinary PC.

\chapter{The hardness spectrum, algebraic witnesses and models}
\label{tx:sec:model-theory}
We first explain the topological classification in the current paper,
then distinguish algebraic witnesses from actual models.
A uniform target invites a semantic question: does every solution of
$\Gamma_n$ satisfy $G_n=0$? Our theorem classifies syntactic algebraic
derivation, not that semantic question. The two are related through
linear witnesses, but the relation is not an identification of the four
types with four classes of ordinary models.
Keep a fixed uniform axiom family $\Gamma_n$ and its fixed target request
over a field $K$ of prime characteristic $p$. Its profile $h_S(n)$ is
the maximum least degree needed for those targets in NS or ordinary PC.

\section{The four types as shapes of a hardness spectrum}
\label{tx:sec:topological-spectrum}
The current paper gives a useful picture of the classification. This
picture applies to \emph{every} aggregate or fixed-orbit uniform request
in the theorem, over an arbitrary field of characteristic $p$. It is
not restricted to Boolean equations or to a particular translation of
first-order logic. Its input is eventual prime-power periodicity of
bounded-degree derivability.

The residue-class structure is already implicit in Riis's 2009
report~\cite[Section~4, pp.~18--19]{RiisLong}. The following reformulates
that structure and extends its reading to the uniform consequences
considered here; it is not a claim to have newly discovered the four
cases. No converse to the Euler-model criterion is needed.

\subsection{A \texorpdfstring{$p$}{p}-adic point is a compatible address in a residue tree}
At level one record a residue modulo $p$; at level two record one modulo
$p^2$ which reduces to the first residue; continue in this way. A
compatible infinite sequence is a $p$-adic integer. For example, in
characteristic two, the successive addresses of the ordinary integer
$5$ are $1\bmod2$, $1\bmod4$, $5\bmod8$, $5\bmod16$, and so on.
Not every infinite address represents a nonnegative ordinary integer.
This distinction is precisely what makes the extra points useful.
For further background see Gouv\^ea~\cite{Gouvea}; the argument
needed here is supplied below.

A neighbourhood specifies a finite part of the address:
$a+p^L\mathbb Z_p$ contains all points with the chosen residue modulo
$p^L$. Such a neighbourhood is both open and closed, or \emph{clopen}:
its complement is the finite union of the other residue classes.
The whole space $\mathbb Z_p$ is compact. One way to see why is to
view its points as infinite paths in the finitely branching residue
tree. If a cover by residue neighbourhoods had no finite subcover,
one could choose a residue at each level whose descendants still had
no finite subcover. The resulting infinite path would be uncovered,
a contradiction. General open covers reduce to these basic ones.

\subsection{Colour the residues which are eventually easy}
Use exactly $N(d)$ and $Q(d)$ from \eqref{tx:eq:NQ}, including the
separate convention $N(0)=n_*$ and $Q(0)=1$. Define
\[
 A_d(h)=\{a\in\{0,\ldots,Q(d)-1\}:\ h(n)\le d
       \text{ for all }n\ge N(d),\ n\equiv a\pmod{Q(d)}\},
\]
\[
 \mathcal E_d(h)=\bigcup_{a\in A_d(h)}(a+Q(d)\mathbb Z_p),\qquad
 \mathcal H_p(h)=\mathbb Z_p\setminus\bigcup_{d\ge0}\mathcal E_d(h).
\]
Here $A_d(h)$ is a \emph{set of residues}, whereas $N(d)$ is a
\emph{size threshold}. The calligraphic $\mathcal E_d(h)$ is a subset
of $\mathbb Z_p$; it is not the set $E(h)$ of recurring degree values.
The sets $\mathcal E_d(h)$ increase with $d$. To check this, pass to
the common finer period and compare sufficiently late positive integer
representatives. Each set is clopen, so the hardness spectrum
$\mathcal H_p(h)$ is closed.

A point lies \emph{outside} this spectrum precisely when it has a
residue neighbourhood on which all sufficiently large integer sizes
have one uniform finite degree bound. This equivalent description
shows that choosing later valid onsets or finer valid prime-power
periods does not change the spectrum. Finitely many exceptional sizes
also make no difference. We have not defined a polynomial system at
a noninteger $p$-adic size: the point records compatible information
about ordinary finite sizes.

\begin{proposition}[The four shapes]\label{tx:prop:topological-spectrum}
For the profiles in the four-way theorem,
\begin{center}
\begin{tabular}{@{}ll@{}}
\toprule
Type & Shape of $\mathcal H_p(h)$\\
\midrule
$\mathrm T_1$ & Empty\\
$\mathrm T_2$ & All of $\mathbb Z_p$\\
$\mathrm T_3$ & Nonempty, proper and clopen\\
$\mathrm T_4$ & Closed but not open\\
\bottomrule
\end{tabular}
\end{center}
For the same uniform request, NS and ordinary PC have the same
hardness spectrum, even if their exact degree profiles differ.
\end{proposition}
\begin{proof}
If the spectrum is empty, the increasing open sets $\mathcal E_d$
cover compact $\mathbb Z_p$. A finite subcover is contained in one
$\mathcal E_D$, so that set is the whole space. This gives an eventual
degree bound $D$, exactly T1. Conversely such a bound makes the
spectrum empty. The spectrum is the whole space exactly when every
$\mathcal E_d$ is empty. This means that no finite degree recurs,
exactly T2.

In T3 the classification theorem gives a fixed nonempty proper union
of easy residue classes, with a uniform bound $D$. No bounded-degree
residue neighbourhood can intersect the linearly hard region: a common
refinement would contain arbitrarily large integers satisfying both
incompatible bounds. Thus $\bigcup_d\mathcal E_d=\mathcal E_D$,
and the spectrum is nonempty, proper and clopen.
Conversely, for a nonempty proper clopen spectrum its complement is
compact. The same finite-subcover argument makes the complement equal
to some $\mathcal E_D$. Any recurring finite value $e$ occupies a
residue neighbourhood on a tail, by truncated periodicity. That
neighbourhood belongs to $\mathcal E_e\subseteq\mathcal E_D$, whence
$e\le D$. At least one finite value recurs, and the nonempty spectrum
excludes T1. This is T3. The remaining closed-but-not-open possibility
is therefore exactly T4.

Finally $h_{\PC}(n)\le h_{\NS}(n)\le F(h_{\PC}(n))$, with $F$
independent of $n$. A residue neighbourhood is eventually bounded for
one system if and only if it is eventually bounded for the other.
This proves equality of their spectra; it does not prove equality of
$\mathcal E_d$ for a fixed $d$.
\end{proof}

\subsection{A worked binary example: hard at the limit, finite at every size}
Take the abstract profile $h(n)=2^{v_2(n)}$ for positive integers $n$.
Thus odd sizes have value one, sizes twice an odd number have value
two, and sizes four times an odd number have value four. This is the
toy profile already used in Chapter~\ref{tx:ch:profiles}, not an
assertion that a particular polynomial family has exactly these degrees.
The easy regions are
\begin{center}
\begin{tabular}{@{}rll@{}}
\toprule
Budget $d$ & $\mathcal E_d(h)$ & Still outside this easy region\\
\midrule
$0$ & $\varnothing$ & $\mathbb Z_2$\\
$1$ & $1+2\mathbb Z_2$ & $2\mathbb Z_2$\\
$2,3$ & $\mathbb Z_2\setminus4\mathbb Z_2$ & $4\mathbb Z_2$\\
$4,\ldots,7$ & $\mathbb Z_2\setminus8\mathbb Z_2$ & $8\mathbb Z_2$\\
$8,\ldots,15$ & $\mathbb Z_2\setminus16\mathbb Z_2$ & $16\mathbb Z_2$\\
\bottomrule
\end{tabular}
\end{center}
Consequently $\mathcal H_2(h)=\bigcap_L2^L\mathbb Z_2=\{0\}$.
Every positive size has finite value, yet there is no bound valid on
any neighbourhood of this limiting address. For every $L$, the
neighbourhood $2^L\mathbb Z_2$ contains the arbitrarily large positive
sizes $2^j$ with values $2^j$. It also contains the arbitrarily large
sizes $2^L(2k+1)$ of the fixed finite value $2^L$.
Thus each neighbourhood of zero meets the easy region but is not
uniformly easy. The singleton $\{0\}$ is closed and not open.
This is the T4 picture. The point $0\in\mathbb Z_2$ is an address,
not a new instance with an empty underlying universe.

For comparison, the abstract profiles $h(n)=1$ and $h(n)=n$ give
empty and full spectra. The profile which is $1$ on odd sizes and
$n$ on even sizes gives $2\mathbb Z_2$, the T3 shape. These examples
separate the four pictures without claiming a realisation theorem
for arbitrary closed subsets of $\mathbb Z_p$.

\subsection{What topology explains, and what still needs the quantitative proof}
In T3, looking at a single finite number of residue digits suffices
to separate the easy and hard regions. In T4 it never suffices near
a boundary point. Compactness is the step turning local finite bounds
into one global bound in T1, or one bound on the entire easy part in T3.
It does not compute that bound.

Topology alone does not give linear lower bounds. For example,
$g(n)=1+v_2(n)$ has the same singleton spectrum as the preceding toy
profile, but $g(n)\le1+\log_2 n$. Its threshold periods grow
exponentially with the degree budget, rather than obeying the linear
period bound proved here. Thus it is not a counterexample to the
paper's theorem. The estimates $N(d)=O(d)$ and $Q(d)=O(d)$ are what
produce the linear gap and finite linear peaks. Keep the qualitative
picture and that quantitative input distinct.

\section{A quotient remembers what remains unproved}
For a fixed degree budget $d$, the ambient polynomial space must also
contain the target. We therefore allow polynomials up to degree
$\max\{d,D_G\}$, where $D_G$ is the fixed target-degree bound.
Form its quotient by $I_d(n)$ for NS or by $C_d(n)$ for PC, the consequence spaces
of Proposition~\ref{tx:prop:closure}. A target is derivable exactly when
its quotient class is zero. If the class is nonzero, elementary linear
algebra separates it from zero by a functional. This is a completely
abstract description of the failure to prove a consequence.

For target one the separating functional is normalised by $\lambda(1)=1$.
For a general target it must instead satisfy $\lambda(G)=1$.
That small change matters: a witness for failure to derive $x$ need not
be normalised on one at all. Nor is it generally evaluation at a model.
The two examples below make the distinction explicit.

\section{Dual witnesses for failure to derive a target}
For refutations, the design viewpoint of Buss~\cite{BussDesign}
organises such linear witnesses. We state the elementary separation
criterion for a prescribed target, whose normalisation is on that
target rather than necessarily on the constant $1$.
For either system write $W_d(n)$ for its degree-$d$ consequence space:
$W_d(n)=C_d(n)$ for PC and $W_d(n)=I_d(n)$ for NS. Let
$\delta=D_G$ be the fixed upper bound on target degree. We view
$W_d(n)$ inside $P_{\le\max\{d,\delta\}}(n)$, so the ambient space
also contains every requested target even when its degree exceeds $d$.

The \emph{dual} of a vector space is its space of linear functionals.
A functional \emph{annihilates} a subspace when it is zero on every
vector of that subspace. Thus the two conditions below say that the
functional ignores all proved consequences but detects the chosen
target. The normalisation on that target fixes the scale of the
functional; it does not make the functional multiplicative.

\begin{definition}[A degree-bounded target witness]
\label{tx:def:consistency-witness}
For a requested target $G_{b,n}$, a degree-$d$ witness is a $K$-linear
functional
\[
 \lambda:P_{\le\max\{d,D_G\}}(n)\longrightarrow K,\qquad
 \lambda(G_{b,n})=1,\qquad\lambda(W_d(n))=0.
\]
For an all-target request, a witness consists of an index $b$ and such
a functional. For a single fixed orbit, one may use any available
representative. With target $1$, this reduces to the usual normalisation
$\lambda(1)=1$ for a bounded-degree non-refutability witness.
\end{definition}

\begin{proposition}[Dual form of nonderivability]
\label{tx:prop:dual-witness}
A degree-$d$ target witness exists exactly when the target has no
derivation of degree at most $d$ in the chosen system. For an all-target
request it exists exactly when $h(n)>d$.
\end{proposition}
\begin{proof}
If $G_{b,n}\in W_d(n)$, a functional annihilating $W_d(n)$ cannot
take value one on it. Conversely its class in the finite-dimensional
quotient $P_{\le\max\{d,D_G\}}(n)/W_d(n)$ is nonzero whenever
it is not in $W_d(n)$. Extend that class to a basis, assign it value
one and assign the other basis vectors zero. Composing with the
quotient map gives the required functional. If the target degree is
greater than $d$, it automatically lies outside $W_d(n)$; the enlarged
ambient space also covers this degree-floor obstruction. Failure of
all-target derivability is failure for at least one target.
\end{proof}

A $K$-valued solution $a$ of $\Gamma_n$ which violates $G_{b,n}=0$ gives a
witness by $\lambda(f)=f(a)/G_{b,n}(a)$. For target one this is ordinary
evaluation, with its multiplicative property. A general witness need
not be multiplicative, positive or normalised on one. It records a
linear separation from all permitted bounded-degree consequences,
not necessarily a truth assignment to the original variables. If the solution is instead defined over an extension field $E/K$, the scaled evaluation initially takes values in $E$. Choose a $K$-linear map $\pi:E\to K$ with $\pi(1)=1$, by extending $1$ to a $K$-basis of $E$. Composing the scaled evaluation with $\pi$ gives a $K$-valued separating functional: it still vanishes on every consequence and takes value one on the target. This descent uses linearity, not multiplicativity of $\pi$.

\begin{example}[Semantic entailment need not be derivation]
The equation $x^2=0$ over a field forces $x=0$ at every solution.
Nevertheless, neither NS nor PC derives $x$ from this axiom: every
consequence belongs to the ideal $(x^2)$. For $d\ge2$, the functional
on $P_{\le d}$ taking value one on $x$ and zero on all other monomials
annihilates the degree-$d$ consequence space. This is a target witness
although no actual solution violates $x=0$. The example explains why
our minimum degree is a syntactic ideal-membership resource.
\end{example}

\begin{example}[Linear consistency without truth-value evaluation]
Over $\mathbb F_3$, take $x^2-x=0$, target one and degree budget two.
The PC consequence space is $K(x^2-x)$: no nonzero vector in that span
has degree at most one, so no multiplication is usable. The functional
\[
 \lambda(1)=1,\qquad\lambda(x)=\lambda(x^2)=2
\]
is a witness but is not evaluation: $\lambda(x^2)=2$ whereas
$\lambda(x)^2=1$. Boolean equations impose the intended truth-value
semantics on actual solutions, without making every linear witness an
actual evaluation.
\end{example}

Annihilating only axiom multiples gives an NS witness. For PC the
functional must annihilate the full closure, including polynomials
obtained by cancellation followed by new multiplication. This is the
dual form of the distinction between $I_d$ and $C_d$.

\section{Four patterns of bounded-degree algebraic information}
Let $h(n)$ be a fixed-orbit or all-target degree profile in the chosen
system. By Proposition~\ref{tx:prop:dual-witness}, a degree-$d$ witness
exists precisely when $h(n)>d$.
Recall that $E(h)$ consists of the finite values occurring at arbitrarily
large ranks. The four alternatives classify eventual boundedness and
whether $E(h)$ is empty, finite nonempty or infinite; they do not classify
the sets of roots of $\Gamma_n$.

\begin{proposition}[A dual reading of the classification]
\label{tx:prop:semantic-types}
The four types have the following mutually exclusive interpretations.
\begin{enumerate}[label=\textup{T\arabic*.}]
\item Some fixed $D$ suffices to exclude degree-$D$ witnesses at every
      sufficiently large size.
\item For every fixed $d$, witnesses exist at every sufficiently large
      size.
\item There are a fixed $D$, a $p$-power period and a nonempty proper
      easy residue set. No degree-$D$ witness exists on easy residues
      at large sizes. On hard residues, witnesses exist at every fixed
      degree for all sufficiently large sizes in those residues.
\item Infinitely many distinct finite degrees $e$ recur, each at
      arbitrarily large sizes. At an occurrence with $e\ge1$, a degree-$(e-1)$ witness
      exists while a degree-$e$ witness does not. There is also a fixed
      degree at which witnesses fail at arbitrarily large sizes. No
      fixed periodic partition separates a uniformly bounded-degree
      part from a part with positive linear hardness.
\end{enumerate}
\end{proposition}
\begin{proof}
T1 follows directly from $h(n)>d$ being equivalent to witness
existence. In T2 the linear bound eventually exceeds any fixed $d$;
conversely eventual witness existence at every fixed degree excludes
every recurring finite value. For T3 use the easy/hard residue split
of Theorem~\ref{tx:thm:classification}. Conversely the displayed split
makes the recurring values a nonempty subset of $\{0,\ldots,D\}$
and leaves the profile unbounded. In T4 use each recurring $e$ for
the two adjacent-degree assertions, and use the least recurring value
for infinitely many failures at a fixed degree. The impossibility of
a fixed periodic bounded/linear split is
Proposition~\ref{tx:prop:no-periodic-split}. Infinitely many recurring
finite values characterise T4, completing the converse.
\end{proof}

For each fixed $d$, the existence of a witness is eventually determined
by a residue modulo a $p$-power. T3 needs only one finite degree window
to distinguish all easy residues. In T4, larger windows keep revealing
new recurring finite levels. The linear onset bounds the last possible
isolated low degree, while the short period locates finite linear peaks.
For indexed targets, the separating index may vary with size within
the fixed requested system; this is not a change to the request itself.

The \emph{finite-model spectrum} of a sentence is the set of positive
integers $n$ for which it has a model with $n$ elements. It records
existence of solutions, without recording a polynomial proof budget.

The refutation examples in Section~\ref{tx:sec:examples} all have no
finite nonempty models yet exhibit all four types. Their identical empty
finite-model spectra cannot distinguish the cases. The fresh-target
construction in Section~\ref{tx:sec:target-examples} also shows that a
nonconstant consequence can retain these algebraic distinctions. The
classification is a property of uniform algebraic reasoning about the
chosen equations and targets, rather than a satisfiability classification.

\section{The connection with Euler structures}
A first-order structure consists of a domain and interpretations of
its relation and function symbols. A set is \emph{definable} when a
formula in that structure describes it; here the language and allowed
parameters are those required by the cited criterion. For example,
a unary relation defines a subset of the domain, and a function graph
defines a subset of its square. A definable map has a definable graph;
its fibre over $y$ is the set of inputs mapped to $y$. The index-sort
hypotheses below require the sets indexing the relevant variables and
sums to be available in that definable setting as well.

Kraj\'i\v cek's Euler structures \cite{KrajEuler} equip definable sets
with a characteristic $\chi$, normalised on points and compatible with
unions, products and definable bijections. The strong version also
respects definable maps with fibres of constant characteristic.
Here $M$ denotes the underlying first-order structure, and $\chi(M)$
the Euler characteristic assigned to its domain as a definable set.
This is a counting rule, not the field characteristic $p$. On a finite
set of three points, additivity and point normalisation give $\chi=3$
in its coefficient ring. For finite sets, the fibre rule is the familiar
identity $|X|=c|Y|$ when a map $X\to Y$ has exactly $c$ points in each
fibre. An Euler structure asks for the corresponding rules on the
permitted definable sets, which can be infinite.
We recall an external criterion with
its definability hypotheses; it is not an additional hypothesis of our
finite-dimensional transfer theorem.
For the degree criterion the coefficient ring is
$\mathbb Z/p^\nu\mathbb Z$, with $\nu$ depending on the fixed proof
degree and indexing bounds. Its reduction modulo $p$ defines summation
of finite-valued definable functions:
\[
 \sum_{x\in X}^{\chi}a(x)
       =\sum_{c\in\mathbb F_p}c\,\overline\chi(a^{-1}(c)).
\]
Under the required definability and index-sort hypotheses, a definable
solution of the axioms violating a uniform target $g=0$ excludes
bounded-degree NS and PC derivations of that target at sufficiently
large sizes $n\equiv\chi(M)\pmod{p^\nu}$
\cite[Theorem 6.1]{KrajEuler}. Taking $g=1$ gives its refutation case.
Thus uniform target membership already belongs to this historical
framework. A merely $\mathbb F_p$-valued characteristic need not
supply the required higher residue, and a weak Euler characteristic
omits the fibre rule. Kraj\'i\v cek and Scanlon \cite{KrajScanlon}
develop the associated Grothendieck-ring viewpoint.

The terminology records these counting identities algebraically.
The \emph{Grothendieck ring} has symbols $[X]$ for definable sets,
identifies definably bijective sets, imposes
$[X]=[Y]+[X\setminus Y]$ for definable $Y\subseteq X$, and uses
$[X][Z]=[X\times Z]$ for multiplication. Passing to the ring allows
formal differences of these symbols. An additive, multiplicative
characteristic assigns values to them while respecting those relations.
The strong fibre rule is an additional requirement, not something
supplied by the ring terminology alone.

This criterion supplies semantic obstructions rather than an
identification of T1--T4 with four model classes. Our classification
instead follows from exact finite-dimensional target and consequence
transport, with explicit linear onset. No Euler structure is an
assumption of the template theorem. An infinite model alone does not
supply the summation operation needed for totality equations, and a
linear separating witness need not extend to a multiplicative
interpretation in an Euler structure. Boolean axioms ensure truth-value
semantics of the first-order compiler; bounded supports and stable
coefficient maps provide the hypotheses of the algebraic transfer.

\section{What the model-theoretic language does not replace}
The finite-dimensional target theorem requires neither an infinite model
nor an Euler structure. Such structures can explain particular residue
obstructions, but the proof of the four types works directly with finite
consequence spaces and their actual maps. Boolean equations are essential
for the intended truth-value interpretation of the first-order encoding,
not for the general support-bounded transfer theorem.

A useful test of understanding is to say which statement is being proved:
that a target vanishes at all roots, that its polynomial lies in an ideal,
or that it has a derivation within one particular degree budget. These
are three different questions. The book's profile concerns the last.

\chapter{No algorithm can distinguish the types}\label{tx:sec:undecidability}
Already in the special case of the fixed target $G_n=1$, the
classification cannot be turned into an algorithm
which determines the case from a sentence. This remains true under the
promise that every finite instance is inconsistent. Throughout this section,
the input is a well-formed normalised sentence and the request is a
refutation of its translated axioms, with the translation of
Chapter~\ref{tx:sec:systems}; the characteristic $p$ is fixed.
For a sentence $\psi$, write $d_S(\psi,n)$ for the least degree of
deriving $1$ from its size-$n$ translation. This is its refutation
profile. Its type is the alternative in Theorem~\ref{tx:thm:classification}:
eventually bounded (T1), no recurring finite level (T2), not eventually bounded with
a finite nonempty set of recurring levels (T3), or infinitely many
recurring finite levels (T4).

\section{Knowing the alternatives need not identify the alternative}
The classification gives four exhaustive structural patterns. It does
not give an algorithm deciding which pattern belongs to a presented
uniform family. The reduction below makes this difference sharp even
for the constant target one, where every finite instance is inconsistent.

Think of a program whose halting is unknown. A uniform computation-prefix
sentence has models of every finite size if the program never halts.
If it does halt, a fixed finite prefix supplies a bounded-degree
contradiction at all large sizes. We use this switch inside a controlled
join of two example families. The join has one type in the nonhalting
case and the other type in the halting case. An algorithm distinguishing
those types would therefore decide halting.

Three algebraic points need proof, and each is supplied below. The
halting prefix yields an actual bounded-degree NS certificate, not just
a semantic contradiction. Adjoining a satisfiable block preserves the
other branch's degree. The guarded disjunction keeps both branches
visible with at most an additive degree-one overhead. These are the
connections between the computational reduction and our exact proof
systems.

An algorithm receives a finite string encoding a sentence: its symbols,
clauses and function arities can all be read effectively. \emph{Partial
computable} allows the algorithm to run forever on some inputs.
The theorem demands termination and a correct answer only on the
\emph{promised} inputs, namely sentences with no finite nonempty model
and type in $A\cup B$. It does not demand a procedure for recognising
that promise. The reduction will produce promised inputs in both
the halting and the nonhalting case, so a separator could not evade
them by failing to halt.

\begin{theorem}\label{tx:thm:undecidability}
Let $A,B$ be disjoint nonempty subsets of $\{\mathrm T_1,\mathrm T_2,
\mathrm T_3,\mathrm T_4\}$. There is no partial computable algorithm which,
on every normalised sentence having no finite nonempty model and having
type in $A\cup B$, halts and correctly distinguishes membership in $A$
from membership in $B$. This holds for NS and for ordinary PC, over every
field of characteristic $p$.
\end{theorem}

In particular, each of the six pairs is undecidable, and so is membership
in any nonempty proper union of types. The last conclusion is the scope
of the historical every-union theorem. The pairwise statement requires the
additional construction below; it does not follow merely by listing all
the undecidable union problems.

\section{A sentence that records a computation prefix}
A \emph{reduction} transforms each finite program description into a
finite sentence description by a fixed effective procedure. It must
not first decide whether the program halts. Calling the transformation
\emph{primitive recursive} states a stronger effectiveness property:
it is built using composition and recursion whose iteration counts
are given as natural-number arguments. The later unbounded search is performed by the
program, not by the procedure producing its sentence.

A counter program has finitely many natural-number registers and instructions
for incrementing a register, branching on zero and decrementing, and halting.
All registers initially contain zero. The number of registers is finite
but need not be two. Write $c_M$ for one more than the largest register index
occurring in $M$, so $c_M\ge1$.

From a program $M$, form a sentence $\theta_M$ with a strict linear order,
least and greatest elements $o,e$, a successor function $S$, counter-value
functions $C_r$, and state predicates $Q_\ell$. Require $S(e)=e$ and require
$S(t)$ to be the immediate successor of $t$ when $t\ne e$. At $o$, all
counters equal $o$ and the state is the initial one. Exactly one state holds
at each element, and halting states are forbidden.

At $t\ne e$, the state determines the transition to $S(t)$. Increment
replaces $C_r(t)$ by $S(C_r(t))$, with the guard $C_r(t)\ne e$.
Decrement branches according as $C_r(t)=o$; on the positive branch, the new
value $v$ satisfies $v<C_r(t)$ and $S(v)=C_r(t)$. Other counters are
unchanged. These are finite universal clauses with flattened function
graphs, and $M\mapsto\theta_M$ is primitive recursive.

\begin{lemma}[Computation switch]\label{tx:lem:switch}
If $M$ never halts, $\theta_M$ has a model of every size $n\ge1$.
If the run of $M$ first reaches a halting state at time $T$, the translated system has an NS refutation of
degree bounded by a constant $C_M$ at every size $n\ge T+1$.
\end{lemma}
\begin{proof}
In the first case, use $[n]=\{0,\ldots,n-1\}$ in its natural order and
record the first $n$ configurations. Every counter value at time $j$ is at
most $j$, so no guarded increment overflows before the last element.
This includes $n=1$.

For the second case put $a_j=S^j(o)$. If the structure has $T+1$ distinct
elements, then $a_j\ne e$ for $j<T$: otherwise immediacy of successor and
linearity of the order leave only $a_0,\ldots,a_j$. Induction along this
finite prefix forces the actual state at $a_j$ and forces
$C_r(a_j)=a_{c_r(j)}$, where $c_r(j)$ is the numerical counter value.
At $a_T$ this gives a forbidden halting state. The argument needs neither
well-ordering nor an induction axiom in the structure.

This local contradiction gives a degree bound by the finite-term
certificate in Lemma~\ref{tx:lem:finite-certificate}. More explicitly, with $c_M$
counters one may take
\[
 m=T+2+2c_M(T+1),\qquad
 C_M=m+\sum_R(m+T+1)^{\operatorname{arity}(R)}
                  +\max_{C\text{ a clause}}|C|.
\]
The term nodes are $o,e,a_1,\ldots,a_T,C_r(a_j),S(C_r(a_j))$;
$T+1$ distinct additional labels supply the size condition. The constant
depends on the program and its halting time, never on $n$.
\end{proof}

For completeness, the reduction needs effective universality, not a
decision procedure for the halting time. Every primitive recursive function
can be compiled into finite counter code: composition uses fresh registers,
and primitive recursion uses a counted loop, preserving the input registers.
Let $H(k,c)$ equal one when the program with code $c$, on its fixed
encoded input, halts within $k$ simulated steps, and zero otherwise.
Bounded simulation computes $H$ by a fixed number of primitive-recursive
transition updates, so this predicate is primitive recursive. A counter program that successively tests
$H(0,c),H(1,c),\ldots$ therefore halts exactly when the coded program does.
Loading $c$ into this fixed search program is primitive recursive. This
reduces the usual undecidable halting problem to the programs just used.
Only effective compilation and the standard undecidability of halting are used; the number of counters is unrestricted but finite.

\section{Finite local contradictions and their certificates}
\label{tx:app:finite-certificate}
The halting switch uses a fixed finite computation prefix even though
function totality is expressed by sums with $n$ summands. The following
elementary certificate construction separates these two sizes.

An acyclic evaluation list is an ordered list of term evaluations
whose arguments have already been named. This is a dependency condition
on the list, not a prohibition on a function taking the same value twice;
constants can be treated as evaluations with no arguments. The proof
sums over every labelling of that finite list. A \emph{falsity indicator}
for a Boolean literal is $1-x$ for a positive literal and $x$ for a
negative literal. The translated clause is the product of its literal
falsity indicators: it vanishes exactly when at least one literal is
true. These products are the actual input polynomials used below.

\begin{lemma}[Finite-term certificate]\label{tx:lem:finite-certificate}
Let $n\ge1$ and let a finite acyclic list of $m$ function evaluations be given, with every
argument supplied by an earlier node. Fix $b$ distinct labels. Suppose that
every assignment of labels to the nodes, and every Boolean assignment to
the relation atoms on the named and fixed labels, either has inconsistent
function values or falsifies a clause instance whose variables lie among the named and fixed
labels and all of whose function evaluations occur in the list. If clauses have length at most $L$, the
graph-and-sum translation has an NS refutation of degree at most
\[
                 m+\sum_R(m+b)^{\operatorname{arity}(R)}+L.
\]
The hypothesis and conclusion concern sizes for which the $b$ distinct
labels exist.
\end{lemma}
\begin{proof}
For a labelling $\alpha$ of the nodes, let $q_\alpha$ be the product of the
$m$ graph variables asserting those function evaluations. Expanding the
totality sum for each node, in the order of the acyclic list, gives
\[
                         \sum_\alpha q_\alpha-1
\]
in the degree-$m$ NS span. This is a telescoping identity: at step $j$, the
product for the first $j-1$ evaluations multiplies the totality equation
for evaluation $j$. Summing over the possible earlier values gives the
next partial sum. Repeated graph variables cause no problem.

If $\alpha$ gives different values to the same function on the same input,
$q_\alpha$ contains a functionality collision and already belongs to the
span within degree $m$. Otherwise, list the distinct relation variables
on the at most $m+b$ named labels. Write this actual set as
$Y_\alpha=\{y_1,\ldots,y_{a_\alpha}\}$, where
$a_\alpha\leq a_{\mathrm{rel}}=\sum_R(m+b)^{\operatorname{arity}(R)}$.
Expand over these distinct variables:
\[
             1=\prod_{i=1}^{a_\alpha}(y_i+(1-y_i))
              =\sum_\beta u_\beta.
\]
Each $u_\beta$ is the Boolean assignment monomial for one partial relational
structure. By hypothesis that structure falsifies an available clause.
Its clause polynomial is a product of at most $L$ falsity indicators.
Write $P=q_\alpha u_\beta$ and the clause polynomial as
$F_1\cdots F_\ell$, where each $F_j$ is the falsity indicator of a
literal false in the chosen partial structure. For each literal,
$P(1-F_j)$ has an explicit expression in the structural NS span:
\begin{itemize}
\item For a positive relation literal, $F_j=1-x$ and $u_\beta=(1-x)u'$.
Thus $P(1-F_j)=q_\alpha u'x(1-x)=-q_\alpha u'(x^2-x)$.
For a negative relation literal, $F_j=x$ and $u_\beta=xu'$; the same
identity with $x(1-x)$ applies.
\item For a negative graph literal false in the partial structure,
its graph variable $x$ is one of the chosen evaluations, so
$q_\alpha=xq'$. Its falsity indicator is $F_j=x$ and
$P(1-F_j)=-q'u_\beta(x^2-x)$.
\item For a positive graph literal false there, the chosen evaluation
has a different output variable $x'$ on the same function input.
Here $q_\alpha=x'q'$, $F_j=1-x$, and
$P(1-F_j)=q'u_\beta x'x$, an explicit functionality-axiom multiple.
\item A false equality or disequality has the constant falsity
indicator $F_j=1$, so $P(1-F_j)=0$.
\end{itemize}
These are all cases of the normalised literals. Each identity has
product degree at most $m+a_{\mathrm{rel}}+1$.
The exact telescoping identity
\[
 P(1-F_1\cdots F_\ell)
   =\sum_{j=1}^{\ell}P(1-F_j)\prod_{i<j}F_i
\]
therefore puts its left side in the NS span within degree
$m+a_{\mathrm{rel}}+\ell$. Since $PF_1\cdots F_\ell$ is a
multiple of the falsified clause axiom, $P$ is in that same span.
Sum first over $\beta$, then over $\alpha$, and use the totality identity
to obtain one.
\end{proof}


A \emph{partial structure} for $\theta_M$ has a domain of labels, all relations defined on it, and
function values defined only on some arguments. It is \emph{locally consistent} if every clause
instance whose variables take values in the domain, and whose function atoms are all defined,
has a true literal.

\begin{lemma}[Local switch]\label{tx:lem:local-switch}
Let the run of $M$ first reach a halting state at time $T$, and let $M$ have $c_M$ counters.
Suppose a locally consistent partial structure defines $o$, $e$, $a_{j+1}=S(a_j)$ for $j<T$
\textup{(}with $a_0=o$\textup{)}, $C_r(a_j)$ and $S(C_r(a_j))$ for $j\le T$, $r<c_M$, and that all these
elements lie in its domain. Then its domain has at most $T$ distinct elements.
\end{lemma}
\begin{proof}
Suppose it contains $T+1$ distinct elements. If $a_j=e$ for some $j<T$, take $j$ least. Every
domain element $x$ satisfies $x=a_0$ or $a_0<x$ (least element), and, if $a_i<x$ with $i<j$, then
$x=a_{i+1}$ or $a_{i+1}<x$ (no element strictly between $a_i$ and $S(a_i)$, and linearity). As
$a_j=e$ is greatest, the domain is $\{a_0,\ldots,a_j\}$, too small. So $a_j\ne e$ for $j<T$, and
the order clauses make $a_0<a_1<\cdots<a_T$. For the counter induction, the initial clauses give the initial
state and $C_r(a_0)=a_0$ for every counter. A counter value at time $j$
is at most $j$, so the notation $a_v$ uses only this named prefix.
Suppose its states and counters agree through time $j<T$.
An increment of value $v$ gives $S(a_v)=a_{v+1}$; because $v\le j<T$,
this successor is among the defined evaluations and its guard is valid.
The zero branch of a decrement is selected exactly when $v=0$.
On a positive decrement, the transition specifies a value $y<a_v$
with $S(y)=a_v$. The immediate-predecessor property forces
$y=a_{v-1}$: if $y<a_{v-1}$ then $a_{v-1}$ lies strictly between
$y$ and $S(y)$, and if $a_{v-1}<y<a_v$ then $y$ lies strictly
between $a_{v-1}$ and $S(a_{v-1})$. Both alternatives contradict
immediacy. Unchanged counters retain their named values. State
uniqueness and the chosen transition identify the next state.
Every clause used has its variables in the domain and its function
atoms in the specified finite evaluation list. This proves the
induction through time $T$. At $j=T$ the state is halting, which a clause forbids.
\end{proof}

\begin{corollary}\label{tx:cor:switch-degree}
In that case, for every $n\ge T+1$ the translation of $\theta_M$ at size $n$ has an NS refutation
of degree at most
\[
 C_M=m+\sum_R(m+T+1)^{\operatorname{arity}(R)}+\max_{C}|C|,\qquad m=T+2+2c_M(T+1).
\]
\end{corollary}
\begin{proof}
Apply Lemma~\ref{tx:lem:finite-certificate} to the $m$ evaluations $o$, $e$, $a_1,\ldots,a_T$,
$C_r(a_j)$, $S(C_r(a_j))$ ($j\le T$, $r<c_M$), each argument supplied by an earlier node, with
$b=T+1$ fixed labels. A consistent labelling and a Boolean assignment give a partial structure of
the kind in Lemma~\ref{tx:lem:local-switch}, whose domain contains the $T+1$ fixed labels; so it is not
locally consistent, which is the hypothesis of Lemma~\ref{tx:lem:finite-certificate}.
\end{proof}

Here $c_M$ is one more than the largest counter index occurring in $M$ (so $c_M\ge1$), and $T$ must be
the first halting time. Sizes $n\le T$ are genuine exceptions: $\theta_M$ can have models there.

\section{A disjunction with controlled degree}
The notation $d_*(\alpha,n)$ below always refers to the original
graph-and-sum translation of $\alpha$, with target $1$; the symbol $*$
stands for either NS or ordinary PC.
A signature is the collection of relation and function symbols.
Disjoint signatures are obtained by renaming the symbols of one branch,
so the two polynomial variable families are independent even though
they use the same universe. \emph{Structural equations} below mean the
Boolean, function-totality and function-functionality equations of the
compiler. The clause equations impose the sentence's further demands.
This distinction lets the inactive branch retain a structural root
while its clause equations are guarded away.

Conjunction is easy to control when one block is satisfiable. In disjoint
signatures, if $\beta$ has a root at size $n$, then
\begin{equation}\label{tx:eq:satisfiable-block}
 d_*(\alpha\wedge\beta,n)=d_*(\alpha,n),
 \qquad *=\NS,\PC.
\end{equation}
One inequality follows by using the refutation of $\alpha$; the other
follows by substituting the root of $\beta$ into every proof line.

We also need a disjunction which does not accidentally simplify either
branch. Introduce a fresh unary relation $Z$. Keep both branches' Boolean,
totality and functionality equations. Replace every clause polynomial $f$
of $\alpha$ by $z_t f$, and every clause polynomial $g$ of $\beta$ by
$(1-z_t)g$, for a fresh universally quantified label $t$. Denote this
normalised sentence by $\alpha\join\beta$.
Thus setting every $z_t$ to one exposes the $\alpha$ clauses, while
setting them to zero exposes the $\beta$ clauses. Keeping the structural
equations makes either specialisation an actual copy of the selected
branch's polynomial system.

The guard $z_t$ is itself a Boolean relation variable. Multiplying a
clause by it makes the clause active when $z_t=1$; multiplying by
$1-z_t$ selects the other value. The lower-bound specialisations set
\emph{all} guards to the same value. For the upper bound one label
$t_0$ suffices, since the universally instantiated inputs include the
clauses guarded by that label. The hypothesis $n\ge1$ guarantees that
such a label exists.

\begin{lemma}[Controlled join]\label{tx:lem:join}
For $n\ge1$, with $\infty$ allowed, each of NS and PC satisfies
\[
 \max\{d_*(\alpha,n),d_*(\beta,n)\}
 \le d_*(\alpha\join\beta,n)
 \le\max\{d_*(\alpha,n),d_*(\beta,n)\}+1.
\]
\end{lemma}
\begin{proof}
Set all selectors to one, and give the inactive block a structural root:
all its functions take one fixed value and all its relations are false.
Its clauses are guarded away, while its structural equations vanish.
The resulting proof refutes $\alpha$ without increasing degree. Selectors
zero give the other lower inequality.

For the upper bound fix a label $t_0$. Simulate each branch refutation
with its fixed guard $c=z_{t_0}$ or $c=1-z_{t_0}$, representing a branch
line $f$ by $cf$. A guarded clause is already an axiom. A structural
input $h$ is retained unguarded, so PC can form $z_{t_0}h$ and,
for the other guard, $h-z_{t_0}h$. Thereafter a linear combination of
branch lines is simulated by the same combination of their guarded
lines, and a step $f\mapsto x_vf$ is simulated by the legal variable
multiplication $cf\mapsto x_v(cf)$. This uses no unguarded intermediate
branch line and remains valid after cancellations. Every line has
degree at most one above the branch budget. The two simulations end
in $z_{t_0}$ and $1-z_{t_0}$, whose sum is one. For NS the same
construction multiplies the two certificates by their guards.
\end{proof}

\section{The reduction}
Let $\psi_1,\psi_2,\psi_3,\psi_4$ be the four sentences of
Section~\ref{tx:sec:examples}. Each has no finite nonempty model. For these particular
witnesses, the dominant type under the maximum operation is ordered by
\[
                 \mathrm T_1\prec\mathrm T_4\prec\mathrm T_3\prec\mathrm T_2.
\]
This is a property of the chosen witnesses, not an ordering assertion about
arbitrary profiles. A bounded profile has no effect on the type of a maximum;
the all-size linear lower bound for $\psi_2$ makes it dominant; and the
maximum of $\psi_4$ and $\psi_3$ is bounded off the multiples of $p$ and
linear on them. An additive error of at most one preserves these conclusions.
More generally, adjoining an eventually bounded profile and allowing this
error preserves each of the four intrinsic types: recurrence transfers
through finite intervals of degrees, and unboundedness transfers directly.

Choose indices $a,b$ with $\mathrm T_a\prec\mathrm T_b$ in this order
and define, using disjoint signatures,
\[
             \Phi_M=\psi_a\join(\psi_b\wedge\theta_M).
\]
If $M$ does not halt, \eqref{tx:eq:satisfiable-block} makes the second branch
have exactly the profile of $\psi_b$. The join has type $\mathrm T_b$. If $M$ halts,
the second branch is eventually bounded by Lemma~\ref{tx:lem:switch}, and the
join has type $\mathrm T_a$. Every $\Phi_M$ is within the all-finite-inconsistent
promise, whether or not $M$ halts. The map is primitive recursive and
uses no information about that alternative. A separator of types $\mathrm T_a,\mathrm T_b$
would decide halting. Choosing $a$ and $b$ on opposite sides of any
disjoint nonempty $A,B$ proves Theorem~\ref{tx:thm:undecidability}.

\section{Where the textbook argument ends}
The general consequence theorem has reduced the target question to a
support-bounded closure calculation, fulfilled its transfer contract by
one common projector, and read off the four profiles from fixed-degree
recurrence. The refutation examples and this reduction show that the
four alternatives are substantial and cannot be algorithmically
separated in general. The exercise solutions which follow return to
small calculations and help check that each logical distinction in this
route has been retained.

\part{Model answers and a reference map}

\chapter{Model answers and a notation guide}\label{tx:ch:solutions}
This chapter gives a model answer to every numbered exercise in the book.
Each question is repeated, followed by the calculation or proof needed
to justify the conclusion. Upper bounds and lower bounds are checked
separately when an exact minimum degree is requested.

The index links to all seventeen answers; the exercise number links
back to its original question. There is also a direct model-answer
link beside every original exercise. The notation guide follows the
answers. These are written pedagogical arguments, not additional
Lean-verified theorem endpoints.

\section*{Index of model answers}
\addcontentsline{toc}{section}{Index of model answers}
\begin{longtable}{@{}lp{0.62\textwidth}r@{}}
\toprule
Exercise & Topic & Answer page\\
\midrule

\ref{tx:ex:zero-constant} & \hyperref[tx:ans:zero-constant]{Empty input and the constant axiom} & \pageref{tx:ans:zero-constant}\\

\ref{tx:ex:small-gap} & \hyperref[tx:ans:small-gap]{A degree-two PC proof and a degree-three NS identity} & \pageref{tx:ans:small-gap}\\

\ref{tx:ex:Boolean-rule} & \hyperref[tx:ans:Boolean-rule]{Why Boolean reduction is not a degree-two PC step} & \pageref{tx:ans:Boolean-rule}\\

\ref{tx:ex:support-count} & \hyperref[tx:ans:support-count]{Monomial support and row-index support} & \pageref{tx:ans:support-count}\\

\ref{tx:ex:orbit-switch} & \hyperref[tx:ans:orbit-switch]{Fixed requests and arbitrary orbit switching} & \pageref{tx:ans:orbit-switch}\\

\ref{tx:ex:stable-floor} & \hyperref[tx:ans:stable-floor]{The safe target-stability rank} & \pageref{tx:ans:stable-floor}\\

\ref{tx:ex:closure-zero} & \hyperref[tx:ans:closure-zero]{Exactly which constants are derivable at degree zero} & \pageref{tx:ans:closure-zero}\\

\ref{tx:ex:closure-witness} & \hyperref[tx:ans:closure-witness]{A semantic consequence outside every PC closure} & \pageref{tx:ans:closure-witness}\\

\ref{tx:ex:dense-count} & \hyperref[tx:ans:dense-count]{Full coefficient lists and the squarefree count} & \pageref{tx:ans:dense-count}\\

\ref{tx:ex:period-arithmetic} & \hyperref[tx:ans:period-arithmetic]{The strict prime-power period and inverse onset} & \pageref{tx:ans:period-arithmetic}\\

\ref{tx:ex:forward-only} & \hyperref[tx:ans:forward-only]{Forward persistence does not preserve exact minima} & \pageref{tx:ans:forward-only}\\

\ref{tx:ex:valuation-profile} & \hyperref[tx:ans:valuation-profile]{Recurring levels, liminf and finite linear peaks} & \pageref{tx:ans:valuation-profile}\\

\ref{tx:ex:zero-T3} & \hyperref[tx:ans:zero-T3]{Where stability at degree zero is essential} & \pageref{tx:ans:zero-T3}\\

\ref{tx:ex:sum-versus-individual} & \hyperref[tx:ans:sum-versus-individual]{The invariant sum and the individual targets} & \pageref{tx:ans:sum-versus-individual}\\

\ref{tx:ex:fresh-small} & \hyperref[tx:ans:fresh-small]{The exact fresh-target degrees over any field} & \pageref{tx:ans:fresh-small}\\

\ref{tx:ex:fresh-window} & \hyperref[tx:ans:fresh-window]{Eligible lower windows and the exceptional size} & \pageref{tx:ans:fresh-window}\\

\ref{tx:ex:fresh-soundness} & \hyperref[tx:ans:fresh-soundness]{Extension-field soundness and the extra Boolean term} & \pageref{tx:ans:fresh-soundness}\\

\bottomrule
\end{longtable}
\clearpage

\section{Proof rules and minimum degrees}

\begin{samepage}

\subsection*{Model answer to Exercise~\ref{tx:ex:zero-constant}}\label{tx:ans:zero-constant}

\begin{quote}\small\textbf{Question.} Find the least NS and PC degrees of the targets $0$, $1$, and $x$ from
$\Gamma=\varnothing$ and from $\Gamma=\{1\}$. Explain the role of the
last-line degree floor.
\end{quote}

\paragraph{Model answer.}
The answer is the same in NS and ordinary PC:
\[
 \begin{array}{c|ccc}
 \Gamma & h_S(\Gamma\vdash0)&h_S(\Gamma\vdash1)&h_S(\Gamma\vdash x)\\
 \hline
 \varnothing&0&\infty&\infty\\
 \{1\}&0&0&1
 \end{array}
 \qquad(S\in\{\NS,\PC\}).
\]
\end{samepage}

From empty input, an NS identity has no nonzero summands. In PC the
only initial line available is zero; linear combinations and variable
multiplication keep it zero. Thus neither system can produce \(1\)
or \(x\). Zero has the empty NS certificate and a one-line PC proof,
both of degree zero.

From \(\{1\}\), the input itself proves \(1\) in degree zero.
For \(x\), NS uses the identity \(x=x\cdot1\); PC lists \(1\) and
then multiplies by \(x\). These proofs have degree one. The last-line
degree floor gives \(h_{\PC}\ge\deg x=1\), and
\(h_{\NS}\ge h_{\PC}\) gives the same lower bound for NS.
Thus the displayed upper bounds are exact. In particular, inconsistency
of the inputs does not make a nonconstant target a degree-zero consequence.


\noindent{\small\hyperref[tx:ex:zero-constant]{Return to Exercise~\ref*{tx:ex:zero-constant}.}}\par


\begin{samepage}

\subsection*{Model answer to Exercise~\ref{tx:ex:small-gap}}\label{tx:ans:small-gap}

\begin{quote}\small\textbf{Question.} For Example~\ref{tx:eg:cancellation}, prove
$h_{\PC}(\Gamma\vdash z)=2$ and
$h_{\NS}(\Gamma\vdash z)=3$. For the NS lower bound, list the monomials
which can occur in a degree-two certificate.
\end{quote}

\paragraph{Model answer.}
Work over \(\mathbb F_2\), with
\[
 f_1=xy+x,\qquad f_2=xy+1,\qquad f_3=xz.
\]
A complete PC derivation is
\[
 \begin{array}{c|l|l}
 \text{line}&\text{polynomial}&\text{justification}\\ \hline
 1&xy+x&f_1\\
 2&xy+1&f_2\\
 3&xz&f_3\\
 4&x+1&1+2\\
 5&xz+z&z\text{ times line }4\\
 6&z&5+3
 \end{array}
\]
Here the notation \(1+2\) means addition of the polynomials on those
two numbered lines. Every line has degree at most two. Multiplication
on line 5 is legal within that budget because line 4 has actual
degree one after cancellation.
\end{samepage}

No nonzero input is usable at degree at most one: all three inputs
have degree two. Starting from zero alone produces only zero.
Therefore \(h_{\PC}(\Gamma\vdash z)=2\).

For NS, expand the identity
\[
 z=zf_1+zf_2+f_3
   =(xyz+xz)+(xyz+z)+xz.
\]
Its product degree is three. In a degree-two certificate, each nonzero
multiplier must be constant, since each input has degree two and
nonzero product degrees add over a field. Such a certificate has form
\[
 af_1+bf_2+cf_3=(a+b)xy+ax+b+cxz.
\]
The only monomials which can occur are \(xy,x,1,xz\). The coefficient
of the monomial \(z\) is always zero, so this expression cannot equal
the target. Hence \(h_{\NS}(\Gamma\vdash z)=3\).


\noindent{\small\hyperref[tx:ex:small-gap]{Return to Exercise~\ref*{tx:ex:small-gap}.}}\par


\begin{samepage}

\subsection*{Model answer to Exercise~\ref{tx:ex:Boolean-rule}}\label{tx:ans:Boolean-rule}

\begin{quote}\small\textbf{Question.} A proposed degree-two PC calculation multiplies the derived line $xy$
by $x$ and replaces $x^2y$ by $xy$, using a Boolean equation. Explain
why this is not an ordinary degree-two step, even if $x^2-x$ is an input.
\end{quote}

\paragraph{Model answer.}
Ordinary PC works with actual polynomials before reduction. The
proposed multiplication produces
\[
                         x\cdot xy=x^2y,
\]
which has degree three. That intermediate line already violates the
degree-two budget.
\end{samepage}

The Boolean input would permit the correction
\[
                 y(x^2-x)=x^2y-xy,
\]
but multiplying that input by \(y\) also produces a degree-three line.
Subtracting the correction would eventually give \(xy\), but the
degree of the proof counts both degree-three intermediate lines.

In the quotient by the Boolean equation the classes of \(x^2y\)
and \(xy\) agree. Ordinary PC has no rule which performs multiplication
and quotient reduction together as one degree-two inference.
This identifies the error in the proposed step; it does not exclude
some different legal degree-two derivation in a system with additional axioms.


\noindent{\small\hyperref[tx:ex:Boolean-rule]{Return to Exercise~\ref*{tx:ex:Boolean-rule}.}}\par


\section{Uniformity and support}

\begin{samepage}

\subsection*{Model answer to Exercise~\ref{tx:ex:support-count}}\label{tx:ans:support-count}

\begin{quote}\small\textbf{Question.} For variables $x_{ij}$, find the label supports and total degrees of
$x_{ij}^3$ and $x_{ij}x_{jk}$ when $i,j,k$ are distinct. Compare the
actual supports with the general bound $2d$. Explain why the row sum
in Example~\ref{tx:eg:row-support} has an index of support one.
\end{quote}

\paragraph{Model answer.}
Each pair-indexed variable uses at most two labels.
With \(i,j,k\) distinct, the requested counts are
\[
 \begin{array}{c|c|c|c}
 \text{monomial}&\text{degree }d&\text{label support}&2d\\ \hline
 x_{ij}^3&3&\{i,j\}\text{ (size }2)&6\\
 x_{ij}x_{jk}&2&\{i,j,k\}\text{ (size }3)&4
 \end{array}
\]
A power repeats the same variable and introduces no extra labels.
In the second monomial, the two variable supports overlap at \(j\).
These are the reasons the general union bound \(2d\) is not an equality.
\end{samepage}

For the row axiom \(f_{i,n}=\sum_{j<n}x_{ij}-1\), the axiom
index is \(i\), with support \(\{i\}\). If a permutation \(\pi\)
fixes \(i\), then
\[
 \pi f_{i,n}=\sum_{j<n}x_{i,\pi(j)}-1=f_{i,n},
\]
because \(\pi\) merely reorders the summands. The whole row sum can
mention every label while still having a one-label axiom index.
Its individual monomials and its axiom index are different basis objects.


\noindent{\small\hyperref[tx:ex:support-count]{Return to Exercise~\ref*{tx:ex:support-count}.}}\par


\begin{samepage}

\subsection*{Model answer to Exercise~\ref{tx:ex:orbit-switch}}\label{tx:ans:orbit-switch}

\begin{quote}\small\textbf{Question.} Take no axioms, a single variable $x$ of empty support, and two target
indices of empty support requesting $0$ and $x$. Determine the two
fixed-orbit profiles and the all-target profile. What happens if one
chooses the first target at sizes in an arbitrary set and the second
elsewhere?
\end{quote}

\paragraph{Model answer.}
Let \(h_0\) be the profile for the fixed target \(0\), \(h_x\) that
for the fixed target \(x\), and \(h_{\mathrm{all}}\) the profile for
the two-target request. At every rank,
\[
            h_0(n)=0,\qquad h_x(n)=\infty,\qquad
            h_{\mathrm{all}}(n)=\max\{0,\infty\}=\infty.
\]
These statements hold separately for NS and PC. The zero certificate
or zero line proves \(0\), whereas empty input cannot derive any
nonzero polynomial.
\end{samepage}

If the selected target is \(0\) at ranks in \(A\) and \(x\) elsewhere,
the selected profile becomes
\[
 h_A(n)=
 \begin{cases}0,&n\in A,\\ \infty,&n\notin A.\end{cases}
\]
Thus its degree-zero threshold can encode an arbitrary set of ranks.
The two available target sorts have stable coefficient rules, but
this rank-dependent selection is neither one fixed global orbit nor
the fixed all-target request. It therefore falls outside the theorem's
request convention. In particular, arbitrary orbit switching cannot
be used to infer another case of the four-profile classification.


\noindent{\small\hyperref[tx:ex:orbit-switch]{Return to Exercise~\ref*{tx:ex:orbit-switch}.}}\par


\begin{samepage}

\subsection*{Model answer to Exercise~\ref{tx:ex:stable-floor}}\label{tx:ans:stable-floor}

\begin{quote}\small\textbf{Question.} Take variables $x_i$, no axioms, and targets $x_i-x_j$ indexed by
ordered distinct pairs. Use $k=1,s=0,t=2,D_G=1,n_0=1$ to compute $n_*$.
At which positive ranks does this request actually contain a target?
Why is the stated stable rank allowed to be larger than necessary?
\end{quote}

\paragraph{Model answer.}
The prescribed parameters give
\[
          n_*=\max\{n_0,1,kD_G+t\}
              =\max\{1,1,1+2\}=3.
\]
At rank \(1\) there is no ordered pair of distinct labels, so the
request is empty and its degree is zero in both systems.
At every rank \(n\ge2\), such pairs exist and every target
\(x_i-x_j\) is a nonzero degree-one polynomial, including in
characteristic two. There are no axioms from which to derive it.
Consequently
\[
       h_S(1)=0,\qquad h_S(n)=\infty\quad(n\ge2).
\]
\end{samepage}

The actual nonempty request and its degree floor stabilise at rank
two. The theorem's rank three is a sufficient uniform bound obtained
by adding a monomial-support bound and an index-support bound.
Here those supports overlap, so the general bound is conservative.
A safe theorem onset need not be the smallest onset for each example.


\noindent{\small\hyperref[tx:ex:stable-floor]{Return to Exercise~\ref*{tx:ex:stable-floor}.}}\par


\section{The closure and its coordinates}

\begin{samepage}

\subsection*{Model answer to Exercise~\ref{tx:ex:closure-zero}}\label{tx:ans:closure-zero}

\begin{quote}\small\textbf{Question.} Describe $I_0$ and $C_0$ for an arbitrary system. Prove directly that a
nonzero constant target has a degree-zero proof exactly when an input
is a nonzero constant.
\end{quote}

\paragraph{Model answer.}
Let \(K\cdot1\) denote the one-dimensional space of constants.
Then
\[
 I_0=C_0=
 \begin{cases}
   \{0\},&\Gamma\text{ contains no nonzero constant},\\
   K\cdot1,&\Gamma\text{ contains a nonzero constant}.
 \end{cases}
\]
At degree zero, no nonconstant input can be listed in PC. In NS,
a nonzero product has degree zero only when both its multiplier
and its input are constants. Zero inputs contribute nothing.
Thus both systems start with the span of the constant inputs.
\end{samepage}

PC multiplication of a nonzero constant by a variable raises its
degree to one, so it contributes no new degree-zero vector. Multiplying
zero still gives zero. Hence the starting span is the whole closure
at degree zero, and it equals the NS span.

If \(c\in\Gamma\) is a nonzero constant and the target is a constant
\(a\ne0\), the identity \(a=(a/c)c\) is an NS certificate of degree
zero and also a legal PC scalar combination. Conversely, a nonzero
degree-zero consequence requires a nonzero constant in that initial
span, so at least one such input must be present.


\noindent{\small\hyperref[tx:ex:closure-zero]{Return to Exercise~\ref*{tx:ex:closure-zero}.}}\par


\begin{samepage}

\subsection*{Model answer to Exercise~\ref{tx:ex:closure-witness}}\label{tx:ans:closure-witness}

\begin{quote}\small\textbf{Question.} With input $\Gamma=\{x^2\}$, show that every vector in every closure
$C_d$ has zero coefficient of $x$. Deduce that $x$ is underivable,
although it vanishes at every field-valued solution of $\Gamma$.
\end{quote}

\paragraph{Model answer.}
Every input belongs to the ideal \((x^2)\). Linear combinations and
variable multiplication preserve that ideal, so an induction on PC
lines shows that every derived polynomial is divisible by \(x^2\).
Its coefficient of the monomial \(x\) is therefore zero.
\end{samepage}

In the one-variable ambient ring \(K[x]\), the closure can be written
explicitly:
\[
 C_d=
 \begin{cases}
 \{0\},&d<2,\\
 \Span_K\{x^2,x^3,\ldots,x^d\},&d\ge2.
 \end{cases}
\]
The forward inclusion follows from divisibility and the degree
cutoff. The reverse inclusion follows by repeatedly multiplying the
input \(x^2\) by \(x\) and then taking linear combinations.
If the ambient ring has additional variables, the divisibility
argument still applies, with more multiples of \(x^2\) in the closure.

The target \(x\) belongs to none of these spaces; its least PC
degree is infinity. It has no NS identity either, since \(x\notin(x^2)\).
On the other hand, a field has no nonzero element whose square is
zero. Every field-valued solution of \(x^2=0\) consequently satisfies
\(x=0\). Vanishing at all roots and algebraic derivability are distinct.


\noindent{\small\hyperref[tx:ex:closure-witness]{Return to Exercise~\ref*{tx:ex:closure-witness}.}}\par


\begin{samepage}

\subsection*{Model answer to Exercise~\ref{tx:ex:dense-count}}\label{tx:ans:dense-count}

\begin{quote}\small\textbf{Question.} List the monomials counted by $B_2(2)$, and compare the full coefficient
list of $x^2$ with its sparse representation. For $v\ge d$, identify
explicitly the $2^d$ entries used in the final lower bound of
Lemma~\ref{tx:lem:dense-size}.
\end{quote}

\paragraph{Model answer.}
With the variable order \(x,y\), one ordering of the degree-at-most-two
monomials is
\[
                       1,\ x,\ y,\ x^2,\ xy,\ y^2.
\]
Thus \(B_2(2)=\binom{2+2}{2}=6\). In that ordering the full coefficient
list of \(x^2\) is \((0,0,0,1,0,0)\). A sparse representation instead
records only the coefficient \(1\) and the exponent pair \((2,0)\).
These representations charge for different data.
\end{samepage}

For \(v\ge d\), choose distinct variables \(x_1,\ldots,x_d\).
Every subset \(S\subseteq\{1,\ldots,d\}\) gives a squarefree monomial
\[
                            \prod_{i\in S}x_i.
\]
Its degree is \(|S|\le d\), so it has an entry in the full coefficient
list. Distinct subsets have different exponent vectors and give
distinct monomials. There are \(2^d\) subsets, including the empty
subset, whose monomial is \(1\). This proves \(B_v(d)\ge2^d\);
the argument also includes \(d=0\), where there is one entry.


\noindent{\small\hyperref[tx:ex:dense-count]{Return to Exercise~\ref*{tx:ex:dense-count}.}}\par


\section{Periodicity, recurrence, and peaks}

\begin{samepage}

\subsection*{Model answer to Exercise~\ref{tx:ex:period-arithmetic}}\label{tx:ans:period-arithmetic}

\begin{quote}\small\textbf{Question.} Let $p=2$, $\kappa=4$, $\sigma=4$, and $n_*=7$. Compute $N(0),Q(0)$,
$N(3),Q(3)$ and $\beta(62)$. Check $N(\beta(62))\le62$ and explain
why the strict inequality defining $Q$ matters.
\end{quote}

\paragraph{Model answer.}
Here \(R(d)=4d+4\). Degree zero uses its separate definition:
\[
                            N(0)=7,\qquad Q(0)=1.
\]
It is not obtained by substituting zero into the positive-degree
period formula.
\end{samepage}

At degree three,
\[
 \begin{aligned}
 R(3)&=4\cdot3+4=16,\\
 N(3)&=\max\{7,2\cdot16+1\}=33,\\
 Q(3)&=\min\{2^a:2^a>16\}=32.
 \end{aligned}
\]
Although \(16\) is a power of two, it equals the support budget and
does not satisfy the strict inequality. Taking it as \(Q(3)\)
would use an unproved transfer step.

For the inverse onset, the numerator is positive, so
\[
 \begin{aligned}
 \beta(62)
   &=\left\lfloor\frac{(62-2\cdot4-1)_+}{2\cdot4}\right\rfloor
     =\left\lfloor\frac{53}{8}\right\rfloor=6,\\
 N(\beta(62))
   &=N(6)=\max\{7,2(4\cdot6+4)+1\}=57\le62.
 \end{aligned}
\]
For comparison, \(N(7)=65>62\). The value \(\beta(62)=6\)
locates the largest integer threshold guaranteed by this displayed
onset formula. It makes no claim that a sharper theorem could not
give a larger admissible threshold.


\noindent{\small\hyperref[tx:ex:period-arithmetic]{Return to Exercise~\ref*{tx:ex:period-arithmetic}.}}\par


\begin{samepage}

\subsection*{Model answer to Exercise~\ref{tx:ex:forward-only}}\label{tx:ans:forward-only}

\begin{quote}\small\textbf{Question.} For $g(1)=2$ and $g(n)=1$ for $n\ge2$, verify forward persistence of
every threshold under $n\mapsto n+1$. Identify a reverse threshold
implication that fails. Explain why exact recurrence of the starting
value cannot be deduced from forward persistence alone.
\end{quote}

\paragraph{Model answer.}
For nonnegative integer thresholds \(d\), all cases are
\[
 \begin{array}{c|c}
 d&\{n\ge1:g(n)\le d\}\\ \hline
 0&\varnothing\\
 1&\{2,3,4,\ldots\}\\
 d\ge2&\{1,2,3,\ldots\}.
 \end{array}
\]
Each set is closed under \(n\mapsto n+1\), proving forward
persistence of every threshold. A predicate can become true at the
next rank while still having this property.
\end{samepage}

At threshold one the reverse implication fails:
\[
                 g(2)=1\le1,\qquad g(1)=2>1.
\]
The starting exact value \(2\) occurs only once, whereas value \(1\)
occurs at every later rank. Forward persistence preserves the upper
bound \(g(n)\le2\) but does not exclude a drop to a smaller minimum.
To preserve the exact value two, one would also need reverse
comparison at threshold one. This is why the recurrence argument
uses two-way transfer at lower thresholds as well as at the chosen budget.


\noindent{\small\hyperref[tx:ex:forward-only]{Return to Exercise~\ref*{tx:ex:forward-only}.}}\par


\begin{samepage}

\subsection*{Model answer to Exercise~\ref{tx:ex:valuation-profile}}\label{tx:ans:valuation-profile}

\begin{quote}\small\textbf{Question.} For $h(n)=p^{v_p(n)}$, determine every recurring value, its liminf,
and a sequence of finite linear peaks. Explain why no eventual linear
lower bound holds at all sizes.
\end{quote}

\paragraph{Model answer.}
Every value of \(h(n)=p^{v_p(n)}\) is a power \(p^a\), with \(a\ge0\).
For each fixed \(a\), the ranks
\[
                     n_j=p^a(1+pj),\qquad j=0,1,2,\ldots
\]
have valuation \(a\), since \(p\nmid1+pj\). They tend to infinity
and all satisfy \(h(n_j)=p^a\). Thus the recurring finite values
are exactly
\[
                         E(h)=\{1,p,p^2,p^3,\ldots\}.
\]
\end{samepage}

Every value is at least one. At the unbounded sequence of ranks
\(n=1+pj\) the value is one, giving
\[
                            \liminf_{n\to\infty}h(n)=1.
\]
For finite linear peaks take \(n_a=p^a\). Then
\[
                          h(n_a)=p^a=n_a\longrightarrow\infty.
\]
Each peak is finite and lies exactly on the line \(h=n\).

There is no eventual all-size lower bound \(h(n)\ge cn\) with \(c>0\).
Given \(c\) and any proposed starting rank, choose \(n=1+pj\) above
that rank and above \(1/c\). Then \(h(n)=1<cn\).
This exercise concerns the displayed abstract sequence. It does not
assert that the normalised weak-bijection system has these exact
degrees or that its rectangle lower bound applies at \(n=p^a\).


\noindent{\small\hyperref[tx:ex:valuation-profile]{Return to Exercise~\ref*{tx:ex:valuation-profile}.}}\par


\begin{samepage}

\subsection*{Model answer to Exercise~\ref{tx:ex:zero-T3}}\label{tx:ans:zero-T3}

\begin{quote}\small\textbf{Question.} Show that T3 cannot have $E(h)=\{0\}$. State exactly where stability of
the degree-zero predicate is used.
\end{quote}

\paragraph{Model answer.}
If \(E(h)=\{0\}\), zero recurs at arbitrarily large ranks, in
particular at some \(n\ge n_*\). The degree-zero threshold is
constant beyond \(n_*\). Since it is true at that rank, it must
be true at every rank \(m\ge n_*\):
\[
                              h(m)\le0.
\]
All degrees are nonnegative, so \(h(m)=0\) throughout that range.
The profile is eventually bounded and is T1, contrary to the
unboundedness required in T3.
\end{samepage}

The essential use of zero-degree stability is the implication
that truth at one eligible rank gives truth at every eligible rank.
Mere periodicity with a larger period would not suffice. For example,
the arbitrary sequence which is zero at even ranks and equals \(n\)
at odd ranks is unbounded and has only zero as an infinitely recurring
finite value. Its zero predicate is not eventually constant, so it
is excluded by the theorem's zero-degree input.

Equivalently, the actual zero-degree period is \(Q(0)=1\).
There is no nonempty proper easy-residue set modulo one from which
to build a T3 split with only the easy value zero.


\noindent{\small\hyperref[tx:ex:zero-T3]{Return to Exercise~\ref*{tx:ex:zero-T3}.}}\par


\section{Consequences and concrete examples}

\begin{samepage}

\subsection*{Model answer to Exercise~\ref{tx:ex:sum-versus-individual}}\label{tx:ans:sum-versus-individual}

\begin{quote}\small\textbf{Question.} For the consistent family of Proposition~\ref{tx:prop:consistent-T3},
compute the two profiles at sizes $1,2,3,4,5,6$ in characteristic three.
Compare the invariant sum request with the indexed request for every
individual variable. Which statements have finite proof degrees?
\end{quote}

\paragraph{Model answer.}
In characteristic three, set
\[
 \Gamma_n=\{x_i-1:i<n\},\qquad G_n=\sum_{i<n}x_i.
\]
Write \(h_{\mathrm{sum}}\) for the invariant target profile and
\(h_{\mathrm{ind}}\) for the request for every target \(x_i\).
The table applies to both NS and ordinary PC:
\[
 \begin{array}{c|cccccc}
 n&1&2&3&4&5&6\\ \hline
 h_{\mathrm{sum}}(n)&\infty&\infty&1&\infty&\infty&1\\
 h_{\mathrm{ind}}(n)&\infty&\infty&\infty&\infty&\infty&\infty
 \end{array}
\]
\end{samepage}

Summing the inputs gives \(G_n-n\). At \(n=3,6\), the scalar \(n\)
is zero in the field, so
\[
                            G_n=\sum_{i<n}(x_i-1).
\]
This is an NS certificate of degree one and gives a PC proof of
the same degree. The target is a nonzero linear polynomial, so its
degree floor proves that the minimum is exactly one.

For the other four ranks, the assignment \(x_i=1\) satisfies all
inputs but has \(G_n=n\ne0\) in the field. Soundness excludes any
derivation. The same assignment has \(x_i=1\ne0\) for every
individual target at every positive rank, excluding the indexed
request in all six cases.
The finite proofs in this table are precisely those for the sum at
ranks three and six. A sum being derivable does not make its summands derivable.


\noindent{\small\hyperref[tx:ex:sum-versus-individual]{Return to Exercise~\ref*{tx:ex:sum-versus-individual}.}}\par


\begin{samepage}

\subsection*{Model answer to Exercise~\ref{tx:ex:fresh-small}}\label{tx:ans:fresh-small}

\begin{quote}\small\textbf{Question.} For $\Gamma=\{x,1-x\}$ over an arbitrary field, write a degree-one PC
proof of the fresh target $y$ and prove that its least NS degree is two.
Do not insert an axiom involving $y$.
\end{quote}

\paragraph{Model answer.}
A legal PC proof over any field is
\[
 \begin{array}{c|l|l}
 1&x&\text{input}\\
 2&1-x&\text{input}\\
 3&1&1+2\\
 4&y&y\text{ times line }3.
 \end{array}
\]
All lines have degree at most one. The target degree is one, so
\(h_{\PC}(\Gamma^\uparrow\vdash y)=1\).
\end{samepage}

For NS,
\[
                            y=yx+y(1-x).
\]
Both nonzero products have degree two, giving an upper bound two.
In a degree-one certificate, each multiplier of the degree-one
inputs must be constant. Any such sum is
\[
                       ax+b(1-x)=(a-b)x+b
                       \qquad(a,b\in K),
\]
which has no \(y\)-term. It cannot equal \(y\).
Thus \(h_{\NS}(\Gamma^\uparrow\vdash y)=2\).

PC first derives the constant \(1\), then multiplies only that last
line. In NS the extra \(y\) is carried in the multipliers of the
original inputs, increasing each product degree. No input involving
\(y\), including \(y^2-y\), has been introduced.


\noindent{\small\hyperref[tx:ex:fresh-small]{Return to Exercise~\ref*{tx:ex:fresh-small}.}}\par


\begin{samepage}

\subsection*{Model answer to Exercise~\ref{tx:ex:fresh-window}}\label{tx:ans:fresh-window}

\begin{quote}\small\textbf{Question.} In characteristic three, apply the fresh-$y$ bounds to the normalised
weak-bijection family at $n=6$ and $n=36$. What can the stated bounds
say at $n=9$? Distinguish the lower-bound size condition from the
all-size upper certificate.
\end{quote}

\paragraph{Model answer.}
Use \(q=3^{v_3(n)}\). The numerical data are
\[
 \begin{array}{c|c|c|c|c}
 n&q&n\ge2q?&h_{\PC}\text{ bounds}&h_{\NS}\text{ bounds}\\ \hline
 6&3&\text{yes}&3\le h_{\PC}\le7&4\le h_{\NS}\le8\\
 36&9&\text{yes}&9\le h_{\PC}\le13&10\le h_{\NS}\le14\\
 9&9&\text{no}&h_{\PC}\le13&h_{\NS}\le14
 \end{array}
\]
Indeed \(6=3\cdot2\) and \(36=3^2\cdot4\), so their valuations
are one and two, respectively. Both satisfy the rectangle size
condition. The original refutation bounds and the fresh-target
identities give
\[
 \begin{aligned}
 q\le d_{\PC}(\widehat B_n)\le q+4
   &\ \Longrightarrow\ q\le h_{\PC}\le q+4,\\
 q\le d_{\NS}(\widehat B_n)\le q+4
   &\ \Longrightarrow\ q+1\le h_{\NS}\le q+5.
 \end{aligned}
\]
Here \(q\ge1\), so the PC maximum with one changes nothing.
\end{samepage}

At \(n=9\), the valuation is still two, but \(9<2q=18\).
The quoted rectangle lower window supplies neither \(h_{\PC}\ge9\)
nor \(h_{\NS}\ge10\). The all-size NS upper certificate, followed
by the fresh-target construction, still supplies the displayed
upper bounds. The target floor gives \(h_{\PC},h_{\NS}\ge1\),
and \(h_{\PC}\le h_{\NS}\) remains valid at every size.
The displayed windows are bounds, not assertions of exact minimum degrees.


\noindent{\small\hyperref[tx:ex:fresh-window]{Return to Exercise~\ref*{tx:ex:fresh-window}.}}\par


\begin{samepage}

\subsection*{Model answer to Exercise~\ref{tx:ex:fresh-soundness}}\label{tx:ans:fresh-soundness}

\begin{quote}\small\textbf{Question.} Suppose $\Gamma$ has a root over some extension field of $K$. Explain
why the fresh target $y$ is underivable from $\Gamma^\uparrow$.
If one instead adds $y^2-y$ as an axiom, identify the extra term which
appears when extracting the coefficient of $y^1$ from an NS identity.
Why does the original coefficient-extraction proof no longer apply
unchanged?
\end{quote}

\paragraph{Model answer.}
Let \(a\) be a root of the old inputs over \(E\supseteq K\).
Extend it to the fresh variable by assigning \(y=1\).
Every old axiom is independent of \(y\), so still evaluates to zero,
whereas the target evaluates to one.
\end{samepage}

An NS identity for \(y\), even with coefficients in \(E\), would
evaluate to the impossible equality \(1=0\). In PC, each input,
linear combination and variable multiple evaluates to zero by
induction, so a last line evaluating to one is impossible.
Thus \(y\) is underivable over \(K\) and over \(E\).

If \(y^2-y\) is instead added as an axiom, a proposed NS identity has
the form
\[
                 y=\sum_f a_f(x,y)f(x)+b(x,y)(y^2-y).
\]
For each old input,
\([y^1](a_f f)=([y^1]a_f)f\). For the new summand,
\[
 [y^1]\bigl(b(x,y)y^2\bigr)=0,\qquad
 [y^1]\bigl(-b(x,y)y\bigr)=-b(x,0).
\]
Taking the coefficient of \(y\) therefore gives
\[
                       1=\sum_f([y^1]a_f)f-b(x,0).
\]
The extra term \(-b(x,0)\) need not be a multiple of an old input.
Hence this is not the old refutation certificate obtained when no
\(y\)-axiom was present. The original one-degree-saving extraction
argument cannot be reused unchanged.

The soundness obstruction above does survive addition of the Boolean
axiom: the extended root with \(y=1\) satisfies \(y^2-y=0\) as well.
The obstruction and the exact-degree certificate argument have
different hypotheses. This calculation does not assert a new
degree identity for inconsistent inputs after adding that axiom.


\noindent{\small\hyperref[tx:ex:fresh-soundness]{Return to Exercise~\ref*{tx:ex:fresh-soundness}.}}\par


\section{Notation guide}
The entries below give the local meaning of frequently reused symbols
and a few terms needed in the later chapters. A degree budget bounds
actual proof lines or NS products, while a support budget bounds labels
in one basis object. Neither is the universe size. Terms such as design
and witness refer to linear tests; they need not describe a solution of
the equations.

\begin{description}[style=nextline,leftmargin=1em,labelindent=0pt,labelwidth=0pt]
\item[$\Gamma_n$] The input polynomials at universe size $n$, all
interpreted as equations equal to zero.
\item[$G_{b,n}$] The requested target polynomial with target index $b$.
The target is not an additional axiom.
\item[$h_S(n)$] The maximum of the least degrees for the requested
targets; infinity means at least one target is underivable. An empty
request has profile zero.
\item[$d_S(\Gamma_n)$] The special minimum degree for deriving $1$.
For a translated sentence $\psi$, this is also written $d_S(\psi,n)$.
\item[$k,s,t$] Bounds on supports of one variable, one axiom index,
and one target index. They do not bound the labels in a whole sum.
\item[$d,D,D_G$] A proposed proof-degree bound, a fixed input-degree
bound in the template, and a fixed target-degree bound, respectively.
In the T3 argument, the locally specified $D=\max E(h)$ instead denotes
the largest recurring degree; in the matching chapter, $D$ is a proof budget.
\item[$P_{n,\le d},P_b(N)$] The ordinary polynomial space of degree at
most $d$, and the tuple space $K[\Inj([b],[N])]$, respectively.
The integer $b$ counts tuple positions, while $d$ bounds polynomial degree.
\item[$I_d,C_d$] The NS span of axiom products of degree at most $d$,
and the stable ordinary-PC consequence closure in $P_{n,\le d}$.
\item[$W,W_N,W_d(n)$] $W$ is a locally specified ambient space or,
in an exact sequence, its quotient term. $W_N$ denotes a degree-cutoff
polynomial ambient, whereas $W_d(n)$ in the model-theory chapter is
the chosen NS or PC consequence subspace. In commutator calculations,
$W=[D,z]$ is instead an invertible operator. The displayed type fixes
which meaning is intended.
\item[$C$] A chosen complement or chain-start space in the layer
construction; also a fresh label set in the quadratic construction
and a column set on a rook board. In the injection certificate,
$C=\sum_j c_j$ is a polynomial. These local meanings are distinct
from the PC closure $C_d$ and the annihilator class $\mathcal C_r$.
\item[$r=R(d),\kappa,\sigma$] The support budget $R(d)=\kappa d+\sigma$.
For indexed targets $(\kappa,\sigma)=(2k,s+t)$; for invariant targets
they are $(k,s)$.
In the rectangular matching system $B_{r,m}$, $r$ instead counts
rows; the rook-forcing proof specifies its row set $R$ separately.
\item[$n,N=n+q$] A retained universe size and the larger size before
a block of labels is deleted.
\item[$z,u,w,e$] The block translation operator, its top power
$u=z^{p-1}$, a common group-algebra witness with $uwu=u$, and the
resulting projector. The witness identity holds on tuple spaces of
length at most $r$. Products and powers of these operators mean
composition; $z^p=0$ means that $p$ successive applications give zero.
In Lemma~\ref{tx:lem:detect}, $u$ instead names an integer support
bound. In $n=p^au$ it is the factor not divisible by $p$.
The scalar grade $e\in\mathbb F_p$ in the layer splitting is also
distinct from the projector $e$.
\item[$\Del,\Lift$] Deletion is the quotient
$\ker z/(\ker z\cap\im z)$. The lift $\Lift=e\iota$ first includes
a smaller-rank polynomial and then applies the common projector.
\item[$V',\mathcal Z_b,T_b$] In $V=P_b(N)$, marginal maps $\partial_i$
sum over one tuple position. The incoming space is
$V'=\sum_i\im\partial_i^{\mathsf T}$, the zero-marginal space is
$\mathcal Z_b=\bigcap_i\ker\partial_i=(V')^\perp$, and the top quotient
is $T_b=V/V'$.
\item[$J_r(N),\mathcal C_r(N)$] The common group-algebra annihilator of
all tuple spaces of length at most $r$, and the modules killed by that
ideal. The latter class is also called $\mathcal M_r$; it is distinct
from the PC consequence space $C_d$.
\item[$J_N$] The actual target map $[b]\mapsto G_{b,N}$, distinct from
the annihilator ideal $J_r(N)$.
\item[$p,q,Q(d)$] The prime characteristic, an admissible block size
$p^a>R(d)$, and, for $d\ge1$, the chosen least such power for the
periodicity bound. The separate zero-degree convention is $Q(0)=1$.
The actual period may be smaller than $Q(d)$.
\item[$n_0,n_*,N(d)$] The axiom stability threshold, the stable rank
including the target-degree floor, and the periodicity onset.
For $d\ge1$, $N(d)=\max\{n_*,2R(d)+1\}$; separately, $N(0)=n_*$.
\item[$E(h)$] The finite degrees recurring at arbitrarily large ranks.
\item[$A_d(h),\mathcal E_d(h),\mathcal H_p(h)$] The eventually easy
residues, their clopen union in $\mathbb Z_p$, and the closed complement
of all such easy regions. These are not the matching-design space
$E_d(N,M)$ or the finite-degree set $E(h)$.
\item[Subquotient; annihilator] A quotient $S'/S$ of an invariant
subspace by a smaller invariant subspace. A linear functional
annihilates $S$ when it vanishes on $S$; $S^\perp$ expresses this
condition using the dot product. The group-algebra annihilator $J_r(N)$
instead consists of operators killing the specified tuple spaces.
\item[Plateau] An equality of consecutive intersections
$I_{d+j}(n)\cap P_{n,\le d}=I_{d+j+1}(n)\cap P_{n,\le d}$.
It supplies closure under the legal degree-$d$ PC multiplications;
it does not say that a truncated consequence space is an ideal.
\item[{$[y^1]a, v_p(n)$}] The coefficient of $y$ in a polynomial $a$,
and the exponent of $p$ in the positive integer $n$, respectively.
A fresh $y$ is an independent indeterminate absent from the old axioms.
\item[$B_{r,m},E_d(N,M),\Lambda_\psi$] The rectangular bijection
system, its degree-$d$ space of matching designs, and the associated
linear functional on ordinary polynomials. Here $E_d$ is unrelated
to the recurring-degree set $E(h)$.
\item[Filtration; extension] Nested degree windows. A design extension
supplies the next window's values while preserving all values in the
previous window.
\item[Definable; $\chi$] Described by a formula in the permitted
structure and language; $\chi$ is the Euler counting characteristic
of its definable sets, not the prime characteristic of the field.
\item[Promise; reduction] The hypotheses on inputs for which an
algorithm must halt correctly; an effective transformation of finite
input descriptions which transfers a decision problem to another.
\end{description}


\section{Connections worth following in the literature}
\label{tx:sec:literature-route}
The following connections have different roles. Some are historical
antecedents, some supply a useful language, and some suggest analogies.
They are not interchangeable external proofs of the present theorem.

\paragraph{Proof complexity and the earlier four-way classification.}
The NS and PC systems originate in the work of Beame, Impagliazzo,
Kraj\'i\v cek, Pitassi and Pudl\'ak~\cite{NSFoundations}, and of Clegg,
Edmonds and Impagliazzo~\cite{PCFoundations}. Riis's LICS 2008 paper
and 2009 report~\cite{Riis,RiisLong} provide the four-way refutation
classification. The present linear quantitative estimates and uniform
consequence treatment should be read against those earlier results.
The tree-resolution gap~\cite{RiisTree} is a related classification
phenomenon, not a theorem about polynomial degree. Its infinite-model
criterion should not be silently transferred to NS or PC.

\paragraph{Representation stability: similar objects, stronger transport needs.}
The characteristic-zero uniform permutation-module work of Riis and
Sitharam~\cite{RiisSitharamModules}, the FI-module framework of Church,
Ellenberg and Farb~\cite{CEF}, and Harman's modular periodicity
\cite{Harman} place fixed-support families in a broader setting.
Chapter~\ref{tx:sec:stability-context} explains the actual comparisons,
including why the uniform sum $\sum_{i<n}x_i$ does not behave as an
FI-submodule under naive inclusion. Hone and coauthors~\cite{HKPSYZ}
provide an explicit categorical-functor perspective. Compatibility with
natural operations already appears in Harman's results. What must be
checked here is preservation and reflection of the specific polynomial
closure and target maps, with the stated numerical range.

\paragraph{Homological algebra: exactness is a substantive step.}
The characteristic-two quotient $\ker z/\im z$ resembles homology.
The short nilpotent chains in Chapter~\ref{tx:ch:prime} explain why
that formula must change for general $p$. The exact-functor vocabulary
in the Stacks Project~\cite{StacksExact} is helpful: a kernel, image or
quotient is information about a map, not just the dimension of an
object. Our deletion is proved exact on the specified class of
representations; arbitrary homology is not exact. Nor does an exact
operation automatically detect every nonzero target obstruction.
The separate survival argument is indispensable.

\paragraph{Combinatorial linear algebra: rook coordinates.}
Bowman, Doty and Martin~\cite{BDM} study integral Schur--Weyl duality
for partition algebras. Their setting gives background for encoding
equivariant maps through finite equality patterns and permutation
actions. Our zero-marginal correction is developed explicitly in the
rook chapter, with a signed calculation and a triangular elimination.
The reference is mathematical context; it does not replace the
required correction identity or its coefficient-by-coefficient check.

\paragraph{Duality and designs: a witness need not be a point.}
Buss's designs~\cite{BussDesign} make the linear-dual approach to NS
lower bounds concrete. Here the target class in a quotient is separated
by a linear functional. For PC, that functional must annihilate the
\emph{whole} consequence closure, not just axiom multiples. The example
$x^2=0$ with target $x$ in Chapter~\ref{tx:sec:model-theory} distinguishes
ordinary ideal membership from vanishing on solutions. This is the
commutative-algebra distinction between an ideal and its radical;
see Lang~\cite{LangAlgebra} for the general algebraic background.

\paragraph{Topology and logic: two distinct meanings of compactness.}
The residue-tree construction in Section~\ref{tx:sec:topological-spectrum}
is the inverse-limit description
$\mathbb Z_p=\varprojlim_L\mathbb Z/p^L\mathbb Z$.
Gouv\^ea~\cite{Gouvea} is an accessible introduction to the surrounding
$p$-adic mathematics. Compactness here glues local \emph{degree bounds}
on congruence neighbourhoods. It does not assert a solution of the
original equations, nor require a nonstandard model or a Boolean
encoding. Kraj\'i\v cek's Euler structures~\cite{KrajEuler} and the
Grothendieck-ring viewpoint of Kraj\'i\v cek--Scanlon~\cite{KrajScanlon}
are a separate semantic route: they replace finite counting by a
characteristic on definable sets, with essential fibre and precision
hypotheses. No general converse to that criterion is asserted here.

\clearpage
\begin{thebibliography}{99}
\small
\setlength{\itemsep}{0pt}\setlength{\parsep}{0pt}
\bibitem{RiisTree}
S. Riis, \emph{A complexity gap for tree resolution},
Computational Complexity \textbf{10}, 179--209, 2001.
\href{https://doi.org/10.1007/s00037-001-8194-y}{doi:10.1007/s00037-001-8194-y}.

\bibitem{BussPitassiInduction}
S. R. Buss and T. Pitassi, \emph{Good degree bounds on Nullstellensatz
refutations of the induction principle}, Journal of Computer and System
Sciences \textbf{57}(2), 162--171, 1998.
\href{https://doi.org/10.1006/jcss.1998.1585}{doi:10.1006/jcss.1998.1585}.
\url{https://mathweb.ucsd.edu/~sbuss/ResearchWeb/modusponens/index.html}.

\bibitem{DantchevRiis}
S. Dantchev and S. Riis, \emph{On relativisation and complexity gap for
resolution-based proof systems}, in Computer Science Logic (CSL 2003),
Lecture Notes in Computer Science \textbf{2803}, 142--154, 2003.
\href{https://doi.org/10.1007/978-3-540-45220-1_14}{doi:10.1007/978-3-540-45220-1\_14}.

\bibitem{DMS}
S. Dantchev, B. Martin and S. Szeider, \emph{Parameterized proof complexity},
in Proceedings of the 48th Annual IEEE Symposium on Foundations of Computer Science
(FOCS 2007), 150--160, 2007.
\href{https://doi.org/10.1109/FOCS.2007.53}{doi:10.1109/FOCS.2007.53}.

\bibitem{DMRank}
S. Dantchev and B. Martin, \emph{Rank complexity gap for Lov\'asz--Schrijver
and Sherali--Adams proof systems}, Computational Complexity \textbf{22},
191--213, 2013.
\href{https://doi.org/10.1007/s00037-012-0049-1}{doi:10.1007/s00037-012-0049-1}.

\bibitem{QBFGap}
O. Beyersdorff, J. Clymo, S. Dantchev and B. Martin, \emph{The Riis complexity
gap for QBF resolution}, Journal on Satisfiability, Boolean Modeling and
Computation \textbf{15}(1), 9--25, 2024.
\href{https://doi.org/10.3233/SAT-231505}{doi:10.3233/SAT-231505}.

\bibitem{Riis}
S. Riis, \emph{On the asymptotic Nullstellensatz and Polynomial Calculus proof complexity}, in
Proceedings of the 23rd Annual IEEE Symposium on Logic in Computer Science
(LICS 2008), pp.~272--283, 2008.
\href{https://doi.org/10.1109/LICS.2008.30}{doi:10.1109/LICS.2008.30}.

\bibitem{RiisLong}
S. Riis, \emph{On the asymptotic Nullstellensatz and Polynomial Calculus proof complexity}, extended preprint,
Electronic Colloquium on Computational Complexity, Report TR09-014, 2009.
\url{https://eccc.weizmann.ac.il/report/2009/014/}.

\bibitem{Krajicek}
J. Kraj\'i\v cek, \emph{On the degree of ideal membership proofs from uniform
families of polynomials over a finite field}, Illinois Journal of
Mathematics \textbf{45}(1), 41--73, 2001.
\href{https://doi.org/10.1215/ijm/1258138254}{doi:10.1215/ijm/1258138254}.

\bibitem{Harman}
N. Harman, \emph{Stability and periodicity in the modular representation
theory of symmetric groups}, preprint, arXiv:1509.06414v3, 2016 (first version 2015).
\url{https://arxiv.org/abs/1509.06414}.

\bibitem{HKPSYZ}
C. Hone, F. Klein, B. Pauwels, A. Sherman, O. Yacobi and V. L. Zhang,
\emph{Semisimplifying categorical Heisenberg actions and periodic equivalences}, preprint,
arXiv:2509.07377v2, 2025.
\url{https://arxiv.org/abs/2509.07377}.

\bibitem{BDM}
C. Bowman, S. Doty and S. Martin, \emph{Integral Schur--Weyl duality for
partition algebras}, Algebraic Combinatorics \textbf{5}(2), 371--399, 2022.
\href{https://doi.org/10.5802/alco.214}{doi:10.5802/alco.214}.
\bibitem{KrajEuler}
J. Kraj\'i\v cek, \emph{Uniform families of polynomial equations over a
finite field and structures admitting an Euler characteristic of definable
sets}, Proceedings of the London Mathematical Society (3)
\textbf{81}(2), 257--284, 2000.
\href{https://doi.org/10.1112/S0024611500012557}{doi:10.1112/S0024611500012557}.

\bibitem{KrajScanlon}
J. Kraj\'i\v cek and T. Scanlon, \emph{Combinatorics with definable sets:
Euler characteristics and Grothendieck rings}, Bulletin of Symbolic Logic
\textbf{6}(3), 311--330, 2000.
\href{https://doi.org/10.2307/421058}{doi:10.2307/421058}.
\bibitem{CEF}
T. Church, J. S. Ellenberg and B. Farb,
\emph{FI-modules and stability for representations of symmetric groups},
Duke Mathematical Journal \textbf{164}(9), 1833--1910, 2015.
\href{https://doi.org/10.1215/00127094-3120274}{doi:10.1215/00127094-3120274}.
\url{https://arxiv.org/abs/1204.4533}.
\bibitem{CompressionCompanion}
S. Riis, \emph{Support compression for symmetric polynomial systems:
Permutation differences and preservation of proof degree},
unpublished companion manuscript and supporting proof sources, 2026.
\bibitem{StacksExact}
The Stacks Project Authors,
Lemma 12.7.2 (Tag 010N) in \emph{Additive functors}, and
Remark 35.4.4 (Tag 08WK) in \emph{Descent for universally injective morphisms}.
\url{https://stacks.math.columbia.edu/tag/010N};
\url{https://stacks.math.columbia.edu/tag/08WK}.
Optional background for the vocabulary of exactness and reflection.

\bibitem{LangAlgebra}
S. Lang, \emph{Algebra}, revised third edition,
Graduate Texts in Mathematics \textbf{211}, Springer, 2002.
\href{https://doi.org/10.1007/978-1-4613-0041-0}{doi:10.1007/978-1-4613-0041-0}.

\bibitem{AlperinLocal}
J. L. Alperin, \emph{Local Representation Theory:
Modular Representations as an Introduction to the Local
Representation Theory of Finite Groups},
Cambridge Studies in Advanced Mathematics \textbf{11},
Cambridge University Press, 1986.
\href{https://doi.org/10.1017/CBO9780511623592}{doi:10.1017/CBO9780511623592}.

\bibitem{Cutland}
N. Cutland, \emph{Computability: An Introduction to Recursive
Function Theory}, Cambridge University Press, 1980.
\href{https://doi.org/10.1017/CBO9781139171496}{doi:10.1017/CBO9781139171496}.


\bibitem{NSFoundations}
P. Beame, R. Impagliazzo, J. Kraj\'i\v cek, T. Pitassi and P. Pudl\'ak,
\emph{Lower bounds on Hilbert's Nullstellensatz and propositional proofs},
Proceedings of the London Mathematical Society (3)
\textbf{73}(1), 1--26, 1996.
\href{https://doi.org/10.1112/plms/s3-73.1.1}{doi:10.1112/plms/s3-73.1.1}.
Preliminary version in FOCS~1994, 794--806.

\bibitem{PCFoundations}
M. Clegg, J. Edmonds and R. Impagliazzo,
\emph{Using the Groebner basis algorithm to find proofs of unsatisfiability},
in Proceedings of the 28th Annual ACM Symposium on Theory of Computing
(STOC~1996), 174--183, 1996.
\href{https://doi.org/10.1145/237814.237860}{doi:10.1145/237814.237860}.

\bibitem{RiisSitharamTautologies}
S. Riis and M. Sitharam,
\emph{Generating hard tautologies using predicate logic and the symmetric group},
Logic Journal of the IGPL \textbf{8}(6), 787--795, 2000.
\href{https://doi.org/10.1093/jigpal/8.6.787}{doi:10.1093/jigpal/8.6.787}.

\bibitem{RiisSitharamModules}
S. Riis and M. Sitharam,
\emph{Uniformly generated submodules of permutation modules:
Over fields of characteristic 0},
Journal of Pure and Applied Algebra \textbf{160}(2--3), 285--318, 2001.
\href{https://doi.org/10.1016/S0022-4049(00)00069-4}{doi:10.1016/S0022-4049(00)00069-4}.

\bibitem{Ajtai}
M. Ajtai, \emph{Symmetric systems of linear equations modulo p},
Electronic Colloquium on Computational Complexity, Report TR94-015, 1994.
\url{https://eccc.weizmann.ac.il/report/1994/015/}.

\bibitem{BussDesign}
S. R. Buss, \emph{Lower bounds on Nullstellensatz proofs via designs},
in P. Beame and S. R. Buss (eds.), \emph{Proof Complexity and Feasible
Arithmetics}, DIMACS Series in Discrete Mathematics and Theoretical
Computer Science \textbf{39}, 59--71, American Mathematical Society, 1998.
\url{https://mathweb.ucsd.edu/~sbuss/ResearchWeb/designs/}.

\bibitem{BeameRiis}
P. Beame and S. Riis, \emph{More on the relative strength of counting
principles}, in P. Beame and S. R. Buss (eds.), \emph{Proof Complexity
and Feasible Arithmetics}, DIMACS Series in Discrete Mathematics and
Theoretical Computer Science \textbf{39}, 13--35,
American Mathematical Society, 1998.
\url{https://homes.cs.washington.edu/~beame/papers/riis.pdf}.

\bibitem{Gouvea}
F. Q. Gouv\^ea, \emph{$p$-adic Numbers: An Introduction}, third edition,
Universitext, Springer, 2020.
\href{https://doi.org/10.1007/978-3-030-47295-5}{doi:10.1007/978-3-030-47295-5}.
Optional background for the residue-tree and compactness viewpoint.

\end{thebibliography}

\end{document}
